% OpenSG-2.0 tutorial -- master document.
% Build:  cd docs/latex && bash build.sh   (pdflatex x2 -> opensg_tutorial.pdf)
\documentclass[11pt,letterpaper,oneside]{report}

\usepackage[margin=1in]{geometry}
\usepackage{amsmath}
\usepackage{amssymb}
\usepackage{graphicx}
\usepackage{booktabs}
\usepackage{longtable}
\usepackage{array}
\usepackage{xcolor}
\usepackage{listings}
\usepackage{fancyvrb}
\usepackage[font=small,labelfont=bf]{caption}
\usepackage[colorlinks=true,linkcolor=blue!55!black,urlcolor=blue!55!black,
            citecolor=blue!55!black]{hyperref}

% Long identifiers and file paths appear inline throughout. Allow a line
% break after each underscore so they wrap instead of overrunning the margin.
\setlength{\emergencystretch}{3em}
\sloppy
\AtBeginDocument{\def\_{\char95\allowbreak}}

\newcommand{\FIGDIR}{../_static}
\graphicspath{{../_static/}}

\definecolor{codebg}{rgb}{0.97,0.97,0.95}
\definecolor{codekw}{rgb}{0.10,0.30,0.65}
\definecolor{codestr}{rgb}{0.55,0.20,0.10}
\definecolor{codecom}{rgb}{0.35,0.45,0.35}

\lstdefinestyle{opensg}{
  language=Python,
  backgroundcolor=\color{codebg},
  basicstyle=\ttfamily\small,
  keywordstyle=\color{codekw}\bfseries,
  stringstyle=\color{codestr},
  commentstyle=\color{codecom}\itshape,
  showstringspaces=false,
  breaklines=true,
  breakatwhitespace=true,
  frame=single,
  rulecolor=\color{black!25},
  framesep=4pt,
  columns=fullflexible,
  keepspaces=true,
  captionpos=b,
}
\lstdefinestyle{shellst}{
  language=bash,
  backgroundcolor=\color{codebg},
  basicstyle=\ttfamily\small,
  keywordstyle=\color{black},
  showstringspaces=false,
  breaklines=true,
  frame=single,
  rulecolor=\color{black!25},
  framesep=4pt,
  columns=fullflexible,
  keepspaces=true,
  captionpos=b,
}
\lstset{style=opensg}

\setcounter{tocdepth}{2}
\setcounter{secnumdepth}{3}

\title{\vspace{-1.5cm}\Huge\bfseries OpenSG-2.0\\[0.35cm]
       \Large Mechanics of Structure Genome homogenization in JAX\\[0.2cm]
       \Large User Tutorial}
\author{Akshat Bagla \\ Purdue University}
\date{}

\begin{document}
\maketitle

\begin{abstract}
\noindent
OpenSG-2.0 computes the macroscopic constitutive law of a heterogeneous
structure -- a beam $6\times6$, a plate ABD, or an equivalent 3-D solid
$6\times6$ -- from a structure gene (SG): the smallest piece of the
structure that carries its heterogeneity. It also runs the calculation
backwards, recovering pointwise 3-D stress inside that SG from macroscopic
loads. Two engines share the theory and the file conventions: \texttt{msg-shell}
works on shell structure genes (a composite cross-section contour, a thin
sheet meshed as a surface, a tapered blade segment), and \texttt{msg-solid}
works on solid structure genes in one, two or three dimensions. Every run has
the same shape: a YAML (or SwiftComp \texttt{.sc}) file goes in, and a timed
output file in the SwiftComp \texttt{.K} layout comes back.

This tutorial covers the architecture of both engines, the input file formats
and sign conventions, the constitutive relations each capability produces, a
worked tutorial for every example shipped with the code, and the full wind
turbine blade pipeline from a windIO file to BeamDyn and back to recovered
wall stress.
\end{abstract}

\tableofcontents

\chapter{Overview and Architecture}
\label{ch:overview}

\section{What OpenSG-2.0 is}
\label{sec:what-is}

OpenSG-2.0 is an implementation of Mechanics of Structure Genome (MSG)
homogenization written in JAX. You give it a \emph{structure gene} (SG) --- the
smallest piece of a structure that carries its heterogeneity --- and it returns
the macroscopic constitutive law of the structure that SG builds: a Timoshenko
beam $6\times6$, a plate ABD (or the shear-refined $8\times8$), or an equivalent
3-D solid $6\times6$. It also runs the same calculation backwards
(\emph{dehomogenization}), rebuilding the pointwise 3-D stress and strain inside
the SG from macroscopic beam or plate loads.

The repository is organised in four layers. Only the top layer knows where files
live on disk.

\begin{table}[htbp]
\centering
\caption{The four layers of OpenSG-2.0. Paths and \texttt{main} blocks exist
only in the top row.}
\label{tab:layers}
\small
\begin{tabular}{@{}llp{7.4cm}@{}}
\toprule
Layer & Package & Role \\
\midrule
user     & \texttt{examples/}        & YAML in $\rightarrow$ \texttt{.out} out; owns all paths \\
msg-shell& \texttt{src/opensg\_shell}& Reissner--Mindlin shell SGs: ring (cross-section), aperiodic segment, 3-D shell cell \\
msg-solid& \texttt{src/opensg\_solid}& solid SGs: plate ABD, beam KKT, 3-D solid law \\
FE core  & \texttt{src/fe\_jax}      & JAX FE architecture: basis/quadrature, periodic assembly map, jitted element kernels, element-by-element (EBE) Chebyshev-CG, sparse direct \\
\bottomrule
\end{tabular}
\end{table}

\subsection{The universal contract}
\label{sec:contract}

Every capability in both engines obeys the same contract: \textbf{a YAML file (or
a SwiftComp \texttt{.sc} file) goes in, and a timed \texttt{.out} file in the
SwiftComp \texttt{.K} layout comes back.} There is exactly one writer for that
file, \texttt{opensg\_solid.sg\_homo.write\_sc\_K}, and it is called from inside
the solver source, not from the example scripts, so no run can forget it. A
minimal session looks like this:

\begin{lstlisting}[style=opensg,caption={The whole contract: the .out file is written as a side effect of the call.},basicstyle=\ttfamily\footnotesize]
from opensg_shell import build_rm_bundle

B = build_rm_bundle("iea_s10_shell.yaml")   # writes iea_s10_shell_Timo.out
B["Timo"]                                    # the Timoshenko 6x6 as a numpy array
\end{lstlisting}

The file that comes back always has the same shape: an \texttt{OpenSG} banner
naming the model and the SG measure $\omega$ on the first line, the effective
stiffness matrix, its inverse (the compliance), an engineering-constants block
for 3-D laws only, and \texttt{Time taken} on the last line. Numbers are printed
in a $19$-character right-justified Fortran field
(\verb|%.7E| mantissa, signed three-digit exponent) so the file drops straight
into existing SwiftComp comparison tooling:

\begin{Verbatim}[frame=single,fontsize=\small,samepage=false]
 OpenSG msg-shell beam model [ext sh2 sh3 twist bend2 bend3]

 The Effective Timoshenko Beam Stiffness Matrix
 --------------------------------------------
     <6 numbers per row>
 ...
 The Effective Timoshenko Beam Compliance Matrix
 --------------------------------------------
 ...
 Time taken: 12.34 sec
\end{Verbatim}

The suffix on the output file names the model, so two different routes run on
the same input YAML do not silently overwrite each other
(Table~\ref{tab:outfiles}).

\begin{table}[htbp]
\centering
\caption{Output file names. There is one file per model; the \texttt{\_Timo} vs
\texttt{\_C3D} suffix disambiguates when two routes can read the same YAML.}
\label{tab:outfiles}
\small
\begin{tabular}{@{}llp{5.2cm}@{}}
\toprule
Model & Entry function & Output file \\
\midrule
msg-solid beam / plate / 3-D & \texttt{plate\_homo\_2d} & \texttt{<base>.out} \\
msg-shell beam (Timoshenko)  & \texttt{build\_rm\_bundle} & \texttt{<yaml>\_Timo.out}, \texttt{<yaml>\_ABDG.out} \\
msg-shell aperiodic segment  & \texttt{segment\_timo\_from\_3dyaml} & \texttt{<yaml>\_Timo.out} \\
msg-shell solid props, cross-section SG & \texttt{build\_solid\_bundle} & \texttt{<yaml>\_C3D.out}, \texttt{<yaml>\_ABDG.out} \\
msg-shell 3-D shell SG       & \texttt{shell\_sg3d} & \texttt{<yaml>\_C3D.out}, \texttt{<yaml>\_ABDG.out} \\
\bottomrule
\end{tabular}
\end{table}

\subsection{One command: \texttt{opensg}}
\label{sec:opensg-cli}

You do not have to know which engine owns a file in order to run it. The
package installs one dispatcher, \texttt{opensg}, alongside the two
engine-specific commands \texttt{opensg\_solid} and \texttt{opensg\_shell};
all three take the same one or two arguments:

\begin{lstlisting}[style=shellst]
opensg <sg.yaml>          # homogenization -- the default
opensg <sg.yaml> H        # homogenization, stated explicitly
opensg <sg.yaml> D        # dehomogenization: homogenize, then recover
\end{lstlisting}

\noindent
\texttt{opensg} (\texttt{src/opensg/cli.py}) is a pure dispatcher and prints
nothing of its own. It resolves the engine with
\texttt{opensg\_solid.sg\_mesh.resolve\_msg} --- the YAML's \texttt{msg:} key
(\texttt{shell} or \texttt{solid}) or, when that key is absent, the mesh
\emph{dialect} (\texttt{nodes}/\texttt{cells}/\texttt{mat\_id} is solid,
\texttt{nodes}/\texttt{elements}/\texttt{sets}/\texttt{sections}/\texttt{elementOrientations}
is shell) --- and hands the untouched argument list to
\texttt{opensg\_shell.cli.main} or \texttt{opensg\_solid.cli.main}. A
\texttt{msg:} that contradicts the dialect is an error, not a silent
override, and a file written before the key existed still runs. The
engine-specific commands remain and behave identically, as do
\texttt{python -m opensg\_solid <yaml>} and
\texttt{python -m opensg\_shell <yaml>}. Everything else --- the macro
model, the classical/shear-refined switch, the macro state for a recovery
--- lives in the YAML header, which is why the command line needs no flags.

Every run prints the same header, so the terminal states what was actually
read before any number appears:

\begin{Verbatim}[frame=single,fontsize=\small,samepage=false]
============================================================================
OpenSG -- a multiscale structural analysis tool based on the Mechanics of
Structure Genome (MSG), developed by the Multiscale Structural Mechanics
Group led by Prof. Wenbin Yu at Purdue University
============================================================================

 input     : /abs/path/to/iea_s10_shell.yaml
 msg       : shell
 SG dim    : 2D
 analysis  : homogenization
 macro model: beam, shear-refined

Timoshenko Beam Stiffness Matrix  [eps11 gam12 gam13 kappa1 kappa2 kappa3]:
[[ 2.76606e+10 ... ]]
Homogenization stored in iea_s10_shell_Timo.out
Time taken: 2.58 sec
\end{Verbatim}

Three lines of that header repay reading. \texttt{msg} is the engine that
was resolved, so a mis-tagged file is visible on the first line of output.
\texttt{SG dim} is the \textbf{ambient} dimension --- the space the SG
occupies, \texttt{node\_span\_dim}, not the intrinsic dimension of the
manifold --- so a Schwarz-P shell surface prints \texttt{3D} while a
cross-section ring and a 2-D solid cell both print \texttt{2D}.
\texttt{macro model} echoes \texttt{n\_model} and \texttt{refined}, and gains
a trailing \texttt{, aperiodic} when the header sets \texttt{aperiodic}. For
a dehomogenization the two closing lines are instead
\texttt{Local field files are computed and stored.} and a single
\texttt{Time taken}, covering the whole run.

\subsection{The cores hold no paths}
\label{sec:no-paths}

Neither \texttt{opensg\_shell} nor \texttt{opensg\_solid} contains a
\texttt{main} block, a hard-coded directory, or an assumption about where you
run from. Every path decision lives in the example script, which names its input
file and calls exactly one entry function on it. That is why the tutorials in
the later chapters can be reduced to two lines of user input and one call: to
run a different cross-section you change the string, not the library.

Two pieces of process-level setup do live in the package \texttt{\_\_init\_\_}
files, and both matter to you:

\begin{itemize}
\item \textbf{Double precision is load-bearing.} Both packages execute
      \texttt{jax.config.update("jax\_enable\_x64", True)} at import. Without
      it JAX silently degrades to float32 and the digits of the $6\times6$ are
      wrong. Import the package (\texttt{import opensg\_shell}) before you
      import anything else that touches JAX.
\item \textbf{A persistent compile cache.} \texttt{opensg\_solid/\_\_init\_\_.py}
      points JAX's compilation cache at \verb|$JAX_COMPILATION_CACHE_DIR|,
      defaulting to \verb|~/.cache/opensg_jax|, with
      \texttt{jax\_persistent\_cache\_min\_compile\_time\_secs} set to $0$. A
      cold run pays each XLA compile once; warm runs load the binaries back. The
      cache is wrapped in a \texttt{try}/\texttt{except} --- it is an
      optimisation, and the engine must run without it. The same file honours
      \verb|$FE_JAX_CORE|, which prepends an external \texttt{fe\_jax} checkout
      to \texttt{sys.path} to override the bundled copy at \texttt{src/fe\_jax}.
\end{itemize}

\section{The theory in one page}
\label{sec:theory-one-page}

Every pipeline in this manual is the same variational statement, specialised. It
is worth reading this section once, because the symbols used here are the
variable names used in the source: \texttt{Dee}, \texttt{Dhe}, \texttt{Dhh},
\texttt{V0}, \texttt{V1}, \texttt{omega}.

\subsection{The structure gene and the strain split}

A structure gene is the smallest piece of the structure that carries the
heterogeneity: a laminate stacking sequence through the thickness (a 1-D SG), a
cross-section contour or a 2-D unit cell (a 2-D SG), a lattice or TPMS unit cell
(a 3-D SG). The displacement field inside the SG is split into a macroscopic
part driven by the macro strains $\bar\varepsilon$ and a fluctuation (warping)
field $w$. In strain form,

\begin{equation}
\varepsilon \;=\; \Gamma_e\,\bar\varepsilon \;+\; \Gamma_h\,w ,
\label{eq:split}
\end{equation}

where $\Gamma_e$ is the operator that embeds the macro strains into the SG and
$\Gamma_h$ is the SG's own strain operator acting on the warping. The two are
not independent: $\Gamma_e$ is $\Gamma_h$ evaluated on the macro embedding, a
consistency requirement written down in
\texttt{Rules/gamma\_e\_gamma\_h\_consistency.md} and enforced by building both
from the same element kernel.

\subsection{The quadratic form}

Substituting~\eqref{eq:split} into the strain energy gives the quadratic form
the code assembles, in exactly the block names used throughout the source:

\begin{equation}
2U \;=\; \bar\varepsilon^{T}\mathbf{D}_{ee}\,\bar\varepsilon
      \;+\; 2\,w^{T}\mathbf{D}_{he}\,\bar\varepsilon
      \;+\; w^{T}\mathbf{D}_{hh}\,w .
\label{eq:energy}
\end{equation}

Read it as three integrals over the SG: $\mathbf{D}_{ee}$ is the macro-macro
block (what you would get with no warping at all, the rule-of-mixtures answer),
$\mathbf{D}_{hh}$ is the SG stiffness seen by the warping alone, and
$\mathbf{D}_{he}$ couples them.

\subsection{The zeroth-order solve}

Minimising~\eqref{eq:energy} over $w$, subject to the SG's boundary treatment
(periodicity, Dirichlet boundary data, or a rigid-body constraint that removes
the null space) gives one linear system with as many right-hand sides as there
are macro strains,

\begin{equation}
\mathbf{D}_{hh}\,\mathbf{V}_0 \;=\; -\,\mathbf{D}_{he},
\label{eq:v0}
\end{equation}

and the effective law follows by substituting the minimiser back:

\begin{equation}
\mathbf{D}_{\mathrm{eff}} \;=\; \mathbf{D}_{ee} + \mathbf{V}_0^{T}\mathbf{D}_{he},
\qquad
\mathbf{C} \;=\; \mathbf{D}_{\mathrm{eff}}/\omega .
\label{eq:deff}
\end{equation}

Here $\omega$ is the SG measure that stays in the model --- the thickness of a
1-D laminate SG, the area of a cross-section SG, the volume of a 3-D solid cell,
the perimeter (or, in three dimensions, the midsurface area) of a shell SG. The
banner of every \texttt{.out} states $\omega$ because a bare matrix is
ambiguous: a whole-matrix ratio that is uniform across all entries (exactly
$0.5000$, say) is a normalisation difference, not a physics discrepancy.

$\mathbf{V}_0$ is stored, not discarded. Its columns are the warping fields per
unit macro strain, and they are what the dehomogenization reads back.

Equation~\eqref{eq:deff} is the \textbf{zeroth-order} answer: the plate ABD, the
Euler--Bernoulli $4\times4$, the equivalent 3-D solid law. It contains no
transverse shear, because the classical macro models it corresponds to have
none.

\subsection{The first-order step: transverse shear}

Recovering transverse shear --- the Reissner--Mindlin $G$ block of a plate, the
$GA_{12}$ and $GA_{13}$ entries of a Timoshenko $6\times6$ --- takes a
\textbf{first-order} step: a second solve for $\mathbf{V}_1$ against a
right-hand side built from $\mathbf{V}_0$, followed by an energy transformation.
Two functions do this, and they are shared by every beam route in the
repository:

\begin{itemize}
\item \texttt{prepare\_v1\_rhs(V0, Dhl, Dll, Dle\_dense, Psi, Dc)} forms the
      first-order right-hand side and the auxiliary blocks
      $\mathbf{V}_0^{T}\mathbf{D}_{ll}\mathbf{V}_0$,
      $\mathbf{D}_{hl}\mathbf{V}_0$ and
      $\mathbf{D}_{hl}^{T}\mathbf{V}_0 + \mathbf{D}_{le}$, projecting the
      right-hand side onto the range of $\mathbf{D}_{hh}$ so the singular
      system stays consistent.
\item \texttt{finalize\_v1\_and\_compute\_deff(V1s\_raw, V0, D\_eff, V0DllV0,
      DhlV0, DhlTV0Dle, Psi, Dc)} projects $\mathbf{V}_1$, symmetrises
      $\mathbf{D}_{\mathrm{eff}}$, and performs the generalized-Timoshenko
      transformation of the energy blocks $(A, B, C, D)$ into the sorted
      $6\times6$.
\end{itemize}

Both live in \texttt{opensg\_shell/fe\_jax/msg\_solver.py} and are
\texttt{@jax.jit}-decorated. The output ordering of the $6\times6$ is
\texttt{[EA, GA12, GA13, GJ, EI2, EI3]}, i.e.\ the strain measures
$[\varepsilon_{11}\;\;\gamma_{12}\;\;\gamma_{13}\;\;\kappa_1\;\;\kappa_2\;\;\kappa_3]$,
which is the VABS diagonal order.

The symmetrisation inside \texttt{finalize\_v1\_and\_compute\_deff} is not
cosmetic. $\mathbf{D}_{\mathrm{eff}}$ is a Hessian and therefore symmetric by
construction, but a Dirichlet-constrained tapered-segment solve produces a small
($2$ to $6$ percent) antisymmetric artifact, because with $\mathbf{V}_0$ pinned
at the end rings $\mathbf{D}_{ee}+\mathbf{V}_0^{T}\mathbf{D}_{he}$ is no longer
the exact symmetric Schur complement. Left alone, that artifact propagates
through the inverse into the antisymmetric shear--bending coupling ($C_{26}$,
$C_{35}$) and inflates it by up to $78$ percent on anisotropic tapers while
leaving the symmetric terms untouched. Enforcing the invariant is a no-op for a
prismatic SG.

\subsection{One factorization for two solves}

Equations~\eqref{eq:v0} and the $\mathbf{V}_1$ system share the same operator.
Every beam route therefore factorizes once and back-substitutes twice. In
\texttt{sg\_homo.ring\_indep} the augmented KKT matrix
\[
\begin{bmatrix}
\mathbf{D}_{hh} & \mathbf{D}_c\\[2pt]
\mathbf{D}_c^{T} & \mathbf{0}
\end{bmatrix}
\begin{Bmatrix}\mathbf{V}\\ \lambda\end{Bmatrix}
=
\begin{Bmatrix}\mathbf{R}\\ \mathbf{0}\end{Bmatrix}
\]
is built once, passed to \texttt{scipy.linalg.lu\_factor}, and then
\texttt{lu\_solve} is called with $\mathbf{R}=-\mathbf{D}_{he}$ for
$\mathbf{V}_0$ and with the \texttt{prepare\_v1\_rhs} output for
$\mathbf{V}_1$. In the tapered-segment route the same idea appears as
\texttt{sg\_assembly.dirichlet\_factor} followed by two
\texttt{dirichlet\_solve\_fac} calls; the second call falls back to a fresh
factorization only if the $\mathbf{V}_1$ Dirichlet set differs from the
$\mathbf{V}_0$ one. In the msg-solid beam engine the equivalent is
\texttt{solve\_fluctuation\_field}, which returns the assembled KKT matrix
alongside $\mathbf{V}_0$ precisely so the caller can reuse it; there the
constraint rows are scaled by $10^{8}$ before assembly to keep the saddle-point
system well conditioned, and empty rows are given a unit diagonal with a zeroed
right-hand side.

\subsection{Dehomogenization: running the chain backwards}

Given macro strains (or macro forces, converted through
$\bar\varepsilon = \mathbf{C}^{-1}\mathbf{F}$), the warping is rebuilt from the
stored $\mathbf{V}_0$ and $\mathbf{V}_1$, equation~\eqref{eq:split} is evaluated
pointwise, and the material law of each element turns that strain into 3-D
stress. Nothing is re-solved: homogenize once, recover as often as you like. For
the shell engine this is a \emph{two-step} recovery --- step 1 rebuilds the six
shell strains plus the two wall transverse shears
$[\varepsilon_{11}\;\varepsilon_{22}\;2\varepsilon_{12}\;\kappa_{11}\;\kappa_{22}\;2\kappa_{12}]$
and $[2\gamma_{13}\;2\gamma_{23}]$ on the ring, and step 2 pushes them through
the through-thickness plate SG to get ply-level 3-D stress. The entry points are
\texttt{opensg\_shell.stress\_at\_points} and
\texttt{opensg\_shell.disp\_at\_points} for the shell engine, and
\texttt{opensg\_solid.sg\_dehom.dehom\_fields} / \texttt{plate\_dehom\_2d} for
the solid engine.

\section{Module map}
\label{sec:module-map}

Both packages are split file-per-concern with the same naming, so a routine you
know in one engine is in the same-named file in the other.

\begin{table}[htbp]
\centering
\caption{\texttt{src/opensg\_shell} --- the MSG shell engine.}
\label{tab:mod-shell}
\small
\begin{tabular}{@{}lp{9.1cm}@{}}
\toprule
File & Concern \\
\midrule
\texttt{sg\_mesh.py} &
SG input loaders and mesh utilities: 1-D ring YAML loaders
(\texttt{load\_ring}, \texttt{load\_ring\_ref}) with explicit laminate-reference
conventions; 3-D segment loading and topological boundary extraction
(\texttt{extract}) into solver \texttt{.npz} bundles; orientation-frame numeric
reports and PNG prechecks. \\
\texttt{sg\_materials.py} &
Section laws: the per-section wall plate ABD $6\times6$ plus the $2\times2$
transverse-shear $G$ (\texttt{\_material\_by\_section}), and the per-station
wall-law emitter \texttt{emit\_station\_abd} / \texttt{load\_station\_abd} that
writes one windIO-style YAML of $8\times8$ Reissner--Mindlin plate laws per
cross-section. \\
\texttt{sg\_assembly.py} &
Shell element operators and assembly: the production 6-DOF drilling element, the
5-DOF general taper element, the prismatic MITC4 element, MITC assumed-strain
tying, the \texttt{compute\_k22} / \texttt{compute\_kg} initial-curvature
corrections, the Dirichlet factorizations
(\texttt{dirichlet\_factor}, \texttt{dirichlet\_solve\_fac}), and the six-block
MSG energy assembly (\texttt{assemble\_segment\_indep},
\texttt{assemble\_constraint}). \\
\texttt{sg\_periodicity.py} &
Periodic local-to-global sparse assembly map for a shell mesh, $6$ DOF per node
$(w_1,w_2,w_3,\omega_1,\omega_2,\omega_3)$ instead of the solid engine's three
translations. The tie is a plain index map; no transformation matrix is formed. \\
\texttt{sg\_homo.py} &
Homogenization drivers: \texttt{ring\_indep} (ring $\rightarrow$ Timoshenko
$6\times6$), \texttt{build\_rm\_bundle} (the packaged bundle the dehom
consumes), \texttt{ring\_solid} / \texttt{build\_solid\_bundle}
(cross-section $\rightarrow$ equivalent 3-D solid), \texttt{shell\_sg3d} (3-D
shell SG $\rightarrow$ \texttt{C3D}), \texttt{segment\_timo\_from\_3dyaml}
(aperiodic / tapered segment), \texttt{elastic\_constants}. \\
\texttt{sg\_dehom.py} &
Reissner--Mindlin two-step dehomogenization and recovery:
\texttt{stress\_at\_points}, \texttt{disp\_at\_points}, built on the same
element operators as the homogenization so it is the energy-consistent adjoint. \\
\texttt{sg\_junction.py} &
Level-2 junction law: local corner and micro-cell solves
(\texttt{corner\_micro\_law} for L/T/X topologies, \texttt{microcell\_law} for a
periodic mini-lattice) that correct the shell lattice for sub-shell junction
physics. \\
\texttt{fe\_jax/} &
The OpenSG-authored \texttt{msg\_*} kernels: \texttt{msg\_materials}
(ABD and $6\times6$ tables, reference shifts), \texttt{msg\_transverse\_shear}
(wall $G$, the $8\times8$ plate law), \texttt{msg\_mesh},
\texttt{msg\_rm}/\texttt{msg\_rm\_timo} (RM operators and the 5-DOF Timoshenko
assembly), \texttt{msg\_solver} (KKT solve, $\mathbf{V}_1$ RHS, Timoshenko
finalization), \texttt{msg\_hermite} (the Hermite $C^1$ Kirchhoff--Love
thin-walled pipeline), \texttt{msg\_dehom}, \texttt{orient\_plot}. \\
\texttt{helper/} &
Format converters (\texttt{ff\_to\_yaml.py}). \\
\bottomrule
\end{tabular}
\end{table}

\begin{table}[htbp]
\centering
\caption{\texttt{src/opensg\_solid} --- the general structure-gene engine.}
\label{tab:mod-solid}
\small
\begin{tabular}{@{}lp{9.1cm}@{}}
\toprule
File & Concern \\
\midrule
\texttt{sg\_materials.py} &
The $6\times6$ stiffness tables, per material and per element:
\texttt{build\_material\_C} and \texttt{get\_heterogeneous\_C\_matrix}, plus the
in-plane rotation. Voigt order everywhere is SwiftComp
$(xx, yy, zz, yz, xz, xy)$. \\
\texttt{sg\_mesh.py} &
SG input parsing and mesh utilities: \texttt{load\_sg\_input} (a SwiftComp
\texttt{.sc} or the helper's YAML), \texttt{\_cell\_basis} (SG dimension plus
nodes-per-element $\rightarrow$ basix cell type and basis degree), and
\texttt{plot\_sg\_mesh}, which draws the actual mesh coloured by material. \\
\texttt{sg\_assembly.py} &
Element kernels, periodic assembly, preconditioners and linear solvers:
\texttt{calculate\_RHS\_and\_Ke\_batch\_periodic},
\texttt{compress\_periodic\_cells\_jax}, \texttt{compute\_block\_inv\_diag},
\texttt{apply\_block\_precond}, \texttt{apply\_chebyshev\_precond},
\texttt{ebe\_jacobian\_product\_periodic},
\texttt{estimate\_max\_eigenvalue}, \texttt{\_sparse\_direct\_solve},
\texttt{solve\_fluctuation\_field},
\texttt{plate\_ladder\_element\_blocks}. \\
\texttt{sg\_homo.py} &
Homogenization drivers and the one public entry \texttt{plate\_homo\_2d}: the
periodic CG pipeline (\texttt{full\_homogenization\_pipeline},
\texttt{\_homo\_direct}) for \texttt{n\_model} 2 and 3, the beam KKT path
(\texttt{\_beam\_homo\_kkt}) for \texttt{n\_model} 1, the Reissner--Mindlin
first-order plate ladder (\texttt{plate\_shear\_ladder}), and the single output
writer \texttt{write\_sc\_K}. \\
\texttt{sg\_dehom.py} &
Dehomogenization: the Gauss-point recovery kernel for \texttt{n\_model} 2 and 3,
the generalized-Timoshenko beam recovery chain for \texttt{n\_model} 1, the
public \texttt{dehom\_fields} and \texttt{plate\_dehom\_2d}, and the Gauss
exports (\texttt{export\_gauss}). The batched recovery is wrapped in a
module-level \texttt{jit} so repeated calls in one process do not re-trace. \\
\texttt{rm\_plate\_1D/} &
The 1-D through-thickness plate stack used by both engines:
\texttt{segment\_plate.plate\_sg\_yaml} (build the through-thickness SG from a
stacking sequence), \texttt{msg\_rm\_plate.rm\_plate\_msg} (the $8\times8$ ABDG
and the warping ladders), \texttt{rm\_dehom} (ply-level recovery),
\texttt{rm\_homo}, \texttt{cli\_rm\_plate}. \\
\bottomrule
\end{tabular}
\end{table}

The shared FE core \texttt{src/fe\_jax} is not an OpenSG concern but a general
JAX finite-element architecture. The pieces the two engines actually import are
\texttt{basis\_quadrature} (\texttt{FiniteElementType}, \texttt{get\_quadrature},
\texttt{eval\_basis\_and\_derivatives}, all backed by basix),
\texttt{setup.mesh\_to\_jax} and
\texttt{setup.mesh\_to\_periodic\_sparse\_assembly\_map},
\texttt{solve\_cg} (a conjugate-gradient solver that returns convergence
information), \texttt{linear\_elasticity}, and \texttt{profiling}.

\section{The pipelines}
\label{sec:pipelines}

Four routes are described here in the order a new user meets them. Each is a
sequence of calls that the entry function makes on your behalf; you only ever
call the entry function. Figure~\ref{fig:pipelines} puts all four side by side:
one input file, one entry function, one \texttt{.out} file per branch.

\begin{figure}[htbp]
\captree{
  \node[funcnode] (f1) {\texttt{sg\_homo.build\_rm\_bundle}\\ \texttt{sg\_mesh.load\_ring\_ref} (it calls\\ \texttt{\_material\_by\_section} and\\ \texttt{compute\_k22}) $\rightarrow$ \texttt{rm\_plate\_msg}\\ $\rightarrow$ \texttt{sg\_homo.ring\_indep}};
  \node[inputnode, left=5mm of f1] (r1) {1-D shell SG ring YAML\\ \texttt{iea\_s10\_shell.yaml}};
  \node[outnode, right=5mm of f1] (o1) {\texttt{iea\_s10\_shell\_Timo.out}\\ \texttt{iea\_s10\_shell\_ABDG.out}\\ the Timoshenko $6\times6$};

  \node[funcnode, below=8mm of f1] (f2) {\texttt{sg\_homo.segment\_timo\_from\_3dyaml}\\ \texttt{sg\_mesh.extract} end rings $\rightarrow$\\ \texttt{ring\_indep} per side $\rightarrow$\\ \texttt{\_boundary\_dirichlet} $\rightarrow$\\ \texttt{dirichlet\_factor}};
  \node[inputnode, left=5mm of f2] (r2) {3-D shell segment YAML\\ \texttt{BAR\_URC\_numEl\_52\_}\\ \texttt{segment\_12.yaml}};
  \node[outnode, right=5mm of f2] (o2) {\texttt{BAR\_URC\_numEl\_52\_}\\ \texttt{segment\_12\_Timo.out}\\ per unit length, $L$ in the banner};

  \node[funcnode, below=8mm of f2] (f3) {\texttt{sg\_homo.shell\_sg3d}\\ one wall law for every element,\\ junction-edge count, periodic or\\ aperiodic boundary, one \texttt{splu}};
  \node[inputnode, left=5mm of f3] (r3) {3-D shell SG YAML\\ \texttt{schwarz\_p\_3Dshell.yaml}};
  \node[outnode, right=5mm of f3] (o3) {\texttt{schwarz\_p\_3Dshell\_C3D.out}\\ \texttt{schwarz\_p\_3Dshell\_ABDG.out}\\ the equivalent solid $6\times6$,\\ per unit-cell volume};

  \node[funcnode, below=8mm of f3] (f4) {\texttt{sg\_homo.plate\_homo\_2d}\\ \texttt{load\_sg\_input} $\rightarrow$ \texttt{\_cell\_basis} $\rightarrow$\\ \texttt{get\_heterogeneous\_C\_matrix} $\rightarrow$\\ \texttt{\_beam\_homo\_kkt} (\texttt{n\_model} 1) or\\ \texttt{\_homo\_direct} (\texttt{n\_model} 2, 3)};
  \node[inputnode, left=5mm of f4] (r4) {solid SG, a SwiftComp \texttt{.sc}\\ or a YAML:\\ \texttt{RHC\_SW\_2UC\_45.yaml}};
  \node[outnode, right=5mm of f4] (o4) {\texttt{RHC\_SW\_2UC\_45.out}\\ \texttt{n\_model} 1 named Euler-\\ Bernoulli Beam or Timoshenko\\ Beam; 2 Classical Plate or\\ Reissner-Mindlin; 3 Cauchy\\ Continuum};

  \draw[capedge] (r1) -- (f1); \draw[capedge] (f1) -- (o1);
  \draw[capedge] (r2) -- (f2); \draw[capedge] (f2) -- (o2);
  \draw[capedge] (r3) -- (f3); \draw[capedge] (f3) -- (o3);
  \draw[capedge] (r4) -- (f4); \draw[capedge] (f4) -- (o4);
}
\caption{The four pipelines, one row each: the three msg-shell routes and the
         msg-solid driver whose \texttt{n\_model} argument selects the macro
         model. Yellow is the file you edit, blue the entry function with the
         calls it makes in order, green the \texttt{.out} file it leaves
         behind. All four \texttt{.out} files are written by the one writer
         \texttt{opensg\_solid.sg\_homo.write\_sc\_K}, which is why they share
         a single layout; the \texttt{\_ABDG.out} companion --- written by
         every route that reduces a layup to a wall plate law, the
         cross-section rings and the 3-D shell SG alike --- is the one
         exception, written by \texttt{write\_abdg\_out}.}
\label{fig:pipelines}
\end{figure}

\subsection{msg-shell: 1-D ring $\rightarrow$ Timoshenko $6\times6$}
\label{sec:pipe-ring}

This is the composite blade cross-section route. The input is a 1-D shell SG
YAML: nodes and two-node elements tracing the wall contour,
\texttt{elementOrientations}, \texttt{sections} carrying a layup per element set,
and \texttt{materials}. The entry function is
\texttt{build\_rm\_bundle(shell\_yaml, ref=None, shear="mitc4\_g23",
g\_source=None)}, and it proceeds as follows.

\begin{enumerate}
\item \texttt{sg\_mesh.load\_ring\_ref} reads the YAML and returns the contour
      arrays --- nodal coordinates \texttt{rx}, connectivity \texttt{cells},
      the per-element section index \texttt{rsub}, the element normals
      \texttt{re3}, the axial index \texttt{ax} and the two cross-section
      indices \texttt{cross}. The laminate reference surface comes from the
      YAML's own \texttt{reference} field (absent $\rightarrow$ \texttt{center},
      the mid-surface), so homogenization and recovery cannot disagree; pass
      \texttt{ref} explicitly only to override. The reference is converted once
      into a thickness fraction \texttt{frac}
      (\texttt{center} $\rightarrow 0.5$, \texttt{oml} $\rightarrow 0.0$,
      \texttt{oml\_flip} and \texttt{iml} $\rightarrow 1.0$) and that single
      number then drives the ring laminate reference, the plate-SG $z$
      reference, the emitted ABD and the recovery depth conversion.
\item \texttt{sg\_materials.\_material\_by\_section} builds, per section, the
      wall plate ABD $6\times6$ at the chosen reference and the $2\times2$
      transverse-shear $G$.
\item The wall $G$ of every section is then replaced by
      \texttt{rm\_plate\_msg}, the MSG/VAM second-order-energy (Yu-2002
      least-squares) projection in \texttt{opensg\_solid.rm\_plate\_1D}, run
      at the same \texttt{fraction=frac}. \textbf{This is the only route and
      there is no switch.} The retired \texttt{g\_source} selector is still
      accepted by the signature so existing callers keep working, but any
      value other than \texttt{msg} raises a \texttt{DeprecationWarning} and
      is \emph{ignored} (\texttt{sg\_materials.check\_g\_source}), and an
      emitted ABD cache records \texttt{g\_source: msg}. The
      complementary-energy (Whitney-1973) shear flow survives only as the
      internal numerical fallback on the rare laminate whose MSG $U^{*}$ fit
      is not positive definite --- a guard, not a user-selectable
      alternative. The per-section laws are written to
      \texttt{<yaml>\_ABDG.out}.
\item \texttt{sg\_assembly.compute\_k22} computes the hoop curvature per edge,
      the initial-curvature correction that makes a curved wall stiffer than a
      faceted approximation of it.
\item \texttt{sg\_homo.ring\_indep} assembles the ring as a one-element-deep
      strip of quads: the 6-DOF drilling element
      $(w_1,w_2,w_3,\theta_1,\theta_2,\omega_3)$ with an element-wise Lagrange
      multiplier enforcing the drilling constraint, and MITC tying of
      $\gamma_{23}$ only. (Under span invariance $\gamma_{13}$ is algebraic in
      the directors and carries no fluctuation gradient, so only $\gamma_{23}$
      pairs a differentiated displacement with a rotation.) The six energy
      blocks \texttt{Dhh, Dhe, Dee, Dhl, Dll, Dle} come out of
      \texttt{assemble\_segment\_indep} and the drilling rows
      \texttt{Gc, Gl, Ge} out of \texttt{assemble\_constraint}; all are divided
      by the artificial strip depth $h$.
\item The KKT matrix is factorized once; $\mathbf{V}_0$ and $\mathbf{V}_1$ are
      obtained by two back-substitutions, and
      \texttt{finalize\_v1\_and\_compute\_deff} returns the generalized-Timoshenko
      $6\times6$, symmetrised.
\item \texttt{write\_sc\_K} writes \texttt{<yaml>\_Timo.out}, and the returned
      bundle carries \texttt{Timo}, \texttt{V0}, \texttt{V1}, the contour
      geometry, the strip, the per-element layup names and the layup and
      material databases --- everything the two-step dehomogenization needs.
\end{enumerate}

The shipped example is
\texttt{examples/OpenSG\_shell/1\_get\_beam\_props\_from\_shell\_cross\_section},
whose input is \texttt{iea\_s10\_shell.yaml} (the IEA-22 blade at $r/R=0.2$):

\begin{lstlisting}[style=shellst]
cd examples/OpenSG_shell/1_get_beam_props_from_shell_cross_section
opensg iea_s10_shell.yaml
\end{lstlisting}

\noindent
The YAML carries \texttt{n\_model: 1}, \texttt{refined: 1} and
\texttt{msg: shell}, so that one command runs the Timoshenko route and writes
\texttt{iea\_s10\_shell\_Timo.out} (the Timoshenko stiffness and compliance)
and \texttt{iea\_s10\_shell\_ABDG.out} (the per-section wall laws). The
printed diagonal is $[EA, GA_2, GA_3, GJ, EI_2, EI_3]$. The shipped driver
\texttt{beam\_homo\_shell.py} makes the same call through
\texttt{build\_rm\_bundle} and additionally draws the two figures the CLI
does not --- \texttt{iea\_s10\_mesh.png} (the contour coloured by layup) and
\texttt{iea\_s10\_orient.png} (the compulsory $e_1/e_2/e_3$ orientation
figure, emitted by \texttt{auto\_emit}) --- so run the script when you want
those:

\begin{lstlisting}[style=shellst]
python beam_homo_shell.py
\end{lstlisting}

\subsection{msg-shell: 3-D segment $\rightarrow$ aperiodic Timoshenko $6\times6$}
\label{sec:pipe-segment}

The ring pipeline assumes the beam is prismatic: periodic along the axis, so one
cross-section suffices. A tapered segment is not periodic, so the periodicity
condition has to be replaced. The entry function is
\texttt{segment\_timo\_from\_3dyaml(seg\_yaml, workdir=None,
lam\_space="elem", shear="full")} and the input is a 3-D shell segment YAML
(nodes, quads, sections, materials).

\begin{enumerate}
\item \texttt{sg\_mesh.extract} finds the two end cross-sections
      \emph{topologically}: a mesh edge used by exactly one quad is a free edge,
      and the connected components of the free-edge graph are the end rings.
      Each is written out as a standalone 1-D SG YAML (so you can inspect or
      re-run it on its own) together with an orientation PNG, and the whole
      bundle --- segment coordinates, cells, subdomains, the three element frame
      vectors, and the node-to-segment map --- is saved to
      \texttt{<base>\_boundaries.npz}.
\item Each end ring is solved as a ring in its own right by
      \texttt{ring\_indep}, giving \texttt{C6}, \texttt{V0} and \texttt{V1} per
      side.
\item The full 3-D quad mesh is assembled with the same
      \texttt{assemble\_segment\_indep} and \texttt{assemble\_constraint}, with
      the element hoop curvature \texttt{k22\_e} and geometric curvature
      \texttt{kg\_e} computed from the element centroids and frames. The
      drilling multipliers are appended as extra rows and columns by
      \texttt{\_augment\_drilling}, giving a saddle-point system with zero
      penalty.
\item \texttt{\_boundary\_dirichlet} maps the ring warping fields onto the
      segment's end nodes. This is the boundary condition that replaces
      periodicity: $\mathbf{V}_0 = \mathbf{V}_{0,\mathrm{ring}}$ and
      $\mathbf{V}_1 = \mathbf{V}_{1,\mathrm{ring}}$ on both end sections.
\item \texttt{dirichlet\_factor} factorizes the constrained interior once;
      $\mathbf{V}_0$ and then $\mathbf{V}_1$ are solved against it. The
      Euler--Bernoulli block is
      $A_{EB} = (\mathbf{D}_{ee} + \mathbf{V}_0^{T}\mathbf{D}_{he})/L$ with $L$
      the axial extent of the nodes, and every energy block handed to
      \texttt{finalize\_v1\_and\_compute\_deff} is likewise divided by $L$, so
      the result is a per-unit-length stiffness.
\item \texttt{<yaml>\_Timo.out} is written with the segment length in its
      banner; the returned dict carries \texttt{S6}, \texttt{C6L}, \texttt{C6R},
      \texttt{L} and \texttt{solve\_time}.
\end{enumerate}

Two acceptance checks come with this route. On a \emph{prismatic} segment the
segment $6\times6$ must reproduce its own boundary ring $6\times6$; measured on
an isotropic and a $\pm45$ tube this identity holds to
$\le 5\times10^{-5}$ relative. On a genuinely tapered segment no identity holds,
and the check is instead that each diagonal term of the segment lies between the
two ring values --- the ratio $(S_{ii}-R_{ii})/(L_{ii}-R_{ii})$ printed by the
example must fall in $[0,1]$. The shipped case is a BAR-URC 100\,m blade
segment:

\begin{lstlisting}[style=shellst]
cd examples/OpenSG_shell/5_taper_segment_aperiodic_boundaries
python run_bar_urc_segment.py
\end{lstlisting}

\noindent
which reads \texttt{BAR\_URC\_numEl\_52\_segment\_12.yaml} and writes
\texttt{BAR\_URC\_numEl\_52\_segment\_12\_Timo.out},
\texttt{BAR\_URC\_numEl\_52\_segment\_12\_boundaries.npz}, the two boundary
YAMLs \texttt{boundary\_..\_L.yaml} and \texttt{boundary\_..\_R.yaml}, three
orientation PNGs, and the comparison table \texttt{bar\_urc\_segment12.dat}.

\subsection{msg-shell: 3-D shell SG $\rightarrow$ equivalent solid $C_{3D}$}
\label{sec:pipe-shell3d}

Here the SG is a surface in three dimensions --- a triply periodic minimal
surface, a lattice of sheets --- and the answer wanted is the $6\times6$ of the
equivalent homogeneous 3-D solid. The entry function is
\texttt{shell\_sg3d(yaml\_path, omega=None, drill\_pen=1.0e-3,
g\_source=None, solver="direct", boundary=None, shear="mitc")}.
\texttt{boundary=None} means \emph{read it from the YAML header}, where the
optional \texttt{aperiodic: 1} key is the opt-out --- the same key
\texttt{opensg\_solid.plate\_homo\_2d} reads, so the two CLIs share one
header contract; an explicit argument always wins.

\begin{enumerate}
\item The YAML's nodes, elements and \texttt{elementOrientations} are read; the
      element normal $e_3$ is columns 7 through 9 of the orientation rows.
      Triangles and quads may be mixed: they are assembled as two separate
      element batches (MITC3 and MITC4) scattered into the same global
      system, never padded into one another.
\item The single wall law (ABD at the mid-surface reference plus the $2\times2$
      $G$, again from \texttt{rm\_plate\_msg}, the only $G$ route) is built
      once and applied to every element.
\item Junction lines are counted: any shell edge shared by more than two
      elements. The count appears in the \texttt{.out} banner, because it tells
      you how much of the cell's stiffness comes from junction physics the
      plain shell lattice does not resolve.
\item The boundary treatment. With \texttt{boundary="periodic"}, the default, the
      periodic map ties opposite faces, edges and corners through the sparse
      assembly map, and the three rigid translations are removed by three
      area-weighted Lagrange rows. With \texttt{boundary="aperiodic"} the macro
      (affine) field is mapped onto the boundary instead, prescribing zero
      \emph{translational} fluctuation $w_1=w_2=w_3=0$ on every node of the
      bounding-box faces while the rotational DOFs stay free. Clamping the
      rotations as well over-stiffens the cut edges; freeing them was measured
      to move $C_{11}$ from $+12$ to $+7$ percent. The aperiodic result is a
      kinematic upper bound on a single cell and converges to the periodic one
      as the SG grows.
\item The batched $\Gamma_h$ and $\Gamma_e$ operators are evaluated at the Gauss
      points, contracted into the sparse $\mathbf{K}$, $\mathbf{D}_{he}$ and
      $\mathbf{D}_{ee}$, with a drilling stabilisation of
      \texttt{drill\_pen} $\times\,A_{11}$ on the $\omega_3$ DOF. One
      \texttt{scipy.sparse.linalg.splu} factorization then serves all six macro
      strain columns.
\item $\mathbf{D}_{\mathrm{eff}} = \mathbf{D}_{ee} +
      \mathbf{V}_0^{T}\mathbf{D}_{he}$ is symmetrised and divided by $\omega$.
      The default $\omega$ is the \emph{volume the equivalent continuum
      occupies}, i.e.\ the node bounding-box volume measured from the mesh
      --- the same convention \texttt{opensg\_solid} uses, and the same cell
      the periodic assembly map ties --- so the returned \texttt{C3D}, the
      printed matrix and \texttt{<yaml>\_C3D.out} are one and the same
      matrix, directly comparable with a solid \texttt{.K}. Pass
      \texttt{omega} explicitly (or set the YAML \texttt{omega:} header key)
      only when the continuum occupies something else; the midsurface
      \emph{surface area}, for a per-unit-area wall law, is returned
      alongside as \texttt{surface\_area}.
\item The step-1 record --- the $8\times8$ Reissner--Mindlin wall plate law
      and its compliance, one block per section --- is written to
      \texttt{<yaml>\_ABDG.out} by the same \texttt{write\_abdg\_out} the
      cross-section routes use.
\end{enumerate}

Because the wall law is a Reissner--Mindlin plate law, the sheet thickness is a
parameter of the material model rather than of the mesh: the same surface mesh
serves any thickness, which is what makes a thickness sweep cheap.
Figure~\ref{fig:schwarz} shows the Schwarz-P unit cell that the shipped example
uses.

\begin{figure}[htbp]
\centering
\includegraphics[width=0.55\textwidth]{\FIGDIR/schwarz_p_mesh.png}
\caption{The Schwarz-P triply periodic minimal surface meshed as a 3-D shell
structure gene. The surface carries a Reissner--Mindlin wall law, so sheet
thickness is a material parameter and one mesh serves an entire thickness
sweep. Produced by \texttt{schwarz\_solid\_props.py} in
\texttt{examples/OpenSG\_shell/4\_get\_solid\_props\_from\_shell\_3D\_SG}.}
\label{fig:schwarz}
\end{figure}

\begin{lstlisting}[style=shellst]
cd examples/OpenSG_shell/4_get_solid_props_from_shell_3D_SG
python make_schwarz_yaml.py          # generates schwarz_p_3Dshell.yaml
opensg schwarz_p_3Dshell.yaml        # homogenizes it
\end{lstlisting}

\noindent
The first command is the mesh generator and has no CLI equivalent; the second
is the plain homogenization, which the YAML's own header
(\texttt{n\_model: 3}, \texttt{msg: shell}) fully describes. The run prints
the junction-edge count, the number of DOF and the solve time, then the
$6\times6$; it writes \texttt{schwarz\_p\_3Dshell\_C3D.out} and
\texttt{schwarz\_p\_3Dshell\_ABDG.out}. The driver
\texttt{schwarz\_solid\_props.py} makes the same call through
\texttt{shell\_sg3d} and additionally writes \texttt{schwarz\_p\_mesh.png}.
The companion script
\texttt{compare\_boundary\_modes.py} runs the same cell under both boundary
treatments and tabulates them in \texttt{boundary\_compare.dat}. This route is
validated against its own solid counterpart: for the Schwarz-P cell the shell SG
sits within $-3$ percent on the normal moduli, $-6$ percent on the shears and
$+0.3$ percent on Poisson's ratio relative to the SwiftComp-exact msg-solid
answer.

\subsection{msg-solid: one driver, three macro models}
\label{sec:pipe-solid}

The msg-solid engine has a single public entry,

\begin{lstlisting}[style=opensg,caption={The msg-solid driver signature: n\_model selects which macro strain set is built.},basicstyle=\ttfamily\footnotesize]
from opensg_solid.sg_homo import plate_homo_2d

r = plate_homo_2d(sc_path,               # a .sc or .yaml SG file
                  material_param=None,   # (n_mat, 9) engineering override
                  angles=None,           # deg per material
                  n_model=None,          # 1 beam, 2 plate, 3 3-D solid
                  refined=None,          # 0 classical, 1 shear-refined
                  workdir=None,
                  elem_rotation=None,    # (E, 9) per-element direction cosines
                  solver="direct",       # "direct" | "cg"
                  shear_refined=False,   # legacy alias of refined=1 (plate)
                  plot=True,             # writes <base>_mesh.png
                  boundary=None,         # "periodic" (default) | "aperiodic"
                  density=None,          # (n_mat,) beam mass rho
                  omega=None)            # user SG measure (3-D solid)
\end{lstlisting}

\noindent
and the argument that selects the physics is \texttt{n\_model}: it decides which
macro strain set $\Gamma_e$ is built. \texttt{n\_model}, \texttt{refined},
\texttt{boundary} and \texttt{omega} all default to \texttt{None}, meaning
\emph{take it from the YAML header}, so a self-describing SG file needs no
arguments at all --- that is what makes \texttt{opensg <yaml>} equivalent to
the call. The steps shared by all three macro models are:

\begin{enumerate}
\item \texttt{sg\_mesh.load\_sg\_input} parses the \texttt{.sc} or YAML into a
      dict carrying \texttt{dim} (the SG dimension, 1, 2 or 3), \texttt{nodes},
      \texttt{cells}, \texttt{mat\_id} and the materials. That \texttt{dim} is
      a \emph{parse result}, not something you write: the input contract has
      no \texttt{dim:} key and no \texttt{scale:} key, and both the dimension
      and the SG measure $\omega$ are read from the mesh itself. With
      \texttt{plot=True} it also writes \texttt{<base>\_mesh.png} of the actual
      mesh, elements coloured by material. Mixed element types in one SG are
      rejected --- the solver runs one homogeneous batch.
\item \texttt{sg\_mesh.\_cell\_basis} maps (SG dimension, nodes per element) to
      a basix cell type and basis degree; \texttt{fe\_jax.basis\_quadrature}
      then supplies the quadrature points \texttt{xi\_qp}, weights \texttt{W\_q},
      basis values \texttt{phi\_qn} and parametric gradients
      \texttt{dphi\_dxi\_qnp}. The quadrature degree is $6$ for higher-order
      elements and $2$ for linear ones.
\item \texttt{fe\_jax.setup.mesh\_to\_jax} produces the element-node coordinate
      array \texttt{x\_end} of shape $(E, N, d)$, and
      \texttt{mesh\_to\_periodic\_sparse\_assembly\_map} produces the master
      connectivity and the DOF map that carries periodicity. Periodicity rides
      entirely in the assembly map; no transformation matrix is formed. With
      \texttt{boundary="aperiodic"} the map is the identity and instead every
      bounding-box-face node gets $w=0$ Dirichlet; that option is available for
      the plate and solid models with \texttt{solver="direct"} only.
\item \texttt{sg\_materials.get\_heterogeneous\_C\_matrix} builds the per-element
      $6\times6$ table \texttt{C\_ess} of shape $(E,6,6)$, applying
      \texttt{angles} and \texttt{elem\_rotation} (per-element frames are
      honoured by all three models, not only the beam).
\end{enumerate}

\noindent
From there the routes diverge:

\begin{description}
\item[\texttt{n\_model=1}, beam.] Control goes to \texttt{\_beam\_homo\_kkt}:
      the four Euler--Bernoulli macro modes
      $[\varepsilon_{11}\;\kappa_1\;\kappa_2\;\kappa_3]$ are solved under four
      rigid-body Lagrange constraints, then the first-order chain runs and the
      generalized-Timoshenko reduction produces the $6\times6$. The result dict
      carries \texttt{C\_eff} (the Timoshenko $6\times6$) and
      \texttt{C\_eff\_EB} (the classical $4\times4$ from the same run). The
      output is \texttt{<base>.out}, named \emph{Timoshenko Beam} with
      \texttt{refined: 1} and \emph{Euler-Bernoulli Beam} with
      \texttt{refined: 0}, with $\omega$ in the banner. This path needs an SG
      of dimension $\ge2$.
\item[\texttt{n\_model=2}, plate.] The macro strain set is
      $\bar\varepsilon = [\varepsilon_{11}\;\varepsilon_{22}\;2\varepsilon_{12}\;
      \kappa_{11}\;\kappa_{22}\;\kappa_{12}]$ and the answer is the plate ABD,
      written as \emph{Classical Plate}. With \texttt{refined: 1} (or the
      legacy \texttt{shear\_refined=True}) the first-order warping ladder
      \texttt{plate\_shear\_ladder} runs as well, the result gains
      \texttt{G\_msg} ($2\times2$), \texttt{ABDG} (the $8\times8$ with the
      shear block on the diagonal) and \texttt{A6\_ladder}, and the file is
      named \emph{Reissner-Mindlin}.
\item[\texttt{n\_model=3}, solid.] The macro strain set is the six 3-D strains
      and the answer is the equivalent solid $6\times6$, named \emph{Cauchy
      Continuum} and written with the
      orthotropic-approximated engineering constants appended
      ($E_1$, $E_2$, $E_3$, $G_{12}$, $G_{13}$, $G_{23}$, $\nu_{12}$,
      $\nu_{13}$, $\nu_{23}$).
\end{description}

Models 2 and 3 share the assembly and solve code (\texttt{\_homo\_direct}), so
switching between a plate law and a solid law on the same mesh is a
one-character change. The shipped examples are numbered by capability under
\texttt{examples/OpenSG-solid/}: \texttt{1\_get\_plate\_props\_from\_1DSG}
(a stacking sequence in \texttt{layup\_db.yaml} $\rightarrow$
\texttt{layup\_db\_plate\_homo.out}, run with
\texttt{python 1d\_sg.py layup\_db.yaml}),
\texttt{3\_get\_plate\_props\_from\_2D\_SG},
\texttt{6\_get\_solid\_props\_from\_3D\_SG}, and
\texttt{7\_get\_beam\_props\_from\_SG}, whose script is:

\begin{lstlisting}[style=shellst]
cd examples/OpenSG-solid/7_get_beam_props_from_SG
python beam_homo_sg.py
\end{lstlisting}

\noindent
That one reads \texttt{RHC\_SW\_2UC\_45.yaml} (a two-unit-cell honeycomb
cross-section) with \texttt{n\_model = 1} and three materials overridden through
\texttt{material\_param} and \texttt{angles}, and writes
\texttt{RHC\_SW\_2UC\_45.out} plus \texttt{RHC\_SW\_2UC\_45\_mesh.png}.

\section{GPU readiness}
\label{sec:gpu}

The msg-solid engine has two interchangeable linear solvers, and the choice is
the \texttt{solver} argument of \texttt{plate\_homo\_2d}.

\texttt{solver="direct"} is the default. The jitted element kernels build the
sparse matrix, which is handed as a CSR to \texttt{\_sparse\_direct\_solve} ---
pypardiso (Intel MKL PARDISO) when it is importable, SciPy's SuperLU otherwise.
One factorization serves all macro strain columns. This is the fastest route on
a workstation, and it is the only stage of the calculation that leaves the
accelerator.

\texttt{solver="cg"} is the GPU route. The element kernels
(\texttt{calculate\_RHS\_and\_Ke\_batch\_periodic} and the jitted
\texttt{jacfwd} stiffness), the element-by-element operator
(\texttt{ebe\_jacobian\_product\_periodic}), the block and Chebyshev
preconditioners (\texttt{compute\_block\_inv\_diag},
\texttt{apply\_block\_precond}, \texttt{apply\_chebyshev\_precond},
\texttt{estimate\_max\_eigenvalue}) and the CG solver itself
(\texttt{fe\_jax.solve\_cg}) are pure JAX and run device-resident: no global
matrix is ever formed, so memory scales with the element batch rather than with
the bandwidth of an assembled operator. On a GPU host,
\texttt{plate\_homo\_2d(..., solver="cg")} therefore keeps the entire solve on
the device.

Switching solvers is digit-safe. The effective law
$\mathbf{C}_{\mathrm{eff}}$ is \emph{stationary} in $\mathbf{V}_0$ ---
equation~\eqref{eq:deff} evaluates the energy at its own minimiser --- so an
error $\delta$ in the fluctuation field enters the answer at
$O(\delta^2)$. A solver tolerance loose enough to be visible in $\mathbf{V}_0$
is still invisible in the printed matrix. Both routes are documented to produce
the same digits.

The msg-shell routes are one step behind. Their operators are assembled as
batched einsum contractions in the two-step form
$\mathbf{C}\mathbf{B}$ then $\mathbf{B}^{T}(\mathbf{C}\mathbf{B})$, and those
batches are written against the NumPy/\texttt{jnp} API and are portable to the
device; the sparse factorization (SuperLU for the ring and the 3-D shell SG,
pypardiso for the KKT solve in \texttt{msg\_solver}) is the remaining CPU stage.
For the SG sizes these routes handle --- a cross-section contour is a few
thousand DOF --- that factorization is not the bottleneck.

\chapter{Input Formats and Conventions}
\label{ch:input}

OpenSG-2.0 reads three kinds of structure-gene (SG) description, plus one
laminate shorthand and one legacy converter. Which dialect you write is decided
by the geometry you have, not by the answer you want: a thin-walled cross-section
contour is a 1-D shell SG, a sheet meshed as a surface in space is a 3-D shell
SG, and anything meshed as a filled domain (a laminate through its thickness, a
honeycomb unit cell, a lattice block) is a solid SG. The five dialects are listed
in Table~\ref{tab:input-dialects}.

\begin{table}[htbp]
\centering
\small
\caption{The five input dialects, the package that reads each, and the example
file shipped with the code.}
\label{tab:input-dialects}
\begin{tabular}{@{}lll@{}}
\toprule
dialect & read by & example file \\
\midrule
1-D shell SG (cross-section contour) & \texttt{opensg\_shell}
  & \texttt{iea\_s10\_shell.yaml} \\
3-D shell SG (surface in space) & \texttt{opensg\_shell}
  & \texttt{schwarz\_p\_3Dshell.yaml} \\
solid SG (1-D / 2-D / 3-D mesh) & \texttt{opensg\_solid}
  & \texttt{RHC\_SW\_2UC\_45.yaml} \\
SwiftComp \texttt{.sc} & converted to the solid SG YAML
  & \texttt{RHC\_SW\_2UC\_45.sc} \\
through-thickness layup database & \texttt{opensg\_solid.rm\_plate\_1D}
  & \texttt{layup\_db.yaml} \\
\bottomrule
\end{tabular}
\end{table}

The two shell dialects share one schema --- the same six top-level keys,
differing only in element dimension --- while the solid dialect is a separate,
deliberately minimal schema whose material blocks mirror SwiftComp's.

Every run, in every dialect, writes a timed \texttt{.out} in the SwiftComp
\texttt{.K} layout. The writer is \texttt{opensg\_solid.sg\_homo.write\_sc\_K}
and it always emits, in order: an \texttt{ OpenSG} banner line naming the model,
\texttt{The Effective <Name> Stiffness Matrix}, \texttt{The Effective <Name>
Compliance Matrix}, and a \texttt{ Time taken: <s> sec} footer. The
\texttt{<Name>} infix is \texttt{Timoshenko} for a beam $6\times6$,
\texttt{Classical Plate} for a $6\times6$ ABD, \texttt{Reissner-Mindlin Plate}
for an $8\times8$ ABDG, and empty for the equivalent 3-D solid law --- which
alone is written with \texttt{constants=True} and therefore also prints
\texttt{The Engineering Constants (Approximated as Orthotropic)}. The full
contract is \texttt{Rules/output\_format\_and\_timing.md} at the repository
root.

%========================================================================
\section{The 1-D shell SG YAML}
\label{sec:input-shell1d}

This is the cross-section of a thin-walled beam: the wall midline (or mold line)
meshed with 2-node line segments, each segment carrying a laminate. It is the
input to \texttt{opensg\_shell.build\_rm\_bundle}, which returns the Timoshenko
$6\times6$.

The example \texttt{iea\_s10\_shell.yaml} in
\texttt{examples/OpenSG\_shell/1\_get\_beam\_props\_from\_shell\_cross\_section/}
is station $r/R = 0.2$ of the IEA-22 blade: 307 nodes, 310 line elements, 6
layups, 6 materials. Its top-level keys, in the order OpenSG's own writers emit
them, are listed in Table~\ref{tab:input-shell-keys}.

\begin{table}[htbp]
\centering
\small
\caption{Top-level keys of the 1-D shell SG YAML.}
\label{tab:input-shell-keys}
\begin{tabular}{@{}lp{10.2cm}@{}}
\toprule
key & content \\
\midrule
\texttt{nodes} & one bracketed, space-separated triple per node:
  the two cross-section coordinates then the beam-axis coordinate \\
\texttt{elements} & one bracketed, space-separated pair per line element,
  1-based node ids \\
\texttt{sets} & \texttt{element:} a list of \texttt{\{name, labels\}} groups;
  \texttt{labels} are 1-based element indices \\
\texttt{sections} & one entry per layup, tied to a set through
  \texttt{elementSet}, carrying the \texttt{layup} ply table \\
\texttt{elementOrientations} & one row of nine direction cosines per element,
  a genuine YAML list of floats \\
\texttt{materials} & named orthotropic materials with \texttt{density} and an
  \texttt{elastic} block of \texttt{E}, \texttt{G}, \texttt{nu} triples \\
\texttt{reference} & one scalar naming the surface the wall laws are referenced
  to: \texttt{center}, \texttt{oml}, \texttt{oml\_flip} or \texttt{iml} \\
\bottomrule
\end{tabular}
\end{table}

\subsection{\texttt{nodes}}

One row per node, three coordinates. In the shell dialect the row is written as a
bracketed, \emph{space-separated} triple --- a single YAML scalar inside a flow
sequence --- and not as a comma list:

\begin{lstlisting}[style=shellst,caption={The first two node rows of iea\_s10\_shell.yaml.}]
nodes:
- [4.70852238 0.09722431 0.00000000]
- [4.64563825 0.11238283 0.00000000]
\end{lstlisting}

The space-separated form is required, not merely conventional: the orientation
plotter \texttt{fe\_jax/orient\_plot.py} parses each row with
\texttt{str(row[0]).split()}, so a comma-separated list of floats will not plot.
The loader \texttt{sg\_mesh.load\_ring\_ref} is more forgiving --- its
\texttt{\_row} helper replaces commas by spaces before splitting --- but writing
the space form keeps every consumer happy at once.

The component order is fixed and is the source of most first-run confusion. The
loader sets \texttt{ax = 2} and \texttt{cross = [0, 1]}, meaning the
\textbf{third column is the beam axis} and the first two columns are the
in-plane cross-section pair $(y_2, y_3)$. For a plane cross-section the beam-axis
column is identically zero, exactly as in the excerpt above. If you are writing a
cross-section by hand and your third column is not all zeros, you have written a
segment, not a cross-section.

\subsection{\texttt{elements}}

One row per element, the node ids of a 2-node line segment, 1-based, again
space-separated inside brackets:

\begin{lstlisting}[style=shellst,caption={The first two element rows; element k here is element label k everywhere else in the file.}]
elements:
- [1 2]
- [2 3]
\end{lstlisting}

The loader detects the base with \texttt{if cells.min() == 1: cells = cells - 1},
so a 0-based mesh is also accepted, but every generator in the repository writes
1-based. Element $k$ in this list is element label $k$ in \texttt{sets}, and
row $k$ of \texttt{elementOrientations} belongs to it. The contour need not be a
single loop: the last element of the IEA station is \texttt{- [307 65]}, which
closes a shear web back onto a skin node.

\subsection{\texttt{sets} and \texttt{sections}}

\texttt{sets} groups element labels by name; \texttt{sections} gives each named
group its laminate. The tie between them is the \texttt{elementSet} field, which
must equal a set \texttt{name}:

\begin{lstlisting}[style=shellst,caption={Abridged sets and sections blocks. In the real file layup\_0 lists 25 labels (elements 1-12 and 209-221) and there are six sets and six sections.}]
sets:
  element:
  - name: layup_0
    labels:
    - 1
    - 2
    - 3
sections:
- type: shell
  elementSet: layup_0
  layup:
  - - gelcoat
    - 0.0005
    - 0.0
  - - glass_triax
    - 0.00300294
    - 0.0
  - - glass_uniax
    - 0.02609467
    - 0.0
  - - glass_triax
    - 0.00300294
    - 0.0
\end{lstlisting}

Labels are 1-based element indices into \texttt{elements}.
\texttt{sg\_materials.\_material\_by\_section} builds one ABD per entry of
\texttt{sections} and keys it by the section's \emph{position in the list}; the
loader then maps every element through \texttt{rsub[label - 1] =
section\_index}. An element that appears in no set silently falls to section 0,
so every element should appear in exactly one set. \texttt{type: shell} is
descriptive metadata --- the shell readers key only off \texttt{elementSet} and
\texttt{layup}.

Each \texttt{layup} row is \texttt{[material, thickness, angle\_deg]}: a
material name that must exist in \texttt{materials}, a ply thickness in metres,
and a ply angle in degrees. The list is ordered \textbf{outer face first} and
stacks inward along $e_3$: \texttt{compute\_ABD\_matrix} documents its input as
layer thicknesses bottom to top with the reference surface at that bottom
(outer) face, and \texttt{emit\_station\_abd} records the same order in its
\texttt{convention} string as \texttt{plies OML->IML}. The sum of the row
thicknesses is the wall thickness used for the reference shift; for
\texttt{layup\_0} above that is $0.0005 + 0.00300294 + 0.02609467 + 0.00300294 =
0.03260055$~m, which is the \texttt{thickness} the cached ABD file reports for
that layup.

\subsection{\texttt{elementOrientations}}

One row of nine direction cosines per element, $[e_1(3)\;e_2(3)\;e_3(3)]$, in the
same component order as the node rows. Unlike \texttt{nodes} and
\texttt{elements}, these rows are true YAML lists of floats:

\begin{lstlisting}[style=shellst,caption={One orientation row: e1 = [0,0,1] out of plane, then the in-plane e2 and e3.}]
elementOrientations:
- [0.0, 0.0, 1.0, -0.9721541292349604, 0.2343423756204075, 0.0, -0.2343423756204075,
  -0.9721541292349604, 0.0]
\end{lstlisting}

The meaning of the triad is covered in Section~\ref{sec:input-orientation}.

\subsection{\texttt{materials}}

A list of named materials, each with a \texttt{density} and an \texttt{elastic}
block holding three-component \texttt{E}, \texttt{G} and \texttt{nu}:

\begin{lstlisting}[style=shellst,caption={One material block from iea\_s10\_shell.yaml.}]
materials:
- name: glass_triax
  density: 1940.0
  elastic:
    E:
    - 28211400000.0
    - 16238800000.0
    - 15835500000.0
    G:
    - 8248220000.0
    - 3491240000.0
    - 3491240000.0
    nu:
    - 0.497511
    - 0.18091
    - 0.27481
\end{lstlisting}

The triples are $[E_1, E_2, E_3]$, $[G_{12}, G_{13}, G_{23}]$ and $[\nu_{12},
\nu_{13}, \nu_{23}]$ --- the argument order of
\texttt{fe\_jax.msg\_materials.build\_stiffness\_6x6}, which fills the
orthotropic compliance
\[
S_{11} = 1/E_1,\quad S_{12} = -\nu_{12}/E_1,\quad S_{13} = -\nu_{13}/E_1,\quad
S_{23} = -\nu_{23}/E_2,
\]
\[
S_{44} = 1/G_{23},\quad S_{55} = 1/G_{13},\quad S_{66} = 1/G_{12},
\]
and inverts it to the $6\times6$ stiffness. Units follow the input; the examples
are SI, so moduli are in Pa, thicknesses in m and \texttt{density} in
kg/m$^3$. \texttt{density} feeds the section mass per unit area that
\texttt{emit\_station\_abd} reports as \texttt{mass\_per\_area} (62.892567
kg/m$^2$ for \texttt{layup\_0} of this station); it is not needed for the
stiffness.

\subsection{\texttt{reference}}

\begin{lstlisting}[style=shellst,caption={The whole reference block -- one scalar, the last line of the file.}]
reference: center
\end{lstlisting}

This single scalar names the surface the wall ABD --- and therefore the whole
cross-section model --- is referenced to, and it is the single source of truth.
\texttt{build\_rm\_bundle(shell\_yaml)} reads it when its \texttt{ref} argument
is left \texttt{None} via \texttt{ref = d.get("reference", "center")}, and the
value is immediately turned into a thickness fraction
\texttt{frac = \{"center": 0.5, "oml": 0.0, "oml\_flip": 1.0, "iml": 1.0\}} that
then drives the ring laminate reference, the plate-SG \texttt{z\_ref}, the
emitted ABD file, and the depth conversion used in stress recovery. Absent, it
defaults to \texttt{center}. The four accepted values and the fraction each maps
to are listed in Table~\ref{tab:input-reference}.

\begin{table}[htbp]
\centering
\small
\caption{Values of \texttt{reference} and the thickness fraction each maps to.}
\label{tab:input-reference}
\begin{tabular}{@{}llc@{}}
\toprule
value & meaning & fraction from the outer face \\
\midrule
\texttt{center} & contour is the laminate mid-surface; ABD shifted by $t/2$ & 0.5 \\
\texttt{oml} & contour is the outer mold line; laminate stacks inward, no shift & 0.0 \\
\texttt{oml\_flip} & diagnostic: reference moved the full thickness the other way & 1.0 \\
\texttt{iml} & contour is the inner mold line; laminate stacks outward & 1.0 \\
\bottomrule
\end{tabular}
\end{table}

Mid-surface meshes (tubes, ellipses, generated cross-sections) use
\texttt{center}; airfoil YAMLs written on the OML use \texttt{oml}. Getting this
wrong shifts $\mathbf{B}$ and $\mathbf{D}$ by a parallel-axis term of the wall
thickness and moves every bending and torsion entry of the beam $6\times6$, which
is why it is recorded in the file rather than passed at the call site.

\subsection{Running it}

\begin{lstlisting}[style=shellst]
cd examples/OpenSG_shell/1_get_beam_props_from_shell_cross_section
python beam_homo_shell.py
\end{lstlisting}

The driver sets \texttt{name = "iea\_s10"}, calls
\texttt{build\_rm\_bundle(name + "\_shell.yaml")}, and prints the $6\times6$
with the diagonal labelled \texttt{EA GA2 GA3 GJ EI2 EI3}:

\begin{lstlisting}[style=opensg,caption={The whole compute part of beam\_homo\_shell.py.}]
from opensg_shell import build_rm_bundle, auto_emit

name = "iea_s10"
B = build_rm_bundle(name + "_shell.yaml")
C6 = np.asarray(B["Timo"])
LBL = ["EA", "GA2", "GA3", "GJ", "EI2", "EI3"]
print("diagonal: " + "  ".join("%s=%.5g" % (LBL[i], C6[i, i]) for i in range(6)))

auto_emit(name + "_shell.yaml", out_png=name + "_orient.png")
\end{lstlisting}

Four files come back. \texttt{iea\_s10\_shell\_Timo.out} holds the Timoshenko
$6\times6$ and its compliance under the banner \texttt{ OpenSG msg-shell beam
model [ext sh2 sh3 twist bend2 bend3]}; for this station the diagonal reads
$EA = 2.766\times10^{10}$~N, $GA_2 = 7.186\times10^{8}$~N,
$GA_3 = 4.217\times10^{8}$~N, $GJ = 2.438\times10^{9}$~N\,m$^2$,
$EI_2 = 3.514\times10^{10}$~N\,m$^2$, $EI_3 = 6.813\times10^{10}$~N\,m$^2$, and
the footer records \texttt{ Time taken: 8.02 sec}.
\texttt{iea\_s10\_shell\_ABDG.out} holds one $8\times8$ RM wall law per section,
written by \texttt{sg\_homo.write\_abdg\_out} with the header
\texttt{rows [eps11 eps22 2eps12 | K11 K22 K12+K21 | 2g13 2g23]}, so you can read
off, for instance, the $A_{11} = 1.370\times10^{9}$~N/m and the transverse-shear
$G_{11} = 8.283\times10^{7}$~N/m of \texttt{layup\_0} directly.
\texttt{iea\_s10\_mesh.png} is the contour coloured by layup and
\texttt{iea\_s10\_orient.png} is the orientation check discussed in
Section~\ref{sec:input-orientation}. On the first run \texttt{build\_rm\_bundle}
additionally caches the per-station wall laws as
\texttt{abd/iea\_s10\_shell\_abd.yaml} at the same reference the YAML declared
(here written out as \texttt{reference: mid}); the write is guarded by an
existence test on that path, so later runs reuse the cached file rather than
recompute it.

%========================================================================
\section{The 3-D shell SG YAML}
\label{sec:input-shell3d}

Same six keys, but the mesh is a surface embedded in 3-D: \texttt{nodes} are
genuine three-dimensional points and \texttt{elements} are 3-noded triangles or
4-noded quads. There is no \texttt{reference} key --- \texttt{shell\_sg3d} calls
\texttt{\_material\_by\_section} with \texttt{center\_ref=True} unconditionally,
so every wall law is referenced to its mid-surface. This is the input to
\texttt{opensg\_shell.sg\_homo.shell\_sg3d}, which returns the equivalent 3-D
solid $6\times6$.

\texttt{schwarz\_p\_3Dshell.yaml} is a Schwarz-P triply-periodic minimal-surface
cell: 13536 nodes and 26360 triangles, one aluminium layup at $T = 0.036457$~m
with $E = 69$~GPa and $\nu = 0.3$. Triangles are read and padded to the quad
connectivity by repeating the last node,
\texttt{el = np.array([e + [e[-1]]*(4 - len(e)) for e in el], int) - 1}, so tri
and quad meshes go through one element path; ids are 1-based as in the 1-D
dialect. The mesh itself is shown in Figure~\ref{fig:input-schwarz-mesh}.

\begin{figure}[htbp]
\centering
\includegraphics[width=0.62\textwidth]{\FIGDIR/schwarz_p_mesh.png}
\caption{The Schwarz-P TPMS 3-D shell SG: 13536 nodes, 26360 triangles, one
aluminium layup. The mesh figure is drawn by the driver from the same YAML the
solver reads.}
\label{fig:input-schwarz-mesh}
\end{figure}

\texttt{make\_schwarz\_yaml.py} in the same folder is the canonical way to write
one, and its emission loop shows every row format at once. The frame is built per
triangle from the geometry: $e_3$ is the unit facet normal, $e_2$ is the unit
$p_1 \to p_2$ edge, and $e_1 = e_2 \times e_3$.

\begin{lstlisting}[style=opensg,caption={The emission loop of make\_schwarz\_yaml.py.}]
o = ["nodes:"]
o += ["- [%.12f %.12f %.12f]" % tuple(p) for p in nd]
o.append("elements:")
o += ["- [%d %d %d]" % tuple(e) for e in el]
o.append("elementOrientations:")
for k in range(ne):
    e1 = np.cross(e2[k], e3[k])
    o.append("- [%.6f, %.6f, %.6f, %.6f, %.6f, %.6f, %.6f, %.6f, %.6f]"
             % (*e1, *e2[k], *e3[k]))
o += ["sections:", "- type: shell", "  elementSet: layup_0", "  layup:",
      "  - - alu", "    - %.6e" % T, "    - 0.0",
      "materials:", "- name: alu", "  density: 2700.0", "  elastic:",
      "    E: [%.6e, %.6e, %.6e]" % (E, E, E),
      "    G: [%.6e, %.6e, %.6e]" % (G, G, G),
      "    nu: [%.6f, %.6f, %.6f]" % (nu, nu, nu),
      "sets:", "  element:", "  - name: layup_0", "    labels:"]
o += ["    - %d" % (e+1) for e in range(ne)]
open("schwarz_p_3Dshell.yaml", "w").write("\n".join(o) + "\n")
\end{lstlisting}

Note the asymmetry the loop makes explicit and that trips up hand-written files:
\texttt{nodes} and \texttt{elements} rows are space-separated inside brackets,
\texttt{elementOrientations} rows are comma-separated, and the \texttt{sets}
labels are 1-based. The generator prints the surface area and the relative
density $\text{area} \times T$ so the shell idealisation can be checked against
the solid cell it replaces.

\begin{lstlisting}[style=shellst]
cd examples/OpenSG_shell/4_get_solid_props_from_shell_3D_SG
python make_schwarz_yaml.py
python schwarz_solid_props.py
\end{lstlisting}

The second command writes \texttt{schwarz\_p\_3Dshell\_C3D.out}, whose banner
records the mesh size, the junction-edge count, the boundary treatment and the
bounding-box volume the law is normalized by. For this cell it reads
\texttt{ OpenSG msg-shell equivalent 3D solid (3-D shell SG, 13536 nodes, 26360
elems, 0 junction edges, periodic in 3 dirs, per unit cell 1)}. Because the 3-D
solid law is written with \texttt{constants=True}, the file ends with the
orthotropic-approximated engineering constants --- here
$E_1 = 9.3296\times10^{8}$, $E_2 = 9.3314\times10^{8}$,
$E_3 = 9.3285\times10^{8}$~Pa and $G_{12} = 1.0409\times10^{9}$~Pa --- followed
by \texttt{ Time taken: 66.43 sec}. The near-equality of the three Young's
moduli is the physical check: a Schwarz-P cell is cubic-symmetric, so any
appreciable spread means the mesh or the frames are wrong.

%========================================================================
\section{The solid SG YAML}
\label{sec:input-solid}

The solid dialect describes a meshed \emph{domain} rather than a laminated
\emph{surface}, so it carries no layups and no orientations by default --- the
material blocks are already $6\times6$ stiffnesses or engineering constants.
\texttt{RHC\_SW\_2UC\_45.yaml} is a two-unit-cell honeycomb plate SG with 4251
nodes, 6640 triangles and 3 materials (Figure~\ref{fig:input-rhc-mesh}). The
excerpt below keeps two rows per list and two of the three material blocks; its
keys are collected in Table~\ref{tab:input-solid-keys} and its material types in
Table~\ref{tab:input-mattypes}.

\begin{lstlisting}[style=shellst,caption={Abridged RHC\_SW\_2UC\_45.yaml.}]
dim: 2
scale: 35.248
nodes:
- [17.6240005, 23.3500004, 0.0]
- [-17.6240005, 23.3500004, 0.0]
cells:
- [0, 66, 345]
- [345, 3, 0]
mat_id: [1, 1, 1, 1, 1]
materials:
  1:
    type: 2
    aux: [0.0, 0.0]
    C:
    - [35696.6, 27696.6, 3251.1, 0.0, 0.0, -25368.7]
    - [27696.6, 35696.6, 3251.1, 0.0, 0.0, -25368.7]
    - [3251.1, 3251.1, 8917.8, 0.0, 0.0, -487.1]
    - [0.0, 0.0, 0.0, 3500.0, -500.0, 0.0]
    - [0.0, 0.0, 0.0, -500.0, 3500.0, 0.0]
    - [-25368.7, -25368.7, -487.1, 0.0, 0.0, 27958.5]
  3:
    type: 0
    aux: [0.0, 0.0]
    E: 69000.0
    nu: 0.3
\end{lstlisting}

\begin{table}[htbp]
\centering
\small
\caption{Keys of the solid SG YAML, as consumed by
\texttt{opensg\_solid.sg\_mesh.load\_sg\_input}.}
\label{tab:input-solid-keys}
\begin{tabular}{@{}lp{10.4cm}@{}}
\toprule
key & meaning \\
\midrule
\texttt{dim} & SG spatial dimension, 1, 2 or 3. Only the first \texttt{dim} node
  coordinates are used (\texttt{points = nodes[:, 0:n\_sg]}), so the trailing
  zeros above are padding. \\
\texttt{scale} & the normalization value carried over from the \texttt{.sc}
  trailing line. It is parsed and preserved on the input dict, but the
  homogenizer computes its own SG measure $\omega$ from the mesh and prints that
  in the banner. \\
\texttt{nodes} & comma-separated coordinate rows, always three components. \\
\texttt{cells} & element connectivity, \textbf{0-based} in this dialect (note
  the \texttt{0} in the first row). All elements in one SG must have the same
  node count; a mixed batch raises \texttt{mixed element types in one SG are
  unsupported}. \\
\texttt{mat\_id} & one material id per cell, \textbf{1-based} (the reader forms
  \texttt{mat\_id - 1} internally). \\
\texttt{materials} & a mapping from that id to a material block. \\
\bottomrule
\end{tabular}
\end{table}

Each material block carries a \texttt{type} that selects how the $6\times6$ is
built, matching the SwiftComp material types.

\begin{table}[htbp]
\centering
\small
\caption{Material types in the solid SG YAML and the \texttt{.sc} it mirrors.}
\label{tab:input-mattypes}
\begin{tabular}{@{}clp{7.6cm}@{}}
\toprule
\texttt{type} & fields & interpretation \\
\midrule
0 & \texttt{E}, \texttt{nu} & isotropic; $G = E/[2(1+\nu)]$ \\
1 & \texttt{engineering} & the nine constants
  $E_1\,E_2\,E_3\,G_{12}\,G_{13}\,G_{23}\,\nu_{12}\,\nu_{13}\,\nu_{23}$ \\
2 & \texttt{C} & the $6\times6$ stiffness stored directly --- typically an
  already-rotated ply, so no \texttt{angles} should be supplied with it \\
\bottomrule
\end{tabular}
\end{table}

\texttt{aux} is the auxiliary line of the material block, carried through
verbatim from the \texttt{.sc}; the elastic path does not consume it.

\begin{figure}[htbp]
\centering
\includegraphics[width=0.72\textwidth]{\FIGDIR/rhc_2dsg_mesh.png}
\caption{The RHC honeycomb 2-D solid SG, elements coloured by material id. This
figure is written automatically as \texttt{RHC\_SW\_2UC\_45\_mesh.png} by
\texttt{plot\_sg\_mesh} on every run with \texttt{plot=True} (the default), and
is drawn from the actual parsed mesh.}
\label{fig:input-rhc-mesh}
\end{figure}

The driver overrides the pre-rotated \texttt{type: 2} blocks with engineering
constants plus explicit per-material angles, which is the recommended pattern
when you want the ply angle to appear in exactly one place:

\begin{lstlisting}[style=opensg,caption={plate\_homo\_2dsg.py, the whole user input and the one call.}]
import jax.numpy as jnp
from opensg_solid.sg_homo import plate_homo_2d

n_model = 2                    # 1: Beam; 2: Plate; 3: 3D elastic
name = "RHC_SW_2UC_45"         # reads <name>.yaml

material_param = jnp.array([
    (108e3, 8e3, 8e3, 4e3, 4e3, 3e3, 0.32, 0.32, 0.30),      # ply +45
    (108e3, 8e3, 8e3, 4e3, 4e3, 3e3, 0.32, 0.32, 0.30),      # ply -45
    (69e3, 69e3, 69e3, 26.54e3, 26.54e3, 26.54e3, 0.30, 0.30, 0.30)])
angles = jnp.array([45.0, -45.0, 0.0])   # deg per material; 0.0 = none

r = plate_homo_2d(name + ".yaml", material_param=material_param,
                  angles=angles, n_model=n_model)
print(r["C_eff"])
\end{lstlisting}

\begin{lstlisting}[style=shellst]
cd examples/OpenSG-solid/3_get_plate_props_from_2D_SG
python plate_homo_2dsg.py
\end{lstlisting}

\texttt{n\_model} selects the macro model: 1 beam (Timoshenko), 2 plate, 3
equivalent 3-D solid. The run writes \texttt{RHC\_SW\_2UC\_45.out} ---
\texttt{The Effective Classical Plate Stiffness Matrix}, rows
\texttt{[N11 N22 N12 M11 M22 M12]}, under the banner
\texttt{ OpenSG msg-solid plate model, omega 35.248001, periodic} --- plus
\texttt{RHC\_SW\_2UC\_45\_mesh.png}. The printed $\omega = 35.248001$ is the
in-plane span measured from the mesh; it agrees with the file's
\texttt{scale: 35.248} to six digits, and that agreement is the check that the
mesh and the \texttt{.sc} header describe the same cell. For this run
$A_{11} = 3.3586\times10^{5}$, $D_{11} = 9.0368\times10^{7}$ in the file's own
units, and the footer reads \texttt{ Time taken: 7.83 sec}.

%========================================================================
\section{The SwiftComp \texttt{.sc} input and the converter}
\label{sec:input-sc}

A \texttt{.sc} file can be used directly --- \texttt{load\_sg\_input} dispatches
on the extension and converts --- or converted once and kept as YAML. The
anatomy the reader expects is documented in the
\texttt{opensg\_solid.rm\_plate\_1D.helper.sc\_to\_yaml} module docstring,
implemented in its \texttt{read\_sc}, and summarized in
Table~\ref{tab:input-sc}.

\begin{table}[htbp]
\centering
\small
\caption{Blocks of a SwiftComp \texttt{.sc} file, in file order.}
\label{tab:input-sc}
\begin{tabular}{@{}lp{10.2cm}@{}}
\toprule
block & content \\
\midrule
model header & 1 to 3 short lines (submodel flags and similar) \\
META line & \texttt{dim  n\_nodes  n\_elems  n\_mats} followed by further
  integers; located as the first line holding at least five integers \\
node records & \texttt{n\_nodes} lines of \texttt{id  x [y [z]]} \\
element records & \texttt{n\_elems} lines of \texttt{id  mat\_id  conn}, the
  connectivity zero-padded to a fixed slot count; the padding zeros are not
  nodes \\
material blocks & per material: a header \texttt{mat\_id  mat\_type  n}, an
  auxiliary line such as \texttt{0 0}, then the properties --- type 0 is
  \texttt{E  nu}, type 1 is \texttt{E1 E2 E3 / G12 G13 G23 / v12 v13 v23}, type
  2 is the 21 upper-triangular constants of the $6\times6$ over six lines \\
trailing line & the \texttt{scale} (volume or thickness normalization) \\
\bottomrule
\end{tabular}
\end{table}

In \texttt{RHC\_SW\_2UC\_45.sc} the META line is \texttt{2 4251 6640 3 0 0} --- a
2-D SG with 4251 nodes, 6640 elements and 3 materials --- and the file ends with
the lone line \texttt{35.248}. The file's last triangle, element 6640 of material
3 on nodes 4251/2988/3079, shows the zero padding in full:

\begin{lstlisting}[style=shellst]
6640 3 4251 2988 3079 0 0 0 0 0 0 0 0 0 0 0 0 0 0 0 0 0
\end{lstlisting}

The documented way to run the conversion is the module's own \texttt{convert}:

\begin{lstlisting}[style=opensg,caption={Converting a .sc by hand.}]
from opensg_solid.rm_plate_1D.helper.sc_to_yaml import convert

sc = convert("RHC_SW_2UC_45.sc")
print(sc["dim"], len(sc["nodes"]), len(sc["cells"]), sc["scale"])
\end{lstlisting}

\texttt{convert} writes \textbf{two} files next to the input ---
\texttt{<base>.yaml}, the solid SG YAML documented above, and
\texttt{<base>.msh}, a gmsh v2.2 mesh carrying the material id as both physical
and elementary tag --- and returns the parsed dict. The module defines no
\texttt{\_\_main\_\_} block, so import \texttt{convert} rather than invoking the
module with \texttt{python -m}. Because \texttt{plate\_homo\_2d} performs the
same conversion automatically when handed a \texttt{.sc} path, converting by hand
is only needed when you want to inspect or edit the YAML before homogenizing.

Two failure modes are worth knowing. The tet10 slot convention is remapped
explicitly (four corners, then the six midsides from slots 7--12), and any 3-D
record whose nonzero slot count is not 4, 8 or 10 is refused loudly rather than
converted into phantom nodes.

%========================================================================
\section{The through-thickness \texttt{layup\_db.yaml}}
\label{sec:input-layupdb}

For a plain laminate there is no mesh to write: the 1-D SG is one element stack
through the thickness, and \texttt{layup\_db.yaml} is the whole user input. The
driver \texttt{1d\_sg.py} reads it, generates the SG mesh (\texttt{1dsg.yaml}
plus \texttt{1dsg.png}) through
\texttt{opensg\_solid.rm\_plate\_1D.segment\_plate.plate\_sg\_yaml}, homogenizes
with \texttt{msg\_rm\_plate.rm\_plate\_msg}, and writes
\texttt{layup\_db\_plate\_homo.out}. The generated mesh is shown in
Figure~\ref{fig:input-plate1dsg} and the file's keys in
Table~\ref{tab:input-layupdb}.

\begin{figure}[htbp]
\centering
\includegraphics[width=0.62\textwidth]{\FIGDIR/plate_1dsg.png}
\caption{The generated through-thickness 1-D plate SG, elements coloured by
material: eight thin graphite/epoxy plies bracketing the thick PVC foam core.
The abscissa is $y_3$, running from $-h/2$ to $+h/2$ because
\texttt{fraction: 0.5} put the origin at mid-thickness.}
\label{fig:input-plate1dsg}
\end{figure}

\begin{table}[htbp]
\centering
\small
\caption{Keys of \texttt{layup\_db.yaml} read by \texttt{1d\_sg.py}.}
\label{tab:input-layupdb}
\begin{tabular}{@{}lp{9.8cm}@{}}
\toprule
key & meaning \\
\midrule
\texttt{model} & \texttt{0} prints the classical $6\times6$ ABD, \texttt{1} the
  shear-refined $8\times8$ ABDG. This also names the matrix in the
  \texttt{.out}: \texttt{Classical Plate} or \texttt{Reissner-Mindlin Plate}. \\
\texttt{fraction} & the reference surface as a fraction of the total thickness:
  \texttt{0} = bottom/OML face, \texttt{0.5} = mid-surface, \texttt{1} = top.
  This is the solid-side analogue of the shell YAML's \texttt{reference}. \\
\texttt{mesh.n\_per\_layer} & SG elements per physical ply. \\
\texttt{mesh.elem\_order} & element polynomial order $p$; nodes per element are
  $p+1$, and the engine supports 1 to 4. \texttt{4} --- the five-noded quartic
  --- is the default in the example and represents the whole warping ladder
  exactly ($V_0$ quadratic, $V_1$ cubic, $V_2$ quartic). \\
\texttt{materials} & a mapping from short name to \texttt{E} (3), \texttt{G}
  (3), \texttt{nu} (3) and \texttt{rho}, with an optional \texttt{full\_name}.
  Every number must parse as a float; the example's comment about
  \texttt{128.0e9} is a real PyYAML trap, which is why \texttt{1d\_sg.py}
  coerces with \texttt{float(v)} on read. \\
\texttt{layup} & the stacking sequence, \textbf{bottom ply first}, each entry
  \texttt{\{material, thickness, angle\}}. An optional \texttt{divisions} key is
  read only by the Abaqus 3-D deck generator and ignored by the homogenizer. \\
\bottomrule
\end{tabular}
\end{table}

The Nayak sandwich in the example is a nine-ply $(0/90/0/90/\text{core})_s$ stack
of graphite/epoxy faces on a HEREX C70.130 PVC foam core:

\begin{lstlisting}[style=shellst,caption={The mesh, reference and stacking blocks of layup\_db.yaml.}]
mesh:
  n_per_layer: 1
  elem_order: 4
fraction: 0.5
layup:
  - {material: ge,    thickness: 0.0009525, angle:  0.0}
  - {material: ge,    thickness: 0.0009525, angle: 90.0}
  - {material: ge,    thickness: 0.0009525, angle:  0.0}
  - {material: ge,    thickness: 0.0009525, angle: 90.0}
  - {material: herex, thickness: 0.1447800, angle:  0.0, divisions: 8}
  - {material: ge,    thickness: 0.0009525, angle: 90.0}
  - {material: ge,    thickness: 0.0009525, angle:  0.0}
  - {material: ge,    thickness: 0.0009525, angle: 90.0}
  - {material: ge,    thickness: 0.0009525, angle:  0.0}
\end{lstlisting}

\begin{lstlisting}[style=shellst]
cd examples/OpenSG-solid/1_get_plate_props_from_1DSG
python 1d_sg.py
\end{lstlisting}

\texttt{1d\_sg.py} accepts an alternative database as its one positional
argument, and the output file is named after that database rather than after the
mesh. The \texttt{.out} banner echoes back every number that mattered, which is
the fastest way to confirm the file was read as intended:
\texttt{ OpenSG msg-solid plate model from 1dsg.yaml: 9 plies, h = 0.152400 m,
fraction = 0.5, rho*h = 30.251400 kg/m\^{}2, rows [e11 e22 g12 k11 k22 k12 2g13
2g23]}. With \texttt{model: 1} the matrix printed is the $8\times8$, whose
transverse-shear block for this sandwich is $G_{11} = 7.7010\times10^{6}$ and
$G_{22} = 7.5424\times10^{6}$~N/m, and whose $B$ block is numerically zero
(entries of order $10^{-6}$ against a $10^{8}$ membrane term) as a symmetric
stack requires.

The generated \texttt{1dsg.yaml} records the same decisions in its own header,

\begin{lstlisting}[style=shellst,caption={Header, first node and first elements of the generated 1dsg.yaml (37 node rows and 9 element rows in full).}]
sg: {type: plate_1d, elem_order: 4, n_per_layer: 1, reference_fraction: 0.5,
  thickness: 0.15239999999999998, n_ply: 9}
nodes:
- [-0.07619999999999999]
elements:
- [1, 2, 3, 4, 5]
- [5, 6, 7, 8, 9]
\end{lstlisting}

and \texttt{read\_plate\_sg\_yaml} cross-checks the stored ply thicknesses and
\texttt{reference\_fraction} against the mesh itself, rejecting a hand-edited
file whose header and coordinates disagree instead of silently trusting it.

Note the element rows carefully, because this is the one place where the two
solid routes genuinely differ. \texttt{segment\_plate.plate\_sg\_yaml} numbers
each element's nodes \emph{sequentially}, \texttt{np.arange(p*e, p*e + p + 1)},
which is why the rows read \texttt{[1, 2, 3, 4, 5]} and \texttt{[5, 6, 7, 8, 9]}
--- consecutive integers, adjacent elements sharing their end node.
\texttt{rm\_plate\_msg} then homogenizes that mesh with its own equispaced
Lagrange basis (\texttt{nodes\_xi = np.linspace(-1.0, 1.0, p + 1)},
$\max(3, p+1)$-point Gauss-Legendre), so the \texttt{layup\_db.yaml} route never
touches the general SG engine's element table at all. The ends-first
gmsh/basix line ordering --- \texttt{[seq[0], seq[-1]] + seq[1:-1]} --- belongs
instead to \texttt{opensg\_solid.sg\_mesh.laminate\_to\_sg}, the parity route
that feeds the general engine, and it is that engine's
\texttt{plate\_homo\_2d} which consults \texttt{\_cell\_basis}
(Section~\ref{sec:input-elements}).

%========================================================================
\section{Finite elements used by each route}
\label{sec:input-elements}

The two packages discretize differently, and knowing which element you are
getting explains both the DOF count and the failure modes.

\subsection{Shell route}

The shell SG never assembles a line element directly. The contour is extruded one
quad deep along the beam axis by \texttt{sg\_homo.\_strip}, with extrusion length
$h$ equal to the mean contour element length, giving quads
\texttt{[a, b, m+b, m+a]}; the top node row is then DOF-mapped onto the bottom
row (\texttt{dof\_map = np.concatenate([np.arange(m), np.arange(m)])}) so the
field is exactly span-invariant. That is the operator's own prismatic reduction.
The element is therefore a 4-node bilinear quad integrated at $2\times2$ Gauss
points ($\pm 1/\sqrt{3}$ in each direction), and the 1-D cross-section problem is
recovered as its span-invariant limit. Two element families are built on that
quad, and they differ only in what they do with the drilling rotation
(Table~\ref{tab:input-shell-elements}).

\begin{table}[htbp]
\centering
\small
\caption{The two shell element families and how each handles the drilling
rotation.}
\label{tab:input-shell-elements}
\begin{tabular}{@{}lccp{5.6cm}@{}}
\toprule
entry point & DOF/node & DOF/elem & drilling \\
\midrule
\texttt{sg\_homo.ring\_indep} (production)
  & 6: \texttt{[w1, w2, w3, om1, om2, om3]} & 24
  & independent nodal $\omega_3$, tied back by a finite drilling residual
    (\texttt{NDOF6 = 6}) \\
\texttt{sg\_assembly.ring\_general}
  & 5: \texttt{[w1, w2, w3, om1, om2]} & 20
  & eliminated algebraically through
    $\omega_3 = S/(2C_{33}) - (C_{3b}/C_{33})\,\omega_b$ (\texttt{NDOF = 5}) \\
\bottomrule
\end{tabular}
\end{table}

The independent-drilling element exists because the elimination divides by
$C_{33} = n \cdot b_3$, which vanishes on a flat wall. In the 6-DOF element
$\omega_3$ is a genuine nodal DOF and the in-plane symmetry that defined it is
re-imposed in its finite, undivided form as the drilling residual
\[
DR = C_{33}\,\omega_3 + C_{3b}\,\omega_b - S/2 ,
\]
which equals $C_{33}(\omega_3 - \omega_3^{\text{eliminated}})$ and stays finite
even at $C_{33} = 0$. It is enforced by an element-wise Lagrange multiplier
(\texttt{lam\_space="elem"}, the production setting) or by a penalty. Transverse
shear uses MITC assumed-strain tying in the Dvorkin--Bathe sense: the production
ring scheme is \texttt{shear="mitc4\_g23"}, which ties only $\gamma_{23}$,
because under span invariance $\gamma_{13}$ is algebraic in the directors and
carries no fluctuation gradient.

The 3-D shell SG uses the same quad element, with triangles entering as quads
whose last node is repeated.

\subsection{Solid route}

\texttt{opensg\_solid.sg\_mesh.\_cell\_basis} maps the SG dimension and the node
count per element onto a basix cell type and Lagrange degree. This table
\emph{is} the list of supported elements, and it is consulted by exactly one
function: \texttt{sg\_homo.plate\_homo\_2d}, the general SG engine entry point
(and therefore also by anything that reaches it, including a \texttt{.sc} handed
to it directly and the \texttt{laminate\_to\_sg} parity route). The supported
combinations are given in Table~\ref{tab:input-cellbasis}.

\begin{table}[htbp]
\centering
\small
\caption{The \texttt{\_cell\_basis} table: every element the general solid SG
engine supports.}
\label{tab:input-cellbasis}
\begin{tabular}{@{}cccc@{}}
\toprule
\texttt{dim} & nodes/elem & basix cell & degree \\
\midrule
1 & 2 & \texttt{interval} & 1 \\
1 & 3 & \texttt{interval} & 2 \\
1 & 4 & \texttt{interval} & 3 \\
1 & 5 & \texttt{interval} & 4 \\
2 & 3 & \texttt{triangle} & 1 \\
2 & 4 & \texttt{quadrilateral} & 1 \\
3 & 4 & \texttt{tetrahedron} & 1 \\
3 & 8 & \texttt{hexahedron} & 1 \\
3 & 10 & \texttt{tetrahedron} & 2 \\
\bottomrule
\end{tabular}
\end{table}

Anything else raises \texttt{unsupported <dim>D element with <n> nodes}. Basis
functions are equispaced Lagrange (\texttt{LagrangeVariant.equispaced}) with
quadrature degree 6 for degree $>1$ and 2 otherwise. Two connectivity
conventions are applied on the way in: \texttt{\_to\_basix\_order} permutes quad4
to \texttt{[0, 1, 3, 2]} and hex8 to \texttt{[0, 1, 3, 2, 4, 5, 7, 6]}, while
triangles, tetrahedra and intervals pass through unchanged.

The through-thickness laminate of Section~\ref{sec:input-layupdb} does
\emph{not} go through this table. Its route ---
\texttt{plate\_sg\_yaml} to \texttt{rm\_plate\_msg} --- carries its own
equispaced Lagrange basis and its own sequential node numbering, as described
there. Only \texttt{laminate\_to\_sg}, the alternative route that packages a
laminate as a generic 1-D SG dict for \texttt{plate\_homo\_2d}, writes the
ends-first gmsh/basix line ordering that the \texttt{\_cell\_basis} elements
expect.

%========================================================================
\section{Sign and order conventions}
\label{sec:input-conventions}

Three orderings appear across the codebase. Every one of them is stated in a
module docstring rather than inferred, and every one of them is a common source
of a plausible-looking but wrong answer.

\subsection{Solid (3-D) Voigt order}

The solid order is the SwiftComp order
\[
[\,11\;\;22\;\;33\;\;23\;\;13\;\;12\,]
\]
with \textbf{engineering shears}. The compliance therefore carries
$S_{44} = 1/G_{23}$, $S_{55} = 1/G_{13}$, $S_{66} = 1/G_{12}$, and the shear rows
of a recovered strain are $2\gamma$, not $\gamma$. This is the order of
\texttt{write\_sc\_K}'s output, of \texttt{build\_stiffness\_6x6}, and of every
\texttt{type: 2} material \texttt{C} block in a solid SG YAML.

\subsection{Beam order}

The beam order is
\[
[\,\varepsilon_{11}\;\;\gamma_{12}\;\;\gamma_{13}\;\;\kappa_1\;\;\kappa_2\;\;
\kappa_3\,],
\]
whose diagonal is the familiar VABS ordering
\[
[\,EA\;\;GA_2\;\;GA_3\;\;GJ\;\;EI_2\;\;EI_3\,],
\]
the \texttt{LBL} list used by the shell drivers. \textbf{Axis 1 is the beam
axis}; axes 2 and 3 are the cross-section axes.

The solver indexes global axes as $(1,2,3) = (\text{axial},\,y_2,\,y_3)$ while an
SG YAML writes each 3-vector as $(y_2,\,y_3,\,\text{axial})$. That is why
\texttt{opensg\_solid.sg\_materials.elem\_rotation\_from\_yaml} cycles
\[
[a,\,b,\,c] \;\longrightarrow\; [c,\,a,\,b]
\]
(implemented as \texttt{o[:, :, [2, 0, 1]]}) before a YAML orientation can be
used as a direction-cosine matrix in the solid beam path. Skipping that cycle
points $e_1$ along $y_3$: the run still completes and still looks plausible, but
it destroys the section's symmetry, and a square tube comes out with
$EI_2 \neq EI_3$.

\subsection{Plate order}

The plate order is $[\,N_{11}\;N_{22}\;N_{12}\;M_{11}\;M_{22}\;M_{12}\,]$ for the
classical $6\times6$, conjugate to the strain vector
$[\varepsilon_{11}\;\varepsilon_{22}\;2\varepsilon_{12}\;\kappa_{11}\;
\kappa_{22}\;\kappa_{12}]$. The Reissner--Mindlin $8\times8$ appends the two
transverse shears, giving the rows
\texttt{[e11 e22 g12 k11 k22 k12 2g13 2g23]} that both
\texttt{write\_abdg\_out} and \texttt{1d\_sg.py} print in their banners, and the
block structure
\[
\mathbf{ABDG} =
\begin{bmatrix}
\mathbf{A} & \mathbf{B} & \mathbf{0} \\
\mathbf{B} & \mathbf{D} & \mathbf{0} \\
\mathbf{0} & \mathbf{0} & \mathbf{G}
\end{bmatrix}.
\]

For a plate SG the \textbf{last SG coordinate is the thickness direction} --- the
$y_3$ of the solid Voigt order above. With \texttt{dim: 2} that is node
coordinate index 1, since only the first \texttt{dim} coordinates are read; for a
through-thickness 1-D SG it is the single coordinate. This is why the RHC cell
reports $\omega = 35.248$, the span of coordinate 0 (the in-plane periodic
direction), and not its larger extent in coordinate 1.

\subsection{Changing the reference surface}

Moving the reference surface by $z_0$ maps the membrane strain
$\varepsilon$ to $\varepsilon + z_0\,\kappa$, so with
$\mathbf{T} = \begin{bmatrix} \mathbf{I} & z_0\mathbf{I} \\ \mathbf{0} &
\mathbf{I}\end{bmatrix}$ the plate law transforms as
\[
\mathbf{ABD}_{\text{new}} = \mathbf{T}^{-\mathsf{T}}\,\mathbf{ABD}\,
\mathbf{T}^{-1},
\]
which is exactly what \texttt{fe\_jax.msg\_materials.shift\_abd\_reference(ABD,
z0)} computes. Use it rather than reversing the layup: reversal flips $e_3$ and
breaks agreement with the material orientation. The transverse-shear block
$\mathbf{G}$ is reference-independent, which is why
\texttt{compute\_ABD\_matrix} can shift the $6\times6$ and leave the $8\times8$'s
$\mathbf{G}$ untouched.

%========================================================================
\section{Material orientation}
\label{sec:input-orientation}

Each row of \texttt{elementOrientations} is nine direction cosines forming three
stacked unit vectors,

\begin{lstlisting}[style=shellst]
[ e1x e1y e1z | e2x e2y e2z | e3x e3y e3z ]
\end{lstlisting}

with the roles fixed across both shell dialects:

\begin{itemize}
\item $e_1$ is the \textbf{beam axis} --- the prismatic direction, which for a
  plane cross-section meshed in $(y_2, y_3)$ is out of plane, so
  \texttt{e1 = [0, 0, 1]}.
\item $e_2$ is the \textbf{in-plane ply-flow direction}, the wall tangent along
  the contour.
\item $e_3 = e_1 \times e_2$ is the \textbf{wall normal}, in plane for a
  cross-section SG, and the direction the layup stacks along from the reference
  surface.
\end{itemize}

The triad must be orthonormal and right-handed.
\texttt{opensg\_shell.sg\_mesh.frame\_report(nodes, cells, e1, e2, e3)} checks
this numerically. It \textbf{returns} the pair \texttt{(ok, txt)} and prints
nothing --- callers such as \texttt{sg\_mesh.extract} do the printing --- and its
one-line summary ends in \texttt{OK} or \texttt{CHECK}. Its pass/fail gate is

\[
\bigl|\,\overline{|e_i|} - 1\,\bigr| < 10^{-6}\ (i = 1,2,3), \qquad
\max_{i\neq j}\bigl|e_i \cdot e_j\bigr| < 10^{-6},
\]
\[
\min\bigl(e_1 \times e_2 \cdot e_3\bigr) > 0.99, \qquad
\min\bigl(e_1 \cdot \hat{x}\bigr) > 0.999 .
\]

The first three are universal; the fourth --- $e_1$ axial --- is
cylinder-specific, as the docstring says. The report also prints a fifth number,
the mean of $e_3 \cdot (-\hat{r})$, i.e.\ how nearly $e_3$ is the inward radial
normal, but that value is \textbf{reported and not gated}: it appears in the text
so you can read it, and a poor value alone will not turn an \texttt{OK} into a
\texttt{CHECK}.

Because any mesh handling can renumber nodes and cells, an orientation PNG is a
compulsory deliverable of every OpenSG run. It is emitted by the \emph{example
driver} calling \texttt{opensg\_shell.auto\_emit} --- see the last line of
\texttt{beam\_homo\_shell.py} --- and not by any homogenization entry point:
\texttt{build\_rm\_bundle}, \texttt{shell\_sg3d} and \texttt{plate\_homo\_2d}
never call it. \texttt{auto\_emit} is cached once per \texttt{(shell, solid)}
path per process and swallows any exception with a
\texttt{[orient\_plot] skipped} message, so it never raises and is safe to call
from inside a compute path; a new driver of your own should call it explicitly.

The colour convention of \texttt{fe\_jax/orient\_plot.py} is \textbf{$e_2$ blue,
$e_3$ green}, drawn at element centroids over a light mesh; $e_1$ is out of the
page and so not drawn. (The separate 3-D plotters in \texttt{sg\_mesh.py} use
red/blue/black for $e_1$/$e_2$/$e_3$, since there $e_1$ is visible.) The
canonical example is the IEA-22 $r/R = 0.2$ station of
Figure~\ref{fig:input-orient}.

\begin{figure}[htbp]
\centering
\includegraphics[width=\textwidth]{\FIGDIR/shell_xsec_orient.png}
\caption{Material orientation of the IEA-22 $r/R = 0.2$ shell cross-section:
$e_2$ in blue (in-plane ply direction), $e_3$ in green (wall normal). This is
\texttt{iea\_s10\_orient.png}, written by the \texttt{auto\_emit} call at the end
of \texttt{beam\_homo\_shell.py}.}
\label{fig:input-orient}
\end{figure}

Read it as a physical check, not decoration. On this station the green $e_3$
arrows point inward everywhere on the skin, which is what an outward-first layup
stacked from the mid-surface requires, and the blue $e_2$ arrows are tangent to
every wall --- running consistently around the closed skin loop and along each of
the three shear webs. An $e_3$ that flips between neighbouring skin elements, or
a web whose $e_2$ opposes its neighbours, means the layup is being stacked in two
directions at once: visible here in seconds, and invisible in the resulting
$6\times6$.

\subsection{Where the ply angle lives}

The ply angle is \textbf{not} baked into $e_1$. It lives in the third column of
each \texttt{layup} row (shell dialect) or in the \texttt{angles} argument (solid
dialect), and it rotates the material about $e_3$, away from $e_1$. The rotation
\texttt{rotation\_6x6(theta)} leaves the Voigt-33 row untouched and mixes the
11/22/12 rows, and the reduced in-plane stiffness built by
\texttt{msg\_transverse\_shear.\_ply\_Q\_and\_G} is the textbook
\[
\bar{Q}_{11} = Q_{11}\cos^4\theta + 2(Q_{12} + 2Q_{66})\sin^2\theta\cos^2\theta
  + Q_{22}\sin^4\theta ,
\]
so \texttt{angle: 0.0} places the fibre along $e_1$ --- the beam axis. A
wind-blade unidirectional spar cap therefore carries \texttt{angle: 0.0}, and a
$\pm45$ ply carries \texttt{45.0} and \texttt{-45.0}.

The governing repository rule is
\texttt{Rules/orientation\_e1\_out\_of\_plane.md}, which also records the failure
mode: a YAML whose $e_1$ already contains the fibre rotation (so
$|e_{1z}| < 1$) is a \emph{fibre} frame, and combining it with a non-zero
\texttt{angle} applies the rotation twice. The rule supplies a two-line check
worth running on any YAML you did not write yourself:

\begin{lstlisting}[style=opensg,caption={The orientation precheck from Rules/orientation\_e1\_out\_of\_plane.md.}]
ori = np.array(d["elementOrientations"], float)
e1, e3 = ori[:, 0:3], ori[:, 6:9]
assert np.allclose(np.abs(e1[:, 2]), 1.0)   # e1 out of plane
assert np.allclose(e3[:, 2], 0.0)           # e3 in plane
\end{lstlisting}

If the first assertion fails, the file carries a fibre frame; either switch to
the frame variant built with the axial rotation only, or set the fibre angle to
zero downstream so the rotation is applied exactly once.

\chapter{Constitutive Relations Recovered by OpenSG}
\label{ch:constitutive}

Every OpenSG run answers the same question. Given a structure gene (SG) --- a
one-dimensional through-thickness stack, a two-dimensional cross-section or
cell, or a three-dimensional unit cell --- what is the constitutive law of the
macroscopic model that replaces it? The SG problem is solved once, and its
output is a small dense stiffness matrix relating \emph{generalized forces} to
\emph{generalized strains}. This chapter lists every such law the code
produces, the exact ordering of the strain and force components in each, the
function that computes it, the input file you edit, the command you type, and
the file the answer is written to.

The map is short. Every law in this chapter is one path through
the tree below: an SG of a given dimension, a model choice made either by the
\texttt{n\_model} argument of the solid engine or by which
\texttt{opensg\_shell} entry point is called, the macro law that comes back, and
the \texttt{.out} file it is written to. Solid lines are the default path of
each entry point; dashed lines are the alternative model choices that reuse the
same input file.

\captree{
\node[funcnode, text width=41mm] (f0)
  {\texttt{opensg\_solid.rm\_plate\_1D}\\ \texttt{.msg\_rm\_plate.rm\_plate\_msg}\\
   \texttt{model:} 0 or 1};
\node[inputnode, text width=29mm, left=5mm of f0] (i0)
  {\texttt{layup\_db.yaml}\\ 1-D SG: one laminate};
\node[datanode, text width=25mm, right=5mm of f0] (d0)
  {ABD $6\times6$ or\\ ABDG $8\times8$};
\node[outnode, text width=33mm, right=5mm of d0] (o0)
  {\texttt{<db>\_plate\_homo.out}};
\node[funcnode, text width=41mm, below=of f0] (f1)
  {\texttt{opensg\_solid.sg\_homo}\\ \texttt{.plate\_homo\_2d},
   \texttt{n\_model=1}\\ $\to$ \texttt{\_beam\_homo\_kkt}
   (needs \texttt{n\_sg} $\ge$ 2)};
\node[funcnode, text width=41mm, below=of f1] (f2)
  {\texttt{opensg\_solid.sg\_homo}\\ \texttt{.plate\_homo\_2d},
   \texttt{n\_model=2}\\ optionally \texttt{shear\_refined=True} $\to$
   \texttt{plate\_shear\_ladder}};
\node[funcnode, text width=41mm, below=of f2] (f3)
  {\texttt{opensg\_solid.sg\_homo}\\ \texttt{.plate\_homo\_2d},
   \texttt{n\_model=3}};
\node[inputnode, text width=29mm, left=5mm of f2] (i1)
  {solid SG\\ \texttt{<sg>.sc} or \texttt{<sg>.yaml}\\ \texttt{n\_sg} = 1, 2 or 3};
\node[datanode, text width=25mm, right=5mm of f1] (d1)
  {Timoshenko $6\times6$\\ (+ EB $4\times4$)};
\node[datanode, text width=25mm, right=5mm of f2] (d2)
  {ABD $6\times6$ or\\ ABDG $8\times8$};
\node[datanode, text width=25mm, right=5mm of f3] (d3)
  {3-D solid $6\times6$\\ + eng.\ constants};
\node[outnode, text width=33mm, right=5mm of d1] (o1) {\texttt{<sg>.out}};
\node[outnode, text width=33mm, right=5mm of d2] (o2) {\texttt{<sg>.out}};
\node[outnode, text width=33mm, right=5mm of d3] (o3) {\texttt{<sg>.out}};
\node[funcnode, text width=41mm, below=of f3] (f4)
  {\texttt{opensg\_shell.sg\_homo}\\ \texttt{.build\_rm\_bundle} $\to$
   \texttt{ring\_indep}};
\node[funcnode, text width=41mm, below=of f4] (f5)
  {\texttt{opensg\_shell.sg\_homo}\\ \texttt{.build\_solid\_bundle} $\to$
   \texttt{ring\_solid}};
\node[inputnode, text width=29mm, left=5mm of f4] (i2)
  {1-D shell SG\\ \texttt{iea\_s10\_shell.yaml}\\ cross-section contour};
\node[datanode, text width=25mm, right=5mm of f4] (d4)
  {Timoshenko $6\times6$\\ + wall $8\times8$ per section};
\node[outnode, text width=33mm, right=5mm of d4] (o4)
  {\texttt{<yaml>\_Timo.out}\\ \texttt{<yaml>\_ABDG.out}};
\node[datanode, text width=25mm, right=5mm of f5] (d5) {3-D solid $6\times6$};
\node[outnode, text width=33mm, right=5mm of d5] (o5)
  {\texttt{<yaml>\_C3D.out}};
\node[funcnode, text width=41mm, below=of f5] (f6)
  {\texttt{opensg\_shell.sg\_homo}\\ \texttt{.shell\_sg3d}};
\node[funcnode, text width=41mm, below=of f6] (f7)
  {\texttt{opensg\_shell.sg\_homo}\\ \texttt{.segment\_timo\_from\_3dyaml}};
\node[inputnode, text width=29mm, left=5mm of f6] (i3)
  {3-D shell SG \texttt{.yaml}\\ surface mesh, or\\ tapered segment};
\node[datanode, text width=25mm, right=5mm of f6] (d6) {3-D solid $6\times6$};
\node[outnode, text width=33mm, right=5mm of d6] (o6)
  {\texttt{<yaml>\_C3D.out}};
\node[datanode, text width=25mm, right=5mm of f7] (d7)
  {Timoshenko $6\times6$};
\node[outnode, text width=33mm, right=5mm of d7] (o7)
  {\texttt{<yaml>\_Timo.out}};
\draw[capedge] (i0) -- (f0);  \draw[capedge] (f0) -- (d0);
\draw[capedge] (d0) -- (o0);
\draw[capedge2] (i1) -- (f1); \draw[capedge] (i1) -- (f2);
\draw[capedge2] (i1) -- (f3);
\draw[capedge] (f1) -- (d1);  \draw[capedge] (d1) -- (o1);
\draw[capedge] (f2) -- (d2);  \draw[capedge] (d2) -- (o2);
\draw[capedge] (f3) -- (d3);  \draw[capedge] (d3) -- (o3);
\draw[capedge] (i2) -- (f4);  \draw[capedge2] (i2) -- (f5);
\draw[capedge] (f4) -- (d4);  \draw[capedge] (d4) -- (o4);
\draw[capedge] (f5) -- (d5);  \draw[capedge] (d5) -- (o5);
\draw[capedge] (i3) -- (f6);  \draw[capedge2] (i3) -- (f7);
\draw[capedge] (f6) -- (d6);  \draw[capedge] (d6) -- (o6);
\draw[capedge] (f7) -- (d7);  \draw[capedge] (d7) -- (o7);
}

Every one of those files is printed by \texttt{write\_sc\_K}, except the
per-section wall law \texttt{<yaml>\_ABDG.out}, which
\texttt{write\_abdg\_out} writes block by block in the same number format. In
both cases the title above the matrix names the law, so the tree above and
Table~\ref{tab:outnames} together identify any output file the code produces.

Two conventions hold throughout and are worth internalising before reading any
matrix.

\begin{itemize}
\item \textbf{Units follow the input.} The single writer
  \texttt{write\_sc\_K} performs no unit conversion, so an SG given in metres
  with moduli in pascals returns N, N$\cdot$m and Pa, while an SG given in
  millimetres with moduli in MPa returns the corresponding mixed units (N/mm,
  N$\cdot$mm).
\item \textbf{Shears are engineering shears.} Every off-diagonal strain
  component carries the factor of two ($2\varepsilon_{12}$,
  $2\bar\Gamma_{23}$, $2\gamma_{13}$), which is why each shear stiffness
  appears once, not four times, in the quadratic strain energy.
\end{itemize}

\section{The \texttt{.out} contract}
\label{sec:out-contract}

Every homogenization in OpenSG-2.0 writes a timed text file in the SwiftComp
\texttt{.K} layout through one function,
\texttt{opensg\_solid.sg\_homo.write\_sc\_K}:

\begin{lstlisting}[style=opensg,caption={The single output writer. Every law in this chapter is printed by this call.}]
write_sc_K(path, C, solve_time=None, model="", constants=True, name="")
\end{lstlisting}

The file opens with an \verb| OpenSG <model>| banner, prints the effective
stiffness $\mathbf C$, then its exact inverse $\mathbf S = \mathbf C^{-1}$ (the
effective compliance), and closes with \verb| Time taken: <t> sec|. Numbers are
written by the helper \texttt{\_knum}, a 19-character right-justified
\verb|%.7E| mantissa with a signed three-digit exponent, which is exactly the
SwiftComp \texttt{.K} column format.

The \texttt{name} argument is the only thing that changes the matrix title, and
it identifies which macro law you are looking at (Table~\ref{tab:outnames}).

\begin{table}[htbp]
\centering
\caption{The \texttt{name} argument of \texttt{write\_sc\_K} and the title it
prints. The title is the fastest way to tell which law an \texttt{.out} file
holds.}
\label{tab:outnames}
\footnotesize
\begin{tabular}{lll}
\toprule
\texttt{name} passed to \texttt{write\_sc\_K} & title printed in the \texttt{.out} & macro law \\
\midrule
\texttt{"Timoshenko"} & \texttt{The Effective Timoshenko Stiffness Matrix} & beam $6\times6$ \\
\texttt{"Classical Plate"} & \texttt{The Effective Classical Plate Stiffness Matrix} & plate ABD $6\times6$ \\
\texttt{"Reissner-Mindlin Plate"} & \texttt{The Effective Reissner-Mindlin Plate Stiffness Matrix} & plate ABDG $8\times8$ \\
\texttt{""} (omitted) & \texttt{The Effective Stiffness Matrix} & 3-D solid $6\times6$ \\
\bottomrule
\end{tabular}
\end{table}

Only the last of these also prints
\verb| The Engineering Constants (Approximated as Orthotropic)|. That is the
\texttt{constants=True} branch, and it is meaningful only for a
three-dimensional elasticity law; the beam and plate calls pass
\texttt{constants=False}.

\section{Beam: the Timoshenko $6\times6$}
\label{sec:timo}

The beam macro model carries six generalized strains: one axial extension, two
transverse shears, one twist and two bending curvatures. In code order --- the
order stated in the module header of \texttt{src/opensg\_solid/sg\_homo.py} ---
these are

\begin{equation}
\bar{\boldsymbol\epsilon} \;=\;
\left\{\;\bar\epsilon_{11}\;\;\gamma_{12}\;\;\gamma_{13}\;\;
\kappa_{1}\;\;\kappa_{2}\;\;\kappa_{3}\;\right\}^{\mathsf T},
\end{equation}

and the conjugate generalized forces are the axial force, the two transverse
shear forces, the torque and the two bending moments --- the VABS ordering
$\{F_1\;F_2\;F_3\;M_1\;M_2\;M_3\}$ in which the shipped blade load table
\texttt{ff51\_rmc\_reform.dat} is written (its header reads
\verb|# eta F1 F2 F3 M1 M2 M3 (RM CENTER-ref BeamDyn, VABS order)|). The law is

\widematrix{
\begin{Bmatrix} F_1 \\ F_2 \\ F_3 \\ M_1 \\ M_2 \\ M_3 \end{Bmatrix}
=
\begin{bmatrix}
S_{11} & S_{12} & S_{13} & S_{14} & S_{15} & S_{16} \\
S_{12} & S_{22} & S_{23} & S_{24} & S_{25} & S_{26} \\
S_{13} & S_{23} & S_{33} & S_{34} & S_{35} & S_{36} \\
S_{14} & S_{24} & S_{34} & S_{44} & S_{45} & S_{46} \\
S_{15} & S_{25} & S_{35} & S_{45} & S_{55} & S_{56} \\
S_{16} & S_{26} & S_{36} & S_{46} & S_{56} & S_{66}
\end{bmatrix}
\begin{Bmatrix} \bar\epsilon_{11} \\ \gamma_{12} \\ \gamma_{13} \\
\kappa_{1} \\ \kappa_{2} \\ \kappa_{3} \end{Bmatrix}.
}

The matrix is fully populated in general --- a composite blade section couples
extension to bending and twist --- but its diagonal carries the familiar
engineering meaning

\begin{equation}
\mathrm{diag}(\mathbf S) = \left[\;EA\;\;GA_2\;\;GA_3\;\;GJ\;\;EI_2\;\;EI_3\;\right],
\label{eq:vabsdiag}
\end{equation}

which is how the drivers label their printout: the list
\verb|LBL = ["EA", "GA2", "GA3", "GJ", "EI2", "EI3"]| appears verbatim in
\texttt{examples/OpenSG\_shell/IEA\_blade\_beam/1\_homo\_beam\_props.py}. Read
Eq.~\eqref{eq:vabsdiag} as a naming convention, not as a claim of
decoupling. $S_{44}$ is the torsional stiffness only to the extent that the
twist channel is uncoupled; when $S_{45}$ or $S_{46}$ is large, $GJ$ as
conventionally understood is not $S_{44}$.

\subsection{Running it on a shell cross-section}

The msg-shell ring route homogenizes a cross-section contour with the
constrained six-DOF shell element (independent drilling $\omega_3$ enforced by
element-constant Lagrange multipliers, $\gamma_{23}$ tied by MITC) and takes
its wall transverse shear from the MSG least-squares plate solution. The user
edits one file, the one-dimensional shell SG YAML, which carries the contour
nodes, the line elements, the per-element orientations, the layup sections, the
materials and the \texttt{reference:} field. From
\texttt{examples/OpenSG\_shell/IEA\_blade\_beam}:

\begin{lstlisting}[style=shellst]
python 1_homo_beam_props.py
\end{lstlisting}

The driver's user block is three lines long; only the station name changes:

\begin{lstlisting}[style=opensg,caption={The complete user input of the shell beam driver. build\_rm\_bundle reads iea\_s10\_shell.yaml.}]
name = "iea_s10"                  # IEA-22 station at r/R = 0.20
B = build_rm_bundle(name + "_shell.yaml")
C6 = np.asarray(B["Timo"])
\end{lstlisting}

\texttt{build\_rm\_bundle} writes two files by itself: the SwiftComp-layout
\texttt{<yaml base>\_Timo.out} (here \texttt{iea\_s10\_shell\_Timo.out}) with
banner \verb| OpenSG msg-shell beam model [ext sh2 sh3 twist bend2 bend3]| and
title \texttt{The Effective Timoshenko Stiffness Matrix}, and the per-section
wall law \texttt{<yaml base>\_ABDG.out} described in
Section~\ref{sec:wallabdg}. The driver additionally saves the bare $6\times6$
as \texttt{iea\_s10\_beam\_Timo.out} and the ring mesh figure
\texttt{iea\_s10\_ring\_mesh.png}.

\begin{figure}[htbp]
\centering
\includegraphics[width=0.86\linewidth]{\FIGDIR/shell_xsec_mesh.png}
\caption{The one-dimensional shell SG solved by \texttt{build\_rm\_bundle}: the
IEA-22 blade cross-section at $r/R = 0.20$, 307 contour nodes and 310 two-node
line elements along the wall mid-surface, coloured by \texttt{sets: element}
entry. The Timoshenko $6\times6$ beam law whose diagonal is quoted in this
section is this mesh's answer, written to \texttt{iea\_s10\_shell\_Timo.out}.}
\label{fig:const-shellxsec}
\end{figure}

For the IEA-22 blade station at $r/R = 0.20$ --- the SG of
Figure~\ref{fig:const-shellxsec} --- that $6\times6$ has the diagonal

\widematrix{
\mathrm{diag}(\mathbf S) =
\left[\,2.7661\times10^{10},\;7.1861\times10^{8},\;4.2174\times10^{8},\;
2.4376\times10^{9},\;3.5144\times10^{10},\;6.8130\times10^{10}\,\right]
}

in SI units, so $EA = 27.7$ GN, $EI_3 = 68.1$ GN$\cdot$m$^2$ edgewise against
$EI_2 = 35.1$ GN$\cdot$m$^2$ flapwise, and $GJ = 2.44$ GN$\cdot$m$^2$. The
largest off-diagonal terms in that file are $S_{15} = 1.66864\times10^{9}$ and
$S_{16} = -2.07730\times10^{9}$ (extension--bending, i.e.\ the section
reference axis is not the tension centre) and
$S_{56} = -5.55712\times10^{9}$ (the two bending channels are not principal).
Terms that should vanish by symmetry sit at $10^{-5}$ against diagonals of
$10^{10}$, which is the numerical zero of this solve.

\subsection{Running it on a solid cross-section}

The msg-solid route runs the KKT beam driver: a zeroth-order $V_0$ solve under
four rigid-body Lagrange constraints, then a first-order $V_{1s}$ solve reusing
the same factorization, then the $6\times6$ reduction in
\texttt{finalize\_v1\_and\_compute\_deff}. The entry point is
\texttt{plate\_homo\_2d(..., n\_model=1)}, which requires \texttt{n\_sg >= 2}
(a two- or three-dimensional solid SG). From
\texttt{examples/OpenSG-solid/7\_get\_beam\_props\_from\_SG}:

\begin{lstlisting}[style=shellst]
python beam_homo_sg.py
\end{lstlisting}

\begin{lstlisting}[style=opensg,caption={The user block of beam\_homo\_sg.py. The material rows are (E1,E2,E3,G12,G13,G23,nu12,nu13,nu23) in MPa; angles rotates each material about the SG normal.}]
n_model = 1                    # 1: Beam; 2: Plate; 3: 3D elastic
name = "RHC_SW_2UC_45"         # reads <name>.yaml

material_param = jnp.array([
    (108e3, 8e3, 8e3, 4e3, 4e3, 3e3, 0.32, 0.32, 0.30),      # ply +45
    (108e3, 8e3, 8e3, 4e3, 4e3, 3e3, 0.32, 0.32, 0.30),      # ply -45
    (69e3, 69e3, 69e3, 26.54e3, 26.54e3, 26.54e3, 0.30, 0.30, 0.30)])
angles = jnp.array([45.0, -45.0, 0.0])

r = plate_homo_2d(name + ".yaml", material_param=material_param,
                  angles=angles, n_model=n_model)
\end{lstlisting}

This writes \texttt{RHC\_SW\_2UC\_45.out} with banner
\verb| OpenSG msg-solid beam model, omega 1| and the mesh figure
\texttt{RHC\_SW\_2UC\_45\_mesh.png} (Figure~\ref{fig:rhcmesh}). Its diagonal is

\widematrix{
\mathrm{diag}(\mathbf S) =
\left[\,9.8993\times10^{6},\;1.9075\times10^{6},\;9.3619\times10^{5},\;
2.9775\times10^{8},\;2.1896\times10^{9},\;9.2899\times10^{8}\,\right]
}

in the input units (N and N$\cdot$mm$^2$, since the SG is in millimetres and the
moduli in MPa).

\begin{figure}[htbp]
\centering
\includegraphics[width=0.72\linewidth]{\FIGDIR/rhc_2dsg_mesh.png}
\caption{The two-dimensional honeycomb solid SG \texttt{RHC\_SW\_2UC\_45}, run
here with $n_{\mathrm{model}} = 1$: the Timoshenko $6\times6$ beam law whose
diagonal is quoted in this section is this mesh's answer, written to
\texttt{RHC\_SW\_2UC\_45.out}. The same mesh is the SG of the plate
($n_{\mathrm{model}} = 2$) and solid ($3$) examples. Colours are the three
material sets of the \texttt{.yaml}: the $+45$ and $-45$ composite plies of the
cell walls and the isotropic face sheets.}
\label{fig:rhcmesh}
\end{figure}

\subsection{The classical $4\times4$ that rides along}
\label{sec:ceffeb}

The Timoshenko reduction is built on top of a zeroth-order (Euler--Bernoulli)
solve, and that intermediate law is a genuine classical beam stiffness in its
own right. It drops the two transverse shears:

\begin{equation}
\begin{Bmatrix} F_1 \\ M_1 \\ M_2 \\ M_3 \end{Bmatrix}
=
\mathbf A_{\mathrm{EB}}
\begin{Bmatrix} \bar\epsilon_{11} \\ \kappa_1 \\ \kappa_2 \\ \kappa_3 \end{Bmatrix},
\qquad \mathbf A_{\mathrm{EB}} \in \mathbb R^{4\times4}.
\end{equation}

On the solid route it is returned as \texttt{r["C\_eff\_EB"]} alongside the
Timoshenko \texttt{r["C\_eff"]}. In \texttt{\_beam\_homo\_kkt} it is formed as
$\mathbf C_{\mathrm{eb}} = (D_{ee} + D_1^{V_0})/\omega$ immediately after the
$V_0$ solve, before the first-order ladder runs, and it is then passed straight
into \texttt{finalize\_v1\_and\_compute\_deff} as the $\mathbf A$ block of the
Timoshenko transformation. It is not written to the \texttt{.out} --- the
\texttt{.out} always carries the $6\times6$ --- so a driver that wants it must
print it, as \texttt{beam\_homo\_sg.py} does.

\subsection{Which routes produce the beam law}

\begin{table}[htbp]
\centering
\caption{The three entry points that return a Timoshenko $6\times6$. The
msg-solid entry point lives in \texttt{opensg\_solid.sg\_homo}, the two
msg-shell ones in \texttt{opensg\_shell}.}
\label{tab:beamroutes}
\footnotesize
\begin{tabular}{lll}
\toprule
route & entry point & SG input \\
\midrule
msg-solid beam & \texttt{plate\_homo\_2d(..., n\_model=1)} & 2-D or 3-D solid SG (\texttt{n\_sg >= 2}) \\
msg-shell ring & \texttt{build\_rm\_bundle} (calls \texttt{ring\_indep}) & 1-D shell SG (cross-section contour) \\
msg-shell segment & \texttt{segment\_timo\_from\_3dyaml} & 3-D shell segment YAML (tapered or prismatic) \\
\bottomrule
\end{tabular}
\end{table}

The segment route replaces axial periodicity by Dirichlet data. Each end
cross-section is extracted topologically by \texttt{sg\_mesh.extract}, solved on
its own with \texttt{ring\_indep} to give that ring's $\mathbf C_6$, $V_0$ and
$V_1$, and those warping fields are then imposed on the segment's boundary
nodes; the segment energy is divided by the segment length
$L_z$. It writes \texttt{<yaml base>\_Timo.out} with the same
\texttt{Timoshenko} title and the banner
\verb|msg-shell tapered/aperiodic segment (boundary V0/V1 Dirichlet, L=<Lz>)|.
For a prismatic segment the result must reproduce the boundary ring's own
$6\times6$; that identity is the standing verification of the route.

\section{Plate: classical ABD $6\times6$ and Reissner--Mindlin ABDG $8\times8$}
\label{sec:plate}

\subsection{Classical (Kirchhoff--Love) plate law}

The classical plate model pairs three membrane stress resultants and three
moment resultants with three membrane strains and three curvatures:

\widematrix{
\begin{Bmatrix} N_{11} \\ N_{22} \\ N_{12} \\ M_{11} \\ M_{22} \\ M_{12} \end{Bmatrix}
=
\begin{bmatrix} \mathbf A & \mathbf B \\[2pt] \mathbf B^{\mathsf T} & \mathbf D \end{bmatrix}
\begin{Bmatrix} \varepsilon_{11} \\ \varepsilon_{22} \\ 2\varepsilon_{12} \\
\kappa_{11} \\ \kappa_{22} \\ \kappa_{12} \end{Bmatrix},
\qquad
\mathbf A, \mathbf B, \mathbf D \in \mathbb R^{3\times3}.
}

The row order \verb|[N11 N22 N12 M11 M22 M12]| is stated in the module header
of \texttt{src/opensg\_solid/sg\_homo.py} and echoed by the driver in
\texttt{examples/OpenSG-solid/3\_get\_plate\_props\_from\_2D\_SG}. The
$\mathbf B$ block is the extension--bending coupling; it vanishes only for a
laminate that is symmetric about the chosen reference surface, so it is a
property of the SG \emph{and} of where you put the reference. Run it with

\begin{lstlisting}[style=shellst]
python plate_homo_2dsg.py
\end{lstlisting}

which calls \texttt{plate\_homo\_2d(name + ".yaml", material\_param=...,
angles=..., n\_model=2)} and writes \texttt{RHC\_SW\_2UC\_45.out} with banner
\verb| OpenSG msg-solid plate model, omega 35.248001, periodic| and title
\texttt{The Effective Classical Plate Stiffness Matrix}. Here $\omega =
35.248001$ is the in-plane SG measure (the cell width), and the leading entries
are $A_{11} = 3.3586\times10^{5}$, $A_{22} = 1.1291\times10^{5}$,
$A_{66} = 1.0552\times10^{5}$ N/mm and $D_{11} = 9.0368\times10^{7}$,
$D_{22} = 5.7011\times10^{7}$, $D_{66} = 4.6026\times10^{7}$ N$\cdot$mm. The
$\mathbf B$ block is at the $10^{0}$ level against an $\mathbf A$ of $10^{5}$,
i.e.\ the cell is symmetric about its own mid-surface to solver precision.

\begin{figure}[htbp]
\centering
\includegraphics[width=0.72\linewidth]{\FIGDIR/rhc_2dsg_mesh.png}
\caption{The same two-dimensional honeycomb cell SG \texttt{RHC\_SW\_2UC\_45} as
Figure~\ref{fig:rhcmesh}, run here with $n_{\mathrm{model}} = 2$: the classical
plate ABD $6\times6$ whose leading entries are quoted in this section is this
mesh's answer, written to \texttt{RHC\_SW\_2UC\_45.out} under the title
\texttt{The Effective Classical Plate Stiffness Matrix}. Only the model choice
changes between this run and the beam run of
Section~\ref{sec:timo}; the SG is byte-identical.}
\label{fig:rhcplate}
\end{figure}

\subsection{Reissner--Mindlin plate law}

The shear-refined model adds the transverse-shear pair. The strain vector grows
to eight components and the resultant vector gains the two transverse shear
forces $Q_1$ and $Q_2$:

\widematrix{
\begin{Bmatrix} \mathbf N \\ \mathbf M \\ \mathbf Q \end{Bmatrix}
=
\begin{bmatrix}
\mathbf A & \mathbf B & \mathbf 0 \\[2pt]
\mathbf B^{\mathsf T} & \mathbf D & \mathbf 0 \\[2pt]
\mathbf 0 & \mathbf 0 & \mathbf G
\end{bmatrix}
\begin{Bmatrix} \boldsymbol\varepsilon \\ \boldsymbol\kappa \\ \boldsymbol\gamma \end{Bmatrix},
\qquad
\boldsymbol\gamma = \{\,2\gamma_{13}\;\;2\gamma_{23}\,\}^{\mathsf T},
\qquad
\mathbf Q = \{\,Q_1\;\;Q_2\,\}^{\mathsf T}.
}

Here $\boldsymbol\varepsilon$ and $\boldsymbol\kappa$ are exactly the first and
second triples of the classical law above.

The zero off-diagonal blocks are structural, not an approximation of
convenience. The code assembles the $8\times8$ as \verb|ABDG[:6,:6] = A6| and
\verb|ABDG[6:,6:] = G_msg|, so the in-plane law and the transverse-shear law
come from two separate variational problems and are then stacked. The
$2\times2$ block $\mathbf G$ is obtained as $\mathbf G = \mathbf X^{-1}$ from
the $\mathbf U^*$ least-squares reduction of the first-order warping ladder
(function \texttt{\_rm\_ls\_reduction}), a 78-equation, 27-unknown system solved
by a column-equilibrated truncated-SVD minimum-norm solve with cutoff
\texttt{\_LS\_RCOND = 1e-12}. Two diagnostics come out with it: the relative
residual of the fit, reported as \texttt{r["Ustar\_rel"]}, and the smallest
eigenvalue of $\mathbf X$. If that eigenvalue is not positive the SPD gate
fails and the block is returned as \texttt{None} --- in which case no
$8\times8$ exists and the code falls back to printing the classical
$6\times6$. Always check that you got an $8\times8$ before trusting a
transverse-shear number.

\subsubsection{The 1-D SG route}

\texttt{opensg\_solid.rm\_plate\_1D.msg\_rm\_plate.rm\_plate\_msg} solves the
through-thickness SG of a single laminate and returns \texttt{A6} (the
$6\times6$ ABD), \texttt{G\_msg} and \texttt{ABDG}, with rows
\verb|e11 e22 g12 k11 k22 k12 2g13 2g23|. Which of the two matrices is printed
is selected by \texttt{model:} in the user's \texttt{layup\_db.yaml}:
\texttt{model: 0} prints the classical $6\times6$, \texttt{model: 1} prints the
full $8\times8$. From
\texttt{examples/OpenSG-solid/1\_get\_plate\_props\_from\_1DSG}:

\begin{lstlisting}[style=shellst]
python 1d_sg.py
\end{lstlisting}

The input the user edits is \texttt{layup\_db.yaml} alone. It carries the
reference surface (\texttt{fraction: 0.5} is the mid-surface, \texttt{0} the
bottom or OML face, \texttt{1} the top), the materials with their densities,
the stacking sequence bottom-ply-first, and the SG discretisation
(\texttt{n\_per\_layer: 1} element per physical ply, \texttt{elem\_order: 4} for
the five-noded quartic element that represents the warping ladder exactly).
The script writes the SG mesh \texttt{1dsg.yaml} plus
\texttt{1dsg.png} (Figure~\ref{fig:plate1dsg}) and the law in
\texttt{layup\_db\_plate\_homo.out}.

The shipped case is the Nayak Example-5 sandwich, $(0/90/0/90/\mathrm{core})_s$,
nine plies, total thickness $h = 0.1524$ m, graphite/epoxy faces
($E_1 = 128$ GPa) over a HEREX C70.130 PVC foam core, each face ply
$0.9525$ mm and the core $144.78$ mm. Its \texttt{.out} banner records
\verb|9 plies, h = 0.152400 m, fraction = 0.5, rho*h = 30.251400 kg/m^2| and
the leading entries are

\begin{equation}
A_{11} = A_{22} = 5.49165\times10^{8}\ \mathrm{N/m},\quad
D_{11} = 3.00055\times10^{6},\quad D_{22} = 2.93711\times10^{6}\ \mathrm{N\!\cdot\!m},
\end{equation}
\begin{equation}
G_{11} = 7.70098\times10^{6}\ \mathrm{N/m},\qquad
G_{22} = 7.54237\times10^{6}\ \mathrm{N/m}.
\end{equation}

Read them physically: the thin stiff faces carry the membrane and bending
channels, while the thick compliant core sets the transverse shear. The
$\mathbf B$ block prints at $10^{-6}$ against an $\mathbf A$ of $10^{8}$,
because the stack is symmetric about \texttt{fraction: 0.5}; change that to
\texttt{fraction: 0.0} and $\mathbf B$ becomes a leading-order term.

\begin{figure}[htbp]
\centering
\includegraphics[width=0.62\textwidth]{\FIGDIR/plate_1dsg.png}
\caption{The through-thickness one-dimensional structure gene of the Nayak
Example-5 sandwich written by \texttt{1d\_sg.py}: eight graphite/epoxy face
plies over a HEREX C70.130 foam core, with the mid-surface reference
(\texttt{fraction: 0.5}) marked. This is the SG whose solution is the
Reissner--Mindlin ABDG $8\times8$ whose $A$, $D$ and $G$ entries are quoted in
this section, written by \texttt{rm\_plate\_msg} to
\texttt{layup\_db\_plate\_homo.out}.}
\label{fig:plate1dsg}
\end{figure}

\subsubsection{The 2-D / 3-D SG route}

\texttt{plate\_homo\_2d(..., n\_model=2, shear\_refined=True)} runs
\texttt{plate\_shear\_ladder} on the general SG assembly after the classical
solve, re-integrating the value-product ladder blocks on a quadrature exact to
degree $2p+2$. It adds \texttt{r["G\_msg"]} ($2\times2$),
\texttt{r["A6\_ladder"]}, \texttt{r["X\_shear"]}, \texttt{r["Ustar\_rel"]} and
\texttt{r["ABDG"]} to the result dict, and the \texttt{.out} switches its title
to \texttt{Reissner-Mindlin Plate} and its matrix to the $8\times8$. Passing
\texttt{shear\_refined=True} with any \texttt{n\_model} other than $2$ raises a
\texttt{ValueError}.

Note the derivation status recorded in the docstring of
\texttt{plate\_shear\_ladder}, because it decides how much to trust the
$\mathbf G$ block. For \texttt{n\_sg = 1} this \emph{is} the Yu (2005)
one-dimensional formulation, gated against the closed-form anchors (an
isotropic layer with $\nu = 0$ returns $\tfrac{5}{6}Gh$). For \texttt{n\_sg = 2}
or \texttt{3} it is the SwiftComp-style extension --- the same energy functional
with $\Gamma_h$ carrying the resolved-dimension derivatives, periodic
master--slave coupling, and the $\langle w \rangle = 0$ average constraint
replacing the one-dimensional node constraint. That extension is gated on the
laminate-as-strip equivalence (a laminate meshed as a 2-D or 3-D SG must
reproduce its 1-D $\mathbf G$ under refinement) but is \textbf{not} yet
validated against an external reference for SGs that are heterogeneous
\emph{in plane}. The shear-refined \emph{recovery} is likewise not ported:
plate dehomogenization stays the classical Kirchhoff-mode recovery.

\section{Solid: the equivalent 3-D law}
\label{sec:solid}

When the macro model is ordinary three-dimensional elasticity, the SG returns a
$6\times6$ anisotropic stiffness in the Voigt order $[11\;22\;33\;23\;13\;12]$
with engineering shear strains:

\widematrix{
\begin{Bmatrix}
\sigma_{11} \\ \sigma_{22} \\ \sigma_{33} \\ \sigma_{23} \\ \sigma_{13} \\ \sigma_{12}
\end{Bmatrix}
=
\mathbf C^{\,\mathrm{3D}}
\begin{Bmatrix}
\bar\Gamma_{11} \\ \bar\Gamma_{22} \\ \bar\Gamma_{33} \\
2\bar\Gamma_{23} \\ 2\bar\Gamma_{13} \\ 2\bar\Gamma_{12}
\end{Bmatrix},
\qquad \mathbf C^{\,\mathrm{3D}} \in \mathbb R^{6\times6}.
}

Internally the solid engine stores strain and stress in the SwiftComp order
$(xx, yy, zz, yz, xz, xy)$, which is the same list. On the msg-shell routes
axis 1 is the prismatic (beam) direction and axes 2 and 3 span the
cross-section, as recorded in the module constant
\texttt{opensg\_shell.sg\_homo.GBAR\_ORDER}:
\verb|strains [e11 e22 e33 2e23 2e13 2e12] -> stiffness C11..C66|
\verb|(1=beam axis, 2/3=cross-section axes)|.

\subsection{The engineering-constants block}

This is the only law for which the \texttt{.out} also prints an
engineering-constants block. It is computed from the compliance
$\mathbf S = (\mathbf C^{\,\mathrm{3D}})^{-1}$ by the orthotropic-approximation
formulas coded in \texttt{write\_sc\_K} (and mirrored on the shell side by
\texttt{opensg\_shell.sg\_homo.elastic\_constants}, which also returns
$\mathrm{cond}(\mathbf C)$):

\widematrix{
E_1 = \frac{1}{S_{11}},\quad E_2 = \frac{1}{S_{22}},\quad E_3 = \frac{1}{S_{33}},
\qquad
G_{12} = \frac{1}{S_{66}},\quad G_{13} = \frac{1}{S_{55}},\quad G_{23} = \frac{1}{S_{44}},
}
\widematrix{
\nu_{12} = -\frac{S_{12}}{S_{11}},\qquad
\nu_{13} = -\frac{S_{13}}{S_{11}},\qquad
\nu_{23} = -\frac{S_{23}}{S_{22}}.
}

The word \emph{approximated} in the printed heading is load-bearing. These nine
numbers reproduce $\mathbf C^{\,\mathrm{3D}}$ exactly only if that matrix is
genuinely orthotropic in the SG axes. A cell with normal--shear coupling still
prints them, and they will not rebuild the full matrix. The matrix, not the
constants block, is the result.

\subsection{Which entry points produce it}

\begin{table}[htbp]
\centering
\caption{The three entry points that return an equivalent 3-D solid law, and
the measure each divides the assembled energy by.}
\label{tab:solidroutes}
\footnotesize
\begin{tabular}{lll}
\toprule
entry point & SG & normalisation \\
\midrule
\texttt{plate\_homo\_2d(..., n\_model=3)} & 1-D/2-D/3-D solid SG & $\mathbf C_{\mathrm{eff}} = (\bar{\mathbf D} + \mathbf D_1)/\omega$ \\
\texttt{opensg\_shell.build\_solid\_bundle} & 1-D shell SG (cross-section contour) & $\mathbf C^{\,\mathrm{3D}} = \mathbf D_{\mathrm{eff}}/A_{\mathrm{cell}}$ \\
\texttt{opensg\_shell.sg\_homo.shell\_sg3d} & 3-D shell SG (surface mesh) & $\mathbf D_{\mathrm{eff}}/\omega$, $\omega$ = midsurface area \\
\bottomrule
\end{tabular}
\end{table}

For \texttt{build\_solid\_bundle} the normalising area defaults to the convex
hull of the contour (\texttt{area\_source = "hull"}); pass \texttt{cell\_area}
explicitly to use a different measure, and the banner records which was used.
The isotropic square-tube example in the folder
\texttt{examples/OpenSG\_shell/3\_get\_solid\_props\_from\_shell\_cross\_section/}
\texttt{square\_section\_isotropic}
sets \texttt{CELL\_AREA = 4*a*t\_w}, the wall material area, so that the shell
route and the solid route are normalised identically and can be compared term
by term. Its \texttt{square\_tube\_1Dshell\_C3D.out} returns
$E_1 = 7.0000000\times10^{10}$ Pa and $\nu_{12} = 0.30000$, recovering the wall
material's own $E$ and $\nu$ along the beam axis exactly, which is the sanity
check that fixes the normalisation.

The last row of Table~\ref{tab:solidroutes} carries a subtlety worth stating
plainly. \texttt{shell\_sg3d} \emph{returns} \texttt{C3D} divided by the SG
measure (the midsurface surface area, the three-dimensional analogue of the
plane-section perimeter), but \emph{writes} the \texttt{.out} divided by the
bounding-box cell volume $V_{\mathrm{cell}}$, so the printed moduli compare
directly against a solid SwiftComp \texttt{.K} file. If you read the numbers
out of the return dict and out of the \texttt{.out} and they differ by a
constant factor, that factor is $\omega/V_{\mathrm{cell}}$ and nothing is
wrong.

\subsection{Periodicity, and the aperiodic alternative}
\label{sec:periodicity}

A free SG driven by six macro strains is rank one by construction: every mode
except $\bar\Gamma_{11}$ can be cancelled at zero energy by an affine
fluctuation. Something must forbid those fields. All three routes
\emph{default} to periodicity, and periodicity is the SwiftComp-parity mode:

\begin{itemize}
\item \texttt{build\_solid\_bundle(..., periodic=True)} is the default; it ties
  opposite bounding-box faces through the shell periodic assembly map, and the
  repeated \texttt{dof\_map[dof\_map]} resolves the edge and corner chains, so
  faces, edges and corners are all tied. The tie rides in the assembly map, so
  no constraint rows are added and the rigid kernel drops from four rows to the
  three translations.
\item \texttt{shell\_sg3d(..., boundary="periodic")} is the default; the same
  map is applied on the full three-dimensional coordinates, and the three rigid
  translations are removed by area-weighted Lagrange rows.
\item \texttt{plate\_homo\_2d(..., boundary=None)} resolves to
  \texttt{"periodic"}.
\end{itemize}

Two of these also accept an aperiodic treatment, and it is a genuine
alternative, not a diagnostic. It is the \emph{boundary-solution} treatment:
for a unit cell the boundary solution is the macro (affine) field itself, and
mapping it onto the boundary prescribes zero fluctuation there.

\begin{itemize}
\item \texttt{shell\_sg3d(..., boundary="aperiodic")} clamps
  $w_1 = w_2 = w_3 = 0$ on every node of the bounding-box faces and leaves the
  three rotational DOFs \textbf{natural}. Translations are what the macro solid
  strain prescribes; clamping the rotations as well over-stiffens the cut
  edges, and the in-code note records the measured effect of freeing them
  ($+12\,\%$ down to $+7\,\%$ on $C_{11}$). No Lagrange border is needed,
  because the Dirichlet data also fixes the rigid modes.
\item \texttt{plate\_homo\_2d(..., boundary="aperiodic")} pins all three
  fluctuation DOFs of every bounding-box-face node --- the solid engine carries
  no rotational DOFs, so those three are the whole node --- instead of the usual
  first-node pin, and assembles on the raw mesh rather than the periodic map.
  It is supported for the plate and solid models with
  \texttt{solver="direct"}; requesting it with \texttt{n\_model=1} or
  \texttt{solver="cg"} raises a \texttt{ValueError}.
\end{itemize}

Kinematic (Dirichlet) homogenization is an \emph{upper bound} on a single cell,
and it converges to the periodic answer as the SG grows, because the boundary
layer scales like $1/N$. That convergence is the acceptance check.

The head-to-head is shipped. On the shell side, the script
\texttt{compare\_boundary\_modes.py} in
\texttt{examples/OpenSG\_shell/4\_get\_solid\_props\_from\_shell\_3D\_SG}
runs \texttt{shell\_sg3d} on the same Schwarz-P TPMS cell
(Figure~\ref{fig:schwarz}) in both modes and appends a per-unit-cell block to
\texttt{boundary\_compare.dat}:

\begin{lstlisting}[style=shellst]
python compare_boundary_modes.py
\end{lstlisting}

\begin{table}[htbp]
\centering
\caption{Schwarz-P shell cell, per-unit-cell stiffness in MPa, from
\texttt{boundary\_compare.dat}. The aperiodic run clamps 720 free-face nodes
(81216 DOF, 42.3 s) against 79056 DOF and 66.5 s for periodic.}
\label{tab:bcshell}
\small
\begin{tabular}{lrrrrrr}
\toprule
& \multicolumn{3}{c}{$t = 0.036457$} & \multicolumn{3}{c}{$t = 0.129330$} \\
\cmidrule(lr){2-4}\cmidrule(lr){5-7}
term & aperiodic & periodic & diff \% & aperiodic & periodic & diff \% \\
\midrule
$C_{11}$ & 2375.89 & 2225.31 & $+6.77$ & 9881.00 & 8916.81 & $+10.81$ \\
$C_{22}$ & 2375.92 & 2225.28 & $+6.77$ & 9881.00 & 8916.85 & $+10.81$ \\
$C_{33}$ & 2376.22 & 2225.03 & $+6.80$ & 9873.34 & 8916.17 & $+10.74$ \\
$C_{12}$ & 1537.47 & 1564.87 & $-1.75$ & 4921.09 & 5352.21 & $-8.06$ \\
$C_{44}$ & 1078.28 & 1040.98 & $+3.58$ & 4150.25 & 3961.63 & $+4.76$ \\
$C_{66}$ & 1078.61 & 1040.93 & $+3.62$ & 4151.27 & 3961.51 & $+4.79$ \\
\bottomrule
\end{tabular}
\end{table}

Table~\ref{tab:bcshell} is exactly the expected signature: the direct terms come
out stiffer (a kinematic upper bound), the Poisson terms softer, and the gap
widens with wall thickness because the clamped boundary layer occupies more of
the cell. The solid-side twin, the identically named script in
\texttt{examples/OpenSG-solid/6\_get\_solid\_props\_from\_3D\_SG},
does the same for the tetrahedral TPMS SGs with
\texttt{plate\_homo\_2d(n\_model=3)} and shows a larger boundary layer on the
denser cell: for Sample\_1, $C_{11}$ is $12977.0$ MPa aperiodic against
$10190.2$ MPa periodic ($+27.3\,\%$) with 17610 boundary nodes, while
$C_{12}$ falls $3.94\,\%$. Both runs take about 150 s, so the aperiodic mode
costs nothing extra --- the choice is physical, not computational.

\begin{figure}[htbp]
\centering
\includegraphics[width=0.60\textwidth]{\FIGDIR/schwarz_p_mesh.png}
\caption{The Schwarz-P TPMS three-dimensional shell SG solved by
\texttt{shell\_sg3d}: 13536 nodes, 26360 shell elements, and zero junction
edges (a smooth TPMS has every edge interior, so the junction treatment never
fires). The banner of the \texttt{.out} reports all three counts. This is the
SG behind Table~\ref{tab:bcshell}: every number there is an entry of this
mesh's equivalent 3-D solid $6\times6$, the law \texttt{shell\_sg3d} writes to
\texttt{<yaml>\_C3D.out}.}
\label{fig:schwarz}
\end{figure}

\subsection{A worked solid law: the TPMS Sample\_1 cell}

The reference case for the equivalent 3-D law of this section on the solid
engine is
\texttt{examples/OpenSG-solid/6\_get\_solid\_props\_from\_3D\_SG/sample\_1}:

\begin{lstlisting}[style=shellst]
python run_sample_1.py
\end{lstlisting}

\begin{lstlisting}[style=opensg,caption={run\_sample\_1.py: a 546k-tet TPMS cell of aluminium, periodic in all three directions. The single material row is isotropic E = 69 GPa, nu = 0.3, matched to the SwiftComp .K it is checked against.}]
E, nu = 69.0e9, 0.30
G = E/(2*(1+nu))
r = plate_homo_2d(SC, material_param=jnp.array([(E, E, E, G, G, G,
                                                 nu, nu, nu)]),
                  angles=jnp.array([0.0]), n_model=3, workdir=".", plot=True,
                  boundary="periodic")   # SwiftComp .K digit-parity mode
\end{lstlisting}

This writes \texttt{Sample\_1.out} (banner
\verb| OpenSG msg-solid 3D elastic model, omega 0.3|, where $\omega = 0.3$ is
the relative density of the cell), the mesh figure
\texttt{Sample\_1\_mesh.png}, the per-unit-cell matrix
\texttt{Sample\_1\_msg\_solid.out}, and the comparison table
\texttt{Sample\_1\_compare.dat}. The engineering-constants block of
\texttt{Sample\_1.out} reads

\begin{table}[htbp]
\centering
\caption{The engineering constants printed by \texttt{Sample\_1.out} for the
TPMS Sample\_1 cell (aluminium, $E = 69$ GPa, $\nu = 0.3$, $\omega = 0.3$).}
\label{tab:tpmsconst}
\small
\begin{tabular}{lrlrlr}
\toprule
constant & value (Pa) & constant & value (Pa) & constant & value \\
\midrule
$E_1$ & $2.0547093\times10^{10}$ & $G_{12}$ & $1.5065850\times10^{10}$ & $\nu_{12}$ & $0.35652386$ \\
$E_2$ & $2.0547272\times10^{10}$ & $G_{13}$ & $1.5065455\times10^{10}$ & $\nu_{13}$ & $0.35654253$ \\
$E_3$ & $2.0546693\times10^{10}$ & $G_{23}$ & $1.5065467\times10^{10}$ & $\nu_{23}$ & $0.35653028$ \\
\bottomrule
\end{tabular}
\end{table}

Three things to read out of Table~\ref{tab:tpmsconst}. First, the three Young's
moduli agree to five significant figures although they come from three
independent unit-strain solves on the same cell --- that is the cubic symmetry
of the geometry reappearing in the answer, and it is the cheapest available
verification that the periodic tie caught all faces, edges and corners.
Second, the aluminium is isotropic but the cell is not: $E_1/E = 0.298$ while
$G_{12}/G = 0.568$, so a TPMS at 30\,\% relative density retains far more shear
stiffness than axial stiffness. Third, the off-diagonal terms of the printed
stiffness sit at $10^{5}$ against diagonals of $10^{10}$, so the orthotropic
approximation behind the constants block is legitimate here. The
per-unit-cell comparison in \texttt{Sample\_1\_compare.dat} agrees with the
SwiftComp \texttt{.K} to $\pm 0.000\,\%$ on all nine reported terms
($C_{11} = 10190.158$ MPa for both).

A second, independent validation of the same law is
\texttt{examples/OpenSG-solid/5\_get\_solid\_props\_from\_SG}
(Figure~\ref{fig:udcomp}), a UD fibre/matrix square cell meshed with 801 nodes
and 1536 triangles. It ships a self-contained constant-strain-triangle driver
that solves the same $6\times6$ law under generalized plane strain with periodic
faces and writes \texttt{UDcomp\_2D\_Deff.out} plus
\texttt{UDcomp\_2D\_constants\_vs\_paper.dat}. Its nine constants reproduce the
ASC-23 Table II Model-3 values ($E_1 = 167.2553$ GPa, $E_2 = E_3 = 11.4136$
GPa, $G_{12} = G_{13} = 6.8017$ GPa, $G_{23} = 3.0955$ GPa,
$\nu_{12} = \nu_{13} = 0.31178$, $\nu_{23} = 0.57574$) to within
$0.002\,\%$ on every entry.

\begin{figure}[htbp]
\centering
\includegraphics[width=0.55\textwidth]{\FIGDIR/udcomp_2dsg_mesh.png}
\caption{The UD fibre/matrix two-dimensional solid SG of
\texttt{examples/OpenSG-solid/5\_get\_solid\_props\_from\_SG}, 801 nodes and
1536 constant-strain triangles. The equivalent 3-D solid $6\times6$ of this
section is this mesh's law, written to \texttt{UDcomp\_2D\_Deff.out}; the nine
constants quoted in this section are its orthotropic-approximation block, and
they reproduce the published values to $0.002\,\%$.}
\label{fig:udcomp}
\end{figure}

\section{The wall law: per-section $8\times8$ ABDG}
\label{sec:wallabdg}

The msg-shell routes need a plate law for \textbf{each wall of the
cross-section} before they can solve the shell SG at all, and that intermediate
is exported so it can be inspected independently.
\texttt{opensg\_shell.sg\_homo.write\_abdg\_out} writes one $8\times8$
stiffness and its compliance per layup section to
\texttt{<yaml base>\_ABDG.out}. Both \texttt{build\_rm\_bundle} and
\texttt{build\_solid\_bundle} call it, so this file appears next to every shell
\texttt{.out} without your having to ask for it.

The file header states the row convention verbatim:

\begin{Verbatim}[fontsize=\small,frame=single]
 OpenSG msg-shell wall plate laws, one block per section
 rows [eps11 eps22 2eps12 | K11 K22 K12+K21 | 2g13 2g23]
\end{Verbatim}

so the block structure is the same
$[[\mathbf A,\mathbf B,\mathbf 0],[\mathbf B^{\mathsf T},\mathbf D,\mathbf 0],
[\mathbf 0,\mathbf 0,\mathbf G]]$ as the Reissner--Mindlin plate law of
Section~\ref{sec:plate}, with the shell
operators' twisting measure $K_{12}+K_{21}$ in the sixth slot rather than
$\kappa_{12}$. Each block is preceded by its section identity and full layup,
for example the first block of \texttt{iea\_s10\_shell\_ABDG.out}:

\begin{Verbatim}[fontsize=\tiny]
 section 0: layup_0  layup [['gelcoat', 0.0005, 0.0], ['glass_triax', 0.00300294, 0.0],
                            ['glass_uniax', 0.02609467, 0.0], ['glass_triax', 0.00300294, 0.0]]
\end{Verbatim}

Each layup entry is \texttt{[material name, thickness (m), angle (deg)]}, bottom
ply first. That section returns $A_{11} = 1.3701540\times10^{9}$ N/m,
$A_{22} = 5.5764061\times10^{8}$ N/m, $A_{66} = 1.3539842\times10^{8}$ N/m,
$D_{11} = 1.0897268\times10^{5}$ N$\cdot$m,
$G_{11} = 8.2828865\times10^{7}$ N/m and
$G_{22} = 9.8804474\times10^{7}$ N/m, with a non-zero coupling block
($B_{11} = 3.1172920\times10^{5}$ N) because the stack --- gelcoat, triax,
uniax, triax --- is not symmetric about the mid-surface.

There is no separate mesh figure for these blocks, because there are two meshes
behind them and both are already shown. Section 0 is one element set of the
cross-section SG of Figure~\ref{fig:const-shellxsec}, and the block itself is
the solution of the through-thickness one-dimensional SG of that set's layup,
built by \texttt{rm\_plate\_msg} exactly as in Figure~\ref{fig:plate1dsg}. One
$8\times8$ per section is the entire wall-law export.

Two facts about how these blocks are built matter when comparing them against
hand calculations.

\begin{enumerate}
\item The ABD is shifted by the parallel-axis rule to the reference surface the
  YAML declares. \texttt{reference: center} gives the laminate mid-surface and
  sets $\mathrm{fraction} = 0.5$; \texttt{oml} gives $0.0$, \texttt{iml} and
  \texttt{oml\_flip} give $1.0$. The $\mathbf B$ block depends on that choice,
  and so therefore does every extension--bending coupling downstream of it.
\item With the default \texttt{g\_source="msg"} the transverse-shear block
  $\mathbf G$ is \textbf{replaced} by the MSG least-squares result from
  \texttt{rm\_plate\_msg}, evaluated at the same \texttt{fraction}. It is not a
  shear-correction-factor estimate. The alternative \texttt{g\_source="whitney"}
  selects the complementary-energy shear-flow value instead.
\end{enumerate}

\section{Where each law is produced and written}
\label{sec:lawmap}

Table~\ref{tab:lawmap} is the index for everything above: given a macro model,
it names the function that computes the law, the file the law lands in, and the
title printed inside that file.

\begin{table}[htbp]
\centering
\caption{Macro model, producing function, output file, and the matrix title
printed in the \texttt{.out}. Function names are given without their module
prefix: \texttt{plate\_homo\_2d}, \texttt{\_beam\_homo\_kkt},
\texttt{plate\_shear\_ladder} and \texttt{write\_sc\_K} live in
\texttt{opensg\_solid.sg\_homo}, \texttt{rm\_plate\_msg} in
\texttt{opensg\_solid.rm\_plate\_1D.msg\_rm\_plate}, and the rest in
\texttt{opensg\_shell}. \texttt{<sg>} is the basename of the
\texttt{.sc}/\texttt{.yaml} input, \texttt{<yaml>} the basename of the shell
YAML, and \texttt{<db>} the basename of the \texttt{layup\_db.yaml}.}
\label{tab:lawmap}
\scriptsize
\begin{tabular}{p{2.5cm}p{4.6cm}p{3.9cm}p{3.5cm}}
\toprule
macro law & producing function & output file & title in the \texttt{.out} \\
\midrule
Timoshenko beam $6\times6$ & \texttt{plate\_homo\_2d(n\_model=1)} $\rightarrow$ \texttt{\_beam\_homo\_kkt} & \texttt{<sg>.out} & \texttt{The Effective Timoshenko Stiffness Matrix} \\
Timoshenko beam $6\times6$ & \texttt{build\_rm\_bundle} $\rightarrow$ \texttt{ring\_indep} & \texttt{<yaml>\_Timo.out} & \texttt{The Effective Timoshenko Stiffness Matrix} \\
Timoshenko beam $6\times6$ & \texttt{segment\_timo\_from\_3dyaml} & \texttt{<yaml>\_Timo.out} & \texttt{The Effective Timoshenko Stiffness Matrix} \\
Euler--Bernoulli beam $4\times4$ & \texttt{\_beam\_homo\_kkt} $\rightarrow$ \texttt{r["C\_eff\_EB"]} & in-memory only & --- \\
Classical plate ABD $6\times6$ & \texttt{plate\_homo\_2d(n\_model=2)} & \texttt{<sg>.out} & \texttt{The Effective Classical Plate Stiffness Matrix} \\
RM plate ABDG $8\times8$ & \texttt{plate\_homo\_2d(n\_model=2, shear\_refined=True)} $\rightarrow$ \texttt{plate\_shear\_ladder} & \texttt{<sg>.out} & \texttt{The Effective Reissner-Mindlin Plate Stiffness Matrix} \\
RM plate ABDG $8\times8$ & \texttt{rm\_plate\_msg} $\rightarrow$ \texttt{r["ABDG"]} & \texttt{<db>\_plate\_homo.out} & \texttt{The Effective Reissner-Mindlin Plate Stiffness Matrix} \\
3-D solid $6\times6$ & \texttt{plate\_homo\_2d(n\_model=3)} & \texttt{<sg>.out} & \texttt{The Effective Stiffness Matrix} + engineering constants \\
3-D solid $6\times6$ & \texttt{build\_solid\_bundle} $\rightarrow$ \texttt{ring\_solid} & \texttt{<yaml>\_C3D.out} & \texttt{The Effective Stiffness Matrix} + engineering constants \\
3-D solid $6\times6$ & \texttt{sg\_homo.shell\_sg3d} & \texttt{<yaml>\_C3D.out} & \texttt{The Effective Stiffness Matrix} + engineering constants \\
Wall plate law $8\times8$, per section & \texttt{sg\_homo.write\_abdg\_out} & \texttt{<yaml>\_ABDG.out} & \texttt{The Effective Reissner-Mindlin Plate Stiffness Matrix}, one block per section \\
\bottomrule
\end{tabular}
\end{table}

\section{Running the laws backwards: two-step recovery}
\label{sec:recovery}

Homogenization is only half of MSG. Each law above was obtained together with
the fluctuation (warping) fields that produced it, and those fields let the
macro answer be pushed back down to pointwise 3-D stress and strain inside the
SG. The rule is \emph{homogenize once, recover as often as needed}: every field
the recovery needs rides along in the result dict, so no second homogenization
is ever required.

\subsection{The shell chain: beam force to plate resultants to 3-D stress}

The shell chain has two interfaces, and naming them is the clearest way to see
why it works. It is verified end to end in
\texttt{examples/OpenSG\_shell/IEA\_blade\_beam/2\_dehom\_v2\_vabs.py}:

\begin{lstlisting}[style=shellst]
python 2_dehom_v2_vabs.py
\end{lstlisting}

\paragraph{Step 2a --- beam to wall (\texttt{opensg\_shell}).}
The macro beam strain follows from inverting the Timoshenko law of
Section~\ref{sec:timo}, $\mathbf{st} = \mathbf C_6^{-1}\,\mathbf{FF}$, with
\texttt{FF} the six beam force and moment components in VABS order read from
the load table \texttt{ff51\_rmc\_reform.dat}. For the IEA-22 station at
$r/R = 0.20$ that row is

\widematrix{
\mathbf{FF} = \left[\,4.1927\times10^{4},\;8.2674\times10^{3},\;
1.4673\times10^{6},\;1.2191\times10^{6},\;-6.0185\times10^{7},\;
3.0426\times10^{5}\,\right]
}

in N and N$\cdot$m, i.e.\ a flapwise-moment-dominated state. The SG being
recovered is the cross-section of Figure~\ref{fig:const-shellxsec}, the same
mesh that produced the $6\times6$ being inverted here.
\texttt{\_macro\_fields} recombines the ring warping modes $V_0$ and
$V_1$ into the nodal fields $w = V_0\,\mathbf{st}_m + V_1\,\mathbf{st}_{c1}$ and
$w' = V_0\,\mathbf{st}_{c1} + V_1\,\mathbf{st}_{c2}$, and
\texttt{\_rm\_shell\_strain} then evaluates, per ring element, the six shell
strains
$\mathbf s_6 = [\varepsilon_{11}\;\varepsilon_{22}\;2\varepsilon_{12}\;
\kappa_{11}\;\kappa_{22}\;2\kappa_{12}]$ together with the transverse pair
$[2\gamma_{13}\;2\gamma_{23}]$. Multiplying by that wall's own MSG ABD gives
the interface quantity handed to the plate SG:

\begin{equation}
\mathbf F_6 \;=\; \mathbf A_6\,\mathbf s_6
\;=\; [\,N_{11}\;\;N_{22}\;\;N_{12}\;\;M_{11}\;\;M_{22}\;\;M_{12}\,]^{\mathsf T}.
\label{eq:f6}
\end{equation}

These are exactly the resultants of the classical plate law of
Section~\ref{sec:plate}: the beam model has handed each wall a plate problem.
They are written to \texttt{<name>\_elem\_plate\_forces.dat}, one row per ring
element, with columns

\begin{Verbatim}[fontsize=\footnotesize,frame=single]
elem  y2(m) y3(m)  e11 e22 2e12 k11 k22 2k12  N11 N22 N12(N/m) M11 M22 M12(N)
\end{Verbatim}

The first row of \texttt{iea\_s10\_elem\_plate\_forces.dat} reads element 0 at
$(y_2, y_3) = (4.677, 0.105)$ m with $\varepsilon_{11} = 5.427\times10^{-4}$ and
$N_{11} = 7.149\times10^{5}$ N/m --- a trailing-edge panel in tension under the
flapwise moment.

\paragraph{Step 2b --- wall to 3-D (\texttt{opensg\_solid.rm\_plate\_1D}).}
For each distinct layup, \texttt{rm\_plate\_msg} builds the through-thickness SG
once. Then \texttt{msgrm\_strain\_at\_depth} --- or
its vectorized twin \texttt{msgrm\_strain\_at\_depth\_batch}, which is the same
algebra evaluated with batched einsums and is what makes a $10^5$-point section
sweep tractable --- returns the 3-D Voigt strain, the 3-D Voigt stress and the
ply angle at any depth $z$. Passing only the first strain gradients
\texttt{dE1}/\texttt{dE2} selects the first-order recovery (Yu Eq.~63); adding
\texttt{dE11}/\texttt{dE12}/\texttt{dE22} selects the second-order recovery
(Eq.~66), which is what can carry the transverse pair $\sigma_{33}$ and
$\sigma_{13}$ that first order cannot produce. Output is in the ply material
frame by default (\texttt{frame="material"}), since failure criteria live
there; pass \texttt{frame="plate"} for anything that integrates or
differentiates across plies.

\subsection{What the shell chain gets right, and where it does not}

The IEA-22 station at $r/R = 0.20$ is validated against a VABS two-dimensional
solid solution of the same section, sampled at every VABS Gauss point. The RMS
differences, in MPa, web / skin, are recorded in \texttt{iea\_s10\_report.txt}
and reproduced in Table~\ref{tab:dehomrms}.

\begin{table}[htbp]
\centering
\caption{RMS difference against VABS, in MPa, web / skin, for the IEA-22
station at $r/R = 0.20$. The last column is the RMS of the VABS field itself,
so a difference equal to it means the chain returns essentially zero for that
component.}
\label{tab:dehomrms}
\small
\begin{tabular}{lrrr}
\toprule
component & first order & second order (V2) & VABS RMS \\
\midrule
$\sigma_{11}$ & 43.073 / 9.529 & 43.073 / 9.529 & 44.440 / 96.676 \\
$\sigma_{22}$ & 0.612 / 0.482 & 0.612 / 0.482 & 0.704 / 0.998 \\
$\sigma_{12}$ & 6.799 / 1.255 & 6.798 / 1.255 & 22.026 / 6.009 \\
$\sigma_{33}$ & 0.081 / 0.049 & 0.081 / 0.049 & 0.081 / 0.049 \\
$\sigma_{13}$ & 0.394 / 0.192 & 0.394 / 0.192 & 0.394 / 0.193 \\
$\sigma_{23}$ & 0.035 / 0.028 & 0.035 / 0.028 & 0.035 / 0.028 \\
\bottomrule
\end{tabular}
\end{table}

Read the last column together with the others. The dominant $\sigma_{11}$ is
recovered to roughly $10\,\%$ of its own RMS in the skin, and the same holds
along a through-thickness path in the spar cap (38.355 MPa difference against a
VABS path RMS of 285.457 MPa). But the transverse components $\sigma_{33}$,
$\sigma_{13}$ and $\sigma_{23}$ show differences \emph{equal to} the VABS RMS
itself, which means the shell chain returns essentially zero there. That gap in
the transverse pair is the honest open item on this route. It is also why the
second-order drivers exist even though their payload at this near-constant load
state is small: V2 changes $\sigma_{11}$ by $1.95\times10^{-4}$ MPa RMS against
a first-order RMS of $8.851\times10^{1}$ MPa, and its visible effect is on
$\sigma_{33}$, where the first-order field is $3.6\times10^{-13}$ MPa (machine
zero) and V2 produces $5.19\times10^{-5}$ MPa against a VABS RMS of
$5.58\times10^{-2}$ MPa.

The recovered in-plane stress over the full section is shown in
Figure~\ref{fig:dehomstress}: the spar caps in compression above and tension
below under the flapwise moment, with the shear flow visible in $\sigma_{12}$
around the leading edge and the trailing-edge panels.

\begin{figure}[htbp]
\centering
\includegraphics[width=0.95\textwidth]{\FIGDIR/shell_dehom_stress.png}
\caption{Recovered $\sigma_{11}$ and $\sigma_{12}$ over the IEA-22 blade
cross-section at $r/R = 0.20$, plotted at the recovery points around the
contour and the shear webs, from the two-step chain of
Eq.~\eqref{eq:f6} followed by \texttt{msgrm\_strain\_at\_depth\_batch}.}
\label{fig:dehomstress}
\end{figure}

\subsection{The solid chain}

For an SG homogenized through \texttt{plate\_homo\_2d}, recovery is a single
call in \texttt{opensg\_solid.sg\_dehom}. From
\texttt{examples/OpenSG-solid/4\_get\_plate\_dehom\_from\_2DSG}:

\begin{lstlisting}[style=opensg,caption={The plate recovery driver plate\_dehom\_2dsg.py, reduced to its essentials. epsilon\_bar is the user's macro state at the point of interest.}]
from opensg_solid.sg_homo import plate_homo_2d
from opensg_solid.sg_dehom import dehom_fields, export_gauss

r = plate_homo_2d("RHC_SW_2UC_45.yaml", material_param=material_param,
                  angles=angles, n_model=2)
# the MACRO plate strain [e11, e22, 2e12, k11, k22, 2k12]
epsilon_bar = jnp.array([0.0, 0.1, 0.0, 0.1, 0.0, 0.0])
Gam, Sig, U = dehom_fields(r, epsilon_bar)
export_gauss(r, Gam, Sig, "RHC_SW_2UC_45_dehom", U_eqd=U)
\end{lstlisting}

\texttt{epsilon\_bar} is whichever macro state matches the model that was
homogenized: the six plate strains
$[\varepsilon_{11}\;\varepsilon_{22}\;2\varepsilon_{12}\;\kappa_{11}\;
\kappa_{22}\;2\kappa_{12}]$ for \texttt{n\_model=2}, the six macro strains of
Section~\ref{sec:solid} for \texttt{n\_model=3}, or the six Timoshenko beam
strains of Section~\ref{sec:timo} for \texttt{n\_model=1}. The beam case is run by
\texttt{examples/OpenSG-solid/8\_get\_beam\_dehom\_from\_SG/beam\_dehom\_sg.py},
whose shipped state is
\verb|epsilon_bar = jnp.array([0.001, 0.0, 0.0, 0.0, 0.01, 0.0])|, i.e.\ a unit
axial extension of $10^{-3}$ combined with a $\kappa_2$ curvature of $10^{-2}$.
For \texttt{n\_model=1} the valid recovery drivers are the classical-channel
states (extension, twist, bending and their combinations); the shear channels
do not drive the recovery chain.

The call returns, at every element Gauss point, the local strain \texttt{Gamma}
of shape $(E, Q, 6)$, the local stress \texttt{Sigma} of shape $(E, Q, 6)$ ---
both in SwiftComp storage order $(xx, yy, zz, yz, xz, xy)$ --- and the
fluctuation-only displacement \texttt{U} of shape $(E, Q, 3)$.
\texttt{plate\_dehom\_2d} is a two-output view of the same single pass, and
\texttt{gauss\_coords(r)} returns the physical coordinates of the recovery
points as $(E \cdot Q, n_{sg})$.

\texttt{export\_gauss} writes five files from one recovery:
\texttt{<prefix>.txt} (the tab-separated Gauss table, one row per point with
coordinates, six strains and six stresses), \texttt{<prefix>.vtk} (a point
cloud for ParaView carrying all twelve components as scalars), and the OpenSG
dehom files \texttt{<prefix>.SM} (stress), \texttt{<prefix>.EM} (strain) and
\texttt{<prefix>.U} (fluctuation displacement). One convention trap is worth
flagging: the printed component order in the \texttt{.SM}, \texttt{.EM} and
\texttt{.U} files is $11\;22\;33\;12\;13\;23$, not the SwiftComp order. The
reorder happens on write only; the arrays you hold in memory stay in
$(xx, yy, zz, yz, xz, xy)$. In those three files the SG coordinates occupy the
\emph{last} $n_{sg}$ of the $(x, y, z)$ columns, with the leading slots zero, so
that the last SG coordinate --- the plate thickness direction --- always lands
in the $z$ column.

\chapter{Tutorials: the Shell Engine (msg-shell)}
\label{ch:shell}

The \texttt{msg-shell} engine solves structure genes whose heterogeneity lives
in a \emph{surface}: the mid-surface of a thin composite wall. Thickness never
appears in the mesh. It lives inside the wall's Reissner--Mindlin plate law and
inside the through-thickness recovery ladder, so changing a ply schedule is a
change of material data, not a remesh.

Five capabilities ship as runnable examples, all under
\texttt{examples/OpenSG\_shell/}. Every one of them takes a YAML file in and
writes a timed output file in the SwiftComp \texttt{.K} layout.

\begin{table}[htbp]
\centering
\small
\begin{tabular}{llll}
\toprule
example folder & driver & entry point & produces \\
\midrule
\texttt{1\_get\_beam\_props\_...}   & \texttt{beam\_homo\_shell.py}  & \texttt{build\_rm\_bundle}          & Timoshenko $6\times6$ \\
\texttt{2\_get\_beam\_dehom\_...}   & \texttt{beam\_dehom\_shell.py} & \texttt{build\_rm\_bundle} + recovery & 3-D wall stress \\
\texttt{3\_get\_solid\_props\_...}  & four sub-folders               & \texttt{build\_solid\_bundle}       & equivalent solid $6\times6$ \\
\texttt{4\_get\_solid\_props\_...}  & \texttt{schwarz\_solid\_props.py} & \texttt{shell\_sg3d}             & equivalent solid $6\times6$ \\
\texttt{5\_taper\_segment\_...}     & \texttt{run\_taper\_segment.py} & \texttt{segment\_timo\_from\_3dyaml} & tapered-segment $6\times6$ \\
\bottomrule
\end{tabular}
\caption{The five \texttt{msg-shell} capabilities and the example folder that
demonstrates each one. Folder names are abbreviated; the full names appear at
the head of each section below.}
\label{tab:shell-capabilities}
\end{table}

%=====================================================================
\section{Beam properties from a shell cross-section}
\label{sec:shell-beam}

\subsection{What this capability computes}

Given the contour of a thin-walled composite cross-section, this run produces
the $6\times6$ Timoshenko beam stiffness matrix of that section: the constitutive
law a beam solver such as BeamDyn or GEBT needs. It also writes the
$8\times8$ Reissner--Mindlin plate law of every distinct wall region, because
those are the intermediate quantities the beam matrix is assembled from.

The example is \texttt{examples/OpenSG\_shell/1\_get\_beam\_props\_from\_shell\_cross\_section},
and the input is \texttt{iea\_s10\_shell.yaml}: station $r/R = 0.20$ of the
IEA-22 MW reference blade.

\subsection{What a 1-D shell SG is}

A \textbf{1-D shell structure gene} is the contour of the cross-section: a chain
of straight two-node line elements along the wall mid-surface, each element
carrying a laminate stack as material data rather than as geometry. A 2-D solid
cross-section mesh, by contrast, discretizes the wall \emph{through} its
thickness with area elements, so every ply interface must be resolved by the
mesh.

The consequence for the user is size. This station spans about 7.2\,m in $y_2$
and 2.5\,m in $y_3$, has spar caps 87.6\,mm thick and three shear webs, and is
fully described by 307 nodes and 310 line elements. The price is that the walls
are modelled as Reissner--Mindlin plates, so wall-thickness effects enter through
the plate law and the recovery ladder rather than through the mesh.

\subsection{The input file the user edits}

\texttt{build\_rm\_bundle} reads exactly one file. These are its blocks, with the
values found in \texttt{iea\_s10\_shell.yaml}.

\begin{table}[htbp]
\centering
\small
\begin{tabular}{lll}
\toprule
block & content & this station \\
\midrule
\texttt{nodes} & \texttt{[y2 y3 0]} contour coordinates, metres & 307 nodes \\
\texttt{elements} & 1-based two-node line connectivity & 310 elements \\
\texttt{sets: element} & named element sets, one per wall region & 6 sets \\
\texttt{sections} & per set: \texttt{elementSet} plus \texttt{layup} as & 6 sections \\
                  & \texttt{[material, thickness, angle]} triples, outer ply first & \\
\texttt{elementOrientations} & per element, the flattened $3\times3$ frame $[e_1\;e_2\;e_3]$ & 310 rows \\
\texttt{materials} & \texttt{name}, \texttt{density}, \texttt{elastic: \{E, G, nu\}} as 3-vectors & 6 materials \\
\texttt{reference} & the reference surface the whole run is referred to & \texttt{center} \\
\bottomrule
\end{tabular}
\caption{The seven YAML blocks read by \texttt{build\_rm\_bundle}.}
\label{tab:shell-yaml-blocks}
\end{table}

The \texttt{reference} field is the single source of truth. \texttt{build\_rm\_bundle}
reads it and maps it to a thickness fraction (\texttt{center} $\to$ 0.5,
\texttt{oml} $\to$ 0.0, \texttt{iml} and \texttt{oml\_flip} $\to$ 1.0). That one
number is then used consistently for the ring laminate reference, the plate-SG
$z$ origin, the emitted ABD cache and the recovery depth conversion. Pass
\texttt{ref=} explicitly only to override it. Because the reference here is the
mid-surface, quantities that depend on the reference axis, in particular $GJ$ and
the bending--extension couplings, are mid-surface values.

\begin{figure}[htbp]
\centering
\includegraphics[width=0.95\textwidth]{\FIGDIR/shell_xsec_mesh.png}
\caption{The 1-D shell SG of IEA-22 station $r/R = 0.20$: 307 nodes and 310 line
elements along the wall mid-surface, coloured by \texttt{sets: element} entry.
This is the figure written by \texttt{beam\_homo\_shell.py} as
\texttt{iea\_s10\_mesh.png}.}
\label{fig:shell-xsec-mesh}
\end{figure}

Each colour in Figure~\ref{fig:shell-xsec-mesh} is one element set, and each set
carries one laminate stack.

\begin{table}[htbp]
\centering
\footnotesize
\begin{tabular}{llrr}
\toprule
layup & plies, outer ply first & total thickness (m) & elements \\
\midrule
\texttt{layup\_0} & gelcoat, glass\_triax, glass\_uniax, glass\_triax & 0.03260055 & 25 \\
\texttt{layup\_1} & gelcoat, glass\_triax, medium\_density\_foam, glass\_triax & 0.07711403 & 161 \\
\texttt{layup\_2} & gelcoat, glass\_triax, carbon\_uniax (0.08106 m), glass\_triax & 0.08756579 & 13 \\
\texttt{layup\_3} & gelcoat, glass\_triax, glass\_uniax, glass\_triax & 0.02808872 & 8 \\
\texttt{layup\_4} & gelcoat, glass\_triax, carbon\_uniax (0.07707 m), glass\_triax & 0.08358040 & 14 \\
\texttt{layup\_5} & glass\_biax, medium\_density\_foam, glass\_biax & 0.04600000 & 89 \\
\bottomrule
\end{tabular}
\caption{The six laminate stacks of station \texttt{iea\_s10}, as read from the
\texttt{sections} block and echoed in the header of each block of
\texttt{iea\_s10\_shell\_ABDG.out}.}
\label{tab:shell-layups}
\end{table}

\texttt{layup\_2} and \texttt{layup\_4} are the carbon spar caps on the upper and
lower surfaces, \texttt{layup\_1} is the foam-cored skin that covers most of the
contour, \texttt{layup\_5} is the three shear webs, \texttt{layup\_0} is the
trailing-edge reinforcement and \texttt{layup\_3} the leading edge. Every ply
angle in this file is $0^\circ$; the biax and triax fabrics are supplied as
smeared orthotropic materials rather than as individual plies, so the material
frame of the recovery coincides with the wall frame here.

\subsection{Orientation frames}

A wrong \texttt{elementOrientations} block is the most common way to get a
plausible-looking but wrong answer, so the example driver calls
\texttt{auto\_emit} explicitly after the homogenization and writes an orientation
figure. Note that \texttt{auto\_emit} is called \emph{by the driver}, not from
inside \texttt{build\_rm\_bundle}: if you write your own script around
\texttt{build\_rm\_bundle}, add the \texttt{auto\_emit} call yourself.

$e_1$ is the beam axis (out of the page), $e_2$ (blue) is the in-plane ply-flow
direction along the contour, and $e_3$ (green) is the wall normal, which sets the
OML $\to$ IML sense used by the recovery.

\begin{figure}[htbp]
\centering
\includegraphics[width=0.95\textwidth]{\FIGDIR/shell_xsec_orient.png}
\caption{Per-element $e_2$ (in-plane ply direction) and $e_3$ (wall normal) on
the contour, written by \texttt{auto\_emit} as \texttt{iea\_s10\_orient.png}.
Check that the normals point consistently into the section on the skin and
consistently to one side on each web before trusting any stiffness number.}
\label{fig:shell-xsec-orient}
\end{figure}

\subsection{The element and the wall law}

The ring is assembled as a one-quad-deep prismatic strip: the contour is extruded
by the mean element length along the beam axis, and the extruded node row is
DOF-mapped back onto the original row, which enforces the prismatic
(span-invariant) condition exactly. All matrices are then divided by that depth,
so the result is per unit length. The strip actually assembled is returned in the
bundle under the key \texttt{strip} as \texttt{(nodes, quads, h)}.

Each node carries \textbf{six} degrees of freedom,
$[w_1\;w_2\;w_3\;\omega_1\;\omega_2\;\omega_3]$. The drilling rotation
$\omega_3$ is an \emph{independent} field rather than an algebraic function of the
displacements. Making it independent is what allows adjacent walls meeting at a
junction to share a consistent rotation, but it also introduces a spurious mode,
which is removed by a weak drilling constraint carried by one element-constant
Lagrange multiplier per element (\texttt{lam\_space="elem"}). Element-constant is
the inf-sup-stable choice; equal-order nodal multipliers over-constrain as the
mesh is refined.

Transverse shear uses assumed-strain (MITC) tying, and the production ring scheme
(\texttt{shear="mitc4\_g23"}, the default of \texttt{build\_rm\_bundle}) ties
\textbf{only} the $\gamma_{23}$ row. The reason is specific to a prismatic strip:
under span invariance $\gamma_{13}$ carries no fluctuation gradient and is
algebraic in the directors, so tying it would alias the drilling mode instead of
curing shear locking. Only $\gamma_{23}$ pairs a differentiated displacement with
a rotation.

The wall constitutive law is therefore the $8\times8$ Reissner--Mindlin plate law
\begin{equation}
\mathbf{ABDG} =
\begin{bmatrix}
\mathbf{A} & \mathbf{B} & \mathbf{0}\\
\mathbf{B} & \mathbf{D} & \mathbf{0}\\
\mathbf{0} & \mathbf{0} & \mathbf{G}
\end{bmatrix},
\qquad
\bar{\varepsilon} =
\begin{bmatrix}
\varepsilon_{11} & \varepsilon_{22} & 2\varepsilon_{12} &
K_{11} & K_{22} & K_{12}+K_{21} & 2\gamma_{13} & 2\gamma_{23}
\end{bmatrix}^{T},
\label{eq:shell-abdg}
\end{equation}
with exactly that row grouping printed in the header of
\texttt{iea\_s10\_shell\_ABDG.out}:
\texttt{rows [eps11 eps22 2eps12 | K11 K22 K12+K21 | 2g13 2g23]}.

The $2\times2$ transverse-shear block $\mathbf{G}$ has a default,
\texttt{g\_source="msg"}: the MSG/VAM second-order-energy projection (the
Yu--2002 least-squares construction, Eq.~61 of Yu, Hodges \& Volovoi,
\emph{Computers \& Structures} \textbf{81}:439--454, 2003), computed by
\texttt{opensg\_solid.rm\_plate\_1D.msg\_rm\_plate.rm\_plate\_msg}. The
alternative, \texttt{g\_source="whitney"}, is the coupling-aware
complementary-energy shear flow. Table~\ref{tab:shell-gblocks} shows why this is
not a detail on a foam-cored blade section.

\begin{table}[htbp]
\centering
\small
\begin{tabular}{lrr}
\toprule
layup & $G_{11}$ (N/m) & $G_{22}$ (N/m) \\
\midrule
\texttt{layup\_0} & $8.2828865\times10^{7}$ & $9.8804474\times10^{7}$ \\
\texttt{layup\_1} & $4.1457954\times10^{6}$ & $4.1451801\times10^{6}$ \\
\texttt{layup\_2} & $2.8674578\times10^{8}$ & $2.1272143\times10^{8}$ \\
\texttt{layup\_3} & $5.8528099\times10^{7}$ & $8.6668255\times10^{7}$ \\
\texttt{layup\_4} & $2.7301565\times10^{8}$ & $2.0317797\times10^{8}$ \\
\texttt{layup\_5} & $2.4877763\times10^{6}$ & $2.4877763\times10^{6}$ \\
\bottomrule
\end{tabular}
\caption{The MSG wall transverse-shear block $\mathbf{G}$ of each layup, rows 7
and 8 of the $8\times8$ blocks in \texttt{iea\_s10\_shell\_ABDG.out}. The
foam-cored skin \texttt{layup\_1} and the foam-cored webs \texttt{layup\_5} are
roughly seventy and a hundred times more shear-compliant than the carbon cap
\texttt{layup\_2}, so the soft core dominates the section's shear compliance and
must not be assigned by a rule of thumb.}
\label{tab:shell-gblocks}
\end{table}

\subsection{Run it}

The driver has a single user-input line:

\begin{lstlisting}[style=opensg,caption={The user-input block of beam\_homo\_shell.py.}]
############### User Input #################################
name = "iea_s10"               # reads <name>_shell.yaml
############################################################
\end{lstlisting}

\begin{lstlisting}[style=shellst]
cd examples/OpenSG_shell/1_get_beam_props_from_shell_cross_section
python beam_homo_shell.py
\end{lstlisting}

\subsection{The one call}

Everything is behind one function.

\begin{lstlisting}[style=opensg,caption={The minimal script: build the bundle, take the 6x6, emit the orientation figure.}]
import numpy as np
from opensg_shell import build_rm_bundle, auto_emit

B  = build_rm_bundle("iea_s10_shell.yaml")   # RM 6-DOF ring + MSG wall G
C6 = np.asarray(B["Timo"])                   # (6, 6) Timoshenko stiffness
print(B["ref"], B["g_source"])               # "center msg"
print(np.diag(C6))
auto_emit("iea_s10_shell.yaml", out_png="iea_s10_orient.png")
\end{lstlisting}

\texttt{build\_rm\_bundle} returns a \textbf{bundle}, not just a matrix, because
the dehomogenization of Section~\ref{sec:shell-dehom} must reuse the identical
discretization. The bundle carries the keys of
Table~\ref{tab:shell-bundle-keys}.

\begin{table}[htbp]
\centering
\small
\begin{tabular}{ll}
\toprule
key & content \\
\midrule
\texttt{Timo} & the $(6,6)$ Timoshenko stiffness \\
\texttt{V0}, \texttt{V1} & zeroth- and first-order warping fields, $(6m,4)$ each, \\
                         & including the drilling $\omega_3$ boundary values \\
\texttt{corners}, \texttt{red\_cells} & the $(m,2)$ contour coordinates and $(n_{el},2)$ connectivity \\
\texttt{rx3}, \texttt{re3}, \texttt{k22} & 3-D ring coordinates, wall normals and hoop curvature per element \\
\texttt{strip} & the prismatic strip actually assembled, as \texttt{(nodes, quads, h)} \\
\texttt{ax}, \texttt{cross} & the beam-axis index and the two cross-section axis indices \\
\texttt{layup\_per\_elem} & the layup name of each element \\
\texttt{layup\_db}, \texttt{material\_db} & the geometry-free plate-SG databases \\
\texttt{frac}, \texttt{ref}, \texttt{g\_source} & the reference fraction, its name, and the wall-$G$ source \\
\bottomrule
\end{tabular}
\caption{Keys of the dictionary returned by \texttt{build\_rm\_bundle}.}
\label{tab:shell-bundle-keys}
\end{table}

\subsection{What comes out}

\begin{table}[htbp]
\centering
\small
\begin{tabular}{lll}
\toprule
file & written by & content \\
\midrule
\texttt{iea\_s10\_shell\_Timo.out} & \texttt{build\_rm\_bundle} & the effective Timoshenko stiffness and \\
                                   & (via \texttt{write\_sc\_K}) & compliance in the SwiftComp \texttt{.K} layout, \\
                                   & & with an \texttt{OpenSG} banner and a \texttt{Time taken} footer \\
\texttt{iea\_s10\_shell\_ABDG.out} & \texttt{build\_rm\_bundle} & one $8\times8$ RM plate stiffness and compliance \\
                                   & (via \texttt{write\_abdg\_out}) & block per section, each headed by its layup \\
\texttt{abd/iea\_s10\_shell\_abd.yaml} & \texttt{emit\_station\_abd} & per-layup cache of the $8\times8$ wall law plus ply \\
                                   & & lists, thickness and mass per area \\
\texttt{iea\_s10\_mesh.png} & the driver & the contour coloured by layup \\
\texttt{iea\_s10\_orient.png} & \texttt{auto\_emit} & the $e_2$/$e_3$ arrows \\
\bottomrule
\end{tabular}
\caption{Output files of \texttt{beam\_homo\_shell.py}.}
\label{tab:shell-beam-outputs}
\end{table}

Two notes on the ABD cache. It is written \textbf{only if it does not already
exist}, so delete \texttt{abd/iea\_s10\_shell\_abd.yaml} after changing the
layups or the change will not take. And it records \texttt{g\_source: whitney}
and \texttt{reference: mid}, because \texttt{emit\_station\_abd} defaults to the
complementary-energy shear flow. It is a cache for downstream station-level tools
such as shell buckling, not a dump of the ring's own MSG wall law, which lives in
\texttt{\_ABDG.out}.

\subsection{The numbers}

This is the stiffness block of \texttt{iea\_s10\_shell\_Timo.out} exactly as
written, in SI units, on the strain set
$[\varepsilon_{11}\;\gamma_{12}\;\gamma_{13}\;\kappa_1\;\kappa_2\;\kappa_3]$
(extension, the two transverse shears, twist, and the two bending curvatures,
which is the VABS ordering):

\begin{Verbatim}[fontsize=\footnotesize,frame=single,framesep=4pt]
 The Effective Timoshenko Stiffness Matrix
 --------------------------------------------
     2.7660599E+010    -1.7230658E-005     4.6650498E-006     5.4020232E-006     1.6686413E+009    -2.0773040E+009
    -1.7230658E-005     7.1861326E+008     5.7520370E+007    -1.0192044E+008     2.3079953E-006     5.2738749E-006
     4.6650498E-006     5.7520370E+007     4.2174428E+008    -1.6217837E+008     8.4816767E-006     2.7931393E-007
     5.4020232E-006    -1.0192044E+008    -1.6217837E+008     2.4375653E+009     1.6449557E-006    -8.1140553E-006
     1.6686413E+009     2.3079953E-006     8.4816767E-006     1.6449557E-006     3.5143887E+010    -5.5571249E+009
    -2.0773040E+009     5.2738749E-006     2.7931393E-007    -8.1140553E-006    -5.5571249E+009     6.8129724E+010
\end{Verbatim}

The diagonal is what the driver prints.

\begin{table}[htbp]
\centering
\small
\begin{tabular}{llll}
\toprule
term & value & units & meaning \\
\midrule
$EA$    & $2.7660599\times10^{10}$ & N               & axial stiffness \\
$GA_2$  & $7.1861326\times10^{8}$  & N               & transverse shear, chordwise \\
$GA_3$  & $4.2174428\times10^{8}$  & N               & transverse shear, flapwise \\
$GJ$    & $2.4375653\times10^{9}$  & N$\cdot$m$^2$   & torsion about the mid-surface reference \\
$EI_2$  & $3.5143887\times10^{10}$ & N$\cdot$m$^2$   & flapwise bending \\
$EI_3$  & $6.8129724\times10^{10}$ & N$\cdot$m$^2$   & edgewise bending \\
\bottomrule
\end{tabular}
\caption{The Timoshenko diagonal of station \texttt{iea\_s10}, read from
\texttt{iea\_s10\_shell\_Timo.out}.}
\label{tab:shell-timo-diag}
\end{table}

The off-diagonals are the physics that a beam table without couplings would throw
away. The extension--bending entries $1.6686\times10^{9}$ and
$-2.0773\times10^{9}$ N$\cdot$m locate the tension centre away from the
mid-surface reference point; $-5.5571\times10^{9}$ N$\cdot$m$^2$ is the
flap--edge bending coupling, which says the section's principal bending axes are
not the $y_2$/$y_3$ axes; and $-1.0192\times10^{8}$ and $-1.6218\times10^{8}$
N$\cdot$m are the shear--twist couplings that place the shear centre.

Entries printed at $10^{-5}$ against diagonals of $10^{10}$ are exact zeros to
about fifteen digits: extension does not couple to transverse shear or to twist
on this section. That is a useful self-check that the drilling constraint and the
rigid-body kernel are behaving.

The footer of the same file reads \texttt{Time taken: 8.02 sec} for this run.
That covers the whole bundle build: reading the YAML, the six per-layup MSG plate
solves, the ring KKT factorization and the layup-database build.

%=====================================================================
\section{Two-step recovery to 3-D wall stress}
\label{sec:shell-dehom}

\subsection{What this capability computes}

Homogenization gave the beam its stiffness. \emph{Dehomogenization} does the
inverse: given the beam's internal forces at this station, it reconstructs the
3-D stress at every point of every wall, which is what a failure criterion
actually needs. The example is
\texttt{examples/OpenSG\_shell/2\_get\_beam\_dehom\_from\_shell\_cross\_section}
and runs on the same \texttt{iea\_s10\_shell.yaml}.

\subsection{The input the user edits}

\begin{lstlisting}[style=opensg,caption={The user-input block of beam\_dehom\_shell.py.}]
############### User Input #################################
name = "iea_s10"                  # reads <name>_shell.yaml
ff_file = "ff51_rmc_reform.dat"   # beam FF per station (# eta F1 F2 F3 M1 M2 M3)
station_row = 10                  # row 10 -> eta = r/R = 0.20 (iea_s10)
n_depth = 9                       # through-thickness recovery points per element
stress_frame = "material"         # "material" (ply axes) | "plate" (wall axes)
############################################################
\end{lstlisting}

\texttt{ff51\_rmc\_reform.dat} is a plain table whose header line reads
\texttt{\# eta F1 F2 F3 M1 M2 M3 (RM CENTER-ref BeamDyn, VABS order)}: one row per
spanwise station, the first column $\eta = r/R$ and the remaining six the beam
force and moment resultants, in the same VABS order and the same SI units (N,
N$\cdot$m) as the $6\times6$. Row 10 is
\begin{equation}
FF = \begin{bmatrix}
4.1927\times10^{4} & 8.2674\times10^{3} & 1.4673\times10^{6} &
1.2191\times10^{6} & -6.0185\times10^{7} & 3.0426\times10^{5}
\end{bmatrix},
\label{eq:shell-ff}
\end{equation}
a 60\,MN$\cdot$m flapwise moment with a 1.47\,MN flapwise shear, which is what
dominates everything that follows. The forces must be referred to the same axis
as the stiffness: this table is center-referenced, matching
\texttt{reference: center} in the YAML.

\subsection{Run it}

\begin{lstlisting}[style=shellst]
cd examples/OpenSG_shell/2_get_beam_dehom_from_shell_cross_section
python beam_dehom_shell.py
\end{lstlisting}

\subsection{The two-step chain}

\textbf{Step 1, along the contour.} The macro beam strain is
$\mathrm{st} = \mathbf{C}_6^{-1}FF$. That strain is fed through
\texttt{\_macro\_fields}, which recombines the bundle's warping fields into the
nodal warping $w = V_0\,\mathrm{st}_m + V_1\,\mathrm{st}_{c1}$ and its axial
derivative $w' = V_0\,\mathrm{st}_{c1} + V_1\,\mathrm{st}_{c2}$, and then through
\texttt{\_rm\_shell\_strain}, which applies the \emph{same} element operators
used to assemble the stiffness. The result is, per element, the six shell strains
$[\varepsilon_{11}\;\varepsilon_{22}\;2\varepsilon_{12}\;K_{11}\;K_{22}\;2K_{12}]$
at mid-arc. Because it reuses \texttt{quad\_ops\_indep}, step 1 is the exact
energy-consistent adjoint of the homogenization rather than a separate
post-processing model.

The recovery also needs the \emph{gradients} of those strains. The span gradient
$\partial/\partial y_1$ comes from the operators directly; the arc gradient
$\partial/\partial y_2$ is a finite difference of nodally averaged values, and
the averaging is \textbf{layup-boundary aware}. Two element values are averaged
at a node only when both elements carry the same layup, so no difference is ever
taken across a section change where the strains legitimately jump.

\textbf{Step 2, through the thickness.} For each of the 6 layups,
\texttt{rm\_plate\_msg} builds a 1-D MSG-RM plate structure gene (the warping
ladder plus the MSG $\mathbf{G}$), and \texttt{msgrm\_strain\_at\_depth}
evaluates the first-order recovery at a depth $z$ measured from the reference
surface, returning the 3-D Voigt strain, the 3-D Voigt stress
$[\sigma_{11}\;\sigma_{22}\;\sigma_{33}\;\sigma_{23}\;\sigma_{13}\;\sigma_{12}]$
and the ply angle at that depth. Both functions are imported from
\texttt{opensg\_solid.rm\_plate\_1D.msg\_rm\_plate}, shared verbatim with the
plate and solid examples.

The driver evaluates \texttt{n\_depth = 9} ply-interior stations per element, at
$\zeta = (k + \tfrac12)/9$ with $\zeta = 0$ at the OML and $\zeta = 1$ at the
IML, giving $310 \times 9 = 2790$ recovery points.

\begin{lstlisting}[style=opensg,caption={The recovery chain in miniature: one element, one depth.}]
import numpy as np
from opensg_shell import build_rm_bundle, _macro_fields, _rm_shell_strain
from opensg_solid.rm_plate_1D.msg_rm_plate import rm_plate_msg, msgrm_strain_at_depth

B  = build_rm_bundle("iea_s10_shell.yaml")
FF = np.loadtxt("ff51_rmc_reform.dat")[10, 1:]
st, st_m, aA, aB = _macro_fields(B, beam_force_vabs=FF)
s6, s2 = _rm_shell_strain(B, 75, 0.5, st_m, aA, aB)   # element 75, mid-arc

ldb = B["layup_db"]; mdb = B["material_db"]; frac = float(B["frac"])
warp = rm_plate_msg(ldb["layup_2"]["thick"], ldb["layup_2"]["angles"],
                    ldb["layup_2"]["mat_names"], mdb, fraction=frac)
h = float(sum(ldb["layup_2"]["thick"]))
Gam, Sig, ply = msgrm_strain_at_depth(warp, (0.056 - frac) * h,
                                      np.asarray(s6, float), frame="material")
print(Sig / 1e6)      # MPa, Voigt [S11 S22 S33 S23 S13 S12]
\end{lstlisting}

\subsection{What comes out}

\begin{table}[htbp]
\centering
\small
\begin{tabular}{ll}
\toprule
file & content \\
\midrule
\texttt{iea\_s10\_dehom\_shell.txt} & 2790 rows, one per recovery point: \texttt{elem}, \texttt{y2}, \texttt{y3}, \\
 & \texttt{z} from the reference surface, \texttt{zeta}, \texttt{ply} angle in degrees, \\
 & then the six stresses in MPa. The header repeats the frame, $\eta$ and the full \texttt{FF}. \\
\texttt{iea\_s10\_dehom\_stress.png} & the section stress cloud, $\sigma_{11}$ and $\sigma_{12}$ \\
\texttt{iea\_s10\_cap\_tt.png} & $\sigma_{11}$ and $\sigma_{12}$ against $\zeta$ along one through-thickness \\
 & line at the arc-centre element of the peak-stress cap \\
\bottomrule
\end{tabular}
\caption{Output files of \texttt{beam\_dehom\_shell.py}.}
\label{tab:shell-dehom-outputs}
\end{table}

The driver reports the peak as

\begin{Verbatim}[fontsize=\small,frame=single,framesep=4pt]
peak |sigma11| = 318.18 MPa at (y2, y3) = (-0.677, 1.311), zeta = 0.06  [material frame]
\end{Verbatim}

\subsection{Reading the result}

\begin{figure}[htbp]
\centering
\includegraphics[width=0.82\textwidth]{\FIGDIR/shell_dehom_stress.png}
\caption{Recovered $\sigma_{11}$ (top) and $\sigma_{12}$ (bottom) over the
cross-section, all 2790 recovery points, in MPa. Points sit slightly off the
contour because each one is displaced from the mid-surface by its own depth $z$
along the wall normal, which is why the cap points reach $y_3 = 1.311$\,m while
the contour stops at $1.273$\,m.}
\label{fig:shell-dehom-stress}
\end{figure}

The axial stress is confined almost entirely to the two spar caps, exactly as a
flap-dominated load state should be.

\begin{table}[htbp]
\centering
\small
\begin{tabular}{lll}
\toprule
quantity & value & where \\
\midrule
minimum $\sigma_{11}$ & $-318.177$ MPa & element 75, \texttt{layup\_2}, $(y_2,y_3) = (-0.677, 1.311)$, $\zeta = 0.056$ \\
maximum $\sigma_{11}$ & $+314.461$ MPa & element 150, \texttt{layup\_4}, $(y_2,y_3) = (-0.145, -1.249)$, $\zeta = 0.056$ \\
maximum $|\sigma_{12}|$ & $5.573$ MPa & element 11, \texttt{layup\_0}, $(y_2,y_3) = (3.895, 0.335)$ \\
\bottomrule
\end{tabular}
\caption{Extremes of \texttt{iea\_s10\_dehom\_shell.txt} under the load state of
Eq.~\eqref{eq:shell-ff}.}
\label{tab:shell-dehom-extremes}
\end{table}

The upper cap is in compression and the lower cap in tension under this
flap-dominated state, at nearly equal magnitude, and both extremes fall at the
outermost recovery station $\zeta = 0.056$, which for these layups lies just
inside the carbon uniaxial ply. That is where a maximum-stress check on the
carbon belongs. The in-plane shear tells a different story: it peaks not in the
caps but in the trailing-edge and leading-edge regions, where the closed cell
carries the torque as shear flow.

\subsection{What this route recovers, and what it does not}

The repository's \texttt{examples/OpenSG\_shell/IEA\_blade\_beam} benchmark folder
compares this same station's recovery against a VABS 2-D solid run of the
identical section (\texttt{iea\_s10.sg.SM}), and records the result in
\texttt{iea\_s10\_report.txt} as rms differences alongside the rms magnitude of
the VABS field itself.

\begin{table}[htbp]
\centering
\small
\begin{tabular}{llrr}
\toprule
component & region & rms difference (MPa) & rms of the VABS field (MPa) \\
\midrule
$\sigma_{11}$ & skin & 9.529 & 96.676 \\
$\sigma_{11}$ & web & 43.073 & 44.440 \\
$\sigma_{12}$ & web & 6.799 & 22.026 \\
$\sigma_{11}$ & spar-cap through-thickness path & 38.355 & 285.457 \\
$\sigma_{11}$ & cap-left junction path & 1.265 & 286.837 \\
\bottomrule
\end{tabular}
\caption{Two-step shell recovery against a VABS 2-D solid run of the same
section, from \texttt{iea\_s10\_report.txt}.}
\label{tab:shell-dehom-vabs}
\end{table}

The in-plane picture is sound where the load is: about 10\,\% rms on the skin,
13\,\% along the cap thickness, and 0.4\,\% along the reference junction path.
The web $\sigma_{11}$ difference is the same size as the field itself, which is
honest to state plainly: the webs carry almost no axial stress at this state, so
the comparison there is a ratio of two small numbers.

Two components are \textbf{not} delivered by this route and should not be used
from the output file:

\begin{itemize}
\item $\sigma_{33}$ is machine-zero in the first-order recovery (rms
  $3.6\times10^{-13}$ MPa against a VABS rms of $5.6\times10^{-2}$ MPa). This is
  visible directly in the \texttt{S33} column of
  \texttt{iea\_s10\_dehom\_shell.txt}, whose entries are of order $10^{-13}$
  (median $1.96\times10^{-13}$, maximum $9.18\times10^{-13}$).
\item $\sigma_{13}$ is recovered at rms $6.3\times10^{-3}$ MPa against a VABS rms
  of $2.4\times10^{-1}$ MPa. This is structural, not a bug: on a prismatic ring
  the $\gamma_{13}$ row is untied by construction, so it carries the flat-wall
  drilling mode and cannot serve as a recovery target. The companion diagnostic
  \texttt{iea\_s10\_q\_consistency.txt} quantifies this. The tied $\gamma_{23}$
  channel closes its equilibrium identity
  $\int \sigma_{23}\,dz = (\mathbf{G}_{\mathrm{msg}}\,s_2)_2$ to 1.05\,\% median
  error, while the untied $\gamma_{13}$ channel does not, giving a $\sigma_{13}$
  target that is invalid by a factor of order 100. A physical $\sigma_{13}$ needs
  the surface-quad route of Sections~\ref{sec:shell-sg3d}
  and~\ref{sec:shell-segment}, where both rows are tied, not a rescale of this
  ring.
\end{itemize}

The same benchmark folder also records that adding the second-order recovery rung
changes $\sigma_{11}$ by an rms of $1.95\times10^{-4}$ MPa against a first-order
rms of 88.51 MPa, so the first-order chain used by \texttt{beam\_dehom\_shell.py}
is the right default for this class of section.

%=====================================================================
\section{Equivalent 3-D solid from a cross-section SG}
\label{sec:shell-solid-xsec}

\subsection{What this capability computes}

The same 1-D shell structure gene can be homogenized to an \textbf{equivalent 3-D
solid} $6\times6$ instead of a beam $6\times6$: the anisotropic continuum
stiffness of a cell that is tiled periodically in the two transverse directions.
This is the MSG thin-walled (MSG-TW) equivalent-continuum problem, and it is what
you want for a cellular solid, a lattice, or any prismatic microstructure whose
walls are thin.

The strain order is the solid Voigt order
\begin{equation}
\bar{\Gamma} = \begin{bmatrix}
\bar\Gamma_{11} & \bar\Gamma_{22} & \bar\Gamma_{33} &
2\bar\Gamma_{23} & 2\bar\Gamma_{13} & 2\bar\Gamma_{12}
\end{bmatrix}^{T},
\label{eq:shell-gbar}
\end{equation}
with axis 1 along the prismatic direction. That string is returned in the bundle
under the key \texttt{order} and exported as \texttt{GBAR\_ORDER}.

Where the beam problem drives the SG with four beam strains, the solid problem
drives it with the six macro strains, so \textbf{periodicity is always on}:
opposite edges and the corners of the 2-D cell are tied
(\texttt{periodic=True} is the default of \texttt{build\_solid\_bundle}). The wall
law is the same Reissner--Mindlin $8\times8$ of Eq.~\eqref{eq:shell-abdg}, so the
walls carry transverse shear that a classical (CLPT/Kirchhoff) wall model
discards. That distinction is the entire point of the benchmarks below.

\subsection{The driver}

\begin{lstlisting}[style=opensg,caption={The one call for equivalent-solid properties from a cross-section SG.}]
from opensg_shell import build_solid_bundle

b = build_solid_bundle("cell_1Dshell.yaml", cell_area=A_cell)
C3D = b["C3D"]            # 6x6 equivalent solid stiffness, Pa
b["cell_area"], b["area_source"]   # the normalizing area and where it came from
b["solve_time"]           # seconds; also written into <yaml>_C3D.out
\end{lstlisting}

\texttt{cell\_area} is the one argument you must think about: it is the measure
$\omega$ the energy is divided by. Pass it explicitly (then
\texttt{area\_source} reads \texttt{"user"}); if you omit it, the convex-hull
area of the contour is used instead and \texttt{area\_source} reads
\texttt{"hull"}. \texttt{build\_solid\_bundle} writes
\texttt{<yaml>\_C3D.out} by default, in the SwiftComp \texttt{.K} layout
(effective stiffness, effective compliance, orthotropic-approximated engineering
constants) with an \texttt{OpenSG} banner and a \texttt{Time taken} footer, plus
\texttt{<yaml>\_ABDG.out} for the wall laws.

\subsection{Validation: Deo \& Yu (2023)}

The reference is A. Deo and W. Yu, \emph{Equivalent 3D properties of thin-walled
composite structures using mechanics of structure genome}, \emph{Mechanics of
Advanced Materials and Structures} \textbf{30}(9), 2023, 1737--1748. Two columns
are quoted from that paper: \textbf{Deo TW}, its own thin-walled model with a
classical wall law, and the reference \textbf{SwiftComp} column, which is the
paper's MSG solid result and the equivalent of an OpenSG 2-D-solid run. All
values below are MPa, and \texttt{ours \%} and \texttt{Deo TW \%} are both
measured against the SwiftComp column. Every table is the literal content of the
\texttt{*\_compare.dat} file written by the driver.

Four sub-folders live under
\texttt{examples/OpenSG\_shell/3\_get\_solid\_props\_from\_shell\_cross\_section}.
Each is a two-command run: a generator that writes the 1-D shell YAML, then the
driver that homogenizes it.

\begin{lstlisting}[style=shellst]
cd examples/OpenSG_shell/3_get_solid_props_from_shell_cross_section/deo_cellular_solid
python make_deo_cellular_yamls.py
python deo_cellular_solid_props.py
\end{lstlisting}

The generator is where the geometry is edited. For the cellular solid it is a
five-line block at the top of \texttt{make\_deo\_cellular\_yamls.py}:
\texttt{l, h, t = 0.10, 0.10, 0.005}, aluminium \texttt{E, nu = 68.9e9, 0.33},
and \texttt{nseg = 10} elements per wall of length $l$. The driver then computes
the cell area itself from the wall angle,
\texttt{area = (2*l*np.cos(th))*(2*(h + l*np.sin(th)))}, and passes it as
\texttt{cell\_area}.

\subsubsection{Cellular solid, $\theta = +15^\circ$ (paper Table 1)}

\begin{table}[htbp]
\centering
\small
\begin{tabular}{lrrrrr}
\toprule
term & OpenSG-TW (RM) & Deo MSG-TW & SwiftComp & ours \% & Deo TW \% \\
\midrule
$C_{11}$ & 4733.84 & 4736.90 & 4678.90 & $+1.17$ & $+1.24$ \\
$C_{12}$ & 1086.01 & 1089.40 & 1105.50 & $-1.76$ & $-1.46$ \\
$C_{13}$ &  380.63 &  381.81 &  386.88 & $-1.61$ & $-1.31$ \\
$C_{22}$ & 2444.24 & 2446.39 & 2488.90 & $-1.79$ & $-1.71$ \\
$C_{23}$ &  846.69 &  847.44 &  860.89 & $-1.65$ & $-1.56$ \\
$C_{33}$ &  306.74 &  306.99 &  311.48 & $-1.52$ & $-1.44$ \\
$\mathbf{C_{44}}$ & \textbf{4.20} & 4.32 & 4.19 & $\mathbf{+0.18}$ & $+3.09$ \\
$C_{55}$ &  562.61 &  564.15 &  573.11 & $-1.83$ & $-1.56$ \\
$C_{66}$ &  993.78 &  997.52 & 1000.10 & $-0.63$ & $-0.26$ \\
\bottomrule
\end{tabular}
\caption{Deo cellular solid at $\theta = +15^\circ$, MPa, from
\texttt{deo\_cellular\_compare.dat}. Aluminium $E = 68.9$ GPa, $\nu = 0.33$;
$l = h = 10$ cm, $t = 5$ mm. Solve 1.5 s.}
\label{tab:deo-cell-p15}
\end{table}

\subsubsection{Re-entrant cellular solid, $\theta = -15^\circ$ (paper Table 2)}

The auxetic lattice, with the negative $C_{13}$ and $C_{23}$ couplings.

\begin{table}[htbp]
\centering
\small
\begin{tabular}{lrrrrr}
\toprule
term & OpenSG-TW (RM) & Deo MSG-TW & SwiftComp & ours \% & Deo TW \% \\
\midrule
$C_{11}$ & 7505.26 & 7507.40 & 7352.90 & $+2.07$ & $+2.10$ \\
$C_{12}$ & 1090.52 & 1094.10 & 1075.20 & $+1.42$ & $+1.76$ \\
$C_{13}$ & $-219.81$ & $-220.53$ & $-213.94$ & $+2.74$ & $+3.08$ \\
$C_{22}$ & 4151.29 & 4154.90 & 4080.00 & $+1.75$ & $+1.84$ \\
$C_{23}$ & $-846.69$ & $-847.45$ & $-821.77$ & $+3.03$ & $+3.12$ \\
$C_{33}$ &  180.61 &  180.75 &  173.48 & $+4.11$ & $+4.19$ \\
$\mathbf{C_{44}}$ & \textbf{2.47} & 2.68 & 2.47 & $\mathbf{+0.17}$ & $+8.38$ \\
$C_{55}$ &  331.26 &  332.17 &  346.01 & $-4.26$ & $-4.00$ \\
$C_{66}$ & 1687.82 & 1694.20 & 1706.40 & $-1.09$ & $-0.71$ \\
\bottomrule
\end{tabular}
\caption{Deo re-entrant cellular solid at $\theta = -15^\circ$, MPa, from
\texttt{deo\_cellular\_compare.dat}. Solve 0.2 s.}
\label{tab:deo-cell-m15}
\end{table}

$C_{44}$ is the in-plane shear channel, the one the paper flags as its largest
error and attributes to the classical wall model dropping the segment transverse
shear, with a first-order-shear wall model named as future work. That is exactly
the RM wall law used here: the $+3.09$\,\% and $+8.38$\,\% errors drop to
$+0.18$\,\% and $+0.17$\,\%. Every other term tracks the paper's own thin-walled
column to a few tenths of a percent, so the two independent MSG-TW
implementations agree where they share physics.

\subsubsection{Hierarchical square (paper Table 3)}

Run from \texttt{deo\_hierarchical\_square}. The generator
\texttt{make\_deo\_square\_yaml.py} carries \texttt{R, r, t = 0.10, 0.05, 0.005}
and \texttt{nseg = 10}, the same aluminium.

\begin{table}[htbp]
\centering
\small
\begin{tabular}{lrrrrr}
\toprule
term & OpenSG-TW (RM) & Deo MSG-TW & SwiftComp & ours \% & Deo TW \% \\
\midrule
$C_{11}$ & 5183.48 & 5186.39 & 5099.03 & $+1.66$ & $+1.71$ \\
$C_{12}$ &   24.22 &   24.83 &   26.75 & $-9.47$ & $-7.18$ \\
$C_{13}$ &   24.22 &   24.83 &   26.75 & $-9.47$ & $-7.18$ \\
$C_{22}$ &   46.23 &   47.13 &   50.46 & $-8.39$ & $-6.60$ \\
$C_{23}$ &   27.15 &   27.94 &   30.59 & $-11.23$ & $-8.66$ \\
$C_{33}$ &   46.23 &   47.13 &   50.45 & $-8.37$ & $-6.58$ \\
$C_{44}$ &    3.20 &    3.22 &    3.39 & $-5.60$ & $-5.01$ \\
$C_{55}$ &  647.56 &  650.00 &  670.92 & $-3.48$ & $-3.12$ \\
$C_{66}$ &  647.56 &  650.00 &  670.93 & $-3.48$ & $-3.12$ \\
\bottomrule
\end{tabular}
\caption{Deo hierarchical square, MPa, from \texttt{deo\_square\_compare.dat}.
$R = 10$ cm, $r = 5$ cm, $t = 5$ mm. Solve 1.8 s.}
\label{tab:deo-hier}
\end{table}

This lattice is \textbf{bending-dominated}: its transverse channels are
$t^3$-scale ($C_{22}\approx 50$ MPa against $C_{11}\approx 5100$ MPa), so the
load path is wall bending, not wall stretching. Two consequences are worth
recording. First, the RM result is mesh-converged and the offset is physics:
refining from 10 to 40 elements per segment moves $C_{23}$ by 1\,\%, and running
the same model with shear-rigid walls (the Kirchhoff limit) reproduces the
paper's classical column to 0.0--0.3\,\%. The gap between the two columns is
therefore the wall transverse-shear compliance, which is real at the $t/l = 0.1$
ligaments; element order is not involved. Second, both thin-walled models miss
the same junction stiffness. The residual against SwiftComp is the finite
rotational stiffness of the joint regions, which no midline model carries. The
classical model's smaller apparent error is a partial cancellation, artificial
shear rigidity offsetting missing joint stiffness, and the RM model removes the
first error, exposing the second.

\subsubsection{Composite square (paper Fig.~13 / Table 4)}

Run from \texttt{deo\_composite\_square}. The cell is $25.4\times25.4$ mm, so the
driver passes \texttt{cell\_area = L2*L2} with \texttt{L2 = 25.4e-3}; horizontal
segments sit at $y_3=\pm5.84$ mm with one vertical segment through the centre,
and every wall is a $[15]_8$ laminate of $t = 1.016$ mm with $E_1 = 141.96$,
$E_2 = E_3 = 9.79$, $G = 6.136$ GPa, $\nu = 0.42$.

\begin{table}[htbp]
\centering
\small
\begin{tabular}{lrrrrr}
\toprule
term & OpenSG-TW (RM) & Deo MSG-TW & SwiftComp & ours \% & Deo TW \% \\
\midrule
$C_{11}$ & 15263.83 & 16133.90 & 15518.11 & $-1.64$ & $+3.97$ \\
$C_{12}$ &   936.74 &  1016.82 &  1008.76 & $-7.14$ & $+0.80$ \\
$C_{13}$ &   468.37 &   468.38 &   459.01 & $+2.04$ & $+2.04$ \\
$C_{15}$ & $+1191.85$ & $-1192.81$ & $-1135.27$ & sign & $+5.07$ \\
$C_{16}$ &  2383.70 &  2589.74 &  2535.23 & $-5.98$ & $+2.15$ \\
$C_{22}$ &   906.11 &   983.56 &   988.88 & $-8.37$ & $-0.54$ \\
$C_{23}$ &     0.00 &     0.00 &    20.44 & $-100$ & $-100$ \\
$C_{26}$ &   292.25 &   317.09 &   311.80 & $-6.27$ & $+1.70$ \\
$C_{33}$ &   453.06 &   453.06 &   456.25 & $-0.70$ & $-0.70$ \\
$C_{35}$ & $+146.13$ & $-146.06$ & $-140.04$ & sign & $+4.30$ \\
$C_{44}$ &     0.94 &     1.08 &     1.16 & $-18.57$ & $-6.90$ \\
$C_{55}$ &   547.31 &   547.32 &   545.43 & $+0.35$ & $+0.35$ \\
$C_{66}$ &  1094.63 &  1188.20 &  1188.30 & $-7.88$ & $-0.01$ \\
\bottomrule
\end{tabular}
\caption{Deo composite square, MPa, from
\texttt{deo\_composite\_square\_compare.dat}. Solve 1.8 s. The \texttt{.dat} file
prints the two sign-flipped rows as $-204.98$ and $-204.35$ percent, which is the
arithmetic of a sign reversal, not a magnitude error.}
\label{tab:deo-comp}
\end{table}

$C_{13}$, $C_{33}$ and $C_{55}$ reproduce the paper's thin-walled column exactly,
and $C_{11}$ improves on it against SwiftComp. Two open items are recorded
honestly. $C_{15}$ and $C_{35}$ match the published magnitudes to 0.1\,\% with
the \emph{opposite sign}, a ply-angle handedness convention still to be
reconciled, and the likely cause of the $-6$ to $-8$\,\% group. And $C_{23}$ is
zero in both thin-walled models, for the reason below.

\subsubsection{A fifth folder: the isotropic square tube}

\texttt{square\_section\_isotropic} is the simplest case in the set and the best
place to start. It reads \texttt{square\_tube\_1Dshell.yaml} and its user-input
block is four lines: \texttt{a, t\_w = 1.0, 0.03} and
\texttt{CELL\_AREA = 4*a*t\_w}, which is the \emph{wall material area}, the same
measure the solid engines use, so the two routes are directly comparable. Run it
with \texttt{python solid\_props\_msg\_shell.py}; it writes
\texttt{square\_solid\_msg\_shell.out} with the $6\times6$ and the nine
engineering constants obtained from \texttt{elastic\_constants(C)}.

\subsection{Why $C_{23}$ vanishes on axis-aligned lattices}

For a wall at angle $\varphi$ the only $\Gamma_e$ entries carrying weight in
\emph{both} the $\bar\Gamma_{22}$ and $\bar\Gamma_{33}$ columns are the membrane
$\varepsilon_{22}$ row ($\cos^2\varphi$, $\sin^2\varphi$) and the
transverse-shear $2\gamma_{23}$ row ($-\sin\varphi\cos\varphi$,
$+\sin\varphi\cos\varphi$), giving a per-length integrand
\begin{equation}
(A_{11}-G_{\mathrm{msg}})\,\sin^2\varphi\cos^2\varphi .
\label{eq:shell-c23}
\end{equation}
On a $0/90$ lattice every wall has $\sin^2\varphi\cos^2\varphi = 0$, so this
$O(t)$ membrane path vanishes identically and the whole of $C_{23}$ is the
$O(t^2)$ \textbf{junction} contribution: the wall-overlap blocks, which have zero
measure on a midline mesh. Inclined-wall cells are unaffected, which is why the
cellular cases above carry $C_{23}$ correctly and only the square lattices expose
the junction term. Optional junction corrections (\texttt{junction="census"},
\texttt{"micro"}, \texttt{"microcell"}) recover it for stretch-dominated cells;
they do \emph{not} apply to bending-dominated lattices such as the hierarchical
square, and they are off by default.

%=====================================================================
\section{Equivalent 3-D solid from a 3-D shell SG}
\label{sec:shell-sg3d}

\subsection{What this capability computes}

A structure gene periodic in \textbf{all three} directions -- a triply-periodic
minimal surface (TPMS) unit cell, a lattice block, any 3-D microstructure whose
material is a thin sheet -- homogenizes to the same equivalent 3-D solid
$6\times6$ of Eq.~\eqref{eq:shell-gbar}. \texttt{shell\_sg3d} solves it when the
microstructure is a thin sheet meshed as a \emph{surface}: the wall carries the
RM $8\times8$ law, so thickness is a parameter rather than a remesh.

It reuses the cross-section $\Gamma_e$ / $\Gamma_h$ operators unchanged (they are
geometry-general) and adds only the 3-D environment: sparse assembly, the
three-direction periodic map, drilling by penalty on the element-constant
residual, and a three-translation kernel. Shell edges shared by more than two
elements are detected as \textbf{junction lines} and reported as
\texttt{n\_junction\_edges}; a smooth TPMS has none.

The example is
\texttt{examples/OpenSG\_shell/4\_get\_solid\_props\_from\_shell\_3D\_SG}.

\begin{figure}[htbp]
\centering
\includegraphics[width=0.62\textwidth]{\FIGDIR/schwarz_p_mesh.png}
\caption{The Schwarz-P TPMS unit cell as a 3-D shell SG: 13\,536 nodes and
26\,360 triangular shell elements on the midsurface, written by
\texttt{schwarz\_solid\_props.py} as \texttt{schwarz\_p\_mesh.png}. The cell is
the unit cube, and the wall thickness is a number in the YAML, not geometry.}
\label{fig:schwarz-mesh}
\end{figure}

\subsection{The input the user edits}

The mesh comes from a gmsh 2.2 triangle file, and
\texttt{make\_schwarz\_yaml.py} converts it. Its user block is four lines:

\begin{lstlisting}[style=opensg,caption={The user block of make\_schwarz\_yaml.py: mesh path, thickness, material.}]
MSH = ("../../../tests/08072026_square_shell_mesh/TPMS_SC_files/"
       "schwarz_p_D2_shell.msh")
T = 0.036457           # Sample_2, volume-consistent
E, nu = 69.0e9, 0.30
G = E/(2*(1+nu))
\end{lstlisting}

The converter computes each element's frame from the triangle geometry
($e_3$ = unit normal, $e_2$ along the first edge, $e_1 = e_2\times e_3$), writes
one \texttt{sections} entry with the single ply \texttt{[alu, T, 0.0]}, and
reports the surface area and the resulting relative density
\texttt{area.sum()*T}.

\subsection{Run it}

\begin{lstlisting}[style=shellst]
cd examples/OpenSG_shell/4_get_solid_props_from_shell_3D_SG
python make_schwarz_yaml.py
python schwarz_solid_props.py
python compare_boundary_modes.py
\end{lstlisting}

\begin{lstlisting}[style=opensg,caption={The two boundary treatments of shell\_sg3d.}]
from opensg_shell.sg_homo import shell_sg3d

r = shell_sg3d("schwarz_p_3Dshell.yaml")                        # boundary="periodic" (default)
r = shell_sg3d("schwarz_p_3Dshell.yaml", boundary="aperiodic")  # non-periodic SGs (explicit)
r["C3D"], r["n_junction_edges"], r["ndof"], r["solve_time"]
\end{lstlisting}

The SG measure $\omega$ defaults to the \textbf{midsurface area}, the 3-D
analogue of $\omega = \text{perimeter}$ for a plane-section shell SG, and can be
overridden with the \texttt{omega=} argument. The \texttt{.out} file is written
per unit-cell volume so its moduli compare directly with a solid \texttt{.K}.

\subsection{What comes out}

\texttt{shell\_sg3d} writes \texttt{schwarz\_p\_3Dshell\_C3D.out}, and the driver
writes \texttt{schwarz\_p\_mesh.png}. The banner line of the \texttt{.out} records
the whole configuration, and the default periodic run at $t = 0.036457$ gives:

\begin{Verbatim}[fontsize=\footnotesize,frame=single,framesep=4pt]
 OpenSG msg-shell equivalent 3D solid (3-D shell SG, 13536 nodes, 26360 elems,
 0 junction edges, periodic in 3 dirs, per unit cell 1)

 The Effective Stiffness Matrix
 --------------------------------------------
     2.2253086E+009     1.5648739E+009     1.5649503E+009  ...
     1.5648739E+009     2.2252770E+009     1.5647791E+009  ...
     1.5649503E+009     1.5647791E+009     2.2250347E+009  ...

 The Engineering Constants (Approximated as Orthotropic)
 ----------------------------------------------------------
  E1  =     9.3296450E+008
  E2  =     9.3313789E+008
  E3  =     9.3284790E+008
  G12 =     1.0409258E+009
  G13 =     1.0410061E+009
  G23 =     1.0409794E+009
  nu12=     4.1277916E-001
  nu13=     4.1304617E-001
  nu23=     4.1288357E-001

 Time taken: 66.43 sec
\end{Verbatim}

Read the numbers as a health check first. The three normal terms agree to five
digits ($2.2253$, $2.2253$, $2.2250$ GPa) and the three shear terms likewise
($1.04098$, $1.04101$, $1.04093$ GPa), so the effective law is cubic to five
digits, as the Schwarz-P symmetry demands. Every symmetry-forbidden coupling sits
at $10^{3}$--$10^{5}$ Pa against diagonals of $10^{9}$ Pa, four to six orders
below. That is the health check on the all-directions periodic tie: if the tie
map were wrong, those entries would not be small. The repository also records a
head-to-head of this same surface against the msg-solid route on the same
structure, where the shell result runs about 3\,\% below on the normal terms and
about 6\,\% below on the shears, with cubic symmetry intact in both; because the
solid column is digit-identical to SwiftComp, those percentages measure the
thin-shell reduction itself at a surface reaching $t/R \approx 0.2$ near the
saddle necks.

\subsection{Boundary treatment: periodic default versus aperiodic}

\texttt{shell\_sg3d} takes a \texttt{boundary} argument, and the default is
\texttt{"periodic"}: opposite faces, edges and corners of the bounding box are
tied through the sparse assembly map, and the three rigid translations are
removed by area-weighted Lagrange rows. This is the exact treatment for a
periodic medium, and it is what to report.

\texttt{boundary="aperiodic"} is the boundary-solution treatment. For a unit cell
the boundary solution is the macro (affine) field itself, so mapping it onto the
boundary nodes prescribes zero \textbf{translational} warping fluctuation
($w_1=w_2=w_3=0$, Dirichlet) on every node of the bounding-box faces, with the
rotational DOFs left natural. That forbids the affine fluctuations which make a
\emph{free} aperiodic SG rank-one, fixes the rigid modes without a Lagrange
border, and is a kinematic upper bound. Clamping the rotations as well would
roughly double the offset.

\texttt{compare\_boundary\_modes.py} runs both on the current YAML, copies each
per-mode \texttt{.out} to \texttt{schwarz\_t<T>\_<mode>.out}, and appends a block
to \texttt{boundary\_compare.dat}. Table~\ref{tab:schwarz-boundary} is that file.

\begin{table}[htbp]
\centering
\small
\begin{tabular}{lrrr@{\hskip 2em}rrr}
\toprule
 & \multicolumn{3}{c}{$t = 0.036457$} & \multicolumn{3}{c}{$t = 0.129330$} \\
\cmidrule(lr){2-4}\cmidrule(lr){5-7}
term & aperiodic & periodic & \%diff & aperiodic & periodic & \%diff \\
\midrule
$C_{11}$ & 2375.89 & 2225.31 & $+6.767$ & 9881.00 & 8916.81 & $+10.813$ \\
$C_{22}$ & 2375.92 & 2225.28 & $+6.770$ & 9881.00 & 8916.85 & $+10.813$ \\
$C_{33}$ & 2376.22 & 2225.03 & $+6.795$ & 9873.34 & 8916.17 & $+10.735$ \\
$C_{12}$ & 1537.47 & 1564.87 & $-1.751$ & 4921.09 & 5352.21 & $-8.055$ \\
$C_{13}$ & 1539.24 & 1564.95 & $-1.643$ & 4931.10 & 5352.71 & $-7.877$ \\
$C_{23}$ & 1538.99 & 1564.78 & $-1.648$ & 4931.00 & 5352.55 & $-7.876$ \\
$C_{44}$ & 1078.28 & 1040.98 & $+3.583$ & 4150.25 & 3961.63 & $+4.761$ \\
$C_{55}$ & 1078.32 & 1041.01 & $+3.585$ & 4150.26 & 3961.63 & $+4.761$ \\
$C_{66}$ & 1078.61 & 1040.93 & $+3.620$ & 4151.27 & 3961.51 & $+4.790$ \\
\midrule
ndof     & \multicolumn{1}{c}{81216} & \multicolumn{1}{c}{79056} & &
           \multicolumn{1}{c}{81216} & \multicolumn{1}{c}{79056} & \\
solve    & \multicolumn{1}{c}{42.3 s} & \multicolumn{1}{c}{66.5 s} & &
           \multicolumn{1}{c}{42.2 s} & \multicolumn{1}{c}{65.4 s} & \\
\bottomrule
\end{tabular}
\caption{Aperiodic versus periodic on the same Schwarz-P YAML at two wall
thicknesses, per unit cell, MPa. The literal content of
\texttt{boundary\_compare.dat}. In both blocks 720 free-face nodes of 13\,536 are
clamped by the aperiodic treatment.}
\label{tab:schwarz-boundary}
\end{table}

The offset is the clamped-face boundary layer of a \textbf{single} cell: the
expected kinematic bound, not an error in either path. It grows with wall
thickness, from $+6.8$\,\% on $C_{11}$ at $t = 0.036457$ to $+10.8$\,\% at
$t = 0.129330$, because a thicker wall makes the clamped skin a larger fraction of
the cell. Aperiodic is about 35\,\% \emph{faster} here, because the Dirichlet
rows leave the factorization and there is no tie map or Lagrange border, despite
the slightly larger nominal DOF count. Report periodic for a periodic medium;
request \texttt{"aperiodic"} only for SGs that genuinely are not periodic.

%=====================================================================
\section{Tapered and aperiodic segments}
\label{sec:shell-segment}

\subsection{What this capability computes}

Everything so far has assumed the beam is prismatic: the classical msg-shell
analysis meshes only the cross-section contour and enforces periodicity along the
beam axis. A \textbf{tapered} segment has no such periodicity, because its two end
cross-sections differ. \texttt{opensg\_shell.sg\_homo.segment\_timo\_from\_3dyaml}
handles this by replacing the axial periodic condition with Dirichlet data taken
from the two ends. The input is no longer a 1-D ring but a genuine 3-D surface
mesh of a slice of the beam.

The pipeline has four stages:

\begin{enumerate}
\item \textbf{Boundary extraction.} The two end cross-sections are found
  \emph{topologically}: a mesh edge used by exactly one quad is a free edge, and
  each connected component of the free-edge graph with at least three nodes is one
  end ring. The axis is the direction between the two most widely separated
  components, and left/right is decided by the axial coordinate. Each ring is
  written out as a standalone 1-D YAML together with the ring-node $\to$
  segment-node map.
\item \textbf{Boundary solves.} Each ring is homogenized on its own with
  \texttt{ring\_indep}, giving its Timoshenko $6\times6$ and the warping fields
  $V_0$ (Euler--Bernoulli) and $V_1$ (Timoshenko), all six DOFs per node
  including the drilling rotation $\omega_3$.
\item \textbf{Mapping.} The ring $V_0$/$V_1$ are transferred onto the segment's
  boundary nodes as Dirichlet data, replacing the periodic boundary condition and
  fixing the segment's rigid-body modes.
\item \textbf{Segment solve.} One LU factorization of the constrained system
  serves both orders,
  \begin{equation}
  \mathbf{D}_{hh}\,V_0 = -\mathbf{D}_{he}, \qquad
  \mathbf{D}_{hh}\,V_1 = b(V_0),
  \label{eq:seg-orders}
  \end{equation}
  with $V_0 = V_0^{\text{ring}}$ and $V_1 = V_1^{\text{ring}}$ on both ends. The
  generalized-Timoshenko finalization divides the energy by the segment length
  $L$, which is measured as the extent of the nodes along the axis.
\end{enumerate}

The example is
\texttt{examples/OpenSG\_shell/5\_taper\_segment\_aperiodic\_boundaries}.

\subsection{Seeing the capability}

\begin{figure}[htbp]
\centering
\includegraphics[width=0.88\textwidth]{\FIGDIR/aperiodic_segment_bar_urc.png}
\caption{The BAR-URC blade segment the example runs, coloured by its per-element
material frames. This is the surface on which the segment problem is solved: a
3-D shell mesh, not a cross-section.}
\label{fig:seg-bar-urc}
\end{figure}

\begin{figure}[htbp]
\centering
\begin{minipage}{0.48\textwidth}
\centering
\includegraphics[width=\textwidth]{\FIGDIR/aperiodic_ring_L.png}
\end{minipage}\hfill
\begin{minipage}{0.48\textwidth}
\centering
\includegraphics[width=\textwidth]{\FIGDIR/aperiodic_ring_R.png}
\end{minipage}
\caption{The two end cross-sections extracted automatically from the free-edge
graph of the segment mesh, left ring and right ring, each written as a standalone
1-D SG and solved on its own. On a tapered segment the two rings differ, and
their two $6\times6$ matrices bracket the segment result.}
\label{fig:seg-rings}
\end{figure}

\begin{figure}[htbp]
\centering
\includegraphics[width=0.72\textwidth]{\FIGDIR/aperiodic_prismatic_segment.png}
\caption{The prismatic circular cylinder used for the acceptance test. The same
free-edge extraction runs on it, and here both ends are identical by
construction, so the segment $6\times6$ must reproduce the ring $6\times6$.}
\label{fig:seg-prismatic}
\end{figure}

\subsection{The input the user edits}

For the validation case, \texttt{make\_cylinder\_segment.py} generates the
surface mesh. Its parameter block is the whole user interface:

\begin{lstlisting}[style=opensg,caption={The parameter block of make\_cylinder\_segment.py.}]
R  = 1.0            # wall mid-surface radius [m]
NC = 160            # circumferential quads
NL = 3              # axial quads over the segment (>=2 so an interior ring exists)
L  = 1.5            # segment length [m]
HR = [0.05, 0.1, 0.2]   # wall/radius ratios: thin, moderate, thick

ISO = {"name": "iso", "density": 1800.0,
       "E": [70e9, 70e9, 70e9], "G": [26.923e9]*3, "nu": [0.3, 0.3, 0.3]}
ANI = {"name": "ani", "density": 1800.0,
       "E": [37e9, 9e9, 9e9], "G": [4e9, 4e9, 4e9], "nu": [0.3, 0.3, 0.3]}
\end{lstlisting}

It writes six files into a \texttt{meshes/} sub-folder, named
\texttt{seg\_<tag>\_hR<ratio>.yaml}: an isotropic single-ply tube and a balanced
$[+45/-45]$ tube at each of the three thickness ratios. The beam axis is $x$, the
cross-section plane is $(y,z)$, the reference surface is the wall mid-surface,
and each element stores the nine-component triad $[e_1\,e_2\,e_3]$ with $e_1$
axial, $e_2$ the hoop tangent and $e_3$ the inward normal.

\subsection{Run it}

\begin{lstlisting}[style=shellst]
cd examples/OpenSG_shell/5_taper_segment_aperiodic_boundaries
python make_cylinder_segment.py
python run_taper_segment.py
python run_bar_urc_segment.py
\end{lstlisting}

The whole pipeline is one call on the 3-D shell YAML. There are no boundary files
to prepare, because the ends are found from the mesh topology. Note the path: the
cylinder meshes live in the \texttt{meshes/} sub-folder, so the argument carries
that prefix.

\begin{lstlisting}[style=opensg,caption={The segment call, with the meshes/ path used by run\_taper\_segment.py.}]
from opensg_shell.sg_homo import segment_timo_from_3dyaml

r = segment_timo_from_3dyaml("meshes/seg_iso_hR0.1.yaml")
r["S6"]             # segment Timoshenko 6x6
r["C6L"], r["C6R"]  # the two boundary-ring 6x6s
r["L"]              # segment length used to normalize the energy
r["solve_time"]     # seconds
\end{lstlisting}

\subsection{What comes out}

Each run writes, alongside the input YAML, the two boundary 1-D YAMLs
(\texttt{boundary\_<tag>\_L.yaml} and \texttt{\_R.yaml}), the
\texttt{<tag>\_boundaries.npz} bundle, three orientation PNGs
(\texttt{orient\_<tag>\_segment.png}, \texttt{orient\_<tag>\_ring\_L.png},
\texttt{orient\_<tag>\_ring\_R.png}), and the timed \texttt{<yaml>\_Timo.out} in
the SwiftComp layout. For the isotropic cylinder that file opens with

\begin{Verbatim}[fontsize=\small,frame=single,framesep=4pt]
 OpenSG msg-shell tapered/aperiodic segment (boundary V0/V1 Dirichlet, L=1.5)
\end{Verbatim}

\noindent
and closes with \texttt{Time taken: 2.51 sec}.

One caveat on the figures. The extraction is called with \texttt{plot="auto"},
which means an orientation PNG is drawn \emph{only if that file is not already on
disk}. If you regenerate or edit a mesh and rerun, delete the three
\texttt{orient\_*.png} files for that tag first, or you will be looking at the
frames of the previous mesh.

\subsection{Acceptance test: the prismatic identity}

For a \textbf{prismatic} segment the taper terms vanish, so the segment
$6\times6$ must equal the boundary ring's own $6\times6$. That identity is the
acceptance test of the whole boundary extraction plus $V_0$/$V_1$ mapping, and
\texttt{run\_taper\_segment.py} measures it on the prismatic circular tube
($R=1$, $t/R=0.1$, $L=1.5$) in both the isotropic and the $\pm45$ layup, writing
\texttt{seg\_ring\_identity.dat}.

\begin{table}[htbp]
\centering
\small
\begin{tabular}{lrrrl}
\toprule
term & segment $S_6$ & boundary ring $C_6$ & \% error & case \\
\midrule
$EA$   & $4.397947\times10^{10}$ & $4.397947\times10^{10}$ & $0.0000$ & \\
$GA_2$ & $8.473281\times10^{9}$  & $8.473258\times10^{9}$  & $0.0003$ & \texttt{iso\_hR0.1} \\
$GA_3$ & $8.473281\times10^{9}$  & $8.473258\times10^{9}$  & $0.0003$ & $L=1.5$ \\
$GJ$   & $1.694078\times10^{10}$ & $1.694078\times10^{10}$ & $0.0000$ & solve 2.5 s \\
$EI_2$ & $2.200273\times10^{10}$ & $2.200269\times10^{10}$ & $0.0002$ & \\
$EI_3$ & $2.200273\times10^{10}$ & $2.200269\times10^{10}$ & $0.0002$ & \\
\midrule
$EA$   & $7.706160\times10^{9}$  & $7.706160\times10^{9}$  & $0.0000$ & \\
$GA_2$ & $3.262832\times10^{9}$  & $3.262807\times10^{9}$  & $0.0007$ & \texttt{aniso\_hR0.1} \\
$GA_3$ & $3.262832\times10^{9}$  & $3.262807\times10^{9}$  & $0.0007$ & $L=1.5$ \\
$GJ$   & $6.526891\times10^{9}$  & $6.526891\times10^{9}$  & $0.0000$ & solve 1.5 s \\
$EI_2$ & $3.855225\times10^{9}$  & $3.855204\times10^{9}$  & $0.0006$ & \\
$EI_3$ & $3.855225\times10^{9}$  & $3.855204\times10^{9}$  & $0.0006$ & \\
\bottomrule
\end{tabular}
\caption{The prismatic identity, diagonal terms, from
\texttt{seg\_ring\_identity.dat}. The worst diagonal error is $0.0003$\,\% (on
$GA$) for the isotropic tube and $0.0007$\,\% for the $\pm45$ layup. The worst
error anywhere in the $6\times6$ is $0.0052$\,\% and it falls on the
$GA_2$--$EI_2$ shear--bending coupling of the anisotropic case, whose ring value
is $-3.684772\times10^{7}$ against a segment value of $-3.684965\times10^{7}$.}
\label{tab:seg-identity}
\end{table}

The \texttt{\% error} matrix in the file reports insignificant ring terms as
zero: entries whose magnitude is below $10^{-6}$ of the largest are divided by
infinity rather than by numerical noise, which is why the file shows exact zeros
where the ring value is $10^{-5}$ against diagonals of $10^{10}$.

\subsection{Tapered demonstration: BAR-URC segment 12}

\texttt{run\_bar\_urc\_segment.py} runs the same pipeline on a real tapered
segment of the 100 m BAR-URC blade,
\texttt{BAR\_URC\_numEl\_52\_segment\_12.yaml}. Here no identity holds. Each
boundary ring has its own $6\times6$, and the segment result is the true
length-averaged Timoshenko stiffness of the tapered piece. The physical check is
that every segment diagonal falls \emph{between} its two ring values, which the
driver reports as the ratio $(S-R)/(L-R)$; a value in $[0,1]$ means the segment
lies between its two ends.

\begin{table}[htbp]
\centering
\small
\begin{tabular}{lrrrr}
\toprule
term & segment $S$ & ring L & ring R & ratio \\
\midrule
$EA$   & $1.487716\times10^{10}$ & $1.498819\times10^{10}$ & $1.483824\times10^{10}$ & 0.260 \\
$GA_2$ & $6.098229\times10^{8}$  & $6.248039\times10^{8}$  & $6.030954\times10^{8}$  & 0.310 \\
$GA_3$ & $1.398945\times10^{8}$  & $1.457255\times10^{8}$  & $1.345727\times10^{8}$  & 0.477 \\
$GJ$   & $3.353538\times10^{8}$  & $3.684346\times10^{8}$  & $3.057355\times10^{8}$  & 0.472 \\
$EI_2$ & $2.996686\times10^{9}$  & $3.244829\times10^{9}$  & $2.765385\times10^{9}$  & 0.482 \\
$EI_3$ & $8.971707\times10^{9}$  & $9.620926\times10^{9}$  & $8.397279\times10^{9}$  & 0.469 \\
\bottomrule
\end{tabular}
\caption{BAR-URC tapered shell segment 12, from \texttt{bar\_urc\_segment12.dat}.
Segment length $L = 3.45293$ m, solve 15.9 s. All six ratios lie in $[0,1]$, so
every segment diagonal is bracketed by its two boundary rings.}
\label{tab:seg-bar-urc}
\end{table}

Read the ratios as a taper diagnostic. The bending and torsion terms sit near
0.47--0.48, close to the arithmetic mean of the two ends, which is what a nearly
linear taper in those channels produces. $EA$ and $GA_2$ sit at 0.26 and 0.31,
noticeably below the mean, which says the axial and chordwise-shear stiffness of
this segment is closer to its thin end than a linear interpolation of the two
ring values would suggest. A station-by-station beam table built by interpolating
ring properties would therefore over-predict those two terms on this segment.

\subsection{Shear-scheme rule for segments}

\texttt{segment\_timo\_from\_3dyaml} defaults to \texttt{shear="full"}
($2\times2$ Gauss), and that is the production setting to keep: MITC tying
aliases the algebraic drilling content on flat-walled and webbed sections. The
tied variants exposed by the element (\texttt{"mitc4\_wonly"}, which ties both
rows at the Dvorkin--Bathe points but keeps every rotation column at its
full-integration value; \texttt{"mitc4\_g23"}, which ties only the $\gamma_{23}$
row; and \texttt{"mitc4\_both"}, naive full tying) are ablation options for
studying the element, not recommended production settings for segments. Note the
contrast with the ring of Section~\ref{sec:shell-beam}, whose production default
\emph{is} \texttt{"mitc4\_g23"}: the ring is a one-quad-deep prismatic strip and
the segment is a genuine surface mesh, and the two geometries do not want the
same treatment.

\chapter{Tutorials: the Solid Engine (msg-solid)}
\label{ch:solid}

The \texttt{msg-solid} engine reduces a \emph{solid} structure gene (SG) --- a
mesh of ordinary continuum elements in one, two or three dimensions --- to a
macroscopic constitutive law, and then runs the same construction backwards to
recover the pointwise three-dimensional strain and stress inside that gene.
Everything in this chapter is reached through two entry points.

\begin{table}[htbp]
\centering
\footnotesize
\begin{tabular}{llp{4.8cm}}
\toprule
entry point & module & structure gene it reads \\
\midrule
\texttt{rm\_plate\_msg} & \texttt{opensg\_solid.rm\_plate\_1D.msg\_rm\_plate}
  & a 1-D line mesh through a wall thickness \\
\texttt{plate\_homo\_2d} & \texttt{opensg\_solid.sg\_homo}
  & a 1-D, 2-D or 3-D solid SG \\
\bottomrule
\end{tabular}
\caption{The two homogenization entry points of \texttt{msg-solid}.}
\label{tab:solid-entry}
\end{table}

The first is the dedicated through-thickness (laminate) route: the wall is a
stack of plies, nothing varies in-plane, so the gene is one-dimensional and the
result is a plate law. The second is the general route. Its \texttt{n\_model}
argument selects which macro model the same SG is reduced to.

\begin{table}[htbp]
\centering
\footnotesize
\begin{tabular}{clll}
\toprule
\texttt{n\_model} & macro model & row and column order & title infix \\
\midrule
1 & beam & \texttt{[eps11 gam12 gam13 k1 k2 k3]} & \texttt{Timoshenko Beam} \\
2 & plate & \texttt{[N11 N22 N12 M11 M22 M12]} & \texttt{Classical Plate} \\
3 & 3-D elastic & Voigt \texttt{[11 22 33 23 13 12]} & \texttt{Cauchy Continuum} \\
\bottomrule
\end{tabular}
\caption{\texttt{n\_model} on \texttt{plate\_homo\_2d}. The infix goes into
\texttt{The Effective <infix> Stiffness Matrix} and is always the macro law's
console title, so the file and the terminal read the same. Two of the five
titles belong to the refinement switch rather than to \texttt{n\_model}: the
classical beam $4\times4$ is \texttt{Euler-Bernoulli Beam}, and the plate route
becomes \texttt{Reissner-Mindlin} when the shear-refined $8\times8$ is
produced.}
\label{tab:solid-nmodel}
\end{table}

Every \texttt{plate\_homo\_2d} call writes a timed SwiftComp-\texttt{.K}-layout
report --- effective stiffness, effective compliance, and, for the 3-D solid law
only, an engineering-constants block --- opening with an \texttt{ OpenSG} banner
and closing with a \texttt{ Time taken: } line in seconds. Unless you pass
\texttt{plot=False} it also writes a \texttt{<base>\_mesh.png} with the elements
coloured by material. One rule governs the names of both: inside
\texttt{plate\_homo\_2d} the output base is
\texttt{base = os.path.splitext(sc\_path)[0]}, so the report is
\emph{always} \texttt{<input basename>.out}, regardless of what the script calls
the case.

The eight runnable examples live under \texttt{examples/OpenSG-solid}. Each
runner is meant to be executed from inside its own folder, because the SG file
is named relative to the working directory.

\begin{table}[htbp]
\centering
\footnotesize
\begin{tabular}{lp{3.2cm}p{5.4cm}}
\toprule
folder & script & what it does \\
\midrule
\texttt{1\_get\_plate\_props\_from\_1DSG} & \texttt{1d\_sg.py} & layup $\rightarrow$ 1-D SG $\rightarrow$ plate law \\
\texttt{2\_get\_plate\_stress\_from\_1DSG} & \texttt{rm\_dehom.py}, \texttt{rm\_dehom\_ff.py} & plate resultants $\rightarrow$ 3-D ply stress \\
\texttt{3\_get\_plate\_props\_from\_2D\_SG} & \texttt{plate\_homo\_2dsg.py} & 2-D cell SG $\rightarrow$ plate ABD \\
\texttt{4\_get\_plate\_dehom\_from\_2DSG} & \texttt{plate\_dehom\_2dsg.py} & macro plate strain $\rightarrow$ fields in the cell \\
\texttt{5\_get\_solid\_props\_from\_SG} & \texttt{solid\_homo\_2dsg.py} & 2-D cell SG $\rightarrow$ equivalent 3-D solid \\
\texttt{6\_get\_solid\_props\_from\_3D\_SG} & \texttt{tpms\_solid\_props.py} and others & 3-D SG $\rightarrow$ equivalent 3-D solid \\
\texttt{7\_get\_beam\_props\_from\_SG} & \texttt{beam\_homo\_sg.py} & 2-D cross-section SG $\rightarrow$ Timoshenko $6\times6$ \\
\texttt{8\_get\_beam\_dehom\_from\_SG} & \texttt{beam\_dehom\_sg.py} & macro beam strain $\rightarrow$ fields in the section \\
\bottomrule
\end{tabular}
\caption{The eight \texttt{msg-solid} example folders.}
\label{tab:solid-examples}
\end{table}

A note on what you will find in those folders before you run anything. Some of
the files named as outputs below are regenerated on every run and are simply not
checked in; where that is the case it is said explicitly. Nothing in this
chapter depends on a pre-computed file being present.

\section{Plate properties from a through-thickness 1-D SG}
\label{sec:solid-plate1d}

\subsection{What it computes}

The structure gene is the wall discretised through its own thickness: one
element per ply, quartic by default, with each node carrying the three warping
(fluctuation) degrees of freedom. Minimising the 3-D strain energy over that
line mesh order by order in $h/\ell$ gives the warping ladder --- the
zeroth-order columns $V_0$ and the first-gradient columns $V_{11}$, $V_{12}$ ---
from which the classical plate stiffness follows as
%
\begin{equation}
\mathbf{A}_6 = \big\langle (\mathbf{\Gamma}_\varepsilon
 + \mathbf{\Gamma}_h \mathbf{V}_0)^{\mathsf T}\,
 \mathbf{C}(z)\,
 (\mathbf{\Gamma}_\varepsilon + \mathbf{\Gamma}_h \mathbf{V}_0) \big\rangle ,
\label{eq:solid-A6}
\end{equation}
%
where $\langle\cdot\rangle$ is the through-thickness integral and
$\mathbf{C}(z)$ the ply stiffness rotated by its fibre angle. With
\texttt{model: 1} the run additionally identifies the $2\times2$ MSG
transverse-shear block $\mathbf{G}$ from the second-order energy and assembles
the $8\times8$ law
%
\begin{equation}
\mathbf{ABDG} =
\begin{bmatrix} \mathbf{A}_6 & \mathbf{0} \\ \mathbf{0} & \mathbf{G} \end{bmatrix},
\label{eq:solid-abdg}
\end{equation}
%
in which the transverse shears are uncoupled from the membrane and bending
measures. Row and column order is
\texttt{[e11 e22 g12 k11 k22 k12 2g13 2g23]}, and that string is printed in the
banner of every \texttt{.out} so a file can always be read without the manual.

The whole capability is one file you edit and five functions that run. The SG
mesh is written to disk and read back before the solve, so the mesh file --- not
the database --- is what the homogenization actually sees.

\captree{
  \node[inputnode] (y) {\texttt{layup\_db.yaml}\\
    \texttt{model}, \texttt{fraction}, \texttt{mesh},\\
    \texttt{materials}, \texttt{layup}};
  \node[funcnode, below=of y] (f1) {\texttt{segment\_plate.plate\_sg\_yaml}\\
    the through-thickness node grid,\\
    \texttt{n\_per\_layer} elements per ply};
  \node[outnode, right=of f1] (o1) {\texttt{1dsg.yaml}\\ \texttt{1dsg.png}};
  \node[funcnode, below=of f1] (f2)
    {\texttt{segment\_plate.read\_plate\_sg\_yaml}\\
     the mesh back as thick, angles,\\ mat\_names, material\_db};
  \node[funcnode, below=of f2] (f3) {\texttt{msg\_rm\_plate.rm\_plate\_msg}\\
    warping ladder \texttt{V0}, \texttt{V11}, \texttt{V12}};
  \node[datanode, below=of f3] (d1) {\texttt{r["A6"]} $6\times6$ ABD\\
    \texttt{r["ABDG"]} $8\times8$ at \texttt{model: 1}};
  \node[funcnode, below=of d1] (f4) {\texttt{sg\_homo.write\_sc\_K}\\
    stiffness, compliance, banner, timing};
  \node[outnode, right=of f4] (o2) {\texttt{layup\_db\_plate\_homo.out}};
  \draw[capedge] (y) -- (f1);
  \draw[capedge] (f1) -- (o1);
  \draw[capedge] (o1) -- (f2);
  \draw[capedge] (f2) -- (f3);
  \draw[capedge] (f3) -- (d1);
  \draw[capedge] (d1) -- (f4);
  \draw[capedge] (f4) -- (o2);
}

\subsection{The input you edit}

Everything lives in \texttt{layup\_db.yaml}. The script reads that file and
nothing else, so a new laminate means editing the YAML, never the code.

\begin{table}[htbp]
\centering
\footnotesize
\begin{tabular}{lp{7.2cm}p{2.6cm}}
\toprule
key & what it sets & value in the shipped file \\
\midrule
\texttt{model} & \texttt{0} = classical $6\times6$ ABD, \texttt{1} = $8\times8$ ABDG & \texttt{1} \\
\texttt{fraction} & reference surface as a fraction of thickness & \texttt{0.5} \\
\texttt{mesh: n\_per\_layer} & SG elements per physical ply & \texttt{1} \\
\texttt{mesh: elem\_order} & element order; \texttt{4} = 5-noded quartic & \texttt{4} \\
\texttt{materials} & per material \texttt{E: [E1,E2,E3]}, \texttt{G: [G12,G13,G23]}, \texttt{nu: [nu12,nu13,nu23]}, \texttt{rho} & \texttt{ge}, \texttt{herex} \\
\texttt{layup} & one entry per ply, bottom ply first & 9 plies \\
\bottomrule
\end{tabular}
\caption{The six keys \texttt{1d\_sg.py} reads from \texttt{layup\_db.yaml}.}
\label{tab:solid-layupdb}
\end{table}

For \texttt{fraction}, \texttt{0} is the bottom (OML) face, \texttt{0.5} the
mid-surface and \texttt{1} the top face. It is a property of the \emph{mesh},
not of the homogenization call: the generated node coordinates are measured from
the reference plane, so they run from $-\text{fraction}\cdot h$ to
$(1-\text{fraction})\cdot h$ and $x_3 = 0$ \emph{is} the reference. The $B$
block of the ABD depends on this choice, so homogenize, analyze and recover with
the same \texttt{fraction}.

The shipped case is a $[0/90/0/90/\text{core}]_s$ sandwich: eight
graphite/epoxy plies of $0.9525$~mm ($E_1 = 128$~GPa, $E_2 = E_3 = 11$~GPa,
$G_{12}=G_{13}=4.48$~GPa, $G_{23}=1.53$~GPa, $\rho = 1500$~kg/m$^3$) around a
$144.78$~mm HEREX C70.130 PVC foam core ($E = 103.63$~MPa, $G = 50$~MPa,
$\nu = 0.32$, $\rho = 130$~kg/m$^3$), total $h = 0.1524$~m.

The same \texttt{layup\_db.yaml} also carries a \texttt{V2} flag, a
\texttt{plate:} block and a \texttt{dehom:} block. Those are read by the
recovery scripts of Section~\ref{sec:solid-dehom1d}; \texttt{1d\_sg.py} itself
uses only the six keys in Table~\ref{tab:solid-layupdb}.

\subsection{Run it}

\begin{lstlisting}[style=shellst]
cd examples/OpenSG-solid/1_get_plate_props_from_1DSG
python 1d_sg.py
\end{lstlisting}

A different database file may be passed as the single positional argument:

\begin{lstlisting}[style=shellst]
python 1d_sg.py my_other_layup.yaml
\end{lstlisting}

\subsection{What comes out}

\begin{table}[htbp]
\centering
\small
\begin{tabular}{lp{9.6cm}}
\toprule
file & content \\
\midrule
\texttt{1dsg.yaml} & the generated 1-D SG mesh: an \texttt{sg:} header
  (\texttt{type: plate\_1d}, \texttt{elem\_order: 4}, \texttt{n\_per\_layer: 1},
  \texttt{reference\_fraction: 0.5}, \texttt{thickness}, \texttt{n\_ply}), 37
  nodes given as single $x_3$ coordinates from $-0.0762$ to $+0.0762$~m, 9
  five-noded elements, the material blocks with densities, and one element set
  plus one \texttt{sections:} entry per ply carrying its material, angle and
  thickness \\
\texttt{1dsg.png} & the mesh figure, plies coloured by material \\
\texttt{layup\_db\_plate\_homo.out} & the plate law, named
  \texttt{<layup\_db stem>\_plate\_homo.out} \\
\bottomrule
\end{tabular}
\caption{Outputs of \texttt{1d\_sg.py}.}
\label{tab:solid-ex1out}
\end{table}

The \texttt{.out} opens with a banner that records the whole run, then the
matrix whose title follows \texttt{model} --- \texttt{Reissner-Mindlin}
here, \texttt{Classical Plate} at \texttt{model: 0}:

\begin{lstlisting}[style=shellst,basicstyle=\ttfamily\tiny,breaklines=true]
 OpenSG msg-solid plate model from 1dsg.yaml: 9 plies, h = 0.152400 m, fraction = 0.5, rho*h = 30.251400 kg/m^2, rows [e11 e22 g12 k11 k22 k12 2g13 2g23]

 The Effective Reissner-Mindlin Stiffness Matrix
 --------------------------------------------
     5.4916502E+008     2.6417019E+007     1.5524611E-010     3.6182957E-006     3.6122933E-006     1.0244011E-024     0.0000000E+000     0.0000000E+000
\end{lstlisting}

The effective compliance follows in the same layout and the file closes with
\texttt{ Time taken: 0.82 sec}. Engineering constants are printed only for the
3-D solid law, never for a plate law, because a plate stiffness is not a
material stiffness.

\subsection{What the numbers mean}

The eight diagonal entries of the shipped run, in SI (the code is unit-agnostic
and simply follows the input, which is metres and pascals here):

\begin{table}[htbp]
\centering
\small
\begin{tabular}{lll}
\toprule
term & value & units \\
\midrule
$A_{11}$ & $5.4916502\times10^{8}$ & N/m \\
$A_{22}$ & $5.4916502\times10^{8}$ & N/m \\
$A_{66}$ & $4.1376600\times10^{7}$ & N/m \\
$D_{11}$ & $3.0005458\times10^{6}$ & N$\cdot$m \\
$D_{22}$ & $2.9371144\times10^{6}$ & N$\cdot$m \\
$D_{66}$ & $2.0111708\times10^{5}$ & N$\cdot$m \\
$G_{11}$ & $7.7009772\times10^{6}$ & N/m \\
$G_{22}$ & $7.5423712\times10^{6}$ & N/m \\
\bottomrule
\end{tabular}
\caption{Diagonal of the $8\times8$ plate law of the shipped sandwich.}
\label{tab:solid-ex1diag}
\end{table}

Three checks are worth making on any run of your own, and all three are visible
in this one.

\begin{itemize}
\item \textbf{The section mass.} The banner's $\rho h = 30.2514$~kg/m$^2$ equals
  $8 \times 0.0009525 \times 1500 + 0.14478 \times 130
  = 11.43 + 18.8214$~kg/m$^2$, so the layup that reached the solver is the
  layup you wrote.
\item \textbf{The $B$ block.} The stack is symmetric about the mid-surface and
  \texttt{fraction: 0.5} puts the reference plane there, so membrane-bending
  coupling must vanish. It does: the largest entry of the $B$ block is
  $3.6183\times10^{-6}$ against an $A_{11}$ of $5.4917\times10^{8}$, fourteen
  orders down.
\item \textbf{The transverse-shear block.} $G_{11} = 7.7010\times10^{6}$~N/m
  sits $6.4\,\%$ above the crude core-only estimate
  $G_{\rm core} h_{\rm core} = 50\times10^{6} \times 0.14478
  = 7.239\times10^{6}$~N/m, which is the expected magnitude for a sandwich whose
  shear is carried almost entirely by the foam. $A_{11}$, by contrast, is set by
  the stiff faces.
\end{itemize}

\begin{figure}[htbp]
\centering
\includegraphics[width=0.5\textwidth]{\FIGDIR/plate_1dsg.png}
\caption{The through-thickness 1-D structure gene of the sandwich laminate, as
written by the run to \texttt{1dsg.png}. The gene is a line, drawn with
thickness for legibility: the thin graphite/epoxy face plies are at top and
bottom, the foam core between them, and the dashed line marks the reference
surface at mid-thickness.}
\label{fig:solid-plate1dsg}
\end{figure}

\subsection{What it is validated against}

\texttt{src/opensg\_solid/rm\_plate\_1D/tests/test\_msg\_rm\_plate.py} gates the
construction on analytic anchors and on an independent plain-NumPy re-assembly
of the same SG matrices. For an isotropic plate with $\nu = 0$ the
transverse-shear block comes out exactly $\tfrac56 G h$, the textbook
shear-correction factor, with the least-squares residual $U^*$ at machine zero
(\texttt{test\_isotropic\_nu0\_gives\_five\_sixths}); $\mathbf{A}_6$ equals the
classical lamination-theory \texttt{msg\_materials.compute\_ABD\_matrix} at the
matching reference plane (\texttt{test\_A6\_matches\_compute\_ABD});
$\mathbf{A}_6(\text{fraction})$ equals the offset transform of
$\mathbf{A}_6(0)$ (\texttt{test\_reference\_plane\_transform}); and the warping
fields satisfy the Euler-Lagrange equations and the gauge condition
$\psi^{\mathsf T} H V = 0$ for $V_0$, $V_{11}$ and $V_{12}$ directly
(\texttt{test\_euler\_lagrange\_residuals}, \texttt{test\_H\_identities\_and\_gauge}).
The formulation is that of Yu, Hodges and Volovoi, \emph{Computers \& Structures}
81 (2003) 439--454.

\section{Through-thickness stress recovery}
\label{sec:solid-dehom1d}

\subsection{What it computes}

Homogenization threw information away; recovery puts it back. Given the section
resultants at a point of the plate, the recovery evaluates the 3-D strain,
stress, warping and displacement at every depth of the gene. The construction is
linear in \textbf{42 drivers},
%
\begin{equation}
\mathbf{D} = [\,\mathbf{E}_6,\ \mathbf{E}_{6,1},\ \mathbf{E}_{6,2},\
 \mathbf{E}_{6,11},\ \mathbf{E}_{6,12},\ \mathbf{E}_{6,22},\ \mathbf{qt}_6\,],
\label{eq:solid-42}
\end{equation}
%
so the code builds three operators once per depth from 42 unit evaluations ---
$\mathbf{T}_E,\mathbf{T}_S \in \mathbb{R}^{6\times42}$ and
$\mathbf{T}_W \in \mathbb{R}^{3\times42}$ --- after which every station is one
matrix-vector product, \texttt{jax.vmap}-ed over the whole field. The plate
strains come from the section law itself,
$\mathbf{E}_6 = \mathbf{A}_6^{-1}\mathbf{F}_6$, so a recovery never bypasses the
homogenization.

Both drivers of this section --- the single station of \texttt{rm\_dehom.py} and
the whole field of \texttt{rm\_dehom\_ff.py} --- walk the same core functions in
\texttt{opensg\_solid.rm\_plate\_1D.rm\_dehom}. They differ only in where the
42-driver vector comes from.

\captree{
  \node[inputnode] (i1) {\texttt{1dsg.yaml}\\
    the 1-D SG of Section~\ref{sec:solid-plate1d}};
  \node[funcnode, below=of i1] (f1) {\texttt{rm\_dehom.load\_sg}\\
    \texttt{read\_plate\_sg\_yaml}, then\\
    \texttt{rm\_plate\_msg} re-run;\\
    \texttt{S6} = inverse of \texttt{A6}};
  \node[inputnode, right=of f1] (i2) {\texttt{F6}, \texttt{Q2v} in the USER\\
    INPUT block, \emph{or} a\\ \texttt{.ff} station table plus the\\
    \texttt{dehom:} block of \texttt{layup\_db.yaml}};
  \node[funcnode, below=of f1] (f2) {\texttt{rm\_dehom.z\_stations}\\
    \texttt{rm\_dehom.build\_ops}\\
    \texttt{TE}, \texttt{TS} $(n_z,6,42)$, \texttt{TW} $(n_z,3,42)$};
  \node[funcnode, below=of f2] (f3) {\texttt{rm\_dehom.assemble\_driver},\\
    or \texttt{np.gradient} on a lattice /\\
    \texttt{rm\_dehom.mls\_gradients} on a cloud};
  \node[datanode, below=of f3] (d1) {\texttt{Eps}, \texttt{Sig}, \texttt{Warp}
    per depth\\ then the \texttt{Q1}/\texttt{Q2} rescale};
  \node[funcnode, below=of d1] (f4)
    {\texttt{rm\_dehom.material\_rotations}\\ laminate frame to ply frame};
  \node[funcnode, below=of f4] (f5) {\texttt{rm\_dehom.write\_field}\\
    \texttt{rm\_dehom.write\_vtk}};
  \node[outnode, right=of f5] (o1) {\texttt{1dsg\_dehom.SM/.EM/.U/.vtk}\\
    or \texttt{<ff stem>\_ff.SM/.EM/.U}\\ and \texttt{<ff stem>\_ff\_gauss.vtk}};
  \draw[capedge] (i1) -- (f1);
  \draw[capedge] (f1) -- (f2);
  \draw[capedge] (f2) -- (f3);
  \draw[capedge] (i2) -- (f3);
  \draw[capedge] (f3) -- (d1);
  \draw[capedge] (d1) -- (f4);
  \draw[capedge] (f4) -- (f5);
  \draw[capedge] (f5) -- (o1);
}

Which driver feeds which component is the thing to understand before choosing an
input.

\begin{table}[htbp]
\centering
\footnotesize
\begin{tabular}{p{2.6cm}p{3.6cm}p{4.2cm}p{3.0cm}}
\toprule
driver & where it comes from & what it produces & without it \\
\midrule
$\mathbf{E}_6$ & $\mathbf{A}_6^{-1}\mathbf{F}_6$ at the station
  & $\sigma_{11}$, $\sigma_{22}$, $\sigma_{12}$ layer by layer & nothing \\
$\mathbf{E}_{6,1}$, $\mathbf{E}_{6,2}$ & in-plane gradients of the resultants
  & $\sigma_{13}$, $\sigma_{23}$ (the $V_{11}/V_{12}$ warping)
  & $\sigma_{13}=\sigma_{23}\equiv0$ \\
$\mathbf{E}_{6,11}$, $\mathbf{E}_{6,12}$, $\mathbf{E}_{6,22}$ & second gradients
  & the $\sigma_{33}$ interior build-up
  & $\sigma_{33}$ keeps only its face terms \\
$\mathbf{qt}_6$ & the \emph{applied} surface-pressure ladder
  \texttt{[q q,1 q,2 q,11 q,12 q,22]}
  & the face traction $\sigma_{33} = -q$ & $\sigma_{33}=0$ at that face \\
\bottomrule
\end{tabular}
\caption{The four driver groups of the through-thickness recovery.}
\label{tab:solid-drivers}
\end{table}

$Q_1$ and $Q_2$ are deliberately not part of $\mathbf{F}_6$: there is no shear
slot among the 42 drivers, because in the asymptotic construction the transverse
shears are produced by the gradient brackets. They enter only as the per-station
rescale targets $\int\sigma_{13}\,dz = Q_1$ and $\int\sigma_{23}\,dz = Q_2$, so
the profile \emph{shape} comes from the gradients while the carried force is
exactly the one you supplied.

\subsection{Driver A: one station, resultants typed into the script}

\texttt{rm\_dehom.py} recovers a single station. The input is the
\texttt{USER INPUT} block at the bottom of the file: \texttt{SG\_YAML} (the mesh
written in Section~\ref{sec:solid-plate1d}), \texttt{F6} and \texttt{Q2v}, plus
the optional gradient and load ladders. \texttt{F6} and \texttt{Q2v} ship as
\texttt{None} and the script raises a \texttt{ValueError} rather than run, by
design --- a built-in number would silently masquerade as a result. The block
below is the file with the two mandatory entries filled in with an example
loading; the values are yours to choose.

\begin{lstlisting}[style=opensg,caption={The USER INPUT block of rm\_dehom.py, with the two mandatory entries supplied.}]
SG_YAML = os.path.join(HERE, "1dsg.yaml")
F6  = [1.0e6, 0.0, 0.0, 5.0e4, 0.0, 0.0]  # [N11, N22, N12, M11, M22, M12]
Q2v = [0.0, 0.0]                          # [Q1, Q2]  [N/m]
DE1 = DE2 = None                          # d(E6)/dx, d(E6)/dy
D11 = D12 = D22 = None                    # second gradients
QT6 = None                                # [q q,1 q,2 q,11 q,12 q,22]
LATTICE = "gauss"                         # "gauss" (default) or "nodes"
\end{lstlisting}

\begin{lstlisting}[style=shellst]
cd examples/OpenSG-solid/2_get_plate_stress_from_1DSG
python rm_dehom.py
\end{lstlisting}

It writes \texttt{1dsg\_dehom.SM}, \texttt{.EM}, \texttt{.U} and \texttt{.vtk}
--- the VTK is part of the default output, not an extra step --- and prints
\texttt{max|S11|} through \texttt{max|S23|}. Because the script refuses to run
until you supply \texttt{F6} and \texttt{Q2v}, these four files are produced by
your own run and are not checked into the folder.

With every optional input left at \texttt{None} this is the honest classical
recovery: $\sigma_{11}, \sigma_{22}, \sigma_{12}$ exact layer by layer,
$\sigma_{13} = \sigma_{23} = 0$ exactly, $\sigma_{33}$ at machine zero. One
station has no neighbours, so it cannot supply gradients, and the code does not
invent them. \texttt{LATTICE} selects the output depths: \texttt{"gauss"} (the
default) puts them at the Gauss abscissae of each element, \texttt{"nodes"} at
the SG nodes.

\subsection{Driver B: a whole field from a \texttt{.ff} station table}

\texttt{rm\_dehom\_ff.py} is the general driver and the one to use for real
work. A \texttt{.ff} file is the plate analogue of a VABS \texttt{.glb}:
whitespace separated, \texttt{\#} comments, one row per in-plane station, in any
row order and on any station geometry. The shipped table was collected from an
Abaqus S8R shell field, but any solver can write it --- the driver contains no
Abaqus. Its first three lines state the format:

\begin{lstlisting}[style=shellst]
# plate dehom FF station table (collected from ex5_shell_S8R_field.dat)
# x[m] y[m] N11 N22 N12[N/m] M11 M22 M12[N] Q1 Q2[N/m] u1 u2 u3[m] wx wy[-]
# 40 x 40 rectangular station lattice, x-fastest
\end{lstlisting}

Ten columns are required (\texttt{x y N11 N22 N12 M11 M22 M12 Q1 Q2}); the
optional columns 11--15 add the mid-surface displacement $u_1, u_2, u_3$ and
slopes $w_x, w_y$, and without them the \texttt{.U} file carries the warping
alone. The loader accepts exactly 10 or 15 columns and raises otherwise.
\textbf{Every gradient driver is computed from the stations themselves}, by one
of two paths chosen automatically: exact tensor-grid finite differences
(\texttt{np.gradient} on the true non-uniform spacings) when the coordinates
form a complete rectangular lattice, or a meshfree KD-tree patch fit
($k = 12$ nearest neighbours, distance-weighted local quadratic least squares,
first and second derivatives from one batched solve) for any other cloud. The
flag \texttt{--scatter} forces the meshfree path.

Unlike the single-station driver, \texttt{rm\_dehom\_ff.py} has no
\texttt{LATTICE} variable at all: it always calls
\texttt{z\_stations(sg, "gauss")} and therefore always reports at the element
Gauss depths. \texttt{LATTICE = "nodes"} exists only in \texttt{rm\_dehom.py}.

The applied face pressure is \emph{not} in the table --- loads are input, not
solver output --- so it lives in \texttt{layup\_db.yaml}:

\begin{lstlisting}[style=opensg,caption={The dehom: block of layup\_db.yaml.}]
dehom:
  ff: Abaqus_results/ex5_shell_S8R_field.ff
  loads:
    top: {q0: 68.9476e6, shape: sinsin}   # uniform | sin_x | sin_y | sinsin
    bottom: none
\end{lstlisting}

Each \texttt{shape} fixes the whole six-term ladder analytically, differentiated
with \texttt{a} and \texttt{b} from the \texttt{plate:} block ($a = b = 1.524$~m
in the shipped file). Bottom-face loads raise \texttt{NotImplementedError}: the
42-driver operator block carries the top ladder only.

\begin{lstlisting}[style=shellst]
cd examples/OpenSG-solid/2_get_plate_stress_from_1DSG
python rm_dehom_ff.py ex5_shell_S8R_field.ff
\end{lstlisting}

Pass the table explicitly, as above. The \texttt{dehom: ff:} key in the shipped
\texttt{layup\_db.yaml} points at
\texttt{Abaqus\_results/ex5\_shell\_S8R\_field.ff}, and that folder does not
exist in this example: the table sits at the top level of the folder as
\texttt{ex5\_shell\_S8R\_field.ff}. Leaving the path to the key therefore trips
the driver's own guard, which reads \texttt{no .ff table: pass it on the command
line or set dehom: ff in} followed by the database path. Either give the file
name on the command line, as above, or correct the key. \texttt{--db} points at a database elsewhere. The run prints its chosen
gradient path (\texttt{lattice FD (40 x 40)} for this table), the operator
build, the \texttt{jax.vmap} over stations, and the peak of each stress
component.

\subsection{What comes out}

Outputs are named after the table, \texttt{<ff stem>\_ff}.

\begin{table}[htbp]
\centering
\footnotesize
\begin{tabular}{lp{8.0cm}}
\toprule
file & content \\
\midrule
\texttt{ex5\_shell\_S8R\_field\_ff.SM} & 3-D stress: columns \texttt{x y z} then
  \texttt{S11 S22 S33 S12 S13 S23} in Pa \\
\texttt{ex5\_shell\_S8R\_field\_ff.EM} & 3-D strain, same layout, \texttt{E11}
  through \texttt{E23} \\
\texttt{ex5\_shell\_S8R\_field\_ff.U} & 3-D displacement \texttt{U1 U2 U3} in m \\
\texttt{ex5\_shell\_S8R\_field\_ff\_gauss.vtk} & ASCII
  \texttt{STRUCTURED\_GRID}, \texttt{DIMENSIONS 40 40 45}, 72\,000 points,
  written only on a lattice \\
\bottomrule
\end{tabular}
\caption{Outputs of \texttt{rm\_dehom\_ff.py} on the shipped table.}
\label{tab:solid-ex2out}
\end{table}

The three text files share a two-line header naming the recovery and its station
set. For this run the second line reads \texttt{1600 .ff stations
[lattice FD (40 x 40)] x the 45 SG-element Gauss depths; ALL gradient drivers
from the stations; S/E in the PLY MATERIAL frame (U global); z = 0 is the
reference surface}. Four conventions matter when you read the files.

\begin{itemize}
\item $z$ is in the \textbf{SG frame}: $z = 0$ is the reference surface, not the
  mid-plane, unless \texttt{fraction} put them at the same place.
\item The depths are the Gauss abscissae of each five-noded element --- five per
  ply here, 45 in all --- so every output point is strictly inside one ply, and
  a material interface, where stress is genuinely double-valued, never appears
  in a file.
\item Stress and strain are written in the \textbf{ply material frame} of each
  depth, which is where failure criteria live. The chain computes in the
  laminate frame --- it must, since the gradients and the $Q$-rescale integrate
  across plies --- and rotates only at write time. The displacement \texttt{.U}
  stays global.
\item Component order in the files is \texttt{11 22 33 12 13 23}, reordered from
  the internal SG Voigt order $(11, 22, 33, 23, 13, 12)$ on write.
\end{itemize}

The \texttt{\_ff\_gauss.vtk} is the file to open when you want to look at the
result rather than read it. Figure~\ref{fig:solid-pv1dsg} is the shipped table
recovered and rendered straight out of it.

\begin{figure}[htbp]
\centering
\includegraphics[width=0.47\linewidth]{\FIGDIR/pv_plate_1dsg_s11.png}\hfill
\includegraphics[width=0.47\linewidth]{\FIGDIR/pv_plate_1dsg_s13.png}
\caption{ParaView renders of the recovered field on the SG, produced from the
shipped \texttt{ex5\_shell\_S8R\_field\_ff\_gauss.vtk} by
\texttt{docs/latex/render\_pv.py}. Left: the \texttt{S11} array, $\sigma_{11}$
in Pa, colour scale $\pm5.23\times10^{10}$; right: the \texttt{S13} array,
$\sigma_{13}$ in Pa, colour scale $\pm8.38\times10^{8}$. The grid points
\emph{are} the recovery stations --- the $40\times40$ in-plane lattice times the
45 SG-element Gauss depths --- so each coloured value is one recovered
Gauss-point value, drawn with flat shading and never averaged across an element
boundary. The camera looks down the thickness axis, so the visible face is the
outermost Gauss depth. $\sigma_{11}$ there follows the sin-sin pressure of the
\texttt{dehom: loads:} block and peaks at the plate centre; $\sigma_{13}$ is
near zero on the full-field scale, because the transverse shear must vanish at
the traction-free face, and changes sign across mid-span in $x$ following the
$\cos(\pi x/a)$ gradient of that same load.}
\label{fig:solid-pv1dsg}
\end{figure}

\subsection{What it is validated against}

The accuracy table in
\texttt{src/opensg\_solid/rm\_plate\_1D/USER\_GUIDE.md} records the recovery
errors against reference elasticity solutions, falling at $\mathcal{O}(S^{-2})$
in the slenderness $S$.

\begin{table}[htbp]
\centering
\small
\begin{tabular}{lccl}
\toprule
component (chain) & $S = 4$ & $S = 10$ & thin \\
\midrule
$\sigma_{13}$ (cross-ply) & 21\,\% & 4.4\,\% & 0.18\,\% at $S = 50$ \\
$\sigma_{33}$ equilibrium & 5.5\,\% & 1.2\,\% & 0.05\,\% \\
$\sigma_{11}$ second-order & 27\,\% & 2.0\,\% & 0.04\,\% at $S = 64$ \\
$\sigma_{12}$ angle-ply & 18\,\% & --- & 0.07\,\% at $S = 64$ \\
$U_1$ & 30\,\% & 1.4\,\% & 0.02\,\% at $S = 64$ \\
\bottomrule
\end{tabular}
\caption{Recovery accuracy against reference elasticity solutions, from
\texttt{USER\_GUIDE.md}.}
\label{tab:solid-accuracy}
\end{table}

Soft-core sandwiches need a larger $S$ for the in-plane second-order recovery;
their $\sigma_{13}$, $\sigma_{33}$ and displacements stay accurate at $S = 10$.
Separately, \texttt{dehom\_README.md} in the example folder records that the
finite-difference gradient path and the analytically differentiated known-mode
path were run side by side on the same field and agree to better than
$1\,\%$ on every component in the interior (four integration points trimmed per
edge); the $11$--$15\,\%$ whole-field figure is entirely the one-sided stencil
band about two elements wide at the boundary. That same README derives the full
chain equation by equation and states which term is Yu's and which is this
project's. The second-order recovery gate itself is
\texttt{examples/OpenSG-solid/CS\_OpeNSG\_exampels/static/validate\_v2.py},
which drives four Pagano cylindrical-bending cases with exact closed-form
harmonic drivers; its reference \texttt{.dat} profiles are produced by the
per-case \texttt{run\_case.py} and are not checked in.

\section{Plate properties from a 2-D cell SG}
\label{sec:solid-plate2d}

\subsection{What it computes}

When the wall has in-plane microstructure --- a honeycomb core, a corrugation, a
truss --- the gene is a 2-D mesh of one periodic cell, and \texttt{plate\_homo\_2d}
homogenizes it to the plate ABD with \texttt{n\_model=2}. Opposite faces, edges
and corners are tied through the sparse periodic map;
\texttt{boundary="periodic"} is the default and is the correct treatment for a
periodic medium.

Everything below the runner happens inside the single \texttt{plate\_homo\_2d}
call, in this order.

\captree{
  \node[inputnode] (i1) {\texttt{RHC\_SW\_2UC\_45.yaml}\\
    the 2-D cell SG (a SwiftComp\\ \texttt{.sc} is converted once by\\
    \texttt{helper.sc\_to\_yaml.convert})};
  \node[funcnode, below=of i1] (f1) {\texttt{sg\_homo.plate\_homo\_2d}\\
    $\rightarrow$ \texttt{sg\_mesh.load\_sg\_input}\\
    the parsed dict: \texttt{dim} (inferred\\ from the mesh),
    \texttt{nodes}, \texttt{cells},\\ \texttt{mat\_id}, \texttt{materials}};
  \node[funcnode, right=of f1] (p1) {\texttt{sg\_mesh.plot\_sg\_mesh}};
  \node[funcnode, below=of f1] (f2) {\texttt{fe\_jax.setup.}\\
    \texttt{mesh\_to\_periodic\_sparse\_assembly\_map}\\
    ties opposite faces, edges, corners};
  \node[outnode, right=of f2] (op) {\texttt{RHC\_SW\_2UC\_45\_mesh.png}};
  \node[funcnode, below=of f2] (f3)
    {\texttt{sg\_materials.get\_heterogeneous\_C\_matrix}\\
     \texttt{C\_ess} $(E,6,6)$, rebuilt from\\
     \texttt{material\_param} and \texttt{angles}};
  \node[inputnode, right=of f3] (i2) {\texttt{n\_model = 2}, \texttt{name},\\
    \texttt{material\_param}, \texttt{angles}\\ in \texttt{plate\_homo\_2dsg.py}};
  \node[funcnode, below=of f3] (f4) {\texttt{sg\_homo.\_homo\_direct}\\
    six unit macro states, one sparse\\
    factorization $\rightarrow$ \texttt{V0}, \texttt{C\_eff}, \texttt{omega}};
  \node[funcnode, right=of f4] (f5) {\texttt{sg\_homo.plate\_shear\_ladder}\\
    only at \texttt{shear\_refined=True}\\
    $\rightarrow$ \texttt{G\_msg}, \texttt{ABDG}};
  \node[datanode, below=of f4] (d1) {\texttt{r["C\_eff"]} $6\times6$ ABD\\
    \texttt{[N11 N22 N12 M11 M22 M12]}};
  \node[funcnode, below=of d1] (f6) {\texttt{sg\_homo.write\_sc\_K}};
  \node[outnode, right=of f6] (o1) {\texttt{RHC\_SW\_2UC\_45.out}};
  \draw[capedge] (i1) -- (f1);
  \draw[capedge] (f1) -- (p1);
  \draw[capedge] (p1) -- (op);
  \draw[capedge] (f1) -- (f2);
  \draw[capedge] (f2) -- (f3);
  \draw[capedge] (i2) -- (f3);
  \draw[capedge] (f3) -- (f4);
  \draw[capedge2] (f4) -- (f5);
  \draw[capedge] (f4) -- (d1);
  \draw[capedge] (d1) -- (f6);
  \draw[capedge] (f6) -- (o1);
}

\subsection{The example SG and the input you edit}

\texttt{RHC\_SW\_2UC\_45.yaml} is a two-unit-cell composite honeycomb sandwich
section: 4251 nodes, 6640 three-node triangles, three materials. Coordinates
are in millimetres and moduli in MPa. The yaml \emph{is} the whole problem: a
short scalar header above the mesh blocks says what to run, and the mesh and
material blocks say what to run it on.

\begin{lstlisting}[style=opensg,caption={The header of a solid SG yaml. Every key is defaulted; a headerless mesh still runs.}]
n_model: 2      # 1 = beam, 2 = plate, 3 = solid -- the macro model
refined: 0      # 0 = classical (plate ABD / beam EB); 1 = shear-refined
analysis: H     # H = homogenization | D = dehomogenization
\end{lstlisting}

\begin{table}[htbp]
\centering
\footnotesize
\begin{tabular}{llp{7.6cm}}
\toprule
key & values & meaning \\
\midrule
\texttt{msg} & \texttt{shell}, \texttt{solid}
  & which engine owns the file; omit and the mesh dialect decides \\
\texttt{n\_model} & \texttt{1}, \texttt{2}, \texttt{3}
  & macro model: 1 beam, 2 plate, 3 3-D solid \\
\texttt{refined} & \texttt{0}, \texttt{1}
  & 0 classical (plate ABD $6\times6$, beam EB $4\times4$), 1 shear-refined
    (plate ABDG $8\times8$, beam Timoshenko $6\times6$); ignored for the solid
    model \\
\texttt{analysis} & \texttt{H}, \texttt{D}
  & optional; the command-line argument wins \\
\texttt{epsilon\_bar} & 6 floats & the macro state for \texttt{D} \\
\texttt{omega} & float
  & optional user SG measure, overriding the one measured from the mesh \\
\texttt{aperiodic} & \texttt{1}
  & optional; zero fluctuation on every bounding-box-face node instead of
    periodicity \\
\bottomrule
\end{tabular}
\caption{The header contract of the solid SG yaml. \textbf{There is no
\texttt{dim:} key and no \texttt{scale:} key.} The SG dimension is read from
the mesh --- the element node count, with ties broken by the leading
node-coordinate columns that actually vary --- and the SG measure $\omega$ is
measured from the mesh. A legacy \texttt{dim:}/\texttt{scale:} left in an old
file is still tolerated by the loader, but neither is part of the input
contract and neither should be written into a new file. (The
\texttt{junction:} family of keys belongs to \texttt{msg-shell} recovery and is
not read on this route.)}
\label{tab:solid-header}
\end{table}

The \texttt{material\_param} and \texttt{angles} arguments of
\texttt{plate\_homo\_2d} are an \emph{override} of the material blocks stored
in the file, kept for scripting; examples 4 and 7 drive the same cell that way:

\begin{lstlisting}[style=opensg,caption={The material override block, as examples 4 and 7 carry it.}]
n_model = 2                    # 1: Beam; 2: Plate; 3: 3D elastic
name = "RHC_SW_2UC_45"         # reads <name>.yaml

material_param = jnp.array([
    (108e3, 8e3, 8e3, 4e3, 4e3, 3e3, 0.32, 0.32, 0.30),      # ply +45
    (108e3, 8e3, 8e3, 4e3, 4e3, 3e3, 0.32, 0.32, 0.30),      # ply -45
    (69e3, 69e3, 69e3, 26.54e3, 26.54e3, 26.54e3, 0.30, 0.30, 0.30)])
angles = jnp.array([45.0, -45.0, 0.0])   # deg per material; 0.0 = none
\end{lstlisting}

\texttt{material\_param} is one row of engineering constants per material,
$(E_1, E_2, E_3, G_{12}, G_{13}, G_{23}, \nu_{12}, \nu_{13}, \nu_{23})$, and
\texttt{angles} one rotation per material in degrees. Passing them
\emph{overrides} the material blocks stored in the SG file. The shipped
\texttt{RHC\_SW\_2UC\_45.yaml} keeps the same data in exactly one place
instead: all three materials are \texttt{type: 1}, the nine engineering
constants together with their ply rotation --- \texttt{angle: 45.0} and
\texttt{angle: -45.0} for the two face laminae, and the isotropic aluminium
core ($E = 69000$~MPa, $G = 26540$~MPa, $\nu = 0.3$) at \texttt{angle: 0.0}.
A file that stores its plies \emph{pre-rotated} instead, as a \texttt{type: 2}
full $6\times6$ $\mathbf{C}$ block, must not be given \texttt{angles} as well,
or the plies are rotated twice. The same numbers are recorded next to the
runner in \texttt{mat\_override.yaml} for reference; omit both arguments to use
the file's own blocks unchanged.

The solver reads \texttt{<name>.yaml}. A SwiftComp \texttt{.sc} is converted
once with the packaged helper. That helper is a library function, not a
command-line tool: the module \texttt{sc\_to\_yaml} has no \texttt{\_\_main\_\_}
block, so
\texttt{python -m opensg\_solid.helper.sc\_to\_yaml file.sc}
does nothing. Call
\texttt{convert} instead, which writes \texttt{<base>.yaml} and
\texttt{<base>.msh} and returns the parsed dictionary:

\begin{lstlisting}[style=opensg,caption={Converting a SwiftComp .sc once.}]
from opensg_solid.helper.sc_to_yaml import convert

sc = convert("RHC_SW_2UC_45.sc")   # writes RHC_SW_2UC_45.yaml + .msh
\end{lstlisting}

\subsection{Run it}

The yaml says what to do, so the run is one command with one argument:

\begin{lstlisting}[style=shellst]
cd examples/OpenSG-solid/3_get_plate_props_from_2D_SG
opensg RHC_SW_2UC_45.yaml
\end{lstlisting}

\texttt{opensg} is the unified entry point: it resolves the engine from the
yaml's \texttt{msg:} key --- or, absent that key, from the mesh dialect --- and
hands the arguments on unchanged. \texttt{opensg\_solid RHC\_SW\_2UC\_45.yaml}
and \texttt{python -m opensg\_solid RHC\_SW\_2UC\_45.yaml} are the same run
through the engine-specific entry point.
\texttt{plate\_homo\_2dsg.py} in the folder is the identical call through the
Python API, \texttt{plate\_homo\_2d(name + ".yaml")}, for driving it from your
own script.

\subsection{What comes out}

The run writes \texttt{RHC\_SW\_2UC\_45.out}, the timed \texttt{Classical Plate}
stiffness and compliance in the \texttt{.K} layout, and
\texttt{RHC\_SW\_2UC\_45\_mesh.png}, the cell mesh with elements coloured by
material. It also prints the same $6\times6$ to the terminal, under the macro
law's own title. The \texttt{.out} banner records the SG measure and the
boundary treatment:

\begin{lstlisting}[style=shellst,basicstyle=\ttfamily\tiny]
 OpenSG msg-solid plate model, omega 35.248001, periodic

 The Effective Classical Plate Stiffness Matrix
 --------------------------------------------
     3.3585973E+005     7.7500043E+004    -1.4366892E+002     4.1488404E-001     1.4861902E+000     3.6602858E-002
\end{lstlisting}

closing with \texttt{ Time taken: 7.83 sec} for this 4251-node, 6640-triangle
cell.

\subsection{What the numbers mean}

$\omega = 35.248001$ is the cell's in-plane period, and it is directly checkable
against the mesh: the nodes run from $y_1 = -17.624$ to $+17.624$~mm. It is
\emph{measured}, never declared --- the yaml carries no normalization key,
because the homogenizer always measures its own SG. For a cell converted from
SwiftComp the measurement is also checkable against the source: the trailing
line of \texttt{RHC\_SW\_2UC\_45.sc} reads \texttt{35.248}, and the measured
$\omega$ agrees with it to the digits that line carries, which is the check
that the mesh and the \texttt{.sc} header describe the same cell. The
diagonal, in the mesh's own units:

\begin{table}[htbp]
\centering
\small
\begin{tabular}{lll}
\toprule
term & value & units \\
\midrule
$A_{11}$ & $3.3585973\times10^{5}$ & N/mm \\
$A_{22}$ & $1.1291139\times10^{5}$ & N/mm \\
$A_{66}$ & $1.0551654\times10^{5}$ & N/mm \\
$D_{11}$ & $9.0367521\times10^{7}$ & N$\cdot$mm \\
$D_{22}$ & $5.7011050\times10^{7}$ & N$\cdot$mm \\
$D_{66}$ & $4.6026320\times10^{7}$ & N$\cdot$mm \\
\bottomrule
\end{tabular}
\caption{Diagonal of the honeycomb cell's plate ABD.}
\label{tab:solid-ex3diag}
\end{table}

The membrane law is strongly anisotropic, $A_{11}/A_{22} = 2.97$, while the
bending law is much closer to isotropic, $D_{11}/D_{22} = 1.59$: the two
directions see the cell very differently in stretching and much less
differently in bending. The largest entry of the whole $B$ block is $4.954$
against $A$ terms of order $10^{5}$ and $D$ terms of order $10^{7}$, which is
what a cell symmetric about its mid-plane should give and is the first health
check to make on a run of your own.

\begin{figure}[htbp]
\centering
\includegraphics[width=0.78\textwidth]{\FIGDIR/rhc_2dsg_mesh.png}
\caption{The two-unit-cell honeycomb sandwich structure gene, elements coloured
by material, as written to \texttt{RHC\_SW\_2UC\_45\_mesh.png}. The
$\pm45^\circ$ face laminae (materials 1 and 2) are the thin horizontal bands at
$y_2 = \pm23$~mm and the isotropic core web (material 3) folds between them over
the in-plane period. The through-thickness coordinate is $y_2$, spanning
$\pm23.35$~mm.}
\label{fig:solid-rhcmesh}
\end{figure}

\subsection{The Reissner-Mindlin option}

Setting \texttt{refined: 1} in the yaml header --- or passing
\texttt{refined=1} to \texttt{plate\_homo\_2d}, of which
\texttt{shear\_refined=True} is the legacy plate-only spelling --- additionally
runs the RM first-order warping ladder and returns \texttt{r["G\_msg"]} (the
$2\times2$ transverse-shear block), \texttt{r["ABDG"]} (the same $8\times8$
block form as Eq.~\eqref{eq:solid-abdg}), \texttt{r["A6\_ladder"]} and
\texttt{r["Ustar\_rel"]}. When the ladder produces an SPD fit the \texttt{.out}
is written from the $8\times8$ and its matrices are titled
\texttt{Reissner-Mindlin} instead of \texttt{Classical Plate}; otherwise
the run falls back to the $6\times6$. The same header switch serves the beam
route: \texttt{n\_model: 1} with \texttt{refined: 0} reports the classical
\texttt{Euler-Bernoulli Beam} $4\times4$, and with \texttt{refined: 1} the
\texttt{Timoshenko Beam} $6\times6$. The solid law has no refined variant.

\begin{lstlisting}[style=opensg,caption={Requesting the 8x8 from a 2-D cell SG.}]
from opensg_solid.sg_homo import plate_homo_2d

r = plate_homo_2d("RHC_SW_2UC_45.yaml", material_param=material_param,
                  angles=angles, n_model=2, shear_refined=True)
print(r["ABDG"])
\end{lstlisting}

\subsection{The 1-D gene in the same folder, and a naming trap}

The folder also ships \texttt{Plate\_1D\_SG\_2UC\_45.yaml}, a \textbf{1-D}
structure gene --- two-node line elements through a stack, its dimension read
from the mesh rather than declared in the header --- which runs through the
identical command by naming that file instead.
\texttt{plate\_homo\_2d} reads 1-D, 2-D and 3-D genes from the same code path,
so a stack of plies can be homogenized either here or through the dedicated
through-thickness route of Section~\ref{sec:solid-plate1d}.

Be careful about the output name. \texttt{plate\_homo\_2d} always derives its
report path from the \emph{input} basename, so setting
\texttt{name = "Plate\_1D\_SG\_2UC\_45"} writes
\texttt{Plate\_1D\_SG\_2UC\_45.out} --- the \texttt{.K}-layout report --- and
nothing else. The file \texttt{Plate\_1D\_SG\_2UC\_45\_plate\_ABD.out} sitting
in that folder is not produced by \texttt{plate\_homo\_2dsg.py} at all: it is an
\texttt{np.savetxt} leftover from an older script style, a bare $6\times6$ under
a single comment line (\texttt{\# plate 6x6 [N11 N22 N12 M11 M22 M12] from the
1D SG Plate\_1D\_SG\_2UC\_45.sc}) with no banner, no compliance block and no
timing. Its contents are still a valid record of that gene ---
$A_{11} = A_{22} = 1.65452926\times10^{6}$,
$A_{66} = 6.42527474\times10^{5}$,
$D_{11} = D_{22} = 7.86727254\times10^{7}$,
$D_{66} = 3.12496023\times10^{7}$ in N/mm and N$\cdot$mm --- but do not expect
your run to reproduce that file name or that layout.

\section{Plate recovery from a 2-D SG}
\label{sec:solid-platedehom2d}

\subsection{What it computes}

Given a converged plate solution you know the macro plate strain at the point of
interest. Feeding it back into the homogenized cell recovers the local 3-D
strain and stress at \textbf{every Gauss point of the gene}, which is where a
honeycomb actually fails --- in the web, not in the smeared plate.
\texttt{dehom\_fields} returns the strain, the stress and the fluctuation
displacement in one pass, using the warping columns $V_0$ that the
homogenization already solved. Nothing is re-solved.

That reuse is the whole point of the tree below: the homogenization node runs
once and its result dictionary carries the warping modes, the element stiffness
table and the quadrature data straight into the recovery.

\captree{
  \node[inputnode] (i1) {\texttt{RHC\_SW\_2UC\_45.yaml}\\ the 2-D cell SG};
  \node[funcnode, below=of i1] (f1) {\texttt{sg\_homo.plate\_homo\_2d}\\
    \texttt{n\_model=2}, the Section~\ref{sec:solid-plate2d} solve};
  \node[funcnode, right=of f1] (f2) {\texttt{sg\_homo.write\_sc\_K},\\
    \texttt{sg\_mesh.plot\_sg\_mesh},\\
    \texttt{np.savetxt} in the runner};
  \node[datanode, below=of f1] (d1) {\texttt{r["C\_eff"]}, \texttt{r["V0"]},\\
    \texttt{C\_ess}, \texttt{x\_end}, \texttt{phi\_qn},\\
    \texttt{dphi\_dxi\_qnp}, \texttt{periodic\_cells\_en}};
  \node[outnode, right=of d1] (o1) {\texttt{RHC\_SW\_2UC\_45.out}\\
    \texttt{RHC\_SW\_2UC\_45\_mesh.png}\\
    \texttt{RHC\_SW\_2UC\_45\_plate\_ABD.out}};
  \node[funcnode, below=of d1] (f3) {\texttt{sg\_dehom.dehom\_fields}\\
    \texttt{\_dehom\_batch} and \texttt{\_u\_at\_gauss}\\
    $\rightarrow$ \texttt{Gam}, \texttt{Sig} $(E,Q,6)$, \texttt{U} $(E,Q,3)$};
  \node[inputnode, right=of f3] (i2) {\texttt{epsilon\_bar}\\
    \texttt{[e11 e22 2e12 k11 k22 2k12]}\\ in \texttt{plate\_dehom\_2dsg.py}};
  \node[funcnode, below=of f3] (f4) {\texttt{sg\_dehom.export\_gauss}\\
    \texttt{sg\_dehom.gauss\_coords},
    \texttt{\_write\_field}};
  \node[outnode, below=of f4] (o2) {\texttt{RHC\_SW\_2UC\_45\_dehom.txt}\\
    \texttt{.vtk}, \texttt{.SM}, \texttt{.EM}, \texttt{.U}};
  \draw[capedge] (i1) -- (f1);
  \draw[capedge] (f1) -- (f2);
  \draw[capedge] (f2) -- (o1);
  \draw[capedge] (f1) -- (d1);
  \draw[capedge] (d1) -- (f3);
  \draw[capedge] (i2) -- (f3);
  \draw[capedge] (f3) -- (f4);
  \draw[capedge] (f4) -- (o2);
}

\subsection{The input you edit}

The recovery needs one thing beyond the homogenization: the macro plate strain,
in the order \texttt{[e11, e22, 2e12, k11, k22, 2k12]}. On the command-line
route it is a header key of the SG yaml, read when the analysis is \texttt{D}:

\begin{lstlisting}[style=opensg,caption={The macro plate strain driving the recovery, in the yaml header.}]
epsilon_bar: [0.0, 0.1, 0.0, 0.1, 0.0, 0.0]
\end{lstlisting}

\texttt{plate\_dehom\_2dsg.py} sets the same six numbers in code, as
\texttt{epsilon\_bar = jnp.array([0.0, 0.1, 0.0, 0.1, 0.0, 0.0])}, alongside
the \texttt{name}, \texttt{material\_param} and \texttt{angles} of
Section~\ref{sec:solid-plate2d}.

The shipped state combines a transverse extension $\varepsilon_{22} = 0.1$ with
a bending curvature $\kappa_{11} = 0.1$~mm$^{-1}$. This vector is yours: it
comes from your plate solution at the station you care about, and there is no
default that would be meaningful.

\subsection{Run it}

The trailing \texttt{D} asks for dehomogenization: the run homogenizes first,
then recovers, and reports one \texttt{Time taken} covering the whole thing.

\begin{lstlisting}[style=shellst]
cd examples/OpenSG-solid/4_get_plate_dehom_from_2DSG
opensg RHC_SW_2UC_45.yaml D
\end{lstlisting}

\texttt{python plate\_dehom\_2dsg.py} is the same chain through the Python API,
and additionally writes the bare $6\times6$
\texttt{RHC\_SW\_2UC\_45\_plate\_ABD.out} through \texttt{np.savetxt}.

\subsection{What comes out}

\begin{table}[htbp]
\centering
\footnotesize
\begin{tabular}{lp{8.2cm}}
\toprule
file & content \\
\midrule
\texttt{RHC\_SW\_2UC\_45.out} & the timed \texttt{.K}-layout
  \texttt{Classical Plate} report written by \texttt{plate\_homo\_2d} inside the
  script; regenerated on each run \\
\texttt{RHC\_SW\_2UC\_45\_plate\_ABD.out} & the plain-text $6\times6$ ABD via
  \texttt{np.savetxt}, header \texttt{plate 6x6 [N11 N22 N12 M11 M22 M12] from
  the 2D SG RHC\_SW\_2UC\_45} \\
\texttt{RHC\_SW\_2UC\_45\_dehom.txt} & one row per Gauss point:
  \texttt{GP\_Index}, \texttt{Coord\_X}, \texttt{Coord\_Y}, then
  \texttt{Eps\_xx Eps\_yy Eps\_zz Eps\_yz Eps\_xz Eps\_xy} and \texttt{Sig\_xx}
  through \texttt{Sig\_xy} --- SwiftComp component order; one header line plus
  19\,920 rows \\
\texttt{RHC\_SW\_2UC\_45\_dehom.vtk} & ASCII \texttt{UNSTRUCTURED\_GRID} point
  cloud of the same 19\,920 points, each strain and stress component as its own
  \texttt{SCALARS} array, ready for ParaView \\
\texttt{RHC\_SW\_2UC\_45\_dehom.SM} / \texttt{.EM} / \texttt{.U} & the same
  fields in the \texttt{rm\_plate\_1D} layout: \texttt{x y z} then
  \texttt{S11 S22 S33 S12 S13 S23}, \texttt{E11} through \texttt{E23}, and
  \texttt{U1 U2 U3} \\
\texttt{RHC\_SW\_2UC\_45\_mesh.png} & the mesh figure, from the
  \texttt{plate\_homo\_2d} call the script makes first \\
\bottomrule
\end{tabular}
\caption{Outputs of \texttt{plate\_dehom\_2dsg.py}.}
\label{tab:solid-ex4out}
\end{table}

The 19\,920 rows are 6640 triangles times three quadrature points. The script
prints the ABD diagonal and the peak of each stress component.

Two conventions are worth knowing before reading the files. In the
\texttt{.SM}/\texttt{.EM}/\texttt{.U} tables the SG coordinates are placed in
the \textbf{last} slots of the \texttt{x y z} triple, so a 2-D gene leaves the
\texttt{x} column at zero and puts $y_1$ in \texttt{y} and the
through-thickness $y_2$ in \texttt{z}. And \texttt{.U} is the
\textbf{fluctuation-only} displacement --- its own header says so --- because a
unit-state SG recovery has no macro displacement to add.

\subsection{Reading the result}

The first Gauss point of the shipped run sits at
$(y_1, y_2) = (17.54067,\ 23.29375)$~mm, right at the top face of the cell whose
half-thickness is $23.35$~mm, and reports
%
\begin{equation}
\varepsilon_{11} = 2.329375, \qquad
\sigma_{11} = 8.755003\times10^{4}, \quad
\sigma_{22} = 7.102971\times10^{4}, \quad
\sigma_{12} = -6.542276\times10^{4}
\end{equation}
%
in MPa. That $\varepsilon_{11}$ is exactly $\kappa_{11} y_2
= 0.1 \times 23.29375$, the macro part of the strain reproduced to every printed
digit, which is the cheapest possible check that the strain state you passed is
the strain state that arrived. The other five components are what the recovery
adds: $\varepsilon_{22} = 0.2643$ and $\varepsilon_{33} = -0.9187$ are the
fluctuation response of the cell, and $\sigma_{12}$ is large even though the
applied macro shear strain $2\varepsilon_{12}$ is zero, which is consistent with
the off-axis face laminae the mesh figure shows at that height.

Over the whole gene the field maxima are:

\begin{table}[htbp]
\centering
\small
\begin{tabular}{lcc}
\toprule
component & $\max\lvert\sigma\rvert$ (MPa) & $\max\lvert\varepsilon\rvert$ \\
\midrule
$xx$ & $1.641632\times10^{5}$ & $2.329375$ \\
$yy$ & $9.199997\times10^{4}$ & $1.499038$ \\
$zz$ & $2.535821\times10^{4}$ & $1.032184$ \\
$yz$ & $1.083274\times10^{4}$ & $1.153314$ \\
$xz$ & $5.210321\times10^{3}$ & $8.807437\times10^{-1}$ \\
$xy$ & $7.369978\times10^{4}$ & $7.881358\times10^{-1}$ \\
\bottomrule
\end{tabular}
\caption{Component maxima over the 19\,920 Gauss points of the shipped plate
recovery.}
\label{tab:solid-ex4max}
\end{table}

The $xx$ strain maximum is again $\kappa_{11} y_2^{\max} = 2.329375$, as it must
be at the extreme fibre. Local states like the $\sigma_{zz}$ and $\sigma_{yz}$
columns are exactly what a smeared plate model cannot show you, and the reason
to run a dehomogenization at all. The driver here is a demonstration magnitude,
not a service load; scale it to your own macro state before reading the stresses
as failure indices.

\begin{figure}[htbp]
\centering
\includegraphics[width=0.62\linewidth]{\FIGDIR/pv_plate_dehom_sxx.png}
\caption{ParaView render of the recovered field on the SG mesh, produced from
the \texttt{RHC\_SW\_2UC\_45\_dehom.vtk} written by this example, using
\texttt{docs/latex/render\_pv.py}. The array is \texttt{Sig\_xx}, the axial
stress $\sigma_{xx}$ in MPa (units follow the SG input, millimetres and
megapascals), on a symmetric colour scale of $\pm1.64\times10^{5}$~MPa --- the
$xx$ row of Table~\ref{tab:solid-ex4max}. Every one of the 19\,920 points is a
single Gauss point carried by its own vertex cell and coloured flat, so no value
has been averaged across an element boundary. The field varies linearly with the
through-thickness coordinate $y_2$, which is the macro part $\kappa_{11}y_2$ of
the driver, and each face pack splits into two sub-bands of visibly different
$\sigma_{xx}$ at almost the same height, consistent with the $+45^\circ$ and
$-45^\circ$ laminae differing in the sign of their $C_{16}$ coupling.}
\label{fig:solid-pvplatedehom}
\end{figure}

\section{Equivalent 3-D solid from a 2-D SG}
\label{sec:solid-3dfrom2d}

\subsection{What it computes}

\texttt{n\_model=3} reduces an SG to an equivalent \textbf{3-D solid} stiffness:
the $6\times6$ material law of the homogenized continuum, in Voigt order
$[11, 22, 33, 23, 13, 12]$ with engineering shears. On a \textbf{2-D} SG this is
the generalized-plane-strain problem: the cell is periodic in its two in-plane
directions and uniform out of plane, six unit macro strains are applied, the
fluctuation field $\mathbf{V}_0$ is solved for each, and the effective law is
%
\begin{equation}
\mathbf{D}_{\rm eff} =
 \frac{\mathbf{D}_{ee} + \mathbf{V}_0^{\mathsf T}\mathbf{D}_{he}}{\omega},
\label{eq:solid-deff}
\end{equation}
%
with $\omega$ the cell area. This is the classical micromechanics unit-cell
problem, which makes it the natural place to validate the route against
published numbers.

The shipped runner takes the self-contained path on the left of the tree; the
engine path on the right is the same macro model reached through
\texttt{plate\_homo\_2d}, and is what a 2-D SG in the ordinary
\texttt{sg\_mesh.load\_sg\_input} dialect goes through.

\captree{
  \node[inputnode] (i1) {\texttt{UDcomp\_2D.msh}\\ gmsh 2.2, physical tags\\
    naming the matrix and fibre};
  \node[funcnode, below=of i1] (f1) {\texttt{make\_udcomp\_yaml.py}\\
    nodes, 1-based triangles, a constant\\
    \texttt{elementOrientations} frame, the two\\
    materials in MPa, \texttt{sets.element}};
  \node[outnode, below=of f1] (o1) {\texttt{UDcomp\_2D.yaml}};
  \node[funcnode, below=of o1] (f2) {\texttt{solid\_homo\_2dsg.C\_from\_eng}\\
    then the CST assembly:
    \texttt{Ke}, \texttt{Fe},\\ \texttt{Dee}, \texttt{Dhe}, \texttt{omega}};
  \node[inputnode, right=of f2] (i2) {\texttt{n\_model = 3}, \texttt{name},\\
    the \texttt{paper} dictionary\\ in \texttt{solid\_homo\_2dsg.py}};
  \node[funcnode, below=of f2] (f3) {\texttt{solid\_homo\_2dsg.pair}\\
    right to left, top to bottom;\\
    corner chains by \texttt{master = master[master]};\\
    three rigid translations pinned};
  \node[datanode, below=of f3] (d1) {\texttt{V0} $(n_{\rm dof},6)$\\
    $\mathbf{D}_{\rm eff} = (\mathbf{D}_{ee}
      + \mathbf{V}_0^{\mathsf T}\mathbf{D}_{he})/\omega$};
  \node[funcnode, right=of d1] (e1) {\texttt{sg\_homo.plate\_homo\_2d}\\
    \texttt{n\_model=3} on
    \texttt{RHC\_SW\_2UC\_45.yaml}\\ $\rightarrow$
    \texttt{RHC\_SW\_2UC\_45.out}\\ with the engineering-constants block};
  \node[outnode, below=of d1] (o2) {\texttt{UDcomp\_2D\_Deff.out}\\
    \texttt{UDcomp\_2D\_constants\_vs\_paper.dat}\\
    \texttt{UDcomp\_2D\_mesh.png}};
  \draw[capedge] (i1) -- (f1);
  \draw[capedge] (f1) -- (o1);
  \draw[capedge] (o1) -- (f2);
  \draw[capedge] (i2) -- (f2);
  \draw[capedge] (f2) -- (f3);
  \draw[capedge] (f3) -- (d1);
  \draw[capedge2] (d1) -- (e1);
  \draw[capedge] (d1) -- (o2);
}

\subsection{The example SG}

\texttt{examples/OpenSG-solid/5\_get\_solid\_props\_from\_SG} homogenizes a
unidirectional fibre/matrix square unit cell: 801 nodes, 1536 three-node
triangles, 512 matrix elements and 1024 fibre elements, cell area
$\omega = 1$.

\begin{table}[htbp]
\centering
\small
\begin{tabular}{lllll}
\toprule
material & $E_1, E_2, E_3$ (MPa) & $G_{12}, G_{13}, G_{23}$ (MPa)
 & $\nu_{12}, \nu_{13}, \nu_{23}$ & density \\
\midrule
matrix & 4760, 4760, 4760 & 1737.226 each & 0.37, 0.37, 0.37 & 1200.0 \\
fibre & 276000, 19500, 19500 & 70000, 70000, 5735 & 0.28, 0.28, 0.70 & 1800.0 \\
\bottomrule
\end{tabular}
\caption{The two phases of the UD unit cell. The matrix shear modulus is written
by the converter as $4760/(2\cdot1.37)$.}
\label{tab:solid-udmat}
\end{table}

\begin{figure}[htbp]
\centering
\includegraphics[width=0.5\textwidth]{\FIGDIR/udcomp_2dsg_mesh.png}
\caption{The unidirectional composite unit cell written to
\texttt{UDcomp\_2D\_mesh.png}, fibre and matrix elements coloured by phase, with
the SG coordinates labelled $y_2$ and $y_3$.}
\label{fig:solid-udmesh}
\end{figure}

\subsection{Step 1: build the SG yaml from the gmsh mesh}

\texttt{make\_udcomp\_yaml.py} is a one-time converter from
\texttt{UDcomp\_2D.msh} (gmsh 2.2 format) to \texttt{UDcomp\_2D.yaml}. It writes
the nodes as \texttt{- [x y z]} rows, the triangles as 1-based
\texttt{- [n1 n2 n3]} rows, a constant \texttt{elementOrientations} frame
($e_1$ = axial $z$, $e_2$ = $+x$, $e_3$ = $+y$), the two materials as
engineering constants in MPa, and \texttt{sets.element} splitting the mesh into
the matrix and fibre phases by gmsh physical tag. The yaml is already checked
in, so this step is only needed if you remesh.

\begin{lstlisting}[style=shellst]
cd examples/OpenSG-solid/5_get_solid_props_from_SG
python make_udcomp_yaml.py
\end{lstlisting}

\subsection{Step 2: homogenize}

\begin{lstlisting}[style=shellst]
python solid_homo_2dsg.py
\end{lstlisting}

The script is kept rather than replaced by the unified command because it does
more than homogenize: it writes the comparison table against the published
constants and the mesh figure. For a 2-D SG in the ordinary solid dialect the
same macro law is one command, \texttt{opensg <name>.yaml} with
\texttt{n\_model: 3} in the header.

The \texttt{User Input} block is three lines: \texttt{n\_model = 3},
\texttt{name = "UDcomp\_2D"} (which reads \texttt{<name>.yaml}), and the
\texttt{paper} dictionary of reference constants to compare against.

This particular runner is deliberately self-contained. It assembles
constant-strain triangles, the periodic master-slave node map (right to left
matching $y$, top to bottom matching $x$, with corner chains resolved by three
passes of \texttt{master = master[master]}) and the three pinned rigid
translations in plain NumPy, so the entire \texttt{n\_model=3} algebra is
readable in one file. It therefore writes its own \texttt{\_Deff.out} through
\texttt{np.savetxt} rather than the SwiftComp-layout report, and it consumes the
\texttt{nodes}/\texttt{elements}/\texttt{elementOrientations}/\texttt{materials}/\texttt{sets}
yaml dialect that \texttt{make\_udcomp\_yaml.py} emits, not the
\texttt{nodes}/\texttt{cells}/\texttt{mat\_id}/\texttt{materials}
dialect that \texttt{opensg\_solid.sg\_mesh.load\_sg\_input} parses. The engine
path itself is exercised on the honeycomb SG: \texttt{n\_model=3} on
\texttt{RHC\_SW\_2UC\_45.yaml} runs through \texttt{plate\_homo\_2d} and writes
the full SwiftComp \texttt{.out} with its engineering-constants block.

\subsection{What comes out}

\begin{table}[htbp]
\centering
\footnotesize
\begin{tabular}{lp{8.2cm}}
\toprule
file & content \\
\midrule
\texttt{UDcomp\_2D\_Deff.out} & the $6\times6$ $\mathbf{D}_{\rm eff}$ in MPa,
  Voigt order \texttt{[e11 e22 e33 2e23 2e13 2e12]}, with $\omega$ in the
  header \\
\texttt{UDcomp\_2D\_constants\_vs\_paper.dat} & the nine effective engineering
  constants next to the published values and the percentage difference \\
\texttt{UDcomp\_2D\_mesh.png} & the unit cell, elements coloured by phase \\
\bottomrule
\end{tabular}
\caption{Outputs of \texttt{solid\_homo\_2dsg.py}.}
\label{tab:solid-ex5out}
\end{table}

\subsection{What the numbers mean, and what they were validated against}

The nine constants are read off the compliance
$\mathbf{S} = \mathbf{D}_{\rm eff}^{-1}$ as $E_i = 1/S_{ii}$,
$G_{ij} = 1/S_{kk}$ and $\nu_{ij} = -S_{ij}/S_{ii}$, and compared against Table
II Model-3 of the ASC-23 OpenSG-Solid paper --- the same case as the FEniCS
\texttt{3DModel.ipynb} in \texttt{wenbinyugroup/OpenSG}. The shipped
\texttt{UDcomp\_2D\_constants\_vs\_paper.dat} reads:

\begin{table}[htbp]
\centering
\small
\begin{tabular}{lrrr}
\toprule
constant & this code & paper & \% difference \\
\midrule
$E_1$ (GPa) & 167.255337 & 167.255300 & $+0.0000$ \\
$E_2$ (GPa) & 11.413657 & 11.413600 & $+0.0005$ \\
$E_3$ (GPa) & 11.413657 & 11.413600 & $+0.0005$ \\
$G_{23}$ (GPa) & 3.095555 & 3.095500 & $+0.0018$ \\
$G_{13}$ (GPa) & 6.801763 & 6.801700 & $+0.0009$ \\
$G_{12}$ (GPa) & 6.801763 & 6.801700 & $+0.0009$ \\
$\nu_{12}$ & 0.311780 & 0.311780 & $-0.0000$ \\
$\nu_{13}$ & 0.311780 & 0.311780 & $-0.0000$ \\
$\nu_{23}$ & 0.575743 & 0.575740 & $+0.0006$ \\
\bottomrule
\end{tabular}
\caption{\texttt{UDcomp\_2D\_constants\_vs\_paper.dat}, verbatim.}
\label{tab:solid-udvspaper}
\end{table}

Every constant agrees with the published value to within $0.002\,\%$, which is
the round-off of the paper's own printed precision. The route is exact against
its reference, not merely close.

The stiffness itself, from \texttt{UDcomp\_2D\_Deff.out} (MPa, $\omega = 1$):
%
\widematrix{
\mathbf{D}_{\rm eff} =
\begin{bmatrix}
1.726544\times10^{5} & 8.658488\times10^{3} & 8.658488\times10^{3} & 0 & 0 & 0\\
8.658488\times10^{3} & 1.750725\times10^{4} & 1.026390\times10^{4} & 0 & 0 & 0\\
8.658488\times10^{3} & 1.026390\times10^{4} & 1.750725\times10^{4} & 0 & 0 & 0\\
0 & 0 & 0 & 3.095555\times10^{3} & 0 & 0\\
0 & 0 & 0 & 0 & 6.801763\times10^{3} & 0\\
0 & 0 & 0 & 0 & 0 & 6.801763\times10^{3}
\end{bmatrix}}

Two health checks are worth making a habit of. First, the symmetry the geometry
demands is recovered: $D_{22} = D_{33}$ and $D_{55} = D_{66}$ to all printed
digits, and every normal-shear coupling term in the shipped file is either
exactly zero or below $3.4\times10^{-2}$~MPa, six to seven orders under the
diagonal, which is the numerical signature of a correctly tied periodic cell.
Second, the transverse plane is \textbf{not} isotropic: a transversely isotropic
material would satisfy $D_{44} = (D_{22} - D_{23})/2 = 3621.7$~MPa, whereas the
computed $D_{44}$ is $3095.6$~MPa, $14.5\,\%$ lower. A square fibre array is
tetragonal, so $G_{23}$ genuinely has to be homogenized rather than inferred
from $E_2$ and $\nu_{23}$, which is precisely why a unit-cell solve is needed
instead of a rule of mixtures.

\section{Equivalent 3-D solid from a 3-D SG}
\label{sec:solid-3dfrom3d}

\subsection{What it computes}

When the SG is periodic in all three directions --- a triply-periodic minimal
surface (TPMS) unit cell, a lattice block, any genuine 3-D microstructure ---
the same macro model is reached with
\texttt{plate\_homo\_2d(..., n\_model=3, boundary="periodic")} on a tet or hex
mesh. \texttt{"periodic"} is the default on every route;
\texttt{boundary="aperiodic"} maps the boundary solution onto the
bounding-box-face nodes as Dirichlet data (solid elements carry only
translations, so this clamps every degree of freedom of those nodes) and is the
explicit request for SGs that genuinely are not periodic. Aperiodic supports the
plate and solid models with \texttt{solver="direct"} only; asking for it with
\texttt{n\_model=1} or \texttt{solver="cg"} raises.

The tree is the sample runner and the boundary comparison, the two paths every
number in this section comes from. Note the loader branch: a SwiftComp
\texttt{.sc} is converted on the way in, so the \texttt{.yaml} and \texttt{.msh}
appear in the working directory as a side effect of the homogenization.

\captree{
  \node[inputnode] (i1) {\texttt{Sample\_1.sc} or \texttt{Sample\_2.sc}\\
    SwiftComp linear tets on a unit cell};
  \node[inputnode, right=of i1] (i2) {\texttt{n\_model = 3},
    \texttt{E}, \texttt{nu},\\ \texttt{boundary},
    \texttt{workdir} in\\ \texttt{run\_sample\_N.py}};
  \node[funcnode, below=of i1] (f1) {\texttt{sg\_homo.plate\_homo\_2d}\\
    the single entry point};
  \node[funcnode, below=of f1] (f2) {\texttt{sg\_mesh.load\_sg\_input}\\
    $\rightarrow$ \texttt{helper.sc\_to\_yaml.convert}};
  \node[funcnode, right=of f2] (p1) {\texttt{sg\_mesh.plot\_sg\_mesh}};
  \node[funcnode, below=of f2] (f3) {\texttt{fe\_jax.setup.}\\
    \texttt{mesh\_to\_periodic\_sparse\_assembly\_map},\\
    or the bounding-box-face Dirichlet set\\
    at \texttt{boundary="aperiodic"}; then\\
    \texttt{sg\_materials.get\_heterogeneous\_C\_matrix}};
  \node[outnode, right=of f3] (oa) {\texttt{Sample\_N.yaml}\\
    \texttt{Sample\_N.msh}\\ \texttt{Sample\_N\_mesh.png}};
  \node[funcnode, below=of f3] (f4) {\texttt{sg\_homo.\_homo\_direct}\\
    six unit macro states $\rightarrow$\\
    \texttt{C\_eff}, \texttt{V0}, \texttt{omega}};
  \node[funcnode, right=of f4] (fb) {\texttt{compare\_boundary\_modes.py}\\
    both samples under both treatments};
  \node[funcnode, below=of f4] (f5) {\texttt{sg\_homo.write\_sc\_K}, then the\\
    runner reads \texttt{Sample\_N.sc.k}\\ and differences it term by term};
  \node[outnode, right=of f5] (ob) {\texttt{boundary\_compare.dat}\\
    \texttt{Sample\_N\_periodic.out}\\ \texttt{Sample\_N\_aperiodic.out}};
  \node[outnode, below=of f5] (o1) {\texttt{Sample\_N.out}\\
    \texttt{Sample\_N\_msg\_solid.out}\\ \texttt{Sample\_N\_compare.dat}};
  \draw[capedge] (i1) -- (f1);
  \draw[capedge] (i2) -- (f1);
  \draw[capedge] (f1) -- (f2);
  \draw[capedge] (f2) -- (p1);
  \draw[capedge] (p1) -- (oa);
  \draw[capedge] (f2) -- (f3);
  \draw[capedge] (f3) -- (f4);
  \draw[capedge] (f4) -- (f5);
  \draw[capedge] (f5) -- (o1);
  \draw[capedge2] (f4) -- (fb);
  \draw[capedge2] (fb) -- (ob);
}

\subsection{The samples and the runners}

Two TPMS samples are supplied as SwiftComp \texttt{.sc} meshes of linear tets on
a unit cell, in aluminium. Reading the \texttt{.sc} meta line directly gives
their sizes.

\begin{table}[htbp]
\centering
\small
\begin{tabular}{llrrr}
\toprule
sample & surface & nodes & tets & relative density \\
\midrule
\texttt{Sample\_1} & Schwarz-P & 116\,851 & 545\,741 & 0.3000 \\
\texttt{Sample\_2} & Schwarz-P & 54\,874 & 191\,957 & 0.0857 \\
\bottomrule
\end{tabular}
\caption{The two TPMS solid structure genes. The relative density is the
fraction of the cell filled by aluminium; it is \emph{not} the SG measure. Both
genes sit on a unit cell, so for both $\omega = 1$.}
\label{tab:solid-tpms}
\end{table}

\begin{figure}[htbp]
\centering
\includegraphics[width=0.6\linewidth]{\FIGDIR/schwarz_p_mesh.png}
\caption{The Schwarz-P triply-periodic minimal surface on its unit cell, drawn
from the real mesh of \texttt{schwarz\_p\_3Dshell.yaml} by
\texttt{examples/OpenSG\_shell/4\_get\_solid\_props\_from\_shell\_3D\_SG/}
\texttt{schwarz\_solid\_props.py}, with the SG coordinates labelled $y_1$, $y_2$
and $y_3$. This is the surface itself, meshed as a \texttt{msg-shell} SG; the
two solid genes of Table~\ref{tab:solid-tpms} discretise the same surface as a
volume of linear tets, which is why \texttt{tpms\_all\_compare.dat} can put the
two engines side by side at matched thickness.}
\label{fig:solid-schwarzp}
\end{figure}

Four runners sit in \texttt{6\_get\_solid\_props\_from\_3D\_SG}.

\begin{table}[htbp]
\centering
\footnotesize
\begin{tabular}{p{4.4cm}p{9.4cm}}
\toprule
script & what it does \\
\midrule
\texttt{sample\_1/run\_sample\_1.py}, \texttt{sample\_2/run\_sample\_2.py} &
  homogenize one sample at $E = 69$~GPa, $\nu = 0.3$ with
  \texttt{boundary="periodic"} and compare term by term against the vendor's own
  \texttt{Sample\_N.sc.k}; writes \texttt{Sample\_N.out},
  \texttt{Sample\_N\_msg\_solid.out}, \texttt{Sample\_N\_compare.dat} and
  \texttt{Sample\_N\_mesh.png} \\
\texttt{compare\_boundary\_modes.py} & runs both samples under both boundary
  treatments and writes \texttt{boundary\_compare.dat}, plus per-mode copies
  \texttt{Sample\_N\_periodic.out} and \texttt{Sample\_N\_aperiodic.out} \\
\texttt{tpms\_solid\_props.py} & the plain homogenization driver, taking the
  sample number as an optional argument \\
\texttt{sc\_to\_3dyaml.py} & converts the \texttt{.sc} tets to the OpenSG 3-D
  solid yaml and reports the tet-volume check that distinguishes a solid mesh
  from a shell surface \\
\bottomrule
\end{tabular}
\caption{The runners of example 6.}
\label{tab:solid-ex6runners}
\end{table}

\begin{lstlisting}[style=shellst]
cd examples/OpenSG-solid/6_get_solid_props_from_3D_SG/sample_1
python run_sample_1.py
\end{lstlisting}

\begin{lstlisting}[style=shellst]
cd examples/OpenSG-solid/6_get_solid_props_from_3D_SG
python compare_boundary_modes.py
\end{lstlisting}

Both scripts are kept as scripts because both do something the solver does not:
the sample runner differences the result against the vendor's \texttt{.K} and
writes \texttt{Sample\_N\_compare.dat}, and
\texttt{compare\_boundary\_modes.py} runs each sample twice and tabulates the
two treatments against each other. The homogenization on its own is one
command --- \texttt{opensg Sample\_1.yaml}, whose header already carries
\texttt{msg: solid}, \texttt{n\_model: 3} and \texttt{refined: 0}.

One detail to check before comparing numbers across scripts: the two sample
runners and \texttt{compare\_boundary\_modes.py} use $E = 69$~GPa, matched to
the SwiftComp \texttt{.K}, whereas \texttt{tpms\_solid\_props.py} and
\texttt{sc\_to\_3dyaml.py} use $E = 70$~GPa. Only the $69$~GPa runs are
comparable with the benchmark below. The \texttt{Sample\_N\_msg\_solid.out}
files that \texttt{tpms\_solid\_props.py} writes at the top level of the folder
are regenerated on each run.

\subsection{Normalization: the \texttt{.out} is already per unit cell}

For a 3-D solid SG the measure $\omega$ is the \textbf{node bounding-box
volume}, not the summed element (material) volume. Both these genes sit on a
unit cell, so $\omega = 1$ and the \texttt{.out} is written per unit cell ---
exactly the normalization a SwiftComp \texttt{.K} uses.
\texttt{Sample\_1.out} and \texttt{Sample\_1\_periodic.out} both open with
\texttt{ OpenSG msg-solid 3D elastic model, omega 1, periodic} and report
$C_{11} = 1.0190158\times10^{10}$~Pa $= 10190.1580$~MPa, against the
\texttt{1.0190158E+004}~MPa of \texttt{Sample\_1.sc.k}; Sample 2 reports
$2.3031598\times10^{9}$~Pa $= 2303.1598$~MPa against \texttt{2.3031598E+003}. The
two files are therefore \emph{directly} comparable, with no rescale in between,
and the comparison scripts' surviving \texttt{C x omega} line is a no-op at
$\omega = 1$.

This is a change from the older behaviour, in which $\omega$ was the material
volume --- $0.3$ and $0.0857$ for these two cells --- so the \texttt{.out} came
out $1/\text{relative density}$ too stiff and only the comparison table's
rescale made the numbers line up. Any note to that effect is obsolete: the
figures printed in the \texttt{.out} now are the figures in the tables below.
Every comparison table below is per unit cell, in MPa.

\subsection{The SwiftComp digit-parity benchmark}

\begin{table}[htbp]
\centering
\small
\begin{tabular}{lrrrcrrr}
\toprule
& \multicolumn{3}{c}{Sample 1} & & \multicolumn{3}{c}{Sample 2} \\
\cmidrule(lr){2-4}\cmidrule(lr){6-8}
term & msg-solid & SwiftComp & err \% & & msg-solid & SwiftComp & err \% \\
\midrule
$C_{11}$ & 10190.158 & 10190.158 & $-0.000$ & & 2303.160 & 2303.160 & $-0.000$ \\
$C_{12}$ & 5646.069 & 5646.068 & $+0.000$ & & 1612.130 & 1612.130 & $-0.000$ \\
$C_{13}$ & 5646.092 & 5646.091 & $+0.000$ & & 1612.141 & 1612.141 & $+0.000$ \\
$C_{22}$ & 10190.131 & 10190.131 & $+0.000$ & & 2303.067 & 2303.067 & $+0.000$ \\
$C_{23}$ & 5646.012 & 5646.012 & $+0.000$ & & 1612.120 & 1612.120 & $+0.000$ \\
$C_{33}$ & 10189.958 & 10189.958 & $+0.000$ & & 2303.019 & 2303.018 & $+0.000$ \\
$C_{44}$ & 4519.640 & 4519.640 & $-0.000$ & & 1106.734 & 1106.735 & $-0.000$ \\
$C_{55}$ & 4519.637 & 4519.637 & $-0.000$ & & 1106.668 & 1106.668 & $+0.000$ \\
$C_{66}$ & 4519.755 & 4519.755 & $-0.000$ & & 1106.664 & 1106.664 & $-0.000$ \\
\bottomrule
\end{tabular}
\caption{\texttt{Sample\_1\_compare.dat} and \texttt{Sample\_2\_compare.dat},
per unit cell in MPa, at $E = 69$~GPa and $\nu = 0.3$. Sample 1 is 545\,741
tets and solved in 165.9~s; Sample 2 is 191\,957 tets and solved in 61.4~s.}
\label{tab:solid-tpmsparity}
\end{table}

Digit-for-digit on both samples, to seven significant figures on every entry.
The effective law is cubic to five digits, and every symmetry-forbidden coupling
sits four to five orders below the diagonal, which is the health check on the
all-directions periodic tie. The engineering constants that the \texttt{.out}
prints for Sample 2 --- $E_1 = 9.7551133\times10^{8}$~Pa,
$G_{12} = 1.1066643\times10^{9}$~Pa, $\nu_{12} = 0.41174$ --- are per unit
cell in the same way, and only the two moduli carry that normalization;
Poisson's ratio is a ratio.

\subsection{Periodic versus aperiodic}

\texttt{compare\_boundary\_modes.py} runs the same \texttt{.sc} through both
treatments, so the percentage column measures the aperiodic boundary layer of a
single cell directly against the exact periodic answer.

\begin{table}[htbp]
\centering
\small
\begin{tabular}{lrrrcrrr}
\toprule
& \multicolumn{3}{c}{Sample 1} & & \multicolumn{3}{c}{Sample 2} \\
\cmidrule(lr){2-4}\cmidrule(lr){6-8}
term & aperiodic & periodic & \% diff & & aperiodic & periodic & \% diff \\
\midrule
$C_{11}$ & 12977.036 & 10190.158 & $+27.349$ & & 2810.329 & 2303.160 & $+22.021$ \\
$C_{22}$ & 12977.030 & 10190.131 & $+27.349$ & & 2810.363 & 2303.067 & $+22.027$ \\
$C_{33}$ & 12976.815 & 10189.958 & $+27.349$ & & 2810.251 & 2303.019 & $+22.025$ \\
$C_{12}$ & 5423.565 & 5646.069 & $-3.941$ & & 1548.627 & 1612.130 & $-3.939$ \\
$C_{13}$ & 5423.674 & 5646.092 & $-3.939$ & & 1548.954 & 1612.141 & $-3.919$ \\
$C_{23}$ & 5423.151 & 5646.012 & $-3.947$ & & 1549.049 & 1612.120 & $-3.912$ \\
$C_{44}$ & 5545.997 & 4519.640 & $+22.709$ & & 1300.352 & 1106.734 & $+17.494$ \\
$C_{55}$ & 5546.044 & 4519.637 & $+22.710$ & & 1300.151 & 1106.668 & $+17.483$ \\
$C_{66}$ & 5546.133 & 4519.755 & $+22.709$ & & 1300.048 & 1106.664 & $+17.474$ \\
\midrule
solve & 146.3~s & 150.2~s & & & 52.4~s & 51.5~s & \\
\bottomrule
\end{tabular}
\caption{\texttt{boundary\_compare.dat}, per unit cell in MPa. The aperiodic
run clamps 17\,610 boundary nodes on Sample 1 and 4456 on Sample 2.}
\label{tab:solid-boundary}
\end{table}

Read the table this way. Clamping the fluctuation on the cube faces forbids the
affine fluctuation modes, so the cell can only come out stiffer: the normal and
shear diagonals rise by $17$ to $27\,\%$, while the off-diagonal $C_{12}$ family
falls by about $4\,\%$. The bias is large because the clamped set is large ---
$17\,610$ of $116\,851$ nodes, or $15\,\%$, for Sample 1, and $4456$ of
$54\,874$, or $8\,\%$, for Sample 2 --- since the solid sheet meets the cube
faces in two-dimensional patches. On the solid route aperiodic is also
\textbf{not} faster: the periodic tie merges only face pairs, so the factorized
size barely changes, and the two timings are within a few percent of each other
in both directions. Periodic is the exact treatment for a periodic medium and is
therefore the default; aperiodic is the right model only when the SG genuinely
is not periodic.

The folder carries two further comparison tables, both regenerated by their own
scripts: \texttt{tpms\_all\_compare.dat}, which lines the solid results up
against the same Schwarz-P surface homogenized by the \texttt{msg-shell} route
at matched thickness, and \texttt{solid\_vs\_shell\_compare.dat}, which adds the
normalized stiffness $C/(E\rho)$ and the Zener ratio $2C_{44}/(C_{11}-C_{12})$
--- 1.989 for Sample 1, 3.203 for Sample 2, and 3.151 for the shell run at the
matched thickness $t = 0.0365$. Those three rows are the cleanest available
statement that the two engines describe the same structure.

\section{Beam properties from a 2-D SG}
\label{sec:solid-beam}

\subsection{What it computes}

\texttt{n\_model=1} reduces a cross-sectional SG to the Timoshenko beam
constitutive law, relating the axial force, two transverse shear forces, torque
and two bending moments to the six generalized strains
$[\epsilon_{11},\gamma_{12},\gamma_{13},\kappa_1,\kappa_2,\kappa_3]$. Internally
this is the Beam\_solid KKT engine merged into
\texttt{sg\_assembly}/\texttt{sg\_homo}: the zeroth-order warping solve $V_0$ is
run under four rigid-body Lagrange constraints, the $l$-chain first-order solve
$V_{1s}$ reuses the same factorization, and the pair is reduced to the
$6\times6$. The classical Euler-Bernoulli $4\times4$ in
$[\epsilon_{11},\kappa_1,\kappa_2,\kappa_3]$ falls out of the same $V_0$ solve
and rides along in \texttt{r["C\_eff\_EB"]} at no extra cost. The beam route
requires an SG of dimension 2 or 3.

The entry point is the same \texttt{plate\_homo\_2d} as every other route, but
\texttt{n\_model=1} branches to a different engine at the first fork, and that
engine is where the two shear rows come from.

\captree{
  \node[inputnode] (i1) {\texttt{RHC\_SW\_2UC\_45.yaml}\\
    the 2-D cross-sectional SG};
  \node[inputnode, right=of i1] (i2) {\texttt{n\_model = 1}, \texttt{name},\\
    \texttt{material\_param}, \texttt{angles}\\ in \texttt{beam\_homo\_sg.py}};
  \node[funcnode, below=of i1] (f1) {\texttt{sg\_homo.plate\_homo\_2d}\\
    $\rightarrow$ \texttt{sg\_homo.\_beam\_homo\_kkt}};
  \node[funcnode, right=of f1] (p1) {\texttt{sg\_mesh.load\_sg\_input},\\
    \texttt{sg\_mesh.plot\_sg\_mesh}};
  \node[funcnode, below=of f1] (f2)
    {\texttt{sg\_assembly.compress\_periodic\_cells\_jax},\\
     \texttt{sg\_materials.get\_heterogeneous\_C\_matrix}\\
     $\rightarrow$ \texttt{C\_asm}, \texttt{C\_stress}};
  \node[outnode, right=of f2] (oa) {\texttt{RHC\_SW\_2UC\_45\_mesh.png}};
  \node[funcnode, below=of f2] (f3)
    {\texttt{sg\_assembly.assemble\_system\_matrices},\\
     \texttt{sg\_assembly.assemble\_rigid\_body\_ops}\\
     $\rightarrow$ \texttt{Dhh} through \texttt{Dle}, \texttt{omega},\\
     \texttt{dehomo\_data}, \texttt{Psi}, \texttt{Dc}};
  \node[funcnode, below=of f3] (f4)
    {\texttt{sg\_assembly.solve\_fluctuation\_field}\\
     $\rightarrow$ \texttt{V0}, \texttt{C\_eb} under four rigid-body\\
     constraints; then \texttt{sg\_homo.prepare\_v1\_rhs} and\\
     \texttt{sg\_homo.finalize\_v1\_and\_compute\_deff}\\
     reuse the same factorization};
  \node[datanode, below=of f4] (d1) {\texttt{r["C\_eff"]} Timoshenko $6\times6$\\
    \texttt{r["C\_eff\_EB"]} classical $4\times4$\\
    \texttt{r["V0"]}, \texttt{r["V1s"]}};
  \node[funcnode, below=of d1] (f5) {\texttt{sg\_homo.write\_sc\_K}\\
    \texttt{Timoshenko} stiffness and compliance};
  \node[outnode, right=of f5] (o1) {\texttt{RHC\_SW\_2UC\_45.out}};
  \draw[capedge] (i1) -- (f1);
  \draw[capedge] (i2) -- (f1);
  \draw[capedge] (f1) -- (p1);
  \draw[capedge] (p1) -- (oa);
  \draw[capedge] (f1) -- (f2);
  \draw[capedge] (f2) -- (f3);
  \draw[capedge] (f3) -- (f4);
  \draw[capedge] (f4) -- (d1);
  \draw[capedge] (d1) -- (f5);
  \draw[capedge] (f5) -- (o1);
}

\subsection{The example SG and the input you edit}

\texttt{examples/OpenSG-solid/7\_get\_beam\_props\_from\_SG} homogenizes the same
\texttt{RHC\_SW\_2UC\_45.yaml} honeycomb cross-section as
Section~\ref{sec:solid-plate2d}: 4251 nodes, 6640 three-node triangles.
Coordinates are in mm and moduli in MPa, so the resulting stiffnesses are in N
and N$\cdot$mm$^2$.

The runner has no command-line arguments --- the SG file and the materials are
the \texttt{User Input} block at the top of the script, identical to the plate
runner except for \texttt{n\_model}.

\begin{table}[htbp]
\centering
\small
\begin{tabular}{cllll}
\toprule
id & material & $E_1, E_2, E_3$ (MPa) & $G_{12}, G_{13}, G_{23}$ (MPa) & angle \\
\midrule
1 & ply & 108000, 8000, 8000 & 4000, 4000, 3000 & $+45^\circ$ \\
2 & ply & 108000, 8000, 8000 & 4000, 4000, 3000 & $-45^\circ$ \\
3 & aluminium core & 69000, 69000, 69000 & 26540, 26540, 26540 & none \\
\bottomrule
\end{tabular}
\caption{The three materials of the honeycomb section; $\nu_{12}, \nu_{13},
\nu_{23}$ are $0.32, 0.32, 0.30$ for the plies and $0.30$ throughout for the
core.}
\label{tab:solid-beammat}
\end{table}

\begin{figure}[htbp]
\centering
\includegraphics[width=0.62\linewidth]{\FIGDIR/rhc_2dsg_mesh.png}
\caption{The mesh the beam route runs on, written by
\texttt{sg\_mesh.plot\_sg\_mesh} to \texttt{RHC\_SW\_2UC\_45\_mesh.png}. It is
the same 4251-node, 6640-triangle gene as Figure~\ref{fig:solid-rhcmesh},
because the plate and beam routes read the identical file; only
\texttt{n\_model} differs. What changes is the reading: with
\texttt{n\_model=1} the mesh is a \emph{cross-section}, its two coordinates are
$y_2$ and $y_3$, and the beam axis $y_1$ is normal to the page --- which is why
the recovery of Section~\ref{sec:solid-beamdehom} reports its extreme fibre at
$y_3 = 23.29375$~mm where the plate recovery of
Section~\ref{sec:solid-platedehom2d} reports the same point as $y_2$.}
\label{fig:solid-beammesh}
\end{figure}

\begin{lstlisting}[style=opensg,caption={beam\_homo\_sg.py, complete.}]
import numpy as np
import jax.numpy as jnp

from opensg_solid.sg_homo import plate_homo_2d

n_model = 1                    # 1: Beam; 2: Plate; 3: 3D elastic
name = "RHC_SW_2UC_45"         # reads <name>.yaml

material_param = jnp.array([
    (108e3, 8e3, 8e3, 4e3, 4e3, 3e3, 0.32, 0.32, 0.30),
    (108e3, 8e3, 8e3, 4e3, 4e3, 3e3, 0.32, 0.32, 0.30),
    (69e3, 69e3, 69e3, 26.54e3, 26.54e3, 26.54e3, 0.30, 0.30, 0.30)])
angles = jnp.array([45.0, -45.0, 0.0])

r = plate_homo_2d(name + ".yaml", material_param=material_param,
                  angles=angles, n_model=n_model)
np.set_printoptions(precision=5)
print("beam 6x6  [eps11 gam12 gam13 kappa1 kappa2 kappa3]  (Timoshenko):")
print(r["C_eff"])
print("beam 4x4  [eps11 kappa1 kappa2 kappa3]  (EB, same run):")
print(r["C_eff_EB"])
\end{lstlisting}

\subsection{Run it}

The homogenization itself is one command; the yaml in this folder carries
\texttt{n\_model: 1} and \texttt{refined: 1} in its header, so it asks for the
Timoshenko $6\times6$ without any argument:

\begin{lstlisting}[style=shellst]
cd examples/OpenSG-solid/7_get_beam_props_from_SG
opensg RHC_SW_2UC_45.yaml
\end{lstlisting}

\texttt{python beam\_homo\_sg.py} runs the same solve through the Python API
and additionally prints the classical Euler-Bernoulli $4\times4$ that rides
along in \texttt{r["C\_eff\_EB"]}, which the command-line route does not show.

\subsection{What comes out}

The run writes \texttt{RHC\_SW\_2UC\_45.out}, the timed SwiftComp-layout report
containing \texttt{The Effective Timoshenko Beam Stiffness Matrix} followed by
\texttt{The Effective Timoshenko Beam Compliance Matrix}, with banner
\texttt{ OpenSG msg-solid beam model, omega 1} and closing line
\texttt{ Time taken: 8.94 sec}; and \texttt{RHC\_SW\_2UC\_45\_mesh.png}, the SG
mesh coloured by material.

The banner reports $\omega = 1$: for a beam model on a 2-D SG the measure is
unity, so the $6\times6$ is the \textbf{total sectional stiffness}, not a
per-area density. No engineering-constants block is printed, because a beam
stiffness is not a material stiffness.

\subsection{What the numbers mean}

The shipped \texttt{RHC\_SW\_2UC\_45.out} stiffness block reads:

\begin{lstlisting}[style=shellst,basicstyle=\ttfamily\tiny]
     9.8992866E+006     6.9973096E+003     1.2577025E+001    -7.7065215E-001     2.2740495E+000    -8.6717860E-001
     6.9973096E+003     1.9074524E+006     5.9324366E+001    -3.4044009E+001    -1.6298115E+000    -6.0670378E-002
     1.2577025E+001     5.9324366E+001     9.3619028E+005    -8.4787043E+001    -9.3250231E-001     8.5948371E-001
    -7.7065184E-001    -3.4044009E+001    -8.4787043E+001     2.9774561E+008    -3.1862083E+006    -7.3057137E+003
     2.2740506E+000    -1.6298115E+000    -9.3250231E-001    -3.1862083E+006     2.1895745E+009    -9.5964309E+004
    -8.6717875E-001    -6.0670378E-002     8.5948371E-001    -7.3057137E+003    -9.5964309E+004     9.2899273E+008
\end{lstlisting}

\begin{table}[htbp]
\centering
\small
\begin{tabular}{lllp{4.6cm}}
\toprule
entry & conjugate pair & value & reads as \\
\midrule
$C_{11}$ & $\epsilon_{11}$ & $9.8992866\times10^{6}$ N & axial (extensional) stiffness \\
$C_{22}$ & $\gamma_{12}$ & $1.9074524\times10^{6}$ N & transverse shear stiffness, direction 2 \\
$C_{33}$ & $\gamma_{13}$ & $9.3619028\times10^{5}$ N & transverse shear stiffness, direction 3 \\
$C_{44}$ & $\kappa_1$ & $2.9774561\times10^{8}$ N$\cdot$mm$^2$ & torsional stiffness \\
$C_{55}$ & $\kappa_2$ & $2.1895745\times10^{9}$ N$\cdot$mm$^2$ & bending stiffness about axis 2 \\
$C_{66}$ & $\kappa_3$ & $9.2899273\times10^{8}$ N$\cdot$mm$^2$ & bending stiffness about axis 3 \\
$C_{45}$ & $\kappa_1$--$\kappa_2$ & $-3.1862083\times10^{6}$ N$\cdot$mm$^2$ & twist-bend coupling from the $\pm45^\circ$ plies \\
\bottomrule
\end{tabular}
\caption{The Timoshenko $6\times6$ of the honeycomb section, entry by entry.}
\label{tab:solid-timo}
\end{table}

Read the diagonal only as a first orientation. When couplings are present the
engineering stiffnesses are the inverses of the \textbf{compliance} diagonal,
which is why the same \texttt{.out} prints the compliance: here
$1/S_{11} = 9.89926\times10^{6}$ against $C_{11} = 9.89929\times10^{6}$
($0.0003\,\%$ apart) and $1/S_{44} = 2.97741\times10^{8}$ against
$C_{44} = 2.97746\times10^{8}$ ($0.0015\,\%$ apart), so on this section the
couplings shift the engineering values by less than $0.002\,\%$.

The twist-bend term $C_{45}$ is the physically meaningful signature of the
$\pm45^\circ$ layup. Quantify it rather than eyeballing it: at
$3.1862\times10^{6}$~N$\cdot$mm$^2$ it is a factor $93$ below $C_{44}$ and a
factor $690$ below $C_{55}$ --- two to three orders down, not three.

\subsection{The built-in cross-check: Euler-Bernoulli against Timoshenko}

\texttt{r["C\_eff\_EB"]} is the classical $4\times4$ from the same $V_0$ solve,
printed by \texttt{beam\_homo\_sg.py}. Its four entries must reproduce the
classical rows of the Timoshenko $6\times6$, since the two differ only by the
$l$-chain refinement that produces the shear terms. On this SG they agree to at
least five significant figures on every classical diagonal.

\begin{table}[htbp]
\centering
\small
\begin{tabular}{lrr}
\toprule
term & EB $4\times4$ & Timoshenko $6\times6$ \\
\midrule
extension & $9.89926177\times10^{6}$ & $9.89928657\times10^{6}$ \\
torsion & $2.97745363\times10^{8}$ & $2.97745606\times10^{8}$ \\
bending 2 & $2.18957447\times10^{9}$ & $2.18957452\times10^{9}$ \\
bending 3 & $9.28992811\times10^{8}$ & $9.28992729\times10^{8}$ \\
\bottomrule
\end{tabular}
\caption{The two classical laws from one run. The Timoshenko column is
\texttt{RHC\_SW\_2UC\_45\_beam\_Timo.out}; the EB column is the same four
numbers, recorded in \texttt{RHC\_SW\_2UC\_45\_beam\_EB.out} in example 8, a
file left over from an earlier version of that script.}
\label{tab:solid-ebvstimo}
\end{table}

The two transverse shear stiffnesses $C_{22}$ and $C_{33}$ exist \textbf{only}
in the Timoshenko $6\times6$; that is the entire reason to run
\texttt{n\_model=1} rather than read the EB block.

\section{Beam recovery}
\label{sec:solid-beamdehom}

\subsection{What it computes}

Dehomogenization takes a macro beam state and recovers the pointwise 3-D strain
and stress inside the SG, at every element Gauss point. For \texttt{n\_model=1}
this runs the generalized-Timoshenko recovery ladder
(\texttt{sg\_dehom.compute\_fluctuations\_gpu}): a successive-derivative
product-rule chain $\mathbf{F}' = \mathbf{R}\,\mathbf{F}$ seeded with
$\mathbf{F}_{1d} = \mathbf{C}_{\rm eff}^{-1}\bar{\boldsymbol{\epsilon}}$, which
builds the shear-corrected states and the zeroth- and first-order warping
displacement fields. Nothing from the homogenization is recomputed: the result
dictionary \texttt{r} carries the warping modes, the material tables and the
quadrature data, so you homogenize once and recover as often as you like.

The tree shows exactly what rides along. Everything below the dashed data node
is a contraction of quantities that already exist.

\captree{
  \node[inputnode] (i1) {\texttt{RHC\_SW\_2UC\_45.yaml}\\
    the 2-D cross-sectional SG};
  \node[funcnode, below=of i1] (f1) {\texttt{sg\_homo.plate\_homo\_2d}\\
    \texttt{n\_model=1}, the Section~\ref{sec:solid-beam} solve};
  \node[funcnode, right=of f1] (fw) {\texttt{sg\_homo.write\_sc\_K},\\
    \texttt{sg\_mesh.plot\_sg\_mesh},\\
    \texttt{np.savetxt} in the runner};
  \node[datanode, below=of f1] (d1) {\texttt{r["C\_eff"]}, \texttt{r["V0"]},
    \texttt{r["V1s"]},\\ \texttt{r["dehomo\_data"]},
    \texttt{r["C\_stress"]},\\ \texttt{r["reduced\_periodic\_cells"]},\\
    \texttt{r["elem\_rotation"]}, \texttt{r["phi\_qn"]}};
  \node[outnode, right=of d1] (o1) {\texttt{RHC\_SW\_2UC\_45.out}\\
    \texttt{RHC\_SW\_2UC\_45\_beam\_Timo.out}\\
    \texttt{RHC\_SW\_2UC\_45\_mesh.png}};
  \node[funcnode, below=of d1] (f2) {\texttt{sg\_dehom.dehom\_fields}\\
    $\rightarrow$ \texttt{sg\_dehom.compute\_fluctuations\_gpu}\\
    seeded with \texttt{C\_eff} inverse times
    \texttt{epsilon\_bar}\\ $\rightarrow$ \texttt{w\_1}, \texttt{w1s\_1},
    \texttt{w1s\_2}, \texttt{w\_2}, \texttt{st\_m\_final}};
  \node[inputnode, right=of f2] (i2) {\texttt{epsilon\_bar}\\
    \texttt{[eps11 gam12 gam13}\\ \texttt{kappa1 kappa2 kappa3]}\\
    in \texttt{beam\_dehom\_sg.py}};
  \node[funcnode, below=of f2] (f3)
    {\texttt{sg\_dehom.recover\_gauss\_point\_fields}\\
     and \texttt{sg\_dehom.\_u\_at\_gauss} $\rightarrow$\\
     \texttt{Gam}, \texttt{Sig} $(E,Q,6)$, \texttt{U} $(E,Q,3)$};
  \node[funcnode, below=of f3] (f4) {\texttt{sg\_dehom.export\_gauss}\\
    \texttt{sg\_dehom.gauss\_coords},
    \texttt{\_write\_field}};
  \node[outnode, below=of f4] (o2) {\texttt{RHC\_SW\_2UC\_45\_dehom.txt}\\
    \texttt{.vtk}, \texttt{.SM}, \texttt{.EM}, \texttt{.U}};
  \draw[capedge] (i1) -- (f1);
  \draw[capedge] (f1) -- (fw);
  \draw[capedge] (fw) -- (o1);
  \draw[capedge] (f1) -- (d1);
  \draw[capedge] (d1) -- (f2);
  \draw[capedge] (i2) -- (f2);
  \draw[capedge] (f2) -- (f3);
  \draw[capedge] (f3) -- (f4);
  \draw[capedge] (f4) -- (o2);
}

\subsection{The input you edit}

\texttt{examples/OpenSG-solid/8\_get\_beam\_dehom\_from\_SG} runs the same
honeycomb SG. Beyond the material block, the one new user input is
\texttt{epsilon\_bar}, the \textbf{macro beam strain}
$[\epsilon_{11},\gamma_{12},\gamma_{13},\kappa_1,\kappa_2,\kappa_3]$ --- a
strain six-vector, not a load vector. The shipped driver is
$\epsilon_{11} = 0.001$ combined with $\kappa_2 = 0.01$~mm$^{-1}$.

\begin{lstlisting}[style=opensg,caption={The body of beam\_dehom\_sg.py after the shared material block.}]
from opensg_solid.sg_homo import plate_homo_2d
from opensg_solid.sg_dehom import dehom_fields, export_gauss

# the MACRO beam state [eps11, gam12, gam13, kappa1, kappa2, kappa3]
epsilon_bar = jnp.array([0.001, 0.0, 0.0, 0.0, 0.01, 0.0])

r = plate_homo_2d(name + ".yaml", material_param=material_param,
                  angles=angles, n_model=n_model)
np.savetxt(name + "_beam_Timo.out", r["C_eff"], fmt="%16.8e",
           header="beam Timoshenko 6x6 [eps11 gam12 gam13 kappa1 kappa2"
                  " kappa3] (Beam_solid KKT engine) from the %dD SG %s"
                  % (r["n_sg"], name))

Gam, Sig, U = dehom_fields(r, epsilon_bar)
export_gauss(r, Gam, Sig, name + "_dehom", U_eqd=U)
for i, nm in enumerate(("xx", "yy", "zz", "yz", "xz", "xy")):
    print("  max|Sig_%s| = %.5e" % (nm, np.abs(Sig[..., i]).max()))
\end{lstlisting}

Only \textbf{classical-channel} states are valid drivers: extension, twist,
bending and their combinations. The recovery chain is kept verbatim from its
validated reference, and its seed is the compliance-scaled state
$\mathbf{C}_{\rm eff}^{-1}\bar{\boldsymbol{\epsilon}}$, which makes the two
transverse shear entries \texttt{epsilon\_bar[1]} and \texttt{epsilon\_bar[2]}
\textbf{recovery-inert} --- a shear-only driver recovers essentially zero
stress. This is a documented deviation from the refined theory stated in the
docstring of \texttt{compute\_fluctuations\_gpu}, not a switch to flip.

\subsection{Run it}

The yaml in this folder carries the same six numbers as its
\texttt{epsilon\_bar:} header key, so the whole chain --- homogenize, then
recover --- is one command with the \texttt{D} argument:

\begin{lstlisting}[style=shellst]
cd examples/OpenSG-solid/8_get_beam_dehom_from_SG
opensg RHC_SW_2UC_45.yaml D
\end{lstlisting}

\texttt{python beam\_dehom\_sg.py} is the same chain through the Python API,
and additionally writes the bare $6\times6$
\texttt{RHC\_SW\_2UC\_45\_beam\_Timo.out} through \texttt{np.savetxt}.

\subsection{What comes out}

\texttt{dehom\_fields} returns three arrays shaped $(E, Q, \cdot)$ over elements
and Gauss points: \texttt{Gam} $(E,Q,6)$ strain, \texttt{Sig} $(E,Q,6)$ stress,
and \texttt{U} $(E,Q,3)$ the \textbf{fluctuation-only} displacement (for the
beam route, the zeroth- plus first-order warping $w_1 + w_{1s}$; there is no
macro contribution in a unit-state SG recovery). \texttt{export\_gauss} writes
them out.

\begin{table}[htbp]
\centering
\footnotesize
\begin{tabular}{lp{8.2cm}}
\toprule
file & content \\
\midrule
\texttt{RHC\_SW\_2UC\_45.out} & the Timoshenko report written by
  \texttt{plate\_homo\_2d} itself, as in Section~\ref{sec:solid-beam};
  regenerated on each run \\
\texttt{RHC\_SW\_2UC\_45\_beam\_Timo.out} & the bare $6\times6$ written by the
  script through \texttt{np.savetxt}, one commented header line --- convenient
  to \texttt{np.loadtxt} \\
\texttt{RHC\_SW\_2UC\_45\_dehom.txt} & tab-separated Gauss table:
  \texttt{GP\_Index}, \texttt{Coord\_X}, \texttt{Coord\_Y}, six \texttt{Eps\_*},
  six \texttt{Sig\_*}; one header line plus 19\,920 rows (6640 triangles times
  three quadrature points) \\
\texttt{RHC\_SW\_2UC\_45\_dehom.vtk} & the same cloud as an ASCII unstructured
  grid of vertex cells, each strain and stress component a \texttt{SCALARS}
  field --- open in ParaView \\
\texttt{RHC\_SW\_2UC\_45\_dehom.SM} / \texttt{.EM} / \texttt{.U} & the OpenSG
  dehom files: stress, strain and fluctuation displacement in the
  \texttt{rm\_plate\_1D} layout, \texttt{x y z} columns then components printed
  in the order \texttt{11 22 33 12 13 23}, with the SG coordinates in the last
  \texttt{n\_sg} of the \texttt{x, y, z} slots; regenerated on each run \\
\texttt{RHC\_SW\_2UC\_45\_mesh.png} & the mesh figure \\
\bottomrule
\end{tabular}
\caption{Outputs of \texttt{beam\_dehom\_sg.py}.}
\label{tab:solid-ex8out}
\end{table}

The \texttt{.txt}/\texttt{.vtk} pair uses the SwiftComp component order
\texttt{xx yy zz yz xz xy}; the \texttt{.SM}/\texttt{.EM}/\texttt{.U} files
reorder to \texttt{11 22 33 12 13 23} on write only. Units follow the SG input
system throughout, MPa here.

\subsection{What the numbers mean}

The shipped table is the run of the driver above. Its component maxima are:

\begin{table}[htbp]
\centering
\small
\begin{tabular}{lcc}
\toprule
component & $\max\lvert\sigma\rvert$ (MPa) & $\max\lvert\varepsilon\rvert$ \\
\midrule
$xx$ & $1.529334\times10^{4}$ & $2.339375\times10^{-1}$ \\
$yy$ & $3.460157\times10^{3}$ & $2.265195\times10^{-1}$ \\
$zz$ & $5.539976\times10^{2}$ & $8.787251\times10^{-2}$ \\
$yz$ & $6.034080\times10^{2}$ & $1.344338\times10^{-1}$ \\
$xz$ & $1.176515\times10^{3}$ & $1.704428\times10^{-1}$ \\
$xy$ & $2.959409\times10^{3}$ & $1.028684\times10^{-1}$ \\
\bottomrule
\end{tabular}
\caption{Component maxima over the 19\,920 Gauss points of the shipped beam
recovery.}
\label{tab:solid-ex8max}
\end{table}

The axial field is the one to sanity-check first, because beam kinematics
predict it directly: $\epsilon_{xx} \approx \bar{\epsilon}_{11}
+ \kappa_2 y_3$ plus a warping correction. The first row of the shipped
\texttt{.txt} sits at $y_3 = 23.29375$~mm and reports
$\epsilon_{xx} = 2.339375\times10^{-1}$, against
$0.001 + 0.01 \times 23.29375 = 0.2339375$: an exact match, and the same value
is the table maximum, as it must be at the extreme fibre.

The useful content of the recovery is everything \emph{else}: the transverse and
shear components, which exist purely because the cross-section is heterogeneous
and warps, and which no beam theory can supply. The shipped driver is a
demonstration magnitude, not a service load; scale it to your own macro state
before reading the stresses as failure indices.

\begin{figure}[htbp]
\centering
\includegraphics[width=0.62\linewidth]{\FIGDIR/pv_beam_dehom_sxx.png}
\caption{ParaView render of the recovered field on the SG mesh, produced from
the \texttt{RHC\_SW\_2UC\_45\_dehom.vtk} written by this example, using
\texttt{docs/latex/render\_pv.py}. The array is \texttt{Sig\_xx}, the axial
stress $\sigma_{xx}$ in MPa (units follow the SG input, millimetres and
megapascals), on a symmetric colour scale of $\pm1.53\times10^{4}$~MPa --- the
$xx$ row of Table~\ref{tab:solid-ex8max}. Each of the 19\,920 points is one
Gauss point carried by its own vertex cell and coloured flat, so nothing is
averaged across an element boundary. The gene is the same honeycomb as
Figure~\ref{fig:solid-pvplatedehom} and the pattern is again linear in the
through-thickness coordinate, but the peak is an order of magnitude lower,
tracking this driver's extreme-fibre strain of $0.2339$ against the $2.3294$ of
the plate recovery.}
\label{fig:solid-pvbeamdehom}
\end{figure}

\section{Files at a glance}
\label{sec:solid-glance}

\begingroup\footnotesize
\begin{longtable}{cp{4.6cm}p{8.4cm}}
\caption{Every \texttt{msg-solid} capability in one table. \texttt{<name>} is
the SG file's own basename --- the \texttt{name} variable where a runner still
sets one --- and \texttt{<ff stem>} the basename of the \texttt{.ff} table.
\texttt{opensg} is the unified entry point; \texttt{opensg\_solid} and
\texttt{python -m opensg\_solid} are the same run.}
\label{tab:solid-glance} \\
\toprule
ex & you edit, then you run & it writes \\
\midrule
\endfirsthead
\toprule
ex & you edit, then you run & it writes \\
\midrule
\endhead
\bottomrule
\endfoot
1 & \texttt{layup\_db.yaml}; \texttt{python 1d\_sg.py}
  & \texttt{1dsg.yaml}, \texttt{1dsg.png},
    \texttt{layup\_db\_plate\_homo.out} \\
2 & the \texttt{USER INPUT} block; \texttt{python rm\_dehom.py}
  & \texttt{1dsg\_dehom.SM}, \texttt{.EM}, \texttt{.U}, \texttt{.vtk} \\
2 & a \texttt{.ff} table and the \texttt{dehom:} block;
    \texttt{python rm\_dehom\_ff.py ex5\_shell\_S8R\_field.ff}
  & \texttt{<ff stem>\_ff.SM}, \texttt{.EM}, \texttt{.U},
    \texttt{<ff stem>\_ff\_gauss.vtk} \\
3 & the yaml header (\texttt{n\_model}, \texttt{refined}, \texttt{analysis});
    \texttt{opensg <name>.yaml}
  & \texttt{<name>.out}, \texttt{<name>\_mesh.png} \\
4 & the above plus \texttt{epsilon\_bar:};
    \texttt{opensg <name>.yaml D}
  & \texttt{<name>.out},
    \texttt{<name>\_dehom.txt}, \texttt{.vtk}, \texttt{.SM}, \texttt{.EM},
    \texttt{.U}, \texttt{<name>\_mesh.png};
    \texttt{python plate\_dehom\_2dsg.py} adds
    \texttt{<name>\_plate\_ABD.out} \\
5 & \texttt{name} and the \texttt{paper} dictionary;
    \texttt{python solid\_homo\_2dsg.py}
  & \texttt{<name>\_Deff.out},
    \texttt{<name>\_constants\_vs\_paper.dat}, \texttt{<name>\_mesh.png} \\
6 & the \texttt{.sc} path, $E$ and $\nu$; \texttt{python run\_sample\_1.py}
    (the homogenization alone is \texttt{opensg Sample\_1.yaml})
  & \texttt{Sample\_1.out}, \texttt{Sample\_1\_msg\_solid.out},
    \texttt{Sample\_1\_compare.dat}, \texttt{Sample\_1\_mesh.png} \\
6 & nothing; \texttt{python compare\_boundary\_modes.py}
  & \texttt{boundary\_compare.dat}, \texttt{Sample\_1\_periodic.out},
    \texttt{Sample\_1\_aperiodic.out} and the Sample 2 pair \\
7 & the yaml header (\texttt{n\_model: 1}, \texttt{refined: 1});
    \texttt{opensg <name>.yaml}
  & \texttt{<name>.out}, \texttt{<name>\_mesh.png} \\
8 & the above plus \texttt{epsilon\_bar:};
    \texttt{opensg <name>.yaml D}
  & \texttt{<name>.out},
    \texttt{<name>\_dehom.txt}, \texttt{.vtk}, \texttt{.SM}, \texttt{.EM},
    \texttt{.U}, \texttt{<name>\_mesh.png}; \texttt{python beam\_dehom\_sg.py}
    adds \texttt{<name>\_beam\_Timo.out} \\
\end{longtable}
\endgroup

\chapter{Full Blade Pipeline: windIO to BeamDyn and Back}
\label{ch:windio}

A wind-turbine blade is designed as a three-dimensional laminated shell but
analysed, aeroelastically, as a one-dimensional beam. This chapter walks the
whole round trip on the IEA-22 blade: a windIO blade description becomes one
1-D shell structure gene per span station, each station is homogenized to a
Timoshenko $6\times6$, the spanwise table of those matrices is what a 1-D beam
solver needs, the beam run returns sectional force and moment resultants
station by station, and those resultants drive the dehomogenization back down
to three-dimensional stress in the wall.

Three separate tools are involved, and the boundary between them is the single
most important thing to understand before running anything. OpenSG never runs
the blade dynamics. It produces the sectional constitutive law that the beam
solver needs, and it consumes the sectional loads that the beam solver returns.

\begin{table}[htbp]
\centering
\small
\caption{The five stages of the round trip, and which tool owns each one.}
\label{tab:windio:stages}
\begin{tabular}{clll}
\toprule
Stage & What happens & Tool & Driver \\
\midrule
1 & windIO blade $\rightarrow$ per-station 1-D shell SG YAML & OpenSG\_io & \texttt{scripts/opensg\_io.py} \\
2 & 1-D shell SG $\rightarrow$ Timoshenko $6\times6$ & OpenSG-2.0 & \texttt{1\_homo\_beam\_props.py} \\
  &  &  & \texttt{4\_sweep\_51\_stations.py} \\
3 & spanwise $6\times6$ table $\rightarrow$ blade dynamics & BeamDyn / OpenFAST & run by you \\
4 & sectional loads returned as a table & BeamDyn / OpenFAST & \texttt{ff51\_rmc\_reform.dat} \\
5 & loads $+\ 6\times6$ $\rightarrow$ 3-D wall stress & OpenSG-2.0 & \texttt{2\_dehom\_v2\_vabs.py} \\
  &  &  & \texttt{4\_sweep\_51\_stations.py} \\
  & correctness check on transverse shear & OpenSG-2.0 & \texttt{3\_q\_consistency\_check.py} \\
\bottomrule
\end{tabular}
\end{table}

\section{The worked example}
\label{sec:windio:example}

Everything in this chapter runs in the directory
\texttt{examples/OpenSG\_shell/IEA\_blade\_beam/}. Run the scripts from inside
that directory: the package holds no path assumptions, and the example owns
every path it uses.

The three single-station scripts all analyse the same section,
\texttt{iea\_s10}, the IEA-22 station at $r/R = 0.20$. Its 1-D shell SG,
\texttt{iea\_s10\_shell.yaml}, is a ring of \textbf{310 wall elements} carrying
\textbf{6 layups} (\texttt{layup\_0} through \texttt{layup\_5}) built from
\textbf{6 materials} (\texttt{gelcoat}, \texttt{glass\_triax},
\texttt{glass\_uniax}, \texttt{medium\_density\_foam}, \texttt{carbon\_uniax},
\texttt{glass\_biax}). The element count is confirmed by the recovery output:
\texttt{iea\_s10\_elem\_plate\_forces.dat} has two header lines and 310 data
rows, one per ring element. The fourth script,
\texttt{4\_sweep\_51\_stations.py}, repeats the same computation for all 51
stations of the blade.

\section{Stage 1: windIO blade to per-station 1-D shell SG YAMLs}
\label{sec:windio:stage1}

\subsection*{What it does}

\texttt{OpenSG\_io} is the input layer, and it is a \emph{separate repository}:
\url{https://github.com/bagla0/OpenSG_io}, documented at
\url{https://bagla0.github.io/OpenSG_io/}. It reads a windIO blade description
-- auto-detecting windIO v1 and v2 -- interpolates the outer shape, the layup
stacks and the material database to a requested span fraction, and writes the
cross-section in the two forms an MSG analysis can use:

\begin{itemize}
  \item the \textbf{1-D shell SG YAML} that \texttt{opensg\_shell} reads
        directly, which is what stages 2 and 5 of this chapter consume; and
  \item a \textbf{PreVABS XML} that meshes to a 2-D solid cross-section, which
        is the route that produces the independent VABS reference used to
        validate the recovery in stage 5.
\end{itemize}

It also reads PreVABS XML cross-sections and OpenFAST ElastoDyn/BeamDyn blade
data, so published beam properties can be pulled in as a validation reference.
Its own README is explicit that it does not run windIO or OpenFAST; it reads
those formats so you do not have to hand-build the structure gene.

\subsection*{Installing it}

\begin{lstlisting}[style=shellst]
git clone --recurse-submodules https://github.com/bagla0/OpenSG_io
cd OpenSG_io
pip install -r requirements.txt
python scripts/fetch_prevabs.py
\end{lstlisting}

The \texttt{requirements.txt} in that repository lists \texttt{numpy},
\texttt{scipy}, \texttt{pyyaml} and \texttt{meshio} as the core dependencies,
plus \texttt{gmsh}, \texttt{tetgen} and \texttt{pyvista} for the mesh
generators. A windIO \emph{v2} input file additionally needs the
\texttt{windIO} package installed; v1 files do not. The
\texttt{fetch\_prevabs.py} script downloads the PreVABS release binary for your
operating system, which is required only for the 2-D solid route -- the 1-D
shell YAML needed for stages 2 through 5 is written by pure Python.

\subsection*{The command line}

One CLI handles any of the supported inputs. It writes a set of stations in a
single call:

\begin{lstlisting}[style=shellst]
python scripts/opensg_io.py BAR_URC.yaml out --name bar --stations 0.3 0.5 0.7 --solid
\end{lstlisting}

Here \texttt{BAR\_URC.yaml} is the windIO blade, \texttt{out} is the output
directory, \texttt{-{}-name bar} is the file-name stem, \texttt{-{}-stations}
takes the span fractions $r/R$, and \texttt{-{}-solid} additionally emits the
PreVABS XML and runs PreVABS to build the 2-D solid section. Omit
\texttt{-{}-solid} if you only want the 1-D shell YAMLs.

\subsection*{The Python API}

The API is what you script a 51-station blade with, and it is what generated
the station files used in the rest of this chapter:

\begin{lstlisting}[style=opensg,caption={Building one station's 1-D shell SG and its PreVABS twin.}]
from opensg_io import load_blade, build_cross_section, emit_opensg_yaml, emit_prevabs

blade = load_blade("IEA-22-280-RWT.yaml")      # windIO v1 or v2, auto-detected
cs = build_cross_section(blade, r=0.5)         # one span fraction
emit_opensg_yaml(cs, "shell_r050.yaml")        # the 1-D shell SG for opensg_shell
emit_prevabs(cs, "prevabs_r050", name="r050")  # PreVABS XML -> 2-D solid route
\end{lstlisting}

Loop over your station list and you have the \texttt{iea\_s00\_shell.yaml}
through \texttt{iea\_s50\_shell.yaml} directory that stage 2 points at.

\subsection*{What the 1-D shell SG YAML contains}

\texttt{iea\_s10\_shell.yaml} is the reference for the format. Its top-level
keys, in the order they appear in the file, are listed in
Table~\ref{tab:windio:yamlkeys}.

\begin{table}[htbp]
\centering
\small
\caption{Top-level keys of the 1-D shell SG YAML, in file order.}
\label{tab:windio:yamlkeys}
\begin{tabular}{ll}
\toprule
Key & Content \\
\midrule
\texttt{nodes} & contour node coordinates \texttt{[y1, y2, y3]} in metres, one per line \\
\texttt{elements} & two-node line segments along the section contour \\
\texttt{sets} & \texttt{element} sets, one per layup, holding the element labels \\
\texttt{sections} & \texttt{type: shell}, \texttt{elementSet: layup\_k}, and the \texttt{layup} as \\
  & \texttt{[material, thickness (m), angle (deg)]} triples, outermost ply first \\
\texttt{elementOrientations} & one $3\times3$ frame per element, flattened; the first row is $\mathbf{e}_1$, \\
  & out of the section plane, along the beam axis \\
\texttt{materials} & \texttt{density} in kg/m$^3$ and orthotropic \texttt{elastic: \{E, G, nu\}} triples in Pa \\
\texttt{reference} & which surface the section geometry refers to \\
\bottomrule
\end{tabular}
\end{table}

A concrete section block from that file, the spar-cap laminate:

\begin{lstlisting}[style=shellst]
sections:
- type: shell
  elementSet: layup_0
  layup:
  - - gelcoat
    - 0.0005
    - 0.0
  - - glass_triax
    - 0.00300294
    - 0.0
  - - glass_uniax
    - 0.02609467
    - 0.0
  - - glass_triax
    - 0.00300294
    - 0.0
\end{lstlisting}

The \texttt{reference} field is load-bearing and is the single source of truth
for the whole chapter. \texttt{build\_rm\_bundle} reads it when no override is
passed, and it fixes the through-thickness offset \texttt{frac} used by the
ring laminate law, the plate SG's reference plane, the emitted ABD and the
recovery depth conversion alike. The IEA stations all carry
\texttt{reference: center}.

\begin{table}[htbp]
\centering
\small
\caption{The \texttt{reference} field and the through-thickness fraction it sets.}
\label{tab:windio:reference}
\begin{tabular}{lll}
\toprule
\texttt{reference} & \texttt{frac} & Meaning \\
\midrule
\texttt{center} & 0.5 & mid-surface (what the IEA stations use) \\
\texttt{oml} & 0.0 & outer mould line \\
\texttt{oml\_flip} & 1.0 & outer mould line, reversed stacking \\
\texttt{iml} & 1.0 & inner mould line \\
\bottomrule
\end{tabular}
\end{table}

\section{Stage 2: each station to a Timoshenko $6\times6$}
\label{sec:windio:stage2}

\subsection*{What it computes}

The RM shell ring homogenization: a six-degree-of-freedom-per-node shell
element with an element-wise drilling Lagrange multiplier, $\gamma_{23}$ tied
by MITC, and the MSG (Yu 2002 least-squares) wall transverse-shear law
$\mathbf{G}$. The ring warping problem is solved for the four Euler--Bernoulli
modes $\mathbf{V}_0$ and the Timoshenko correction $\mathbf{V}_1$, and the
energy is condensed to the beam $6\times6$ in the VABS strain order
\[
  \boldsymbol{\varepsilon} =
  \begin{bmatrix} \varepsilon_{11} & \gamma_{12} & \gamma_{13} &
                  \kappa_1 & \kappa_2 & \kappa_3 \end{bmatrix}^{\mathsf T},
\]
whose diagonal is $[\,EA\;\;GA_2\;\;GA_3\;\;GJ\;\;EI_2\;\;EI_3\,]$.

\subsection*{What you edit}

One line at the top of \texttt{1\_homo\_beam\_props.py}:

\begin{lstlisting}[style=opensg]
############### User Input #################################
name = "iea_s10"                  # IEA-22 station at r/R = 0.20
############################################################
\end{lstlisting}

\subsection*{The command}

\begin{lstlisting}[style=shellst]
python 1_homo_beam_props.py
\end{lstlisting}

\subsection*{The API, if you drive it yourself}

\begin{lstlisting}[style=opensg,caption={The whole homogenization is one call.}]
import numpy as np
from opensg_shell import build_rm_bundle

B = build_rm_bundle("iea_s10_shell.yaml")
C6 = np.asarray(B["Timo"])        # (6,6) Timoshenko stiffness
print(B["ref"], B["g_source"])    # 'center'  'msg'
\end{lstlisting}

The full signature is
\texttt{build\_rm\_bundle(shell\_yaml, ref=None, shear="mitc4\_g23",
g\_source="msg")}. Passing \texttt{ref=None} means the YAML's own
\texttt{reference} field decides, which is why homogenization and
dehomogenization cannot silently disagree; pass an explicit \texttt{ref} only
to override. \texttt{g\_source} chooses the wall transverse-shear law,
\texttt{"msg"} for the Yu 2002 least-squares projection or \texttt{"whitney"}
for the complementary-energy shear flow.

What comes back is a bundle, not just a matrix. It carries \texttt{Timo}, the
warping modes \texttt{V0} and \texttt{V1}, the ring geometry
(\texttt{corners}, \texttt{red\_cells}, \texttt{rx3}, \texttt{re3},
\texttt{k22}, \texttt{ax}, \texttt{cross}, \texttt{strip}),
\texttt{layup\_per\_elem}, the geometry-free \texttt{layup\_db} and
\texttt{material\_db}, and \texttt{frac}, \texttt{ref}, \texttt{g\_source}.
Stage 5 consumes every one of those entries, so do not throw the bundle away
after reading the matrix out of it.

\subsection*{What comes out}

\begin{table}[htbp]
\centering
\small
\caption{Files written by a single-station homogenization run.}
\label{tab:windio:stage2out}
\begin{tabular}{lll}
\toprule
File & Written by & Content \\
\midrule
\texttt{iea\_s10\_shell\_Timo.out} & \texttt{build\_rm\_bundle} & the timed SwiftComp \texttt{.K} layout: the banner \\
  &  & \texttt{OpenSG msg-shell beam model [ext sh2 sh3} \\
  &  & \texttt{twist bend2 bend3]}, \texttt{The Effective Timoshenko} \\
  &  & \texttt{Stiffness Matrix}, its compliance, and a \\
  &  & \texttt{Time taken} footer \\
\texttt{iea\_s10\_shell\_ABDG.out} & \texttt{build\_rm\_bundle} & one block per section: \texttt{The Effective Reissner-} \\
  &  & \texttt{Mindlin Plate Stiffness Matrix}, the $8\times8$ ABD \\
  &  & plus $\mathbf{G}$, and its compliance \\
\texttt{abd/iea\_s10\_shell\_abd.yaml} & \texttt{build\_rm\_bundle} & the per-station ABD YAML, emitted once and \\
  &  & cached; carries \texttt{mass\_per\_area} per layup \\
\texttt{iea\_s10\_beam\_Timo.out} & the script & the bare $6\times6$ as text, with the strain order \\
  &  & and the reference recorded in the header \\
\texttt{iea\_s10\_ring\_mesh.png} & the script & the real computed ring, elements coloured by layup \\
\bottomrule
\end{tabular}
\end{table}

Two details about the ABD YAML are worth knowing before you reuse it. It is
written by \texttt{emit\_station\_abd} with \texttt{ref="mid"} when the station
is \texttt{center}-referenced, so its \texttt{reference: mid} field means the
same surface as the ring's \texttt{center}. And it is written with the
emitter's default transverse-shear source, so the file records
\texttt{g\_source: whitney}: the $8\times8$ block in that YAML uses the Whitney
shear-flow $\mathbf{G}$, whereas the ring homogenization itself used the MSG
$\mathbf{G}$. The membrane and bending blocks are identical either way; only
the two transverse-shear rows differ. For \texttt{iea\_s10} the file reports
\texttt{n\_layups: 6}, and for \texttt{layup\_0} a total \texttt{thickness} of
$0.03260055$~m and a \texttt{mass\_per\_area} of $62.89$~kg/m$^2$.

\subsection*{What the numbers mean}

Table~\ref{tab:windio:c6} is the content of \texttt{iea\_s10\_beam\_Timo.out},
in SI units. The YAML is in metres and pascals, so $EA$, $GA_2$ and $GA_3$ are
in N and $GJ$, $EI_2$, $EI_3$ in N$\cdot$m$^2$.

\begin{table}[htbp]
\centering
\footnotesize
\caption{The Timoshenko $6\times6$ of IEA-22 station \texttt{iea\_s10} at $r/R = 0.20$,
         in the VABS strain order, from \texttt{iea\_s10\_beam\_Timo.out}.}
\label{tab:windio:c6}
\setlength{\tabcolsep}{4pt}
\begin{tabular}{lrrrrrr}
\toprule
 & $\varepsilon_{11}$ & $\gamma_{12}$ & $\gamma_{13}$ & $\kappa_1$ & $\kappa_2$ & $\kappa_3$ \\
\midrule
$\varepsilon_{11}$ & 2.76606e+10 & $-$1.72307e$-$05 & 4.66505e$-$06 & 5.40202e$-$06 & 1.66864e+09 & $-$2.07730e+09 \\
$\gamma_{12}$ & $-$1.72307e$-$05 & 7.18613e+08 & 5.75204e+07 & $-$1.01920e+08 & 2.30800e$-$06 & 5.27387e$-$06 \\
$\gamma_{13}$ & 4.66505e$-$06 & 5.75204e+07 & 4.21744e+08 & $-$1.62178e+08 & 8.48168e$-$06 & 2.79314e$-$07 \\
$\kappa_1$ & 5.40202e$-$06 & $-$1.01920e+08 & $-$1.62178e+08 & 2.43757e+09 & 1.64496e$-$06 & $-$8.11406e$-$06 \\
$\kappa_2$ & 1.66864e+09 & 2.30800e$-$06 & 8.48168e$-$06 & 1.64496e$-$06 & 3.51439e+10 & $-$5.55712e+09 \\
$\kappa_3$ & $-$2.07730e+09 & 5.27387e$-$06 & 2.79314e$-$07 & $-$8.11406e$-$06 & $-$5.55712e+09 & 6.81297e+10 \\
\bottomrule
\end{tabular}
\end{table}

Read the matrix as physics rather than as six numbers on a diagonal. The
extension--bending entries $1.67\times10^{9}$ and $-2.08\times10^{9}$ say that
the section's tension centre does not sit on the reference axis. The
bending--bending coupling $-5.56\times10^{9}$ couples flap and edge bending, so
the principal bending axes are not the reference axes either. The shear--twist
entries $-1.02\times10^{8}$ and $-1.62\times10^{8}$ locate the shear centre
away from the reference axis as well. The entries printed at $10^{-5}$ against
diagonals of $10^{10}$ are numerically zero: extension does not couple to
transverse shear or to twist in this section, and the solver reproduces that
null result to fifteen orders of magnitude, which is a useful sanity check that
the ring assembled and condensed correctly.

\subsection*{The spanwise table: all 51 stations}

\texttt{4\_sweep\_51\_stations.py} repeats the homogenization for every station
of the blade and, in the same loop, runs the first-order recovery of stage 5.
What you edit is the block at the top:

\begin{lstlisting}[style=opensg]
############### User Input #################################
yaml_dir = os.path.expanduser(
    "~/OpenSG-TW-claude/examples/data/iea_all_stations/shell51/1d_yaml")
ff_file = "ff51_rmc_reform.dat"
n_depth = 9                  # through-thickness recovery points per element
stations = range(51)         # s00 .. s50
############################################################
\end{lstlisting}

\texttt{yaml\_dir} is where \texttt{iea\_s00\_shell.yaml} through
\texttt{iea\_s50\_shell.yaml} live, which is the OpenSG\_io output of stage 1.
Point it at your own directory. The script builds each path as
\texttt{iea\_s\%02d\_shell.yaml} against that directory, and reads the load
table \texttt{ff\_file} from the current directory.

\begin{lstlisting}[style=shellst]
python 4_sweep_51_stations.py
\end{lstlisting}

A station whose YAML is missing prints
\texttt{s\%02d: yaml missing -{}- skipped} and is passed over, and any station
that raises is caught and reported as \texttt{s\%02d: FAILED} so that one bad
section cannot kill a sweep that has already spent minutes on the others. In
the shipped run, \textbf{49 of the 51 stations completed}: both
\texttt{iea51\_beam\_props.dat} and \texttt{iea51\_peak\_stress.dat} carry
$\eta = 0.00$ through $0.94$ in steps of $0.02$ plus $\eta = 1.00$, with
$\eta = 0.96$ and $\eta = 0.98$ absent.

One indexing rule matters here. The sweep takes station $i$'s load vector as
row $i$ of the FF table, \texttt{FF = ffall[i, 1:]}, and its span fraction as
\texttt{eta = ffall[i, 0]}. The row order of the load table must therefore match
the station numbering of the YAML files. In the shipped data, station
\texttt{s10} is row 10 and $\eta = 0.20$.

\texttt{iea51\_beam\_props.dat} holds one row per completed station: $\eta$
followed by the six diagonal terms, written with the header
\texttt{IEA-22 spanwise Timoshenko beam properties (RM shell ring)}. Selected
rows, verbatim from the file:

\begin{table}[htbp]
\centering
\footnotesize
\caption{Selected rows of \texttt{iea51\_beam\_props.dat}. Stiffnesses in N
         ($EA$, $GA_2$, $GA_3$) and N$\cdot$m$^2$ ($GJ$, $EI_2$, $EI_3$).}
\label{tab:windio:props}
\setlength{\tabcolsep}{4.5pt}
\begin{tabular}{lrrrrrr}
\toprule
$\eta = r/R$ & $EA$ & $GA_2$ & $GA_3$ & $GJ$ & $EI_2$ & $EI_3$ \\
\midrule
0.00 & 6.44275512e+10 & 7.96306621e+09 & 8.36135930e+09 & 1.27413308e+11 & 2.97295270e+11 & 2.19213722e+11 \\
0.10 & 2.89212767e+10 & 1.11237808e+09 & 1.04545606e+09 & 1.14245218e+10 & 8.40070906e+10 & 1.21502270e+11 \\
0.20 & 2.76605989e+10 & 7.18613262e+08 & 4.21744284e+08 & 2.43756534e+09 & 3.51438873e+10 & 6.81297238e+10 \\
0.50 & 2.15596413e+10 & 5.34263266e+08 & 2.02128533e+08 & 4.95011412e+08 & 7.26524060e+09 & 1.82791086e+10 \\
0.80 & 9.28038409e+09 & 3.90243799e+08 & 7.82870454e+07 & 6.27349440e+07 & 5.76401363e+08 & 2.07536066e+09 \\
0.94 & 2.64020156e+09 & 1.80834082e+08 & 3.43479097e+07 & 1.35795150e+07 & 6.21667876e+07 & 4.27156667e+08 \\
1.00 & 2.65581077e+08 & 1.56743220e+07 & 1.88135914e+06 & 5.15174420e+04 & 3.04364786e+05 & 1.60428557e+06 \\
\bottomrule
\end{tabular}
\end{table}

The $\eta = 0.20$ row is digit-for-digit the diagonal of the standalone
\texttt{iea\_s10} run of Table~\ref{tab:windio:c6}: $EA$ is
\texttt{2.76605989e+10} in both, $GJ$ is \texttt{2.43756534e+09} in both, and so
on for all six. The sweep and the single-station script are the same
computation, which is the cheapest available check that your sweep is wired
correctly -- run \texttt{1\_homo\_beam\_props.py} on one station and match its
row.

\begin{figure}[htbp]
\centering
\includegraphics[width=0.95\textwidth]{\FIGDIR/iea51_beam_props.png}
\caption{Spanwise Timoshenko diagonal of the IEA-22 blade: the six terms
         $EA$, $GA_2$, $GA_3$, $GJ$, $EI_2$, $EI_3$ on log axes against $r/R$,
         one panel each, produced by \texttt{4\_sweep\_51\_stations.py}.}
\label{fig:windio:beamprops}
\end{figure}

Figure~\ref{fig:windio:beamprops} plots those six columns. Three regimes are
visible. The root collapses fast as the circular root transitions into an
aerofoil: $GJ$ falls from $1.274\times10^{11}$ at $\eta = 0$ to
$1.142\times10^{10}$ N$\cdot$m$^2$ at $\eta = 0.10$, an order of magnitude in
the first tenth of span, and $GA_2$ falls from $7.96\times10^{9}$ to
$1.11\times10^{9}$ N over the same interval. The mid-span then decays smoothly
as the section thins. The tip drops off a cliff. Across the whole blade $GJ$
goes from $1.27\times10^{11}$ to $5.15\times10^{4}$ N$\cdot$m$^2$, six and a
half decades, which is why the figure is drawn on log axes and why a beam
solver must be given the table rather than an average.

\section{Stage 3: handing the table to BeamDyn}
\label{sec:windio:stage3}

Be explicit about what changes hands here. \textbf{OpenSG produces the sectional
$6\times6$ stiffness matrix, and nothing else.} BeamDyn, the geometrically exact
beam module of OpenFAST, is a separate program that you install, configure and
run yourself. OpenSG-2.0 ships no BeamDyn input file and no BeamDyn results, and
this chapter does not attempt to document BeamDyn's input format: for the
meaning, ordering and units of every field in a BeamDyn blade input file, and
for the choice of quadrature and the number of input stations, consult the
OpenFAST BeamDyn documentation, which is the authority on that file.

\subsection*{The bridge that OpenSG\_io provides}

\texttt{OpenSG\_io} contains an \texttt{openfast\_io} module whose job is
exactly this hand-off. Two functions write the OpenFAST side:

\begin{lstlisting}[style=opensg,caption={Writing a BeamDyn blade input from a spanwise table of OpenSG 6x6 matrices.}]
import numpy as np
from opensg_io.openfast_io import timo_to_beamdyn, write_beamdyn_blade

# etas: (n,) span fractions; Ks: list of n OpenSG Timoshenko 6x6 matrices
Kbd = timo_to_beamdyn(Ks[0])                 # one matrix, reordered
write_beamdyn_blade(etas, Ks, M_list=None,
                    out_path="blade_BeamDyn_Blade.dat")
\end{lstlisting}

\texttt{timo\_to\_beamdyn(Ktimo)} reorders an OpenSG Timoshenko $6\times6$,
which is stored in the VABS order $[EA, GA_2, GA_3, GJ, EI_2, EI_3]$, into the
BeamDyn degree-of-freedom order. It applies a symmetric permutation
\texttt{Ktimo[np.ix\_(q, q)]} with \texttt{q = [1, 2, 0, 4, 5, 3]}, which is the
mapping in Table~\ref{tab:windio:bdorder}. Because the permutation is applied to
both rows and columns, every off-diagonal coupling is carried across intact.

\begin{table}[htbp]
\centering
\small
\caption{The reordering performed by \texttt{timo\_to\_beamdyn}: which OpenSG
         diagonal term lands in which BeamDyn slot.}
\label{tab:windio:bdorder}
\begin{tabular}{clc}
\toprule
BeamDyn slot & Meaning & From OpenSG term \\
\midrule
1 & shear\_x & $GA_2$ \\
2 & shear\_y & $GA_3$ \\
3 & axial & $EA$ \\
4 & bend\_x & $EI_2$ \\
5 & bend\_y & $EI_3$ \\
6 & torsion & $GJ$ \\
\bottomrule
\end{tabular}
\end{table}

\texttt{write\_beamdyn\_blade(etas, K\_timo\_list, M\_list=None,
out\_path="blade\_BeamDyn\_Blade.dat", damp=(0.003, 0.002, 0.002, 0.002, 0.003,
0.002))} takes the span fractions and the list of OpenSG matrices, reorders each
one internally with \texttt{timo\_to\_beamdyn}, and writes a BeamDyn blade input
with a \texttt{station\_total} equal to the number of entries you supplied,
\texttt{damp\_type} set to $1$, the six damping coefficients from
\texttt{damp}, and then a stiffness block followed by a mass block for each
station. Verify the resulting file against the BeamDyn documentation for the
OpenFAST version you are running before trusting it.

\subsection*{The mass matrix is yours}

\texttt{write\_beamdyn\_blade} takes \texttt{M\_list} as an \emph{optional}
argument and \textbf{writes a $6\times6$ block of zeros for every station when it
is omitted}. That is a placeholder, not a mass model. Nothing in this pipeline
computes a sectional mass matrix: the only mass quantity OpenSG-2.0 produces
here is \texttt{mass\_per\_area} per laminate in the emitted ABD YAML
(Section~\ref{sec:windio:stage2}). Building the sectional inertia matrix,
assembling the complete BeamDyn input, and running the aeroelastic simulation
are all yours, outside OpenSG.

\subsection*{Use the full matrix, not the table's diagonal}

\texttt{iea51\_beam\_props.dat} records only the six diagonal terms, for
plotting and inspection. The couplings displayed in Table~\ref{tab:windio:c6}
are precisely what a composite blade beam model exists to carry, so feed the
beam solver the full $6\times6$ of each station -- either from that station's
\texttt{<station>\_Timo.out} or, more conveniently, straight from
\texttt{B["Timo"]} in your own driver script, collected into the
\texttt{K\_timo\_list} above.

\subsection*{Reading OpenFAST data back as a reference}

The same module reads in the other direction, which is how you validate an
OpenSG table against a published blade. \texttt{read\_beamdyn\_blade(path)}
parses a BeamDyn blade input and returns \texttt{(etas, K\_list, M\_list)} with
each matrix in BeamDyn order; \texttt{beamdyn\_to\_timo(Kbd)} reorders one back
into the OpenSG order so it can be compared term by term.
\texttt{read\_elastodyn\_blade(path)} parses an ElastoDyn blade file and returns
arrays of \texttt{BlFract}, \texttt{PitchAxis}, \texttt{StrcTwst} in degrees,
\texttt{BMassDen} in kg/m, and \texttt{FlpStff} and \texttt{EdgStff} in
N$\cdot$m$^2$, which map to OpenSG's $EI_2$ (flap) and $EI_3$ (edge);
\texttt{elastodyn\_at(ed, r)} interpolates those to any span fraction.

\section{Stage 4: the loads come back as the FF table}
\label{sec:windio:stage4}

A BeamDyn or GEBT run returns, at every station, the sectional force and moment
resultants the blade actually carries in a given load case. That table is the
input to recovery. \texttt{ff51\_rmc\_reform.dat} is such a table. It has one
comment header line and 51 data rows, one per station. The header and the first
data rows, verbatim:

\begin{lstlisting}[style=shellst]
# eta F1 F2 F3 M1 M2 M3 (RM CENTER-ref BeamDyn, VABS order)
0.00000000e+00 -6.44900592e-06 -2.39422988e-08 2.23607775e+06 1.21251850e+06 -1.12030544e+08 -4.11329552e-07
2.00000000e-02 2.07495361e+03 7.56198425e+01 2.15497000e+06 1.21427150e+06 -1.05899552e+08 2.54692651e+03
4.00000000e-02 4.74173926e+03 1.68900665e+02 2.07209938e+06 1.23365012e+06 -9.99422000e+07 5.32342041e+03
\end{lstlisting}

Column 1 is $\eta = r/R$, running from $0.00$ to $1.00$ in steps of $0.02$.
Columns 2 through 4 are the forces $F_1$, $F_2$, $F_3$ in N; columns 5 through 7
are the moments $M_1$, $M_2$, $M_3$ in N$\cdot$m. The row used by the
single-station scripts is row 10:

\begin{lstlisting}[style=shellst]
2.00000000e-01 4.19268867e+04 8.26739746e+03 1.46733512e+06 1.21914500e+06 -6.01849720e+07 3.04264594e+05
\end{lstlisting}

The frame and the component order stated in that header are not decoration. The
resultants must be expressed about the same reference axis and in the same
component order as the $6\times6$ that produced them, or the round trip does not
close and the recovered stress is silently wrong. Here both are \emph{RM
center-ref} and VABS order, matching \texttt{reference: center} in the station
YAMLs. Note also the $\eta = 0$ row: its axial and in-plane entries print as
$-6.4\times10^{-6}$ and $-2.4\times10^{-8}$ N, a numerically zero root row
rather than a physical load.

A converter turns the table into an annotated YAML and pulls one station out of
it:

\begin{lstlisting}[style=opensg,caption={ff51\_rmc\_reform.yaml: the FF table as a self-describing file.}]
from opensg_shell.helper.ff_to_yaml import convert, station_ff

d = convert("ff51_rmc_reform.dat")                      # writes ff51_rmc_reform.yaml
FF, eta = station_ff("ff51_rmc_reform.yaml", eta=0.20)  # (6,) resultants, and 0.2
\end{lstlisting}

\texttt{convert(dat\_path, out\_base=None, frame="RM center-ref")} reads the
seven-column table, raising if it has fewer than seven columns, and writes
\texttt{<base>.yaml}. \texttt{station\_ff(yaml\_path, eta=None, row=None)}
returns the $(6,)$ resultant vector and the $\eta$ of one station, selected
either by nearest $\eta$ or by row index. The resulting
\texttt{ff51\_rmc\_reform.yaml} records the convention explicitly next to the
data, which is what makes the file readable months later:

\begin{lstlisting}[style=shellst]
source: ff51_rmc_reform.dat
convention:
  order: [F1, F2, F3, M1, M2, M3]
  frame: RM center-ref
  units: {force: N, moment: N.m}
stations:
- row: 0
  eta: 0.0
  FF: [-6.44900592e-06, -2.39422988e-08, 2236077.75, 1212518.5, -112030544.0, -4.11329552e-07]
- row: 10
  eta: 0.2
  FF: [41926.8867, 8267.39746, 1467335.12, 1219145.0, -60184972.0, 304264.594]
\end{lstlisting}

\section{Stage 5: dehomogenization back to 3-D wall stress}
\label{sec:windio:stage5}

\subsection*{What it computes}

A two-step chain, because a blade wall is a laminate inside a section:
\[
  \text{beam } \mathbf{FF}
  \;\xrightarrow[\text{RM shell ring}]{}\;
  \text{per-element plate forces}
  \;\xrightarrow[\text{RM plate SG}]{}\;
  \boldsymbol{\sigma}(z)\ \text{through the wall.}
\]

\emph{Step 5a, the section} (\texttt{opensg\_shell}). The macro beam strain is
$\bar{\boldsymbol{\varepsilon}} = \mathbf{C}_6^{-1}\mathbf{FF}$;
\texttt{\_macro\_fields(B, beam\_force\_vabs=FF)} solves it and recombines the
nodal RM warping fields $\mathbf{w} = \mathbf{V}_0\bar{\boldsymbol{\varepsilon}}
+ \mathbf{V}_1\boldsymbol{\varepsilon}'$ and their derivatives, and
\texttt{\_rm\_shell\_strain(B, e, xi, st\_m, aA, aB)} evaluates, for ring
element \texttt{e} at arc coordinate \texttt{xi}, the six shell strains
\[
  \mathbf{s}_6 = \begin{bmatrix} \varepsilon_{11} & \varepsilon_{22} &
  2\varepsilon_{12} & \kappa_{11} & \kappa_{22} & 2\kappa_{12}\end{bmatrix}.
\]
That wall's own MSG ABD turns them into the \textbf{plate reaction forces}
\[
  \mathbf{F}_6 = \mathbf{A}_6\,\mathbf{s}_6 =
  \begin{bmatrix} N_{11} & N_{22} & N_{12} & M_{11} & M_{22} & M_{12}\end{bmatrix},
\]
the interface quantity handed down to the plate SG.

\emph{Step 5b, the wall} (\texttt{opensg\_solid.rm\_plate\_1D}).
\texttt{rm\_plate\_msg(thick, angles\_deg, mat\_names, material\_db,
fraction=frac)} builds the through-thickness MSG-RM structure gene for each
layup once, and \texttt{msgrm\_strain\_at\_depth} -- or its vectorized twin
\texttt{msgrm\_strain\_at\_depth\_batch}, which evaluates many points of one
laminate in a handful of einsums -- returns the 3-D Voigt strain and stress at
depth $z$. Both default to \texttt{frame="material"}, meaning the result is
rotated into the ply material frame where failure criteria live; pass
\texttt{frame="plate"} to stay in the wall frame.

Recovery needs strain \emph{gradients}, and they come from the ring itself. The
span gradient $\partial\mathbf{E}/\partial x_1$ comes from the macro recovery
ladder, computed per element as \texttt{BDe @ st\_cl1 + BDh @ aB[g]} with the
element operators from \texttt{quad\_ops\_indep}. The arc gradient
$\partial\mathbf{E}/\partial x_2$ comes from nodal differencing of
$\mathbf{s}_6$ between the two elements sharing a node -- but only where those
two elements carry the \emph{same} layup, because a layup change is a physical
jump and must never be differenced across. Passing the second gradients as well
(\texttt{dE11}, \texttt{dE12}, \texttt{dE22}) selects the second-order (V2)
recovery, the rung that can in principle carry the transverse pair
$\sigma_{33}$, $\sigma_{13}$ that first order cannot.

\subsection*{What you edit}

The user block at the top of \texttt{2\_dehom\_v2\_vabs.py}:

\begin{lstlisting}[style=opensg]
############### User Input #################################
name = "iea_s10"                     # IEA-22 station at r/R = 0.20
ff_file = "ff51_rmc_reform.dat"      # BeamDyn FF (RM center-ref, VABS order)
station_row = 10                     # row 10 -> eta = 0.20
vabs_sm = "iea_s10.sg.SM"            # the VABS 2-D solid reference cloud
jpath_file = "iea_s10.lp_sparcap_left_thickness.coords"
############################################################
\end{lstlisting}

\begin{lstlisting}[style=shellst]
python 2_dehom_v2_vabs.py
\end{lstlisting}

\subsection*{What comes out}

\begin{table}[htbp]
\centering
\small
\caption{Files written by \texttt{2\_dehom\_v2\_vabs.py}.}
\label{tab:windio:stage5out}
\begin{tabular}{ll}
\toprule
File & Content \\
\midrule
\texttt{iea\_s10\_report.txt} & the validation report: the FF vector, the Timo diagonal, rms \\
  & differences against VABS split web/skin, the V2 payload block, \\
  & and the through-thickness path tables \\
\texttt{iea\_s10\_elem\_plate\_forces.dat} & 310 rows, one per ring element: element index, \texttt{y2 y3}, the six \\
  & shell strains, and the six plate reaction forces \\
  & $N_{11}, N_{22}, N_{12}$ [N/m] and $M_{11}, M_{22}, M_{12}$ [N] \\
\texttt{iea\_s10\_dehom\_v2.txt} & 166\,302 rows, one per VABS gauss point: \texttt{y2 y3 elem z}, then \\
  & the six recovered V2 stresses and the six VABS stresses, all MPa \\
\texttt{iea\_s10\_dehom\_v2.npz} & the same fields plus \texttt{C6}, \texttt{FF}, \texttt{s6mid}, \texttt{F6mid} and every strain \\
  & gradient, for re-plotting without re-solving \\
\texttt{iea\_s10\_v2\_section.png} & section contours, VABS beside RM V2, for $\sigma_{11}$, $\sigma_{12}$, $\sigma_{13}$, $\sigma_{33}$ \\
\texttt{iea\_s10\_parity.png} & parity clouds, first order and V2 against VABS \\
\texttt{iea\_s10\_thickness\_paths.png} & through-thickness profiles on a spar-cap wall, a web, and the \\
  & station's own reference path \\
\bottomrule
\end{tabular}
\end{table}

The first two data rows of \texttt{iea\_s10\_elem\_plate\_forces.dat} show what
the interface quantity looks like in practice:

\begin{lstlisting}[style=shellst]
# per-element PLATE REACTION FORCES from the RM shell beam model (Timo), eta=0.20
# elem  y2(m) y3(m)  e11 e22 2e12 k11 k22 2k12  N11 N22 N12(N/m) M11 M22 M12(N)
     0  4.677080e+00  1.048036e-01  5.426575e-04 -1.683049e-04  5.685206e-04 ...
\end{lstlisting}

\subsection*{What it was validated against}

\texttt{iea\_s10.sg.SM} is the VABS stress cloud of the \emph{same}
cross-section meshed as a 2-D solid, which is the independent reference: the
recovery is projected onto its gauss points and differenced there, 166\,302
points in all. Every point is assigned to its nearest ring element, a local arc
coordinate and a signed wall depth, and the VABS columns are permuted into the
$(11, 22, 33, 23, 13, 12)$ order and converted from Pa to MPa before the
comparison. Table~\ref{tab:windio:vabs} is copied from
\texttt{iea\_s10\_report.txt}, with the VABS rms alongside so each difference
has a scale.

\begin{table}[htbp]
\centering
\small
\caption{rms difference against VABS in MPa, web / skin, from
         \texttt{iea\_s10\_report.txt} at $\eta = 0.20$.}
\label{tab:windio:vabs}
\begin{tabular}{lrrr}
\toprule
Component & first order (web / skin) & V2 (web / skin) & VABS rms (web / skin) \\
\midrule
$\sigma_{11}$ & 43.073 / 9.529 & 43.073 / 9.529 & 44.440 / 96.676 \\
$\sigma_{22}$ & 0.612 / 0.482 & 0.612 / 0.482 & 0.704 / 0.998 \\
$\sigma_{12}$ & 6.799 / 1.255 & 6.798 / 1.255 & 22.026 / 6.009 \\
$\sigma_{33}$ & 0.081 / 0.049 & 0.081 / 0.049 & 0.081 / 0.049 \\
$\sigma_{13}$ & 0.394 / 0.192 & 0.394 / 0.192 & 0.394 / 0.193 \\
$\sigma_{23}$ & 0.035 / 0.028 & 0.035 / 0.028 & 0.035 / 0.028 \\
\bottomrule
\end{tabular}
\end{table}

Read this honestly, because the table says three different things at once.

\emph{The in-plane skin stress is good.} $\sigma_{11}$ in the skin differs by
9.53 MPa rms against a VABS rms of 96.68 MPa, and $\sigma_{12}$ by 1.26 against
6.01. The report also tabulates three through-thickness paths, and those
sharpen the picture considerably. On the station's own reference thickness path,
whose 15 points sit \textbf{0.00 mm} from the nearest VABS gauss point and are
therefore exactly in the reference frame, the $\sigma_{11}$ difference is
\textbf{1.265 MPa against a VABS rms of 286.837 MPa}, better than half a
percent. The two paths generated from the ring geometry sit 2.44 mm and 2.43 mm
from the nearest VABS point, and their larger differences -- 38.355 MPa rms on
the spar cap, 2.299 MPa on the web -- are dominated by that geometric offset
rather than by the recovery itself. The printed nearest-VABS distance is the
frame check: millimetres means the same frame, hundreds of millimetres means you
are comparing against a foreign mesh.

\emph{The web is the weak spot.} For $\sigma_{11}$ in the web the difference,
43.07, is essentially the VABS rms itself, 44.44, which means the web axial
stress is not being reproduced. The same pattern holds for the three transverse
components everywhere: the difference equals the VABS rms to three digits, which
is what you get when the recovered field is very nearly zero. This is the known
transverse recovery gap of the shell-to-plate route, not a tuning problem.

\emph{The second-order rung is inert at this load state.} The report's V2
payload block measures what V2 \emph{changes}, rather than a rounded difference
that could hide a small effect. At this station $\sigma_{11}$ moves by
$1.950\times10^{-4}$ MPa rms against a first-order rms of $8.851\times10^{1}$,
$\sigma_{12}$ by $5.685\times10^{-3}$ against $1.068\times10^{1}$,
$\sigma_{33}$ by $5.190\times10^{-5}$ MPa against a VABS rms of
$5.576\times10^{-2}$, and $\sigma_{13}$ and $\sigma_{23}$ by exactly
$0.000\times10^{0}$. For this constant-force blade state the V2 drivers carry no
payload, and the first-order $\sigma_{33}$ rms is $3.603\times10^{-13}$ MPa,
machine zero.

\subsection*{The correctness check on transverse shear}

\texttt{3\_q\_consistency\_check.py} tests an identity that any recovered
transverse shear must satisfy, wall by wall: the thickness integral of the
recovered stress must equal the wall's own transverse-shear resultant, which the
RM shell already knows.
\[
  I_a = \int_h \sigma_{a3}\,\mathrm{d}z
  \;\;\overset{!}{=}\;\;
  Q_a = \left(\mathbf{G}_{\mathrm{msg}}\,\mathbf{s}_2\right)_a ,
  \qquad
  \mathbf{s}_2 = \begin{bmatrix} 2\gamma_{13} & 2\gamma_{23}\end{bmatrix}.
\]
The plate pipeline enforces this by construction -- shape from the strain
gradients, amplitude from the resultant, $\sigma_{a3}$ rescaled by $Q/I$. The
shell-to-plate chain never uses $\mathbf{s}_2$, so nothing pins the amplitude,
and the ratio $r = I/Q$ measures how far off it is. Consistency means $r = 1$.

\begin{lstlisting}[style=opensg]
############### User Input #################################
name = "iea_s10"
ff_file = "ff51_rmc_reform.dat"
station_row = 10                  # eta = r/R = 0.20
n_z = 41                          # integration points through each wall
############################################################
\end{lstlisting}

\begin{lstlisting}[style=shellst]
python 3_q_consistency_check.py
\end{lstlisting}

The script writes \texttt{iea\_s10\_q\_consistency.txt} and
\texttt{iea\_s10\_q\_consistency.png}, the latter plotting recovered $|I_{13}|$
against wall $|Q_1|$ on log axes with the consistency line, beside the histogram
of $\log_{10}(I_{13}/Q_1)$. The measured result:

\begin{table}[htbp]
\centering
\small
\caption{Q-consistency ratios at $\eta = 0.20$, from
         \texttt{iea\_s10\_q\_consistency.txt}.}
\label{tab:windio:qcheck}
\begin{tabular}{lrrrrl}
\toprule
Ratio & median & mean & min & max & Verdict \\
\midrule
$I_{13}/Q_1$ & 5.228e$-$01 & $-$6.974e+00 & $-$1.299e+03 & 2.260e+01 & target invalid \\
$I_{23}/Q_2$ & 9.895e$-$01 & 1.096e+00 & $-$1.217e+01 & 1.194e+01 & target valid \\
\bottomrule
\end{tabular}
\end{table}

\textbf{$\sigma_{23}$ passes.} The median $I_{23}/Q_2$ is $0.9895$, a
\textbf{1.05\% Q-consistency error}, with a wall $Q_2$ rms of $8.9686$ N/m
against a recovered $I_{23}$ rms of $8.3085$ N/m.

\textbf{$\sigma_{13}$ is not measurable in a cross-section model}, and the
script explains why rather than papering over it. This ring is homogenized with
\texttt{shear="mitc4\_g23"}: $\gamma_{23}$ is tied by MITC while $\gamma_{13}$
is raw and carries the flat-wall drilling mode, so a $Q_1$ built from it is not
a physical wall shear and cannot serve as a recovery target. The fingerprint is
in the magnitudes -- $Q_1$ rms is $9.9064\times10^{2}$ N/m against a $Q_2$ rms
of $8.9686$ N/m, a factor of about 110, which is implausible for a thin-walled
section that carries its load as in-plane shear flow. The ring is a
one-quad-deep prismatic strip, so the $\gamma_{13}$ tying interpolates along a
direction in which the row is already linear and reproduces it exactly: the
script measures \texttt{rms |Q1(g23) - Q1(both)| = 0.0000e+00 N/m} and the same
for $Q_2$, so \texttt{mitc4\_g23} and \texttt{mitc4\_both} are bit-identical
here even though both flags are plumbed and branch. A physical $\gamma_{13}$
needs the surface-quad route, where both rows are genuinely tied, not a rescale
of this ring against a contaminated $Q_1$.

The practical rule that follows: trust $\sigma_{11}$, $\sigma_{22}$ and
$\sigma_{12}$ from this pipeline, treat $\sigma_{23}$ as good to about one
percent in amplitude, and do not use the cross-section route's $\sigma_{13}$ or
$\sigma_{33}$ as a design quantity.

\subsection*{The same recovery, all 51 stations}

\texttt{4\_sweep\_51\_stations.py} runs the \emph{first-order} recovery -- it
passes only the two first gradients to
\texttt{msgrm\_strain\_at\_depth\_batch} -- at \texttt{n\_depth = 9} depths per
wall element. The depths are mid-bin in the normalized thickness coordinate,
$\zeta = (k + \tfrac12)/9$ for $k = 0 \ldots 8$, converted to a physical depth
by $z = (\zeta - \texttt{frac})\,h$ with $h$ the total thickness of that
element's layup and \texttt{frac} taken from the bundle. The peak $|\sigma|$ per
component over the whole section, in the ply material frame, goes to
\texttt{iea51\_peak\_stress.dat}.

\begin{table}[htbp]
\centering
\footnotesize
\caption{Selected rows of \texttt{iea51\_peak\_stress.dat}: peak $|\sigma|$ per
         component over the whole section, ply frame, in MPa. Column order
         follows the file header.}
\label{tab:windio:peak}
\setlength{\tabcolsep}{4.5pt}
\begin{tabular}{lrrrrrr}
\toprule
$\eta = r/R$ & $\sigma_{11}$ & $\sigma_{22}$ & $\sigma_{33}$ & $\sigma_{23}$ & $\sigma_{13}$ & $\sigma_{12}$ \\
\midrule
0.00 & 1.49271417e+02 & 1.53636678e+00 & 4.13972884e$-$13 & 5.39458803e$-$03 & 4.12768841e$-$02 & 2.24186054e+00 \\
0.16 & 3.25793641e+02 & 4.14543750e+00 & 8.97794962e$-$13 & 3.91335971e$-$02 & 3.51941097e$-$02 & 7.01052831e+00 \\
0.20 & 3.18177433e+02 & 3.08795353e+00 & 9.12696123e$-$13 & 1.57149557e$-$02 & 3.45770096e$-$02 & 5.66031217e+00 \\
0.54 & 2.14382707e+02 & 1.62953363e+00 & 8.13510269e$-$13 & 1.74810053e$-$02 & 8.23774625e$-$02 & 5.64234617e+01 \\
0.80 & 1.11316749e+02 & 1.94984342e+00 & 1.93947926e$-$13 & 1.15167225e$-$02 & 5.58043894e$-$02 & 3.85749086e+01 \\
1.00 & 4.88165692e$-$01 & 2.71026227e$-$02 & 2.03726813e$-$15 & 4.70632777e$-$05 & 5.13380478e$-$02 & 7.89278421e+00 \\
\bottomrule
\end{tabular}
\end{table}

Reading the whole file rather than these six rows: the peak axial stress of the
entire blade in this load case is \textbf{325.79 MPa at $\eta = 0.16$}, and the
peak in-plane shear is \textbf{56.42 MPa at $\eta = 0.54$}, where the shear webs
take over from the caps. The $\sigma_{33}$ column never leaves the range
$2\times10^{-15}$ to $2\times10^{-12}$ MPa over all 49 stations -- identically
zero to round-off, the first-order limit noted above -- so do not read it as a
physical through-thickness stress. The $\sigma_{13}$ column, which stays between
$0.033$ and $0.105$ MPa over the whole blade, is subject to the same caveat as in the Q-consistency check and is
not a design quantity either.

\begin{figure}[htbp]
\centering
\includegraphics[width=0.95\textwidth]{\FIGDIR/iea51_sweep.png}
\caption{The two halves of the 51-station sweep. Left: the spanwise peaks of
         $\sigma_{11}$ and $\sigma_{12}$ against $r/R$. Right: the recorded
         per-station wall time, homogenization against dehomogenization.}
\label{fig:windio:sweep}
\end{figure}

Figure~\ref{fig:windio:sweep} shows both halves. The left panel plots the
spanwise peaks of $\sigma_{11}$ and $\sigma_{12}$: the axial peak climbs steeply
out of the root, crests near $\eta \approx 0.16$ and decays to almost nothing at
the tip, while the shear peak rises in distinct steps rather than smoothly, near
$\eta \approx 0.36$ and $\eta \approx 0.46$, where the peak jumps from 4.85 to
25.75 MPa and from 24.40 to 54.30 MPa respectively as a different wall takes
over the shear. The right panel records the per-station wall time of the two
halves separately, so you can see where a sweep's cost actually sits before
scaling it to more stations or more blades. The same arrays are saved to
\texttt{iea51\_sweep.npz} under the keys \texttt{eta}, \texttt{props},
\texttt{peak}, \texttt{t\_homo}, \texttt{t\_dehom} and \texttt{stations}, so a
sweep never has to be repeated just to redraw a figure.

\section{Where each tool ends}
\label{sec:windio:whodoeswhat}

\begin{table}[htbp]
\centering
\small
\caption{Who does what across the three tools in the round trip.}
\label{tab:windio:whodoeswhat}
\begin{tabular}{p{0.30\textwidth}p{0.34\textwidth}p{0.28\textwidth}}
\toprule
OpenSG\_io & OpenSG-2.0 & OpenFAST / BeamDyn \\
\midrule
Reads windIO v1 and v2, PreVABS XML, and OpenFAST ElastoDyn/BeamDyn blade data.
&
Homogenizes each 1-D shell SG to the Timoshenko $6\times6$ and writes
\texttt{<station>\_Timo.out}.
&
Runs the blade dynamics. Installed, configured and run by you, outside OpenSG.
\\[4pt]
Writes the per-station 1-D shell SG YAMLs that \texttt{opensg\_shell} consumes.
&
Writes the RM wall law per section, \texttt{<station>\_ABDG.out}, and the ABD
YAML with \texttt{mass\_per\_area}.
&
Returns the per-station sectional force and moment resultants that come back as
the FF table.
\\[4pt]
Writes the PreVABS XML for the 2-D solid route, which supplies the independent
VABS reference.
&
Recovers the 3-D wall stress from the FF table and validates it against that
VABS cloud.
&
Owns the blade input file format; consult the BeamDyn documentation for it.
\\[4pt]
Bridges the hand-off with \texttt{timo\_to\_beamdyn} and
\texttt{write\_beamdyn\_blade}, and reads OpenFAST data back as a validation
reference.
&
Runs the Q-consistency diagnostic on the recovered transverse shear.
&
Sectional mass matrices are not produced by either OpenSG tool and remain your
responsibility.
\\
\bottomrule
\end{tabular}
\end{table}

The chain closes only because the conventions match at every hand-off: the same
reference surface (\texttt{center}), the same component order (VABS), and the
same units (SI, metres and pascals in, newtons and newton-metres back). Those
three lines are the first things to check when a recovered stress field looks
wrong, and they are checkable without rerunning anything -- the station YAML
records its \texttt{reference}, the FF table records its frame and order in the
header line, and the ABD YAML records the reference it was emitted at.


\end{document}
